\documentclass[11pt,a4paper]{article}
\usepackage[T1]{fontenc}
\usepackage[utf8]{inputenc}
\usepackage{lmodern}
\ifdefined\pdfmapfile
  \pdfmapfile{=lm.map}
  \pdfmapfile{+cm.map}
  \pdfmapfile{+symbols.map}
  \pdfmapfile{+euler.map}
  \pdfmapfile{+rsfs.map}
\fi
\usepackage[a4paper,margin=27mm]{geometry}
\usepackage{amsmath,amssymb,amsthm,mathtools,mathrsfs}
\usepackage{booktabs,tabularx,longtable,array,enumitem}
\usepackage{microtype,xcolor,etoolbox}
\usepackage{hyperref}
\usepackage{xurl}
\hypersetup{colorlinks=true,linkcolor=blue!45!black,citecolor=blue!45!black,
urlcolor=blue!45!black,pdfauthor={Non-Archimedean Bedford--Taylor research record},
pdftitle={Unbounded energy, capacity, and relative Dirichlet theory}}
\usepackage[capitalise,noabbrev]{cleveref}
\numberwithin{equation}{section}
\newtheorem{theorem}{Theorem}[section]
\newtheorem{lemma}[theorem]{Lemma}
\newtheorem{proposition}[theorem]{Proposition}
\newtheorem{corollary}[theorem]{Corollary}
\theoremstyle{definition}
\newtheorem{definition}[theorem]{Definition}
\newtheorem{example}[theorem]{Example}
\theoremstyle{remark}
\newtheorem{remark}[theorem]{Remark}
\newcommand{\Cb}{\mathcal C_b}
\newcommand{\R}{\mathbf R}
\newcommand{\Q}{\mathbf Q}
\newcommand{\Z}{\mathbf Z}
\newcommand{\PSH}{\operatorname{PSH}}
\newcommand{\supp}{\operatorname{supp}}
\newcommand{\dd}{d'd''}
\newcommand{\an}{\mathrm{an}}
\setlist{topsep=4pt,itemsep=2pt,parsep=0pt}
\setlength{\parindent}{1.3em}
\setlength{\parskip}{2pt}
\setlength{\emergencystretch}{2em}
\makeatletter
\patchcmd{\l@section}{1.5em}{2.1em}{}{}
\renewcommand*\l@subsection{\@dottedtocline{2}{1.5em}{3.6em}}
\makeatother


\title{Unbounded Energy, Capacity, and Relative Dirichlet Theory\\[5pt]
\large Eleventh research record in non-Archimedean Bedford--Taylor theory}
\author{Continuation of the exact authorized R8--R10 research chain}
\date{11 November 2026 (UTC)}
\begin{document}
\maketitle
\begin{abstract}
Over the original finite extension of $\mathbf Q_p$ and with its original
absolute value, we develop an unbounded Bedford--Taylor calculus for
arbitrary original symmetric middle-positive separately closed backgrounds.
We prove coefficient-controlled closedness of non-polar products, an exact
no-loss capacity-tail criterion, unbounded comparison and domination,
all-degree power-energy recursion, and graded first-order and energy
stability on every original refined test. Exact level currents yield
coarea, nonsmooth chain formulas, weak maxima and quantitative pole
representatives. On explicit strict analytic boxes we obtain full global
first-energy Green symmetry, the constant-one mixed energy inequality,
the sharp midpoint modulus, and a complete relative strong first-energy
space without a prescribed actual lower barrier. It has comparison,
domination, intrinsic truncation carriers, and bounded finite-mass
approximation. Explicit dominated data have a unique variational
maximizer on this full space. A separate barrier completion gives an
unbounded constrained variational retraction.
In remaining degree one we solve
the complete finite signed diffuse renormalized zero-boundary problem,
including exact level carriers and global total-variation convergence of
clipping curvatures, and its finite-diffuse obstacle problem. Counterexamples
preserve the boundary, integrability,
endpoint and classical-admissibility gates. Frozen results and rediscoveries
are distinguished from new proofs. General nonlinear classical Dirichlet
existence and the analytic PS/PL quotient remain unresolved in their
appropriate, explicitly stated scopes.
\end{abstract}
\tableofcontents
\newpage

% BEGIN REVIEWED MODULE research_r11/scope_and_main_results.tex
\section{Fixed scope, proof dependencies, and the R11 advance}
\label{sec:r11-scope-main-results}

\subsection{The original analytic category is unchanged}

Throughout, $K$ is the fixed original finite extension of
$\mathbf Q_p$, equipped with its original absolute value.
The space $X$ is good, Hausdorff, boundaryless, quasi-smooth,
and pure of dimension $n$.  Opens are ordinary analytic
opens.  We write $D=d'd''$.
Unless a statement explicitly specializes its degree,
$S$ is an arbitrary original symmetric middle-positive,
separately $d'$- and $d''$-closed $(p,p)$-current, and
$q=n-p\geq1$.
It is neither replaced by a smooth or geometric current
nor approximated by a special background class.
Smooth and integration backgrounds appear only in
explicit counterexamples to universal assertions.

All current identities are identities on every original
compactly supported smooth CLD complementary test form.
Every actual finite moment refinement of such a test
is permitted.  Middle, strong, and weak positivity cones
are distinct: rank-one squares and strongly positive
forms serve as coefficient controls, while the actual
positive-current assertions use the middle-positive
background and the correct complementary cone.
No ordinary current equality is automatically an
equality on an additional family of GK delta tests.

The terms classical psh, bounded BT product, original
coefficient-Radon measure, and classical extension to
complete-field curve tests retain their frozen meanings.
The seven bounded BT properties, relative CLN and
bounded energy calculus are \texttt{INTERNAL\_INPUT}.
For a locally bounded classical tuple the notation is
\[
 T_{S,k}(u_1,\ldots,u_k)
              =Du_1\wedge\cdots\wedge Du_k\wedge S.
 \tag{M1}
\]
Every new unbounded object is constructed and justified
from those bounded inputs, with its own stated mass or
energy hypothesis.  Multiplication by an unbounded
potential is never performed before coefficient
integrability has been proved.

\subsection{The exact frozen continuation point}

The immediate baseline is the complete R10 archive,
whose exact SHA256 is
{\small\nolinkurl{4891735a31504e5154d7e027efb2fdf739dbf4645fb963f45a4b4943360b1402}}.
Its 226 listed manifest files and ZIP checks passed.
The archive is included unchanged under
\path{previous/NA_BT_R10_Complete_Archive_2026-11-11.zip}.
Its nested R9 archive has SHA256
{\small\nolinkurl{dafae83ac877f4bd5648ef7939de39f9f5bf114c6e7a711fe3f6732a3f3a68a4}};
the original authorized R8 archive has SHA256
{\small\nolinkurl{378c2ef39f98fb60dca80f3300cb6ec0bf5f0c370a6f93adc95714a8cd5c3254}}.
The complete earlier chain is retained through those exact archives.

R11 does not rewrite the matroidal, two-fan, polytopal,
real-product, actual quadratic-cover, intrinsic PS
curvature, or cap-order results already frozen in that
chain.  The internal research papers R, B and A2 keep
their authorized \texttt{INTERNAL\_INPUT} status.
A2's completed geometric adelic nefness is not reopened.
Targeted source lookbacks in this round concern only
a specific dependency or novelty claim; they are not
new audits of the old inputs.

Each formal result below prints its status.
\begin{center}
\begin{tabular}{ll}
\texttt{EXTERNAL\_KNOWN}&\texttt{INTERNAL\_INPUT}\\
\texttt{DIRECT\_COROLLARY}&\texttt{NEW\_PROVED}\\
\texttt{CONDITIONAL}&\texttt{OPEN}\\
\texttt{FALSE}&\texttt{FAILED\_ATTEMPT}
\end{tabular}
\end{center}
The last status preserves discarded routes and valid rediscoveries.
A new definition is distinguished from a theorem.
The word new means a proved R11 delta relative to the
frozen research chain, not a claim of priority over
the entire mathematical literature.

\subsection{Three domain levels and three distinct boundaries}

The new results have different domain hypotheses.
They must not be merged into an unconditional global
theorem by dropping those hypotheses.
\begin{enumerate}
\item The coefficient-controlled non-polar closedness
and general local mixed-energy criteria are ordinary-local
statements on the allowed analytic space.
\item The all-degree capacity controls require the
specified ordinary domain to have compact closure
inside the interior of one original affinoid.
The analytic compensator theorem proves why this
hypothesis controls every original refined test.
\item Global finite-mass tester approximation, full
$E^1$ Green calculus and the explicit simultaneous
pole representatives use a strict analytic Laurent
box, with the original $S$ defined on a neighborhood
of its closed box.  Arbitrary currents on arbitrary
opens are not asserted to extend across their boundary.
\end{enumerate}

For a negative extended classical potential $u$,
the compact-sublevel condition is
\[
 \mathrm{BG}(u,0):
       \quad\overline{\{u<-t\}}\Subset V
                              \quad\hbox{for every }t>0.
 \tag{M2}
\]
More generally, the comparison theorems use exactly
the printed compact positive-margin gaps.
This is an ordinary boundary condition.
It differs both from finite truncation loss and from
the zero-boundary condition in a relative energy
completion.  The last is specified by the closure
of the original compact smooth core and, for
unbounded measure data, by clipping membership and
vanishing radiation.

\subsection{Non-polar products and the exact no-loss cone}

For an actual negative classical tuple, put
$u_{i,j}=\max(u_i,-j)$ and
$E_j=\bigcap_i\{u_i>-j\}$.
The finite-level restrictions
\[
 {\bf1}_{E_j}\,
       T_{S,k}(u_{1,j},\ldots,u_{k,j})
 \tag{M3}
\]
form a consistent increasing family of positive
coefficient measures.  Theorem
\ref{thm:r11-closed-nonpolar-product}
proves that their locally coefficient-finite limit
is separately closed on every original test.
The proof controls derivatives of bounded cutoffs
by a positive energy identity.  It does not assert
that an individual restricted current is closed.
The pole union need not be ordinary closed.

For one potential, write
$\mu_{u,j}=T_{S,q}(u_j,\ldots,u_j)$,
$\mu_u$ for its non-polar top measure, and
$\nu_{u,j}={\bf1}_{\{u\leq-j\}}\mu_{u,j}$.
The no-loss boundary cone is
\[
 \mathcal F_{S,\partial}(V)
   =\{u\leq0:\ u\text{ extended classical psh},
       \ \mathrm{BG}(u,0),\ \nu_{u,j}(V)\longrightarrow0\}.
 \tag{M4}
\]
The exact capacity-tail criterion is
\[
 u\in\mathcal F_{S,\partial}(V)
 \quad\Longleftrightarrow\quad
       t^q C_S(\{u<-t\};V)\longrightarrow0
 \tag{M5}
\]
within the printed boundary class.
Here the raw capacity is computed from the entire
normalized bounded classical tester family on the
same domain; its ordinary-open outer regularization
is denoted $C_S^*$.

The new comparison theorem needs no loss on the lower
potential and the stated compact-gap condition.
It yields the corresponding domination and uniqueness
modulo outer relative capacity.
Every compact-sublevel classical potential, even
before imposing no loss, has the proved
$O(t^{-q})$ deep-sublevel capacity bound and is
locally quasicontinuous for the relative outer
capacity.  General raw/outer equality is not
an input to these results.

\subsection{Power classes and the all-degree calculus}

For $\sigma>0$ the power class is
\[
 \mathcal E^\sigma_{S,\partial}(V)
     =\{u\in\mathcal F_{S,\partial}(V):
                  e_\sigma(u):=\int_V(-u)^\sigma\,d\mu_u<\infty\}.
 \tag{M6}
\]
The capacity integral
\[
 J_\sigma(u)=
       \int_0^\infty t^{q+\sigma-1}
                        C_S(\{u<-t\};V)\,dt
 \tag{M7}
\]
characterizes this class with the stated boundary
hypothesis.  It controls mixed energies, exact
truncation tails, and nonmonotone stability in the
full localized capacity topology (PE16).
Finite non-polar energy alone, without no loss,
does not define this class.

On the one-affinoid domain, every finite tuple in a
positive power class has closed positive original
mixed products in every degree.  The complete
recursive equation holds for every $\sigma>0$:
\[
                D(u_iN_\alpha)=N_{\alpha+e_i},
                   \qquad |\alpha|\leq q-1.
 \tag{M8}
\]
All original coefficient multipliers are justified
before differentiation.  With $k=|\alpha|$, the
natural exponent budgets are
\[
\begin{array}{c|c}
\text{object with weight growth }W^a&
                    \text{fixed-family integrability budget}\\ \hline
 N_\alpha&a+k\leq q+\sigma\\
 d'(u_iN_\alpha),\ d''(u_iN_\alpha)&a+k+1\leq q+\sigma\\
 \Gamma_{N_\alpha}(u_i,u_l)&a+k+2\leq q+\sigma .
\end{array}
 \tag{M9}
\]
Energy terms are asserted only within their printed
degree range.  The endpoint budget is valid for a
fixed family, and for varying families under
uniform tail integrability.
A uniform bound of power energies alone gives
strict-exponent stability, not unconditional
endpoint continuity.
The explicit shell example sends the critical
weighted top measures to a nonzero rigid-point
mass although the potentials converge in the
specified localized capacity to zero.

On the strict boxes, the finite-mass boundary core
computes every raw Borel capacity exactly.
The new full global $E^1$ Green symmetry then gives
the optimal constant-one inequality
\[
 \int_V(-u_0)M(u_1,\ldots,u_q)
              \leq\prod_{i=0}^{q}e_1(u_i)^{1/(q+1)} .
 \tag{M10}
\]
It also gives global Aubin--Mabuchi variation and
strict concavity on the relative capacity quotient.
This is an integrated energy argument using the
positive energy matrix.  It does not assume a
pointwise Alexandrov--Fenchel inequality for $S$.

\subsection{Relative completions and the sharp strong geometry}

On those same strict boxes, put
\[
 \mathscr I_S(u,v)=\int_V(u-v)(d\mu_v-d\mu_u),\qquad
 \mathscr A_S(u)=-\frac{e_1(u)}{q+1}.
 \tag{M12}
\]
The new first-energy tester estimate gives quantitative
control of every local original outer capacity by
$\mathscr I_S(u,v)^{2^{-q}}$ on bounded-energy sets.
The sharp uniform midpoint inequality is
\[
 \mathscr A_S\left(\frac{u+v}{2}\right)
       -\frac{\mathscr A_S(u)+\mathscr A_S(v)}2
      \geq\frac{1-2^{-q}}{2(q+1)}\mathscr I_S(u,v).
 \tag{M13}
\]
Its constant is optimal on every same-background class
containing a potential of positive first energy.
The proof is the all-degree finite binomial argument,
not an extrapolation from a finite numerical check.

The barrier completion first constructs Borel limits
above an actual $E^\sigma$ lower barrier, retaining
the original closed currents and their complete
recursive and graded calculus. Its comparison and
domination hold for the full normalized tester capacity;
strict Borel maximum locality and the exact nonpolar
truncation carriers are proved after the construction.
A completed envelope and an unbounded barrier-constrained
variational retraction are then defined without ordinary
upper regularization.

There is also a larger strong first-energy completion
$\widehat{\mathcal E}^1_{\rm str}$, defined by actual
$\mathscr I_S$-Cauchy rows without a prescribed common lower barrier.
The new uniform deep-tail and global boundary-tightness
proofs construct its actual Borel representatives, every
original current and critical multiplier, and its global
mixed energies. The resulting space is complete and satisfies
\[
 \widehat{\mathscr I}_S(u_j,u)\longrightarrow0
 \quad\Longleftrightarrow\quad
 \begin{cases}
 u_j\longrightarrow u\text{ in all local truncated outer capacities},\\
 \widehat e_1(u_j)\longrightarrow\widehat e_1(u).
 \end{cases}
 \tag{M14}
\]
Completeness uses an explicitly proved uniform continuity
modulus for the discrepancy, not an assumed triangle
inequality. The limit is not asserted to be an ambient
classical psh function or to have compact negative sublevels.
Its boundary class is the printed strong approximation class.

The characterization can also be made intrinsic: every bounded
truncation must be an admissible same-domain bounded capacitary
limit, and
\[
                  \int_0^\infty t^q
                        C_S^*(\{u<-t\};V)\,dt<\infty.
 \tag{M16}
\]
These two conditions are equivalent to membership in the
strong first-energy space. The proof first constructs legal
finite-mass barrier stages and only then the unbounded
current; it does not assume the sought product.
Comparison, domination, strict maximum locality and intrinsic
nonpolar carriers hold throughout this space. Increasing
sequences converge strongly; decreasing sequences do so under
the explicit uniform first-energy bound.

Every nonempty family in this strong space has a relative
Perron least upper bound, obtained by minimizing energy over
its finite maxima. No common actual lower barrier is needed.
The countable selected increasing row preserves a common
source lower bound and Borel upper constraints. This gives
an envelope with a retained subsolution inequality, not
an unrestricted source equality.

For $0\leq\mu\leq A\mu_\psi$ with one actual
$\psi\in\mathcal E^1$, the full strong-space functional
$-\widehat e(u)/(q+1)-\int u\,d\mu$ has a unique relative
maximizer, even when $\mu(V)=\infty$. Its exact output is
\[
 \int(v-u)\,d(\widehat\mu_u-\mu)\leq0
                 \quad(v\in\widehat{\mathcal E}^1_{\rm str}),
 \qquad
 \widehat e(u)\leq A^{(q+1)/q}e_1(\psi).
 \tag{M17}
\]
The energy bound is attained when $\mu=A\mu_\psi$.

This enlargement does not solve the unrestricted source
equation. In the explicit remaining-degree-one example,
the bounded first energies of any defining row give local
$L^1$ bounds on its convex Gauss traces through a patched
quadratic tester. Convexity then gives a finite local uniform
trace limit. The frozen compact-source rigidity survives
the completion and rules out the nonzero compact datum
(DO9) throughout $\widehat{\mathcal E}^1_{\rm str}$.
Thus the distinction between convergence, a variational
maximum, and a source equation is mathematically necessary.

\subsection{Level currents, poles, and the relative linear problem}

The new level-current calculus separates exact
equality-level mass from strict maximum locality.
It proves a nonsmooth scalar chain rule, coarea,
positive equality-level defects for maxima, and
finite subsolution maximum closure.
Below top degree, a positive level current generally
fails to be separately closed; an explicit original
torus calculation identifies its defect.
Such a current is never inserted as a closed BT
background without a separate proof.

For a compact-sublevel pole set on a strict box,
logarithmic attenuation followed by protected-core
gluing constructs one actual pole representative
in every positive power class.
The function is independent of the background
and works simultaneously for every allowed
extending original $S$, with background-dependent
constants.  Its exponential capacity and coefficient
tails are quantified.
This concerns the precisely specified family of
complete pole sets.  It is not an unrestricted
identification of relative capacity-null sets with
ambient pluripolar or negligible sets.

When $q=1$, the analytic affinoid compensator
removes the old algebraic-domain restriction from
the faithful original energy form.
For every finite signed energy-diffuse datum $\mu$,
the relative zero-boundary problem has a unique
renormalized solution $u$, satisfying, on every
original compact smooth scalar test,
\[
 \begin{split}
 D(uS)&=\mu,\\
 D(\Phi(u)S)&=\Phi'(\widetilde u)\mu
                  +\Phi''(\widetilde u)\Gamma_S(u),\\
 \|D(T_ku\,S)-\mu\|_{\rm TV}
                       &\leq2|\mu|\{|\widetilde u|\geq k\}\longrightarrow0.
 \end{split}
 \tag{M11}
\]
The chain formula has the printed bounded-derivative
and compact-second-derivative hypotheses.
The compatible unweighted energy of an arbitrary
measure-data solution is only sigma-finite in
general; a new explicit example proves that it
need not be locally finite.
Exact level carriers and coarea remain valid on
bounded scalar strips.

For total-variation convergent finite signed diffuse data,
every fixed clipping of the solution converges strongly
in the original part energy norm. Its energy measures and
the stated smooth scalar curvatures converge in global
total variation. Curvature after a fixed hard clipping
need not have that stability: the explicit original
annulus example moves nonzero level atoms.

The finite signed diffuse obstacle theorem uses
$A=-D(\,\cdot\,S)$ and an allowed obstacle $\psi\in H$
with finite signed curvature $\tau=A\psi$ and finite-energy
variation. It gives a unique renormalized pair
\[
 Aw=\mu+\rho,\qquad w\geq\psi,\qquad
 \rho\{\widetilde w>\widetilde\psi\}=0,\qquad
 0\leq\rho\leq(\tau-\mu)^+.
 \tag{M15}
\]
The original finite-diffuse data class is unrestricted
within those stated hypotheses, while the obstacle and
part-boundary assumptions remain explicit.

Only one additional external measure-theoretic
input is needed at the measure-data extension:
the finite-zero-energy nest lemma for a smooth
measure of a transient quasi-regular Dirichlet
form, in Klimsiak--Rozkosz,
\emph{On the structure of bounded smooth measures
associated with a quasi-regular Dirichlet form},
arXiv:1410.4927v2, Lemma 2.1.
Its hypotheses are checked in the original
faithful form.  It does not supply an original
CLD current equation, a renormalized chain,
or an ambient classical psh solution; those
distinctions are maintained in the proofs.

\subsection{What the enlarged theory still does not assert}

The new theorems give a substantial unbounded
arbitrary-background part of the classical BT
program, together with a complete, precisely
specified relative linear measure-data problem.
They do not complete the unrestricted nonlinear
Monge--Ampere Dirichlet theory.
In fact the frozen compact-source obstruction
extends to the new unbounded power classes:
finite energy-dual control, even domination by
one normalized smooth tester, does not ensure
an ambient classical solution for arbitrary $S$.
An energy solution and a classical psh solution
are different admissible classes.

In remaining degree one the strong completed classical-approximation
cone is identified isometrically with a closed cone in the
original part Hilbert space. Its exact finite-dual source
criterion is membership of the Riesz vector in that cone,
or the equivalent full polar-cone inequalities. The source
projection and its quantitative defect are direct Hilbert
consequences. Although original smooth compact directions
are differences of finite-mass cone members, this does not
make both signs feasible at a specified maximizer.

The full analytic PS/PL relation quotient GQ,
low-real-degree signed CTO, $W_{\rm loc}$ density,
and the auxiliary-field quotient remain open.
The relative capacity quotient proved for
actual power-class potentials does not identify
finite analytic delta presentations.
The closing section lists the precise remaining
gates, false shortcuts, and dependency boundaries.

% END REVIEWED MODULE research_r11/scope_and_main_results.tex


% BEGIN REVIEWED MODULE research_r11/dirichlet/analytic_affinoid_energy_domain.tex
\section{Faithful relative energy on arbitrary original affinoid neighborhoods}
\label{sec:r11-affinoid-energy}

Throughout this section the field is the fixed finite extension
$K/\mathbf Q_p$, with its original absolute value.  Let $X$ be good,
Hausdorff, boundaryless, quasi-smooth, and pure of dimension $n\geq1$.
Choose a compact $K$-affinoid domain $B\subset X$ and an ordinary open
set $\Omega$ such that
\[
                 \overline\Omega\subset\operatorname{Int}_X(B).
 \tag{AE1}
\]
The closure is compact.  We impose no algebraizability hypothesis.
Fix an arbitrary original symmetric middle-positive separately
$d'$- and $d''$-closed $(n-1,n-1)$-current $S$ on $\Omega$.
All coefficients, restrictions and tests are original CLD objects.
The frozen bounded BT calculus, coefficient-Radon calculus, compact
cutoffs, and exact countable original moment reduction are
\texttt{INTERNAL\_INPUT}.  In particular, the last input is the
proved R10 Lemma \texttt{r10-exact-countable-moment-reduction},
not a density assertion for background currents.

\begin{lemma}[Exact replacement by quotients of functions on one affinoid]
\label{lem:r11-affinoid-rational-unit-replacement}
Let $x\in\operatorname{Int}_X(B)$, let $W$ be an ordinary neighborhood
of $x$, and let $u_1,\ldots,u_m$ be analytic units on $W$.
There are a compact Laurent domain $C\subset B\cap W$ with
$x\in\operatorname{Int}_X(C)$ and elements
$b_i,c_i\in\mathcal O(B)$ such that both $b_i$ and $c_i$ are
nonvanishing on $C$ and
\[
              |u_i|=|b_i/c_i|\quad\hbox{at every point of }C.
 \tag{AE2}
\]
The domains and quotients required for this assertion can be chosen
from one countable family depending only on $B$.
Consequently every finite original unit moment, including an
additional actual refinement, has this exact local replacement.
\end{lemma}

\noindent\textbf{Status: \texttt{NEW\_PROVED}} as the
global-affinoid-source strengthening of the frozen countable reduction.

\begin{proof}
Put $A=\mathcal O(B)$.  The Banach algebra $A$ is separable:
a generalized Tate algebra over the separable metric field $K$
has a countable dense polynomial set, and this property passes
to a quotient.  Choose a countable dense subring $A_0\subset A$.
The evaluation functions $|f|$, $f\in A_0$, generate the topology
of $B$, since evaluations of arbitrary elements of $A$ are
uniform limits of these evaluations.

As in the exact countable reduction, finitely many closed Laurent
inequalities in elements of $A_0$, with strictly interior positive
rational bounds at $x$, give $C\subset B\cap W$ whose ordinary
interior contains $x$.  Its defining lower-bound functions are
nonvanishing on $C$.  If these functions are $g_1,\ldots,g_s$,
then the image of
\[
                       A[g_1^{-1},\ldots,g_s^{-1}]
                           \longrightarrow\mathcal O(C)
 \tag{AE3}
\]
is dense.  Indeed the Laurent localization is the quotient of a
generalized Tate algebra obtained by adjoining the bounded upper
variables and the inverses of the $g_j$; polynomial truncations
are dense, and substitution for these generators gives (AE3).
Approximating the finitely many coefficients in $A$ by $A_0$
preserves density.  A polynomial in the inverses is one quotient
$b/c$ with $b\in A_0$ and $c$ a product of the $g_j$.

Choose such a quotient $r_i=b_i/c_i$ so that, in a fixed Banach
norm on $\mathcal O(C)$,
\[
                     \|u_i-r_i\|\,\|u_i^{-1}\|<1.
 \tag{AE4}
\]
Then $r_i=u_i(1+(r_i-u_i)u_i^{-1})$ is a unit by its Neumann
series.  The denominator is already a unit, so $b_i$ is
nonvanishing on $C$.  For each $z\in C$ one has
$|u_i(z)-r_i(z)|<|u_i(z)|$, and the ultrametric inequality
gives (AE2).  No approximation remains in the logarithmic
functions or their CLD differentials on the ordinary interior.
There are countably many possible Laurent domains, denominators
and numerators from $A_0$, proving the last assertion.
\end{proof}

\begin{lemma}[A global analytic compensator]
\label{lem:r11-affinoid-compensator}
Let $H\Subset\Omega$ and let $t_i=\log|r_i|$ be finitely many
original moment functions near $H$, with
$r_i=b_i/c_i$ as in (AE2) on a neighborhood of $H$.
There are an original smooth classical psh function $a$ on
$\operatorname{Int}_X(B)$ and a finite $N>0$ such that
\[
 0\leq a\leq1\ \hbox{on }\Omega,\qquad
 Da\geq N^{-1}\sum_i d't_i\wedge d''t_i\ \hbox{near }H,
 \qquad d'a\wedge d''a\leq2a\,Da .
 \tag{AE5}
\]
All inequalities are inequalities of original positive $(1,1)$-forms.
\end{lemma}

\noindent\textbf{Status: \texttt{DIRECT\_COROLLARY}} of the
new exact replacement and the explicit R4 compensator construction.
The following proof records why analytic functions on $B$ suffice.

\begin{proof}
Choose a smooth nonnegative $\kappa$ of integral one supported in
$(-1,0)$ and set
\[
 \zeta(s)=\int_{-\infty}^{s}(s-r)\kappa(r)\,dr,\qquad F=\zeta^2.
\]
Then $F$ is smooth, convex and nondecreasing, vanishes for $s\leq-1$,
has $F''=2$ for $s\geq0$, and $(F')^2\leq2FF''$.
Choose $A>0$ so that $\log|b_i|+A$ and $\log|c_i|+A$
are positive near $H$.  The functions
$F(\log|b_i|+A)$ and $F(\log|c_i|+A)$, extended by zero
near the corresponding analytic zero loci, are original smooth
and classical psh on $\operatorname{Int}_X(B)$.
This uses the frozen convex nondecreasing composition rule on
every complete-extension curve, including its constant zero branch.
Their sum $g$ is bounded on the compact $\overline\Omega$ and satisfies
\[
 Dg\geq\sum_i d't_i\wedge d''t_i,\qquad
 d'g\wedge d''g\leq2g\,Dg.
\]
The first inequality follows from $(\xi-\eta)^2\leq
2\xi^2+2\eta^2$; the second is scalar Cauchy--Schwarz applied
to the nonzero summands and then extended by smoothness.
Take $N=\max(1,\sup_{\overline\Omega}g)$ and $a=g/N$.
\end{proof}

\begin{theorem}[A faithful coercive original form without algebraizability]
\label{thm:r11-general-affinoid-faithful-form}
Under (AE1), there is a bounded continuous classical psh function
$0\leq h\leq1$ on $\Omega$ for which
\[
                      \lambda=D(hS)
 \tag{AE6}
\]
is a positive Radon measure with
\[
 \lambda(E)=0\ \Longleftrightarrow\ {\bf1}_E S=0
 \quad\hbox{for every Borel }E\subset\Omega.
 \tag{AE7}
\]
For every original compact smooth complementary $(1,1)$-form
$\eta$ one has $|\mu_{S,\eta}|\leq C_\eta\lambda$.
Writing $F_S=\operatorname{supp}\lambda=\operatorname{supp}S$, the
form on restrictions of real original smooth compact functions
\[
 \mathcal E(f,g)=\int_\Omega
       \tfrac12(d'f\wedge d''g+d'g\wedge d''f)\wedge S
 \tag{AE8}
\]
is closable on $L^2(F_S,\lambda)$.  Its closure
$(\mathcal E,\mathscr D)$ is a regular, strongly local,
symmetric Dirichlet form, and
\[
 \mathcal E(f,f)\geq\int_{F_S}\frac{f^2}{2-h}\,d\lambda
                         \geq\tfrac12\|f\|_{L^2(\lambda)}^2 .
 \tag{AE9}
\]
For every original compact smooth real $g$,
\[
         D(gS)=b_g\lambda,\quad b_g\in L^\infty(\lambda),
 \qquad \Gamma_S(g)\leq C_g\lambda ,
 \tag{AE10}
\]
with compactly supported densities.  All assertions include every
original moment refinement.  They hold for arbitrary original $S$;
no special-current approximation occurs.
\end{theorem}

\noindent\textbf{Status: \texttt{NEW\_PROVED}} as the removal of the
algebraic-domain assumption in the frozen faithful form theorem.
The form-theoretic argument is the same proved R4 argument once
the global analytic compensators have been supplied.

\begin{proof}
The countable replacement lemma, compact exhaustions of its
countable unit charts, and original smooth cutoffs give countably
many positive compact controls
$\Theta_l=\chi_l\sum_i d't_{li}\wedge d''t_{li}$, $0\leq\chi_l\leq1$,
cofinal for all original compact complementary coefficients.
This means precisely that every such coefficient variation is
bounded by a finite sum of the measures $\Theta_l\wedge S$
with finite positive constants.  It follows by writing the form
in a finite moment presentation and dominating its symmetric
matrix by a multiple of the positive identity matrix; a finite
partition handles its compact support.

Choose $a_l,N_l$ by (AE5), so that
$D(a_lS)\geq N_l^{-1}\Theta_l\wedge S$.  Put
\[
 h=\sum_{l\geq1}2^{-l}a_l,\qquad
 h_m=\sum_{l=1}^m2^{-l}a_l+2^{-m}.
\]
The sequence $h_m$ decreases uniformly to $h$.  Frozen relative
CLN bounds give locally finite mass to the increasing positive
series $\sum_l2^{-l}D(a_lS)$, and the original uniform product
identity identifies this series with $D(hS)$.  Each control is
bounded by $2^lN_l\lambda$, giving the stated coefficient
domination.  Conversely every $D(a_lS)=Da_l\wedge S$ is a
coefficient measure of $S$ locally.  This proves (AE7), its
support consequence, and (AE10).

For clarity, the analytic extension introduces no extra
form-theoretic premise.  The support $F_S$ is locally compact
and second countable.  Original compact smooth scalars separate
points, by smooth truncations of affinoid evaluation functions
and compact cutoffs.  Stone--Weierstrass gives the required
regular-core density.  Compact Stokes gives
\[
                  \mathcal E(f,g)=-\int f b_g\,d\lambda
                      \quad(f,g\hbox{ in the original core}).
 \tag{AE11}
\]
If $f_j\to0$ in $L^2(\lambda)$ and is energy Cauchy, its abstract
energy-completion limit is orthogonal to every core vector by
(AE10)--(AE11), hence is zero.  This proves closability.
The original chain rule for smooth normal contractions proves
the Markov property after weak compactness and closure.
Multiplication by a fixed original smooth cutoff is bounded
for the closed form norm, by
\[
 \mathcal E(\chi f,\chi f)
   \leq2\|\chi\|_\infty^2\mathcal E(f,f)
                              +2C_\chi\|f\|_2^2.
 \tag{AE12}
\]
If $g$ is constant near the support of $f$, insert such a cutoff
in a core approximation to $f$ and use (AE11) and the zero total
mass of $D(\chi f_jS)$.  This proves strong locality on $F_S$.

Finally let $v_m=2-h_m$, so $1\leq v_m\leq2$.
The original nonnegative square
$\Gamma_S(f)-\Gamma_S(v_m,f^2/v_m)$ and compact Stokes yield
\[
                  \mathcal E(f,f)\geq
                     \int(f^2/v_m)\,D(h_mS).
\]
Increasing measure convergence and uniform convergence of $v_m$
give (AE9) on the core, then on the closed domain.
The estimates $\Gamma_S(a_l)\leq2D(a_lS)$ also show that
$\chi h_m\to\chi h$ in form norm for each original compact
cutoff $\chi$, exactly by the weighted Cauchy--Schwarz estimate
\[
 \Gamma_S(h_k-h_m)\leq2^{1-m}\lambda\qquad(k>m).
 \tag{AE13}
\]
Thus $h$ is locally in $\mathscr D$, a fact needed for the
bounded barrier argument below.  Every step uses exact original
moments and their finite common refinements.
\end{proof}

\begin{corollary}[Part spaces, torsion and all-original Green identities]
\label{cor:r11-affinoid-part-green}
Let $V\Subset\Omega$ be ordinary open, let $F=V\cap F_S$, and let
$H=\mathscr D_0(V)$ be the energy closure of the original compact
smooth scalars supported in $V$.  The part form is regular on $F$
and coercive for $\mathcal E^{1/2}$.  For $f\in L^2(F,\lambda)$
there is a unique $Gf\in H$ satisfying
\[
                  \mathcal E(Gf,\phi)=\int_F f\phi\,d\lambda
                                   \quad(\phi\in H).
 \tag{AE14}
\]
If $f\in L^\infty(F,\lambda)$, then
$|Gf|\leq(2-h)\|f\|_\infty\leq2\|f\|_\infty$ energy-quasi-everywhere.
In particular the torsion vector $z=G1$ satisfies $0\leq z\leq2$.
For every $v\in H$ and original $\phi\in A_c^{0,0}(V;\mathbf R)$,
\[
       \mathcal E(v,\phi)
             =-\langle D(vS),\phi\rangle
             =-\int_Fv\,b_\phi\,d\lambda.
 \tag{AE15}
\]
\end{corollary}

\noindent\textbf{Status: \texttt{DIRECT\_COROLLARY}} of
Theorem~\ref{thm:r11-general-affinoid-faithful-form} and the frozen
R4 part-space barrier proof.

\begin{proof}
The part space is the closure of the stated original core; it
inherits coercivity and is a regular part Dirichlet form.
Riesz representation gives (AE14).  The normal-contraction test
with positive parts gives positivity of $G$.  Choose a compact
cutoff equal to one near $\overline V$ and use its product with
$2-h$; (AE13) puts this barrier in the local form domain and
$\mathcal E(2-h,\phi)=\int\phi\,d\lambda$ on $V$.
Testing with $(Gf-\|f\|_\infty(2-h))^+$ and its negative
counterpart proves the bound, with the support-preserving
approximation proving these tests belong to $H$.
For (AE15), $v$ is locally $L^1(\lambda)$ and therefore $vS$
is defined on every original coefficient by (AE10).
Approximate $v$ in form norm, hence in $L^2(\lambda)$, by
the original core and pass (AE11).
\end{proof}

\paragraph{Scope.}
The construction is available around every point of the permitted
analytic space; one chooses (AE1) inside an affinoid neighborhood.
It does not impose an algebraic origin on either $X$ or $S$.
It does retain the stated one-affinoid relatively compact domain
for a coercive zero-boundary problem.  No global coercivity on an
arbitrary compact boundaryless space, pointwise ordinary boundary
regularity, or ambient classical psh property of $Gf$ is asserted.

% END REVIEWED MODULE research_r11/dirichlet/analytic_affinoid_energy_domain.tex


% BEGIN REVIEWED MODULE research_r11/dirichlet/analytic_laurent_collars.tex
\section{Analytic Laurent collars and bounded relative quasicontinuity}
\label{sec:r11-analytic-Laurent-collars}

The field, absolute value and analytic space have the standing original
hypotheses.  The background in this section may have any degree:
$S$ is an arbitrary original symmetric middle-positive, separately
$d'$- and $d''$-closed $(p,p)$-current, and $q=n-p\geq1$.
The local smooth decreasing regularization of bounded classical psh
functions, on one common ordinary neighborhood for a finite family,
is a frozen \texttt{INTERNAL\_INPUT} from R/B and the R4 calculus.
The compact Green identity used below is the frozen R4 original-current
identity for bounded differences with compact support.  It is local:
local simultaneous smooth rows prove it on each original test stage,
and a finite partition of unity cancels the auxiliary cutoff derivatives.
It does not require a global regularization on the analytic collar.
No new regularization assertion is made here.  What is supplied is an
analytic collar construction, without a polynomial or algebraic
condition on the containing affinoid, and the exact capacity argument
on those collars.  This makes the domain scope of subsequent
quasicontinuity arguments explicit.

\begin{lemma}[Smooth common collars on analytic Laurent boxes]
\label{lem:r11-analytic-Laurent-collar}
Let $x\in X$, let $W$ be an ordinary neighborhood, and let finitely
many locally uniformly bounded decreasing sequences
$a_{i,j}\downarrow a_i$ of bounded classical psh functions be defined
on $W$.  After shrinking about $x$, there are compact Laurent domains
\[
 x\in\operatorname{Int}_X C_1,\qquad
 C_1\subset\operatorname{Int}_X C_{\rm mid},\qquad
 C_{\rm mid}\subset\operatorname{Int}_X C_2\Subset W,
 \tag{AC1}
\]
a smooth classical psh function $b$ on a neighborhood of $C_2$,
a constant $c$, and an original smooth compact cutoff $\chi$ in
$\Omega=\operatorname{Int}_X C_2$, such that
\[
 v_{i,j}=\max\{a_{i,j}+c,b\},\qquad
 v_i=\max\{a_i+c,b\}
 \tag{AC2}
\]
are decreasing bounded classical families, coincide with the
unpatched functions plus $c$ near $C_1$, and all coincide with $b$
on one fixed outer collar containing $C_2\setminus
\operatorname{Int}_X C_{\rm mid}$.  The cutoff is one near
$C_{\rm mid}$ and has its derivative supports in that collar.
If $a_{i,j}$ is continuous, so is $v_{i,j}$.
\end{lemma}

\noindent\textbf{Status: \texttt{NEW\_PROVED}}, as the analytic
Laurent-domain extension of the frozen polynomial collar construction.

\begin{proof}
Choose a compact affinoid neighborhood $B$ with $x\in\operatorname{Int}_X B$
and shrink $W$ inside its ordinary interior.  The topology of $B$ is
generated by the evaluation functions $|f|$, $f\in\mathcal O(B)$.
A finite basic neighborhood, with thresholds moved strictly inward,
therefore gives three nested boxes as in (AC1).  They are defined by
finitely many inequalities
\[
 |f_a|\le R_{a,l},\qquad |g_b|\ge r_{b,l},\qquad
 R_{a,1}<R_{a,{\rm mid}}<R_{a,2},\quad
 r_{b,1}>r_{b,{\rm mid}}>r_{b,2}>0.
 \tag{AC3}
\]
The values at $x$ satisfy all inner inequalities strictly.  A closed
box lies in the original basic open set after the outer thresholds
are chosen strictly within its defining bounds; this ensures
$C_2\Subset W$.  These are generalized Laurent domains, so no
condition on the thresholds belonging to $|K^\times|$ is needed.
All lower-bound functions are nonvanishing near $C_2$.

Choose intermediate thresholds
$R_{a,1}<\tau_a<R_{a,{\rm mid}}$ and
$r_{b,{\rm mid}}<\sigma_b<r_{b,1}$.
Let $F:\mathbf R\to\mathbf R_+$ be smooth, convex and nondecreasing,
zero on $(-\infty,0]$ and strictly positive on $(0,\infty)$.
For example take $F''(s)=e^{-1/s}$ for $s>0$, zero for $s\le0$,
and integrate twice with $F(0)=F'(0)=0$.
Set
\[
 b_0=\sum_a F(\log|f_a|-\log\tau_a)
          +\sum_b F(\log\sigma_b-\log|g_b|).
 \tag{AC4}
\]
Terms containing a zero of $f_a$ are defined as zero there; they
vanish on an entire neighborhood of that zero.  Thus (AC4) is an
original smooth function near $C_2$.  Every active branch is a
convex nondecreasing function of the logarithmic norm of an analytic
unit or its inverse.  The frozen curve composition rule, after every
complete extension, proves that $b_0$ is classical psh.  No inference
from positive smooth curvature to classical psh is being substituted.

The function $b_0$ vanishes near $C_1$.  At every point of
$C_2\setminus\operatorname{Int}_X C_{\rm mid}$ at least one middle
inequality in (AC3) fails strictly or is an equality.  Hence
$b_0\ge\delta>0$ there, where $\delta$ is the minimum of the
finitely many positive threshold values in (AC4).  By continuity,
$b_0\ge\delta/2$ on a neighborhood of that compact set.
Endpoint bounds and a common constant $c$ make all $a_{i,j}+c,a_i+c$
lie in $[1,M]$ near $C_2$.  Multiply $b_0$ by a constant making it
larger than $M+3$ on the outer collar, and call the result $b$.
The maxima (AC2) have all the asserted properties.  Original compact
smooth cutoff existence gives $0\le\chi\le1$, equal to one near
$C_{\rm mid}$ and compactly supported in $\operatorname{Int}_X C_2$;
its derivative supports lie in the outer collar.
\end{proof}

For an ordinary open $\Omega$ write
\[
 C_S(E;\Omega)=
 \sup_{\substack{w\text{ bounded classical psh on }\Omega\\0\le w\le1}}
                  \bigl((Dw)^q\wedge S\bigr)(E),
 \tag{AC5}
\]
with its ordinary-open outer regularization $C_S^*$.
This is the original relative BT capacity; it is not an energy-form
capacity, and no unrestricted outer regularity is assumed.

\begin{proposition}[Capacity convergence on an analytic common collar]
\label{prop:r11-analytic-collar-capacity}
Suppose $v_j\downarrow v$ are bounded classical psh functions on
$\Omega=\operatorname{Int}_X C_2$ as above, uniformly bounded near
$C_2$, and equal to the same smooth $b$ outside one fixed compact
$L\subset\operatorname{Int}_X C_{\rm mid}$.  For $g_j=v_j-v\ge0$
and $\delta>0$ one has
\[
 C_S(\{g_j>\delta\};\Omega)
   \le \frac{A}{\delta}
      \left(\int g_j(Dv)^q\wedge S\right)^{2^{-q}}
                       \longrightarrow0.
 \tag{AC6}
\]
The constant depends on the compact collar, the local bound for $v$,
$q$, and finitely many original positive cutoff evaluations of $S$;
it is independent of $j$ and of the normalized capacity tester.
\end{proposition}

\noindent\textbf{Status: \texttt{DIRECT\_COROLLARY}} of the frozen
bounded Green and energy calculus.  The proof is recorded to verify
that no algebraic hypothesis remains in this consequence.

\begin{proof}
Fix a normalized tester $w$ and put
\[
 I_r=\int g_j(Dw)^r(Dv)^{q-r}\wedge S,\qquad
 Q_r=(Dw)^{r-1}(Dv)^{q-r}\wedge S\quad(1\le r\le q).
\]
Each $Q_r$ is an allowed arbitrary-background product of degree
$(n-1,n-1)$.  The compact support of $g_j$, frozen Green identity
and energy Cauchy--Schwarz give
\[
 I_r^2\le \left(\int\Gamma_{Q_r}(g_j)\right)
                \left(\int\chi\Gamma_{Q_r}(w)\right).
\]
Positivity and the exact energy identity give the one-sided estimate
\[
 0\le\int\Gamma_{Q_r}(g_j)
     =\int g_jD(vQ_r)-\int g_jD(v_jQ_r)\le I_{r-1}.
\]
The second factor is bounded uniformly in $w,j$ by a finite $A_r$:
indeed $0\le w\le1$ and
\[
 \int\chi\Gamma_{Q_r}(w)
 =\tfrac12\langle Q_r,w^2D\chi\rangle
                         -\int\chi wD(wQ_r)
 \le\tfrac12|\langle Q_r,w^2D\chi\rangle|,
\]
whose absolute coefficient control is uniformly bounded by relative
CLN for the fixed bounded $v$ and the normalized $w$.
Iteration gives
$I_q\le A_q^{1/2}A_{q-1}^{1/4}\cdots A_1^{1/2^q}I_0^{1/2^q}$.
Chebyshev's inequality followed by the tester supremum proves (AC6).
The final convergence is bounded convergence against the fixed locally
finite measure $(Dv)^q\wedge S$, since $g_j\downarrow0$ and all
supports lie in one compact set.  Every integration identity and
coefficient estimate uses original tests; added actual moment
coordinates are harmless finite refinements of each tested stage.
\end{proof}

\begin{corollary}[Bounded local relative quasicontinuity on arbitrary $X$]
\label{cor:r11-analytic-bounded-relative-QC}
Let $\Omega$ be any ordinary open set in $X$, let $S$ be as above,
and let finitely many locally bounded classical psh functions $a_i$
be defined on $\Omega$.  For every compact $H\Subset\Omega$ and
$\varepsilon>0$ there is an ordinary open $G\subset\Omega$ with
\[
                   C_S(G;\Omega)<\varepsilon
 \tag{AC7}
\]
such that each $a_i|_{H\setminus G}$ is continuous.
The same conclusion holds for any bounded difference of such
functions.  Local decreasing bounded classical sequences converge
in the relative capacity on smaller compact cores.  These conclusions
hold for arbitrary original middle-positive closed $S$ of every
permitted degree, on every actual moment refinement.
\end{corollary}

\noindent\textbf{Status: \texttt{NEW\_PROVED}} as an explicit
analytic-domain extension of the frozen local norm-box consequence;
the bounded regularization input and the energy estimate are not new.

\begin{proof}
First work near a point of $H$.  Invoke the frozen common local
regularization once, before choosing the boxes, to obtain smooth
classical $a_{i,j}\downarrow a_i$ on a common ordinary neighborhood.
They have common bounds on a compact smaller neighborhood after
finitely many terms are discarded.  Lemma
\ref{lem:r11-analytic-Laurent-collar} gives patched continuous
classical sequences $v_{i,j}\downarrow v_i$ with a common collar.
Their differences have support in one fixed compact subset of
$\operatorname{Int}_X C_{\rm mid}$, because equality with $b$ holds
on a neighborhood of the compact outer collar.  Proposition
\ref{prop:r11-analytic-collar-capacity} applies.

For each $i$, choose indices $j_l$ such that
\[
 C_S(\{v_{i,j_l}-v_i>2^{-l}\};\operatorname{Int}_X C_2)
                                        <\varepsilon_i2^{-l-1}.
\]
These exceptional sets are ordinary open: $v_{i,j_l}$ is continuous
and $v_i$ is upper semicontinuous.  Their union has capacity less
than $\varepsilon_i$ by countable subadditivity.  Outside it the
continuous functions converge uniformly to $v_i$, so $a_i$ is
continuous on the smaller core after removal of the same open set.
A finite union treats the finitely many $i$.

Cover $H$ by finitely many such ordinary smaller cores, with containing
boxes $\Omega_l\subset\Omega$, and assign the error budgets so that
their sum is less than $\varepsilon$.  Each local exceptional set is
ordinary open in $\Omega$.  Domain monotonicity gives
$C_S(G_l;\Omega)\le C_{S|\Omega_l}(G_l;\Omega_l)$.
Their union therefore satisfies (AC7).  The smaller cores give a
relative open cover of $H\setminus G$, on each of which every $a_i$
is continuous.  This proves global continuity on that subspace.
The decreasing-convergence assertion uses the same local collars,
finite cover and domain monotonicity.  The proof does not identify
raw and outer capacities on general Borel sets, and does not assert
an unrestricted pluripolar-null equivalence.
\end{proof}

% END REVIEWED MODULE research_r11/dirichlet/analytic_laurent_collars.tex


% BEGIN REVIEWED MODULE research_r11/comparison/analytic_affinoid_capacity_localization.tex
\section{Analytic affinoid extension of local capacity testers}
\label{sec:r11-analytic-capacity-localization}

Keep only (DG1): the ordinary open domain $V$ has compact
closure contained in $\operatorname{Int}_X(B)$ for one
original compact $K$-affinoid domain $B$. The background $S$
on $V$ remains arbitrary symmetric middle-positive and
separately closed, of degree $(p,p)$, and $q=n-p\geq1$.
No extension of $S$ across $\partial V$ is required here,
and $V$ itself need not be a Laurent box.

The exact frozen predecessor is R2
\texttt{domain\_comparison.tex},
\path{lem:r2-uniform-\allowbreak framed-extension} and
\path{thm:r2-domain-\allowbreak capacity-comparison}. Those statements
use a relatively compact domain in a smooth affine algebraic
variety and a projective framing construction. Smooth
pluriharmonic invariance of bounded arbitrary-$S$ products,
\path{lem:r2-harmonic-\allowbreak frame-invariance}, is retained as
\texttt{INTERNAL\_INPUT}. The new argument below gives an
explicit bounded analytic construction on $B$ itself.
It uses neither a projective compactification nor an
unproved algebraization of the containing analytic domain.

\begin{lemma}[A bounded analytic extension up to a fixed harmonic term]
\label{lem:r11-affinoid-harmonic-tester-extension}
Let $\Omega\subset V$ be ordinary open and let $x\in\Omega$.
There is an ordinary neighborhood $W$ of $x$ with compact
closure in $\Omega$, a constant $A_0>0$, and a fixed
original smooth pluriharmonic function $H_0$ near $\overline W$
such that every bounded classical psh $a$ on $\Omega$ with
$0\leq a\leq1$ has a bounded classical extension tester
$\widetilde a$ on $V$ satisfying
\[
 0\leq\widetilde a\leq1,\qquad
       \widetilde a=A_0^{-1}a+H_0
                         \quad\hbox{near }\overline W.
 \tag{AL1}
\]
The constants and $H_0$ depend only on the chosen analytic
boxes and $B$, not on $a$ or $S$. In every original degree,
\[
 \mathbf1_W T_{S,k}(\widetilde a^{[k]})
       =A_0^{-k}\mathbf1_W T_{S,k}(a^{[k]})
                         \quad(0\leq k\leq q).
 \tag{AL2}
\]
The equality holds as original currents on a neighborhood
of $\overline W$, before the displayed Borel restriction.
The same construction and scale may be applied to finitely
many competitors, giving the corresponding mixed identity.
\end{lemma}

\noindent\textbf{Status: \texttt{NEW\_PROVED}} as the
original affinoid analytic extension of the frozen framed
algebraic tester lemma.

\begin{proof}
The topology of $B$ is generated by $|f|$, $f\in\mathcal O(B)$.
Choose two nested closed Laurent boxes $C_1,C_2$ defined by
one finite list of such functions and thresholds with strict
slack, so that
\[
 x\in\operatorname{Int}_X C_1,\qquad
 C_1\subset U_2\Subset\Omega,
 \quad
 \begin{cases}
 C_2=\{|f_a|\leq R_a,\ |g_b|\geq r_b\},\\
 U_2=\{|f_a|<R_a,\ |g_b|>r_b\},
 \end{cases}
 \tag{AL3}
\]
with $\overline{U_2}\subset C_2\Subset\Omega$.
Both boxes lie in $\operatorname{Int}_X(B)$.
All lower-bound functions are nonvanishing near $C_2$.
Set, on such a neighborhood,
\[
 \rho_2=\max\left(-1,\max_a\log\frac{|f_a|}{R_a},
                           \max_b\log\frac{r_b}{|g_b|}\right),
 \qquad H=\sum_b\log|g_b|.
 \tag{AL4}
\]
Empty lower lists give $H=0$. The function $\rho_2$ is
bounded continuous classical psh there, negative on $U_2$
and zero on its ordinary boundary. The function $H$ is
original smooth pluriharmonic. Strict slack gives
$\rho_2\leq-c<0$ near $C_1$.

The lower logarithms in $\rho_2$ cannot individually be
extended across a zero of $g_b$. To remove precisely this
problem, put $G=\prod_b g_b$ and $G_b=\prod_{c\ne b}g_c$.
On the neighborhood where all $g_b$ are units,
\[
 \rho_2+H
  =\max\left(\log|G|-1,
        \max_a\{\log|f_aG|-\log R_a\},
        \max_b\{\log|G_b|+\log r_b\}\right).
 \tag{AL5}
\]
Every product displayed on the right is an analytic function
on the original $B$; $G_b$ means the printed product, not a
meromorphic quotient at a zero. If the lower list is empty,
$G=1$ and the first branch is the constant $-1$.
Choose a real constant $L$ strictly below the minimum of
$\rho_2+H$ on $C_2$, and define $P$ on $\operatorname{Int}_X B$
to be the maximum of $L$ and all the right-hand branches in
(AL5). It is bounded continuous classical psh: the constant
branch makes every logarithmic zero inactive near its zero
set, and the analytic functions are bounded on the compact
affinoid. It agrees with $\rho_2+H$ on a neighborhood of
$C_2$. Real thresholds and this constant introduce no field
extension or change of the original norm.

Choose $A>0$ with $A\rho_2<-2$ near $C_1$, and for the
given normalized $a$ put
\[
                     b_a=\max(a-2,A\rho_2).
 \tag{AL6}
\]
Near $C_1$ it equals $a-2$. Near the boundary of $U_2$ it
equals $A\rho_2$, because $a-2\leq-1$ whereas
$A\rho_2>-1/2$ on a fixed ordinary collar. On that collar
$b_a+AH=AP$. Therefore the prescription
\[
 F_a=b_a+AH\text{ on }U_2,
       \qquad F_a=AP\text{ outside }U_2
 \tag{AL7}
\]
glues to a classical psh function on $\operatorname{Int}_X B$.
For clarity this is ordinary open gluing: the modification
is contained in the compact inner sublevel
$\{\rho_2\leq-1/(2A)\}\Subset U_2$, and both formulas
agree on an open neighborhood of its outer complement in
the local defining neighborhood.

The functions $F_a$ have uniform upper and lower bounds
independent of $a$. Outside $U_2$ this follows from boundedness
of $AP$ on $B$; inside, one has $-2\leq b_a\leq0$ and
$H$ is bounded on $C_2$. Choose common bounds $m\leq F_a\leq M$
and set $A_0=1+M-m$ and
$\widetilde a=(F_a-m)/A_0$ on $V$.
This lies in $[0,1]$. Near $C_1$ it equals
$A_0^{-1}a+A_0^{-1}(AH-2-m)$, giving (AL1) with
$W=\operatorname{Int}_X C_1$ and the indicated fixed
smooth pluriharmonic $H_0$.

The frozen harmonic invariance lemma and bounded scaling
give (AL2) on the ordinary neighborhood where (AL1) holds.
They use $D(H_0Q)=0$ for every separately closed original
current $Q$, by the original smooth product rule. This
identity is valid on every actual moment refinement.
There is no inference from equality of arbitrary ambient
curvature currents or from a signed PS/PL quotient.
\end{proof}

\begin{theorem}[All local capacities are comparable on compact cores]
\label{thm:r11-affinoid-domain-capacity-comparison}
For every ordinary $\Omega\subset V$ and every ordinary
$W$ with $\overline W\Subset\Omega$, there is a finite
constant $C=C(B,V,\Omega,W,q)$ independent of $S$ such that
\[
 C_S(E;V)\leq C_{S|\Omega}(E;\Omega)
                  \leq C C_S(E;V)
                    \quad(E\subset W\text{ Borel}).
 \tag{AL8}
\]
For every subset $E\subset W$, not necessarily Borel,
\[
 C_S^*(E;V)\leq C_{S|\Omega}^*(E;\Omega)
                  \leq C C_S^*(E;V).
 \tag{AL9}
\]
These estimates compare domains; they do not assert that raw
and outer capacities coincide on a general Borel set.
\end{theorem}

\noindent\textbf{Status: \texttt{DIRECT\_COROLLARY}} of
the new analytic extension, retaining the full tester class.

\begin{proof}
The first raw inequality is restriction of a normalized
competitor on $V$ to $\Omega$. Cover the compact $\overline W$
by finitely many inner cores of the preceding lemma. Let
$A_i$ be their extension scales. For any normalized tester
$a$ on $\Omega$, its separate extensions on these cores give
\[
 \mu_a(E)\leq\sum_i\mu_a(E\cap W_i)
       =\sum_i A_i^q\mu_{\widetilde a_i}(E\cap W_i)
       \leq\left(\sum_i A_i^q\right)C_S(E;V).
\]
Take the supremum over all $a$. Thus $C=\sum_i A_i^q$
works. Every scale is independent of the background.

For outer capacities, any ordinary open subset of $\Omega$
is ordinary open in $V$, so the first inequality follows
from the first raw comparison. Conversely, if $G$ is open
in $V$ and contains $E\subset W$, then $G\cap W$ is open
in $\Omega$ and contains $E$. Apply the raw comparison to
this set, which is contained in $W$, to get
\[
 C_{S|\Omega}^*(E;\Omega)
       \leq C_{S|\Omega}(G\cap W;\Omega)
       \leq C C_S(G\cap W;V)\leq C C_S(G;V).
\]
Taking the infimum over $G$ proves (AL9) directly from the
definition, with no raw/outer interchange.
\end{proof}

\begin{corollary}[The fixed domain detects all local convergence and nullity]
\label{cor:r11-affinoid-local-capacity-topology}
For real Borel functions on $V$, convergence in $C_S(\,\cdot\,;V)$
on every compact core is equivalent to convergence in all
smaller-domain original relative capacities on every compact
core. The same assertion holds with outer capacities.
For extended negative potentials apply this statement to
each of their fixed bounded truncations. Under (DG1), this
is precisely the equivalence of fixed-domain truncated
convergence and Definition~\ref{def:r11-truncated-local-capacity}.
The corresponding local raw and outer null classes also agree.
\end{corollary}

\noindent\textbf{Status: \texttt{DIRECT\_COROLLARY}}.

\begin{proof}
For a compact $K\Subset\Omega\subset V$, first choose an
ordinary $W$ with $K\subset W$ and $\overline W\Subset\Omega$,
then apply (AL8) or (AL9) to the relevant gap subset of $K$.
Conversely cover a given compact core in $V$ by finitely many
smaller ordinary domains and use domain monotonicity and
finite subadditivity. The same argument with zero values
proves local nullity equivalence. It concerns Borel functions
and capacities only; it does not construct a classical psh
representative of a limit.
\end{proof}

\paragraph{Exact advancement.}
The smooth harmonic invariance and the finite-cover capacity
argument are the frozen R2 mechanisms. Formula (AL5) is the
new bounded analytic replacement for their global projective
frame: clearing all lower-bound factors makes a finite list
of actual analytic products on the same affinoid, and the
constant lower cap gives a bounded psh branch everywhere.
The original analytic norm replacement principle used in the
faithful-form construction is consistent with this argument;
here the boxes are chosen directly from $\mathcal O(B)$, so
no quotient approximation or change of field is needed.
All assertions keep an arbitrary original background on $V$.

% END REVIEWED MODULE research_r11/comparison/analytic_affinoid_capacity_localization.tex


% BEGIN REVIEWED MODULE research_r11/capacity/psh_level_disintegration.tex
\section{Positive level currents and an intrinsic energy disintegration}
\label{sec:r11-level-disintegration}

The construction in this section applies to every allowed original
background.  It introduces positive currents carried by the level sets
of a bounded difference of classical psh functions.  They need not be
closed.  Their exact differential defect is retained, and their
Lebesgue integral recovers the energy current.  In particular no
arbitrary-current pullback, auxiliary residue field, or analytic
PS-relation quotient occurs in the construction.

Keep the original finite extension $K/\mathbf Q_p$ and its absolute
value.  Let $X$ be good, Hausdorff, boundaryless, quasi-smooth and pure
of dimension $n$, and let $U\subset X$ be ordinary open.  Every test
below is an original compactly supported smooth CLD test, of its
complementary bidegree.  Let $Q$ be an arbitrary symmetric
middle-positive, separately $d'$- and $d''$-closed current of degree
$(p,p)$, where $0\le p<n$.  Write $D=d'd''$.  A bounded difference
means a function $f=a-b$, with $a,b$ locally bounded classical psh;
all assertions are local, and finitely many such presentations may
be used on a compact test support.

The inputs are the frozen bounded weighted BT calculus, finite
polarization, strict Borel maximum locality, and coefficient-Radon
control on all original tests.  These are
\texttt{INTERNAL\_INPUT}.  In particular $D(fQ)$ is a signed
coefficient-Radon current, defined by linearity from the two bounded
classical potentials.

\subsection{The bounded-difference chain rule}

\begin{lemma}[Smooth scalar calculus on bounded differences]
\label{lem:r11-bounded-difference-chain}
The distribution
\[
 \Gamma_Q(f,g)
 =\frac12\{D(fgQ)-fD(gQ)-gD(fQ)\}
 \tag{LC1}
\]
is a symmetric coefficient-Radon current of degree $(p+1,p+1)$.
It is bilinear, and $\Gamma_Q(f):=\Gamma_Q(f,f)$ is
middle-positive.  On every relatively compact witness stage,
for a smooth real function $F$ on an interval containing the range
of $f$, one has
\[
 D(F(f)Q)=F'(f)D(fQ)+F''(f)\Gamma_Q(f),
 \qquad
 \Gamma_Q(F(f),g)=F'(f)\Gamma_Q(f,g).
 \tag{LC2}
\]
All multipliers use the given actual Borel values.  The formulas
are independent of the psh presentation of $f,g$ and of every
actual moment refinement.
\end{lemma}

\noindent\textbf{Status: \texttt{DIRECT\_COROLLARY}} of the
frozen weighted bounded calculus.  The passage from top-degree
energy measures to lower-degree coefficient currents is included
in the following argument.

\begin{proof}
Take simultaneous smooth decreasing classical rows for the finitely
many psh summands on a compact regularization stage.  The smooth
identity for (LC1) is
\[
 \Gamma_Q(f,g)
 =\tfrac12(d'f\wedge d''g+d'g\wedge d''f)\wedge Q.
\]
For $f=g$ the first factor is strongly positive, so the current is
middle-positive.  The two weighted terms in (LC1) converge by the
frozen weighted theorem, and the first term converges by bounded
convergence against the fixed coefficient measures of $Q$.
Thus (LC1) is the distributional limit of positive currents when
$f=g$.  The original positive-current coefficient theorem makes
it coefficient-Radon; polarization treats $f,g$.  This proof works
in every permitted bidegree, not only in top degree.

For completeness, a smooth function $H$ of finitely many bounded
psh variables is locally a difference of two smooth convex,
coordinatewise nondecreasing functions of those variables.  On a
slightly larger compact real box, add a sufficiently large positive
multiple of $\sum x_i^2$, followed by a sufficiently large linear
multiple of $\sum x_i$.  The resulting function and the added
function are convex and coordinatewise nondecreasing on that box.
Their compositions with the psh variables are bounded classical
psh functions.  Applied to decreasing rows, each such composition
decreases to the actual composition.  This gives, in particular,
bounded-difference presentations for $F(f)$ and all smooth scalar
coefficients needed in the smooth chain identities.

One can therefore pass every weighted term of those identities by
the frozen joint decreasing theorem after the indicated finite
decompositions.  For an energy term multiplied by a smooth scalar
coefficient, expand (LC1): the $D(f^2Q)$ term is a difference of
bounded one-potential BT currents by the same convex decomposition
of $f^2$, and its coefficient, as well as the coefficients of
$fD(fQ)$, has such a decomposition.  Every passage is consequently
a finite combination of the already proved weighted BT passages.
This proves both formulas in (LC2).  They involve only intrinsic
distributional operations and actual coefficient multiplication,
which proves independence of the chosen witnesses and refinements.
\end{proof}

\begin{definition}[The level curvature current]
\label{def:r11-level-current}
For a bounded difference $f$ and $t\in\mathbf R$ put
\[
 \mathsf L^Q_f(t)
 :=D(\max\{f,t\}\,Q)
                 -\mathbf1_{\{f>t\}}D(fQ).
 \tag{LC3}
\]
The first current is well defined: if $f=a-b$, then
$\max\{f,t\}=\max\{a,b+t\}-b$ is again a bounded
difference.  Formula (LC3) is the definition, not an assertion
that a Borel-restricted closed current is closed.
\end{definition}

\begin{theorem}[Positivity, carriers, and the differential defect]
\label{thm:r11-positive-level-current}
For every real $t$, $\mathsf L^Q_f(t)$ is a middle-positive
coefficient-Radon current of degree $(p+1,p+1)$ and is carried by
the actual Borel set $\{f=t\}$.  More precisely,
\[
 \mathbf1_{\{f\ne t\}}\mathsf L^Q_f(t)=0.
 \tag{LC4}
\]
Its separate derivatives are exactly
\[
 \begin{split}
 d'\mathsf L^Q_f(t)
       &=-d'\bigl(\mathbf1_{\{f>t\}}D(fQ)\bigr),\\
 d''\mathsf L^Q_f(t)
       &=-d''\bigl(\mathbf1_{\{f>t\}}D(fQ)\bigr).
 \end{split}
 \tag{LC5}
\]
No vanishing of these derivatives is asserted when $p+1<n$.
The construction is local in the original ordinary topology and
commutes with every actual refinement of the same functions and tests.
\end{theorem}

\noindent\textbf{Status: \texttt{NEW\_PROVED}}.

\begin{proof}
Choose a smooth convex $\psi_\epsilon:\mathbf R\to\mathbf R$
which is zero on $(-\infty,0]$, has derivative in $[0,1]$, and
has derivative one on $[\epsilon,\infty)$.  It can be obtained by
integrating twice a nonnegative smooth function of integral one
supported in $(0,\epsilon)$.  It satisfies
\[
 |\psi_\epsilon(s)-s_+|\le\epsilon,
 \qquad
 \psi_\epsilon'(s)\longrightarrow\mathbf1_{\{s>0\}}
                   \quad\text{at every real }s.
\]
The chain identity (LC2) gives
\[
 D(\psi_\epsilon(f-t)Q)
        -\psi_\epsilon'(f-t)D(fQ)
            =\psi_\epsilon''(f-t)\Gamma_Q(f)\ge0.
 \tag{LC6}
\]
On a compact original test support the first term tends to
$D((f-t)_+Q)$ by uniform convergence and the fixed coefficient
Radon bounds for $Q$.  The second tends to
$\mathbf1_{\{f>t\}}D(fQ)$ by bounded convergence against the
fixed signed coefficient measures of $D(fQ)$.  As $D(tQ)=0$,
the limiting distribution is (LC3).  Positivity is closed under
this convergence on all original tests, proving the claim.

To identify its carrier, use $\max\{f,t\}=\max\{a,b+t\}-b$.
Strict Borel maximum locality gives
\[
 \begin{split}
 \mathbf1_{\{f>t\}}D(\max\{f,t\}Q)
        &=\mathbf1_{\{f>t\}}D(fQ),\\
 \mathbf1_{\{f<t\}}D(\max\{f,t\}Q)&=0.
 \end{split}
\]
This is a coefficient-measure identity; neither strict Borel set
is being declared ordinary open.  Subtracting the second term
of (LC3) proves (LC4).  Applying $d'$ and $d''$ to (LC3),
and using $d'D=d''D=0$, proves (LC5).  Finally every term in
the definition already has the original all-test and refinement
meaning, so no additional descent is needed.
\end{proof}

\subsection{Coarea on all original tests}

\begin{theorem}[Lebesgue disintegration of the energy current]
\label{thm:r11-level-coarea}
For every bounded Borel function $\vartheta:\mathbf R\to\mathbf R$
and every original compact smooth complementary test $\eta$,
\[
 \int_{\mathbf R}\vartheta(t)
                 \langle\mathsf L^Q_f(t),\eta\rangle\,dt
       =\langle\vartheta(f)\Gamma_Q(f),\eta\rangle.
 \tag{LC7}
\]
The integration is local on the test support: levels outside a
compact interval containing its nearby $f$-range give zero.
In particular,
\[
 \int_{\mathbf R}\mathsf L^Q_f(t)\,dt=\Gamma_Q(f),
 \qquad
 \mathbf1_{\{f=t\}}\Gamma_Q(f)=0\quad(t\in\mathbf R).
 \tag{LC8}
\]
The second identity holds for every level, including levels
charged by $D(fQ)$.

For a fixed positive trace on a compact moment stage, the resulting
positive scalar measures satisfy the full kernel formula
\[
 \int_B \mathsf L^Q_f(t)(E)\,dt
       =\Gamma_Q(f)(E\cap\{f\in B\})
 \tag{LC9}
\]
for Borel $B\subset\mathbf R$ and Borel $E$ in that stage.
Here the notation in (LC9) refers to the chosen scalar trace;
finite polarization recovers every original coefficient pairing.
\end{theorem}

\noindent\textbf{Status: \texttt{NEW\_PROVED}}.

\begin{proof}
First take $\vartheta\in C_c^\infty(\mathbf R)$ and set
\[
 F(s)=\int_{\mathbf R}\vartheta(t)(s-t)_+\,dt.
\]
Then $F''=\vartheta$ and
$F'(s)=\int_{t<s}\vartheta(t)\,dt$.  The function $F$ is
smooth.  On a compact test support the integrands in the
potential term of (LC3) are uniformly bounded, while the
restricted-curvature term is controlled by the fixed total
variations of $D(fQ)$.  Fubini and distributional differentiation
therefore give
\[
 \int\vartheta(t)\mathsf L^Q_f(t)\,dt
      =D(F(f)Q)-F'(f)D(fQ)
      =\vartheta(f)\Gamma_Q(f).
\]
The last equality is (LC2).

For a positive compact smooth trace test, the map
$t\mapsto\langle\mathsf L^Q_f(t),\eta\rangle$ is
nonnegative and Borel: the maximum term is continuous in $t$
on each compact support, and the restricted signed-measure term
is Borel by ordinary dominated integration.  Its integral is
finite on every relevant compact interval by the proved formula
with a smooth cutoff equal to one there.  Uniqueness of finite
Radon measures on $\mathbf R$ extends the formula from smooth
$\vartheta$ to bounded Borel $\vartheta$.  A finite collection
of positive trace tests dominates all coefficient variations
on a compact stage, so finite polarization gives (LC7) for
every original $\eta$.

For a fixed scalar trace, weak Borel measurability of the positive
Radon measures $\mathsf L^Q_f(t)$ implies measurability of their
values on Borel sets.  This can also be checked directly by
approximating open-set indicators from a countable continuous
core on a relatively compact affinoid stage and then using the
monotone-class theorem.  Formula (LC7), first for compact
continuous scalar weights on the stage and then by Radon
uniqueness, gives (LC9).  These arguments are local on finitely
many such stages and do not impose global second countability
on $X$.

Taking $B=\{t\}$ in (LC9) proves the zero-level energy assertion,
because Lebesgue measure has no atoms.  Taking $\vartheta=1$
on the local range proves the first identity in (LC8).
\end{proof}

\begin{corollary}[A localized bounded-variation energy density]
\label{cor:r11-energy-bv-density}
Suppose now that $p=n-1$, so the level currents are measures.
For an original smooth compact $\chi\ge0$ define the finite
signed measure on $\mathbf R$
\[
 \nu_{\chi,f,Q}
  :=f_*\bigl(\chi D(fQ)+\Gamma_Q(\chi,f)\bigr).
 \tag{LC10}
\]
Then $\nu_{\chi,f,Q}(\mathbf R)=0$, and for every $t\in\mathbf R$
\[
 k_{\chi,f,Q}(t):=\int\chi\,d\mathsf L^Q_f(t)
       =\nu_{\chi,f,Q}(( -\infty,t])\ge0.
 \tag{LC11}
\]
Thus $k$ is a right-continuous compactly supported BV function,
$\operatorname{Var}(k)\le|\nu_{\chi,f,Q}|(\mathbf R)$, and
\[
 f_*(\chi\Gamma_Q(f))=k_{\chi,f,Q}(t)\,dt,
 \qquad
 \|k_{\chi,f,Q}\|_\infty
                     \le|\nu_{\chi,f,Q}|(\mathbf R).
 \tag{LC12}
\]
\end{corollary}

\noindent\textbf{Status: \texttt{NEW\_PROVED}}.

\begin{proof}
The compact Stokes identity gives
$\int\chi D(fQ)=-\int\Gamma_Q(\chi,f)$, hence total mass
zero in (LC10).  Apply the first-derivative chain identity to the
one-sided smooth truncations used in (LC6).  Their uniform
convergence and the pointwise convergence of their derivatives
give
\[
 \int\chi D(\max\{f,t\}Q)
       =-\int_{\{f>t\}}\Gamma_Q(\chi,f).
\]
All limits here concern the fixed finite mixed energy measure;
its total variation is finite by Cauchy--Schwarz.  Subtracting
the strict curvature term of (LC3) yields
$k(t)=-\nu_{\chi,f,Q}((t,\infty))$, which is (LC11).
The support of $\chi$ and its differential is compact, and
$f$ is bounded on a neighborhood of that support, so $\nu$
and $k$ have the claimed compact support.  The BV and norm
bounds are the elementary cumulative-measure bounds.  Finally
(LC7) identifies the pushforward density with $k$.
\end{proof}

\begin{proposition}[A nonclosed level current on the original two-torus]
\label{prop:r11-nonclosed-level-example}
The conclusion that every positive level current in
Theorem~\ref{thm:r11-positive-level-current} is separately closed is
\texttt{FALSE}, even when the potential is smooth and classical and
the background is the original integration current.  This failure
occurs in every dimension $n\geq2$.
\end{proposition}

\noindent\textbf{Status: \texttt{NEW\_PROVED}} counterexample.

\begin{proof}
First take $X=(\mathbf G_{m,K}^2)^{\mathrm{an}}$, set
$t_i=-\log|T_i|$, and work on the ordinary domain
$U=\{|t_1|<1,\ |t_2|<1\}$.  Put
\[
 Q=[U],\qquad f=t_1^2+t_2,\qquad
 B_1=d't_1\wedge d''t_1.
 \tag{LC13}
\]
The function $f$ is bounded, original smooth and classical psh:
its convex moment expression has positive semidefinite Hessian.
One has $D(fQ)=2B_1\wedge[U]$.  The distributional derivative of
the half-space indicator in the two original moment coordinates
therefore gives
\[
 d'\mathsf L_f^Q(0)
 =-2\delta(t_1^2+t_2)\,d't_2\wedge B_1\wedge[U].
 \tag{LC14}
\]
This can be verified directly by the smooth one-sided functions
from (LC6): their derivatives converge to the strict indicator,
and ordinary real integration by parts in the moment variables
gives the displayed hypersurface distribution.

Choose $\rho\in C_c^\infty((-1,1)^2)$ nonnegative and positive
on a nonempty segment of $\{(s,-s^2):|s|<1/2\}$.
The original smooth form $\rho(t_1,t_2)d''t_2$ has compact
support in $U$: the inverse image of the compact moment support
is a compact polyannulus.  Original tropical calibration of the
two-torus has multiplicity one, so its pairing with (LC14),
up to the fixed orientation convention, has absolute value
\[
 2\int_{-1}^{1}\rho(s,-s^2)\,ds>0.
 \tag{LC15}
\]
Thus $d'\mathsf L_f^Q(0)\ne0$.  By symmetry the double-prime
defect is nonzero as well.  In higher dimensions take the product
with $n-2$ original open annuli, keep $f$ independent of the extra
coordinates, and wedge the test with a positive compact moment
volume form in those coordinates.  Its nonzero pairing is the
product of (LC15) with a strictly positive real integral.  This
uses actual original tests in every dimension; no delta-test
enlargement is used.
\end{proof}

\paragraph{Closedness and the relation to the general target.}
The currents $\mathsf L^Q_f(t)$ provide an actual level
disintegration of energy, rather than a disintegration of an
arbitrary BT curvature measure into closed backgrounds.
Equation (LC5) is indispensable below top degree.  In particular
one cannot insert each level current into the frozen closed-background
calculus and integrate without separately resolving that derivative
defect.  Neither general OPEN-GQ nor general Borel regularity of
the raw relative BT capacity is inferred from this construction.

% END REVIEWED MODULE research_r11/capacity/psh_level_disintegration.tex


% BEGIN REVIEWED MODULE research_r11/capacity/nonsmooth_scalar_curvature.tex
\section{Nonsmooth scalar curvature and the two level currents}
\label{sec:r11-nonsmooth-scalar-curvature}

Keep the original hypotheses and notation of
Section~\ref{sec:r11-level-disintegration}.  Thus $Q$ is any
original symmetric middle-positive separately closed current of
degree $(p,p)$, $p<n$, and $f$ is a locally bounded difference
of classical psh functions.  All identities are identities of
original coefficient-Radon currents on ordinary open sets, tested
against every original compact smooth complementary form.  The
positive level-current theorem of that section is the only new
input to the present extension.

\begin{definition}[Lower and upper slope conventions]
Write
\[
 \begin{split}
 \mathsf L_f^{-}(t)
    &=D(\max(f,t)Q)-\mathbf1_{\{f>t\}}D(fQ),\\
 \mathsf L_f^{+}(t)
    &=D(\max(f,t)Q)-\mathbf1_{\{f\geq t\}}D(fQ).
 \end{split}
 \tag{SC1}
\]
The superscripts indicate the lower and upper values assigned to
the derivative of the scalar hinge at its corner; in particular
$\mathsf L_f^{-}(t)=\mathsf L_f^Q(t)$ from (LC3).
The background $Q$ is suppressed in this section's notation only.
\end{definition}

\begin{proposition}[Two positive level currents and their jump]
\label{prop:r11-two-level-currents}
Both currents in (SC1) are middle-positive and carried by
$\{f=t\}$.  They satisfy
\[
 \mathsf L_f^{+}(t)=\mathsf L_{-f}^{-}(-t),\qquad
 \mathsf L_f^{-}(t)-\mathsf L_f^{+}(t)
                      =\mathbf1_{\{f=t\}}D(fQ).
 \tag{SC2}
\]
For every original compact smooth test $\eta$, the scalar function
$t\mapsto\langle\mathsf L_f^{-}(t),\eta\rangle$ is right
continuous and has bounded variation on each compact level interval.
Its left limit at every $t$ is
$\langle\mathsf L_f^{+}(t),\eta\rangle$.  On any compact real
interval of levels, the local positive coefficient masses of both
families are uniformly bounded.  Both admit the energy coarea
identity (LC7), with the same energy current $\Gamma_Q(f)$.
\end{proposition}

\noindent\textbf{Status: \texttt{NEW\_PROVED}}.

\begin{proof}
The elementary identity
$\max(-f,-t)=\max(f,t)-f-t$ and $D(tQ)=0$ give the first
formula in (SC2); Theorem~\ref{thm:r11-positive-level-current}
applied to $-f$ proves positivity and the carrier of
$\mathsf L_f^{+}(t)$.  Subtracting (SC1) proves the second formula.

For a fixed compact test, the function
$t\mapsto\langle D(\max(f,t)Q),\eta\rangle$ is Lipschitz:
the potentials differ by at most $|t-s|$, and the fixed signed
coefficient measure $\mu_{Q,D\eta}$ has finite variation on the
test support.  Continuity from below and above for the fixed
signed coefficient measures of $D(fQ)$ gives
\[
 \begin{split}
 \lim_{s\downarrow t}\mathbf1_{\{f>s\}}D(fQ)
       &=\mathbf1_{\{f>t\}}D(fQ),\\
 \lim_{s\uparrow t}\mathbf1_{\{f>s\}}D(fQ)
       &=\mathbf1_{\{f\geq t\}}D(fQ),
 \end{split}
\]
in local coefficient total variation.  This proves the stated
one-sided continuity after combining the two terms of (SC1).
The first term is Lipschitz in the level, and the second is a
tail of a fixed finite signed measure after pushforward by $f$;
their sum therefore has bounded variation on compact intervals.

For the uniform mass assertion, fix a compact original moment
stage and a positive trace control $\chi\beta^{n-p-1}$, where
$\beta$ is the sum of the elementary positive moment squares.
The closed current $Q'=\beta^{n-p-1}\wedge Q$ has degree
$(n-1,n-1)$ and
$\mathsf L_f^{\pm,Q'}(t)=\beta^{n-p-1}\wedge
\mathsf L_f^{\pm,Q}(t)$ by the bounded smooth-slot identities.
The bounded-variation density bound (LC12), applied to $f$ and
$-f$, bounds the two positive trace masses uniformly in $t$.
Finite original coefficient domination proves the assertion for
every compact stage, including any actual moment refinement.

Finally the difference in (SC2), integrated in $t$ against a
bounded Borel scalar function, is zero: for each point $x$ the
indicator of the singleton $\{f(x)\}$ has Lebesgue integral
zero, and Fubini applies to each fixed finite coefficient
variation of $D(fQ)$.  Therefore (LC7) for the minus convention
also proves it for the plus convention.
\end{proof}

\begin{theorem}[An exact scalar curvature formula with a real measure]
\label{thm:r11-stieltjes-scalar-curvature}
Let $J\subset\mathbf R$ be an open interval, let $I\Subset J$
contain the range of $f$ near the compact test support, and let
$F$ be a finite convex real function on $J$.  Let $\nu_F=F''$
be its positive distributional second-derivative measure, and
write $F'_-,F'_+$ for its finite one-sided derivatives on $I$.
Then $F(f)$ is locally a bounded difference of classical psh
functions, and
\[
 \boxed{\quad
 D(F(f)Q)
  =F'_-(f)D(fQ)+\int_J\mathsf L_f^{-}(t)\,d\nu_F(t)
  =F'_+(f)D(fQ)+\int_J\mathsf L_f^{+}(t)\,d\nu_F(t).
 \quad}
 \tag{SC3}
\]
The integrals are weak integrals of coefficient-Radon currents;
they are well defined locally, middle-positive, and carried by
the actual inverse image of a carrier of $\nu_F$.
Only levels in a compact interval near the local range contribute.

The same formula holds for a real difference of two finite convex
functions, with its signed second-derivative measure and the
corresponding signed integral.  All multipliers in (SC3) use the
actual given Borel values.  The construction commutes with
ordinary restriction and every actual moment refinement.
\end{theorem}

\noindent\textbf{Status: \texttt{NEW\_PROVED}}.

\begin{proof}
Write $f=a-b$ on the chosen witness stage.  On a compact interval
strictly larger than $I$ and contained in $J$, the convex function
$F$ is Lipschitz.  Extend its restriction affinely past the two
endpoints using its interior one-sided slopes; this gives a global
convex Lipschitz extension $\widehat F$ agreeing with $F$ on $I$.
Choose $L$ greater than the absolute values of its slopes.  On the
compact range box of $(a,b)$ the function
\[
                  (s,t)\longmapsto\widehat F(s-t)+L(s+t)
 \tag{SC4}
\]
is convex and nondecreasing in each variable.  The frozen
convex nondecreasing composition rule makes (SC4), composed
with $(a,b)$, a bounded classical psh function.  Subtracting
$L(a+b)$ gives a bounded-difference presentation of $F(f)$.

Choose $c<\inf I\leq\sup I<d$ with $[c,d]\Subset J$.
The elementary real convex representation, on $I$, is
\[
 F(s)=A+Bs+\int_{(c,d)}(s-t)_+\,d\nu_F(t),\qquad
 F'_-(s)=B+\nu_F((c,s)).
 \tag{SC5}
\]
The constants $A,B$ absorb the affine contribution of levels
outside $(c,d)$.  The measure in (SC5) is finite.  The
potentials $(f-t)_+$ are uniformly bounded on the compact
stage, so Fubini against the fixed coefficient measures of
$Q$ and distributional differentiation give
\[
 D(F(f)Q)=B D(fQ)
           +\int_{(c,d)}D((f-t)_+Q)\,d\nu_F(t).
 \tag{SC6}
\]
The positive level masses are uniformly bounded by
Proposition~\ref{prop:r11-two-level-currents}; their compact
test pairings are Borel.  They consequently have a well-defined
weak integral with local coefficient-Radon control.  Insert
$D((f-t)_+Q)=\mathbf1_{\{f>t\}}D(fQ)+\mathsf L_f^-(t)$
in (SC6), and apply Fubini to the fixed signed coefficient
measures in its first term.  Its multiplier is exactly
$B+\nu_F((c,f))=F'_-(f)$.  This proves the first equality
in (SC3).

Replacing the strict indicator by the weak indicator gives
$B+\nu_F((c,f])=F'_+(f)$ and the plus level current, proving
the second equality.  This also checks explicitly the jump
contribution at a real atom of $\nu_F$: changing slope
conventions changes the curvature multiplier by
$\nu_F(\{f\})D(fQ)$ and changes the level integral by its
negative, through (SC2).

For any Borel carrier $A$ of $\nu_F$, the positive scalar
kernel construction in (LC9), applied to the finite measure
$\nu_F$ in place of Lebesgue measure for measurability only,
makes the weak integral a positive measure kernel.  Each level
current is carried by $\{f=t\}$, so its integral is carried
by $f^{-1}(A)$.  No assertion that this inverse image is
closed is required.  Finite polarization gives the same
carrier assertion for every original coefficient.
Linearity treats a difference of convex scalar functions.
Intrinsic definitions and the already established refinement
invariance of each level current prove the final assertions.
\end{proof}

\begin{corollary}[Original Kato inequalities and the exact defect]
\label{cor:r11-original-kato}
For $f$ as above, put $f_+=\max(f,0)$ and use
$\operatorname{sgn}(0)=0$.  Then
\[
 \begin{split}
 D(f_+Q)-\mathbf1_{\{f>0\}}D(fQ)&=\mathsf L_f^-(0)\geq0,\\
 D(|f|Q)-\operatorname{sgn}(f)D(fQ)
                 &=\mathsf L_f^-(0)+\mathsf L_f^+(0)\geq0.
 \end{split}
 \tag{SC7}
\]
Both defects are positive coefficient-Radon currents carried by
the actual zero set of $f$.  The same identities hold after
every original positive coefficient trace and against every
original complementary test.  Below top degree the defects
need not be closed.
\end{corollary}

\noindent\textbf{Status: \texttt{DIRECT\_COROLLARY}} of the new
two-level theorem.

\begin{proof}
The first identity is (SC1).  Apply it to $-f$ and add, using
$|f|=f_++(-f)_+$ and
$\mathsf L_{-f}^-(0)=\mathsf L_f^+(0)$ from (SC2).
The two strict indicators give $\operatorname{sgn}(f)$ with
the printed value at zero.  Positivity and the carriers were
proved before the identities were added.  Nonclosedness is
allowed by the explicit differential defects in (LC5), and
Proposition~\ref{prop:r11-nonclosed-level-example} prevents a
universal closedness inference.
\end{proof}

\paragraph{What this calculus resolves and what it leaves open.}
Formula (SC3) integrates a real curvature measure against actual
level currents of a bounded psh difference.  It therefore gives
a precise replacement for formal multiplication by a real Dirac
mass in a nonsmooth scalar chain rule.  Its hypotheses do not
define an arbitrary pullback of a cycle through a moment map,
and no relation quotient for analytic delta presentations follows.
Moreover the level currents retain the defects (LC5).  Treating
them as separately closed backgrounds in a second BT iteration
is a \texttt{FAILED\_ATTEMPT}; (LC13)--(LC15) give its concrete
failure already on the original two-torus.


% END REVIEWED MODULE research_r11/capacity/nonsmooth_scalar_curvature.tex


% BEGIN REVIEWED MODULE research_r11/capacity/uniform_coarea_capacity.tex
\section{Uniform coarea capacity for arbitrary backgrounds}
\label{sec:r11-uniform-coarea-capacity}

Fix an original ordinary domain $\Omega$ and an arbitrary
allowed positive closed $(p,p)$-current $S$.  Put $q=n-p\ge1$.
All assertions below are local on compact subsets of $\Omega$,
and the capacity tester class always uses this same domain:
\[
 \mathcal P(\Omega)
  =\{v\in\mathcal C_b(\Omega):0\le v\le1\},
 \qquad
 C_S(E;\Omega)
  =\sup_{v\in\mathcal P(\Omega)}T_{S,q}(v,\ldots,v)(E).
 \tag{EC1}
\]
Its ordinary-open outer regularization is $C_S^*$.
No unrestricted Borel identification of $C_S$ and $C_S^*$
is assumed.  Let $f=a-b$, where $a,b$ are classical psh and
bounded on $\Omega$.  For a normalized tuple
$\mathbf v=(v_1,\ldots,v_{q-1})\in\mathcal P(\Omega)^{q-1}$,
write
\[
 Q_{\mathbf v}=T_{S,q-1}(v_1,\ldots,v_{q-1}).
 \tag{EC2}
\]
When $q=1$ this means $Q_{\varnothing}=S$.
Each $Q_{\mathbf v}$ is an original positive separately closed
$(n-1,n-1)$-current, even though $S$ is otherwise arbitrary.

\begin{definition}[The energy capacity of a bounded difference]
\label{def:r11-energy-layer-capacity}
For Borel $E\subset\Omega$ and arbitrary $A\subset\Omega$ put
\[
 \begin{split}
 \mathsf e_{S,f}(E;\Omega)
    &=\sup_{\mathbf v\in\mathcal P(\Omega)^{q-1}}
                          \Gamma_{Q_{\mathbf v}}(f)(E),\\
 \mathsf e_{S,f}^{\mathrm o}(A;\Omega)
    &=\inf_{\substack{G\supset A\\G\text{ ordinary open in }\Omega}}
                                      \mathsf e_{S,f}(G;\Omega).
 \end{split}
 \tag{EC3}
\]
This is a capacity attached to the energy of the specified $f$;
it is not a redefinition of $C_S$.  It is finite on compact
cores by the frozen relative energy and CLN estimates.
\end{definition}

\begin{lemma}[A uniform bound by the original BT capacity]
\label{lem:r11-energy-layer-BT-bound}
Let $m_a=\sup_\Omega a-\inf_\Omega a$ and define $m_b$
similarly.  Then for all Borel $E$ and arbitrary $A$,
\[
 \begin{split}
 \mathsf e_{S,f}(E;\Omega)&\le B_f C_S(E;\Omega),\\
 \mathsf e_{S,f}^{\mathrm o}(A;\Omega)&\le B_f C_S^*(A;\Omega),
 \qquad
 B_f=\frac{q^q}{q!}(m_a^2+m_b^2).
 \end{split}
 \tag{EC4}
\]
Both capacities in (EC3) are monotone and countably subadditive.
\end{lemma}

\noindent\textbf{Status: \texttt{DIRECT\_COROLLARY}} of mixed
positivity and the bounded energy identity.

\begin{proof}
For $m_a>0$ put $a_0=(a-\inf a)/m_a\in\mathcal P(\Omega)$.
Since $a_0^2$ is normalized classical and $Q_{\mathbf v}$ is
positive closed,
\[
 \Gamma_{Q_{\mathbf v}}(a)
  =m_a^2\left(\tfrac12D(a_0^2Q_{\mathbf v})
                              -a_0D(a_0Q_{\mathbf v})\right)
  \le\tfrac12m_a^2T_{S,q}(a_0^2,\mathbf v).
\]
If $m_a=0$ its energy is zero.  The same applies to $b$.
Positivity of $\Gamma_Q(a+b)$ gives
$\Gamma_Q(a-b)\le2\Gamma_Q(a)+2\Gamma_Q(b)$.
For any normalized tuple $(w_1,\ldots,w_q)$, the expansion of
$T_{S,q}(\sum_iw_i,\ldots,\sum_iw_i)$ contains the mixed
term $q!T_{S,q}(w_1,\ldots,w_q)$ and has all terms positive.
The average $q^{-1}\sum_iw_i$ is normalized, hence
\[
 T_{S,q}(w_1,\ldots,w_q)(E)
                         \le\frac{q^q}{q!}C_S(E;\Omega).
\]
Combining these inequalities proves the first line of (EC4).
Taking ordinary open neighborhoods proves its second line.
The remaining assertions follow by taking suprema of positive
measure inequalities and then ordinary open covers.
\end{proof}

\begin{theorem}[A uniform bounded coarea density]
\label{thm:r11-uniform-coarea-density}
Let $K\Subset\Omega$ be compact.  Choose an original smooth
compact function $0\le\chi\le1$, equal to one on a neighborhood
of $K$.  There is a finite constant $A_{\chi,S,a,b}$ such that,
for every normalized tuple $\mathbf v$, the coarea density of
Corollary~\ref{cor:r11-energy-bv-density} satisfies
\[
 \|k_{\chi,f,Q_{\mathbf v}}\|_{\infty}
 +\operatorname{Var}(k_{\chi,f,Q_{\mathbf v}})
                                      \le A_{\chi,S,a,b}.
 \tag{EC5}
\]
In particular, for every Borel $B\subset\mathbf R$,
\[
 \sup_{\mathbf v}
  \int_{\{f\in B\}}\chi\,d\Gamma_{Q_{\mathbf v}}(f)
                      \le A_{\chi,S,a,b}\,\mathcal L^1(B).
 \tag{EC6}
\]
The constants are independent of the tuple and of actual
moment refinements.  They may depend on the fixed original
compact geometry, the local coefficient masses of $S$, and
the indicated bounds for $a,b$.
\end{theorem}

\noindent\textbf{Status: \texttt{NEW\_PROVED}}.

\begin{proof}
By (LC10)--(LC12) it suffices to bound uniformly
\[
 \int\chi\,|D(fQ_{\mathbf v})|
                  +\int|\Gamma_{Q_{\mathbf v}}(\chi,f)|.
 \tag{EC7}
\]
Because $D(fQ_{\mathbf v})=D(aQ_{\mathbf v})-D(bQ_{\mathbf v})$
is a difference of two positive mixed BT measures, its first
term is bounded by the relative CLN estimate, uniformly in the
normalized tail $\mathbf v$.  For the second term use the
measure-valued Cauchy--Schwarz inequality on a fixed compact
neighborhood of $\operatorname{supp}\chi$:
\[
 \int|\Gamma_{Q_{\mathbf v}}(\chi,f)|
 \le
 \left(\int\Gamma_{Q_{\mathbf v}}(\chi)\right)^{1/2}
 \left(\int_{\operatorname{supp}\chi}
                                 \Gamma_{Q_{\mathbf v}}(f)\right)^{1/2}.
\]
The first factor is a fixed smooth coefficient pairing against
$Q_{\mathbf v}$ and is uniformly bounded by CLN.  The second
is uniformly bounded either by the frozen energy estimate or
by (EC4) and local finiteness of $C_S$ on that compact core.
A slightly larger compact smooth cutoff may be inserted in
both factors.  This proves a common bound for (EC7), and hence
(EC5), after increasing the constant by a factor of two.
Equation (EC6) follows by integrating the uniformly bounded
density.  All applications use the original compact test
sheaf; no chosen tropical pushforward of $S$ is involved.
\end{proof}

\begin{theorem}[Outer coarea control and energy-null level sets]
\label{thm:r11-outer-coarea-capacity}
For each compact $K\Subset\Omega$ there is a finite $A_K$
such that every subset $B\subset\mathbf R$, whether Borel or
not, satisfies
\[
 \boxed{\quad
 \mathsf e_{S,f}^{\mathrm o}
                  (K\cap f^{-1}(B);\Omega)
                 \le A_K\,\mathcal L^{1,*}(B).
 \quad}
 \tag{EC8}
\]
Consequently, every real-null set of values has relatively
energy-outer-null inverse image on every compact core.  In
particular each actual level set $\{f=t\}$ is locally null
for this outer energy capacity, simultaneously for all
normalized mixed tails in (EC2).
\end{theorem}

\noindent\textbf{Status: \texttt{NEW\_PROVED}}.  This is an
ordinary-open outer assertion, not only a zero-mass statement
for each individual energy measure.

\begin{proof}
Take $K\subset V\Subset\Omega$ and a compact smooth $\chi$
equal to one on $\overline V$.  The previous theorem gives
\[
 \mathsf e_{S,f}(V\cap f^{-1}(J);\Omega)
                                  \le A_K\mathcal L^1(J)
 \tag{EC9}
\]
for every real open set $J$.  Bounded relative quasicontinuity
on arbitrary original analytic neighborhoods, proved in
Corollary~\ref{cor:r11-analytic-bounded-relative-QC} by the new
analytic Laurent collar extension of the frozen bounded argument,
applied to $a,b$ and a finite compact cover, gives for each
$\epsilon>0$ an ordinary open $G$
with $C_S(G;\Omega)<\epsilon$ such that $f$ is continuous
on $\overline V\setminus G$.  Domain monotonicity permits all
exceptional sets from the finitely many smaller regularization
stages to be measured in the same capacity domain $\Omega$.

Choose a real open $J\supset B$ with its Lebesgue measure
arbitrarily close to $\mathcal L^{1,*}(B)$; when that outer
measure is infinite there is nothing to prove.  Relative
continuity supplies an ordinary open $H\subset V$ such that
\[
 H\cap(V\setminus G)
       = f^{-1}(J)\cap(V\setminus G).
\]
Then $O=H\cup(G\cap V)$ is an ordinary open neighborhood of
$K\cap f^{-1}(B)$ and
\[
 \begin{split}
 \mathsf e_{S,f}(O;\Omega)
 &\le\mathsf e_{S,f}(V\cap f^{-1}(J);\Omega)
                              +\mathsf e_{S,f}(G;\Omega)\\
 &\le A_K\mathcal L^1(J)+B_f\epsilon.
 \end{split}
\]
Take the infimum over such neighborhoods, then let $\epsilon$
and the real outer-cover error tend to zero.  This proves
(EC8) without any use of Borel outer regularity of $C_S$.
\end{proof}

\begin{proposition}[Scalar-set outer regularity and compact energy densities]
\label{prop:r11-scalar-energy-outer-regularity}
Fix an original compact smooth $\chi\geq0$, and set
\[
 \mathsf e_{\chi,S,f}(E)
   =\sup_{\mathbf v}\int_E\chi\,d\Gamma_{Q_{\mathbf v}}(f),
 \qquad
 c_{\chi,S,f}(B)=\mathsf e_{\chi,S,f}(f^{-1}(B)).
 \tag{EC12}
\]
Write $\mathsf e_{\chi,S,f}^{\mathrm o}$ for ordinary-open outer
regularization on $\Omega$.  For every real Borel set $B$,
\[
 \mathsf e_{\chi,S,f}^{\mathrm o}(f^{-1}(B))
       =\mathsf e_{\chi,S,f}(f^{-1}(B))
       =\sup_{\mathbf v}\int_B k_{\chi,f,Q_{\mathbf v}}(t)\,dt.
 \tag{EC13}
\]
The value capacity $c_{\chi,S,f}$ is continuous on increasing
Borel sequences and on decreasing Borel sequences.  It is inner
regular by compact real sets and outer regular by real open sets.
Its extension to arbitrary subsets by real-open outer covers is
a finite Choquet capacity on the real line.
Moreover the family of scalar measures
$f_*(\chi\Gamma_{Q_{\mathbf v}}(f))$ is precompact in total
variation, with a positive BV density for every limiting measure.
These statements concern scalar inverse images for the fixed $f$,
not arbitrary subsets of the analytic space.
\end{proposition}

\noindent\textbf{Status: \texttt{NEW\_PROVED}}.

\begin{proof}
Equations (LC12) and (EC5) give the last expression in (EC13)
and a finite constant $A$ such that all its densities are supported
in one compact real interval $I$ and bounded by $A$.  Thus
\[
 |c_{\chi,S,f}(B)-c_{\chi,S,f}(B')|
                 \le A\,\mathcal L^1((B\mathbin\triangle B')\cap I)
 \tag{EC14}
\]
for Borel $B,B'$.  Ordinary real measure regularity now proves
inner and outer regularity and continuity on decreasing sequences.
Continuity on increasing sequences also follows by commuting the
two suprema over sequence indices and normalized tails.
For completeness, every arbitrary subset has a Borel hull with
the same outer value, obtained by intersecting open neighborhoods
whose values tend to the infimum.  For increasing arbitrary sets
$A_j$, choose such Borel hulls $B_j$ and replace them by the
increasing Borel sets $H_j=\bigcap_{k\geq j}B_k\supset A_j$.
Monotonicity and Borel continuity from below give the same
continuity for the outer extension.  Continuity from above on
compact sets was already proved.  This verifies the stated
Choquet axioms on arbitrary subsets as well.

To prove the analytic outer identity, choose $V\Subset\Omega$
containing $\operatorname{supp}\chi$.  Given $\epsilon>0$,
bounded relative quasicontinuity supplies an ordinary open $G$
with $C_S(G;\Omega)<\epsilon$ such that $f$ is continuous on
$\overline V\setminus G$.  Choose a real open $J\supset B$
with $\mathcal L^1((J\setminus B)\cap I)<\epsilon$.
There is an ordinary open $H\subset V$ whose intersection
with $V\setminus G$ equals $f^{-1}(J)\cap(V\setminus G)$.
The ordinary open set
\[
 O=H\cup(G\cap V)\cup(\Omega\setminus\operatorname{supp}\chi)
\]
contains $f^{-1}(B)$, and the last summand has zero weighted
energy for every tail.  The bound (EC4) gives
\[
 \mathsf e_{\chi,S,f}(O)
 \le c_{\chi,S,f}(J)
          +\|\chi\|_\infty B_f\epsilon
 \le c_{\chi,S,f}(B)
          +(A+\|\chi\|_\infty B_f)\epsilon.
\]
The opposite inequality follows from monotonicity on every
open neighborhood.  Let $\epsilon$ tend to zero to obtain
(EC13).  This argument uses no outer regularity of $C_S$ on
general Borel sets.

Finally partition $I$ into intervals of mesh at most $h$ and
replace each density by its interval averages.  For a BV
function, the $L^1$ error of this replacement is at most
$h\operatorname{Var}(k)$.  The average vectors lie in a fixed
bounded finite-dimensional box, hence have compact closure.
The common variation bound in (EC5) proves total boundedness
of the original density family in $L^1(I)$.  Any resulting
$L^1$ limit remains nonnegative and retains the variation bound
by the distributional definition of variation.  Equality between
the $L^1$ distance of densities and total variation of their
measures proves the last assertion.
\end{proof}

\begin{theorem}[A one-dimensional Sobolev functional calculus]
\label{thm:r11-sobolev-scalar-calculus}
Let $I$ be a bounded interval containing the range of $f$ on
a neighborhood of a fixed compact core.  For
$F\in W^{1,2}(I)$ use its continuous one-dimensional
representative.  If smooth $F_j\to F$ in $W^{1,2}(I)$,
then $F_j(f)$ converge locally uniformly to $F(f)$ and are
Cauchy in every localized energy of $Q_{\mathbf v}$,
uniformly over normalized tails.  For a compact smooth cutoff
$\chi\ge0$ as above,
\[
 \sup_{\mathbf v}
 \int\chi\,d\Gamma_{Q_{\mathbf v}}
                       (F_j(f)-F_l(f))
       \le A_{\chi,S,a,b}\|F'_j-F'_l\|_{L^2(I)}^2.
 \tag{EC10}
\]
The resulting energy measures are independent of the smooth
approximations and satisfy
\[
 \Gamma_{Q_{\mathbf v}}(F(f),G(f))
           =F'(f)G'(f)\Gamma_{Q_{\mathbf v}}(f)
 \tag{EC11}
\]
for $F,G\in W^{1,2}(I)$.  Changing derivatives on real
Lebesgue-null sets has no effect.  The convergence of the
localized energy measures is in total variation, uniformly
over the same tails.

This constructs a relative finite-energy composition.  It
does not assert that a general $W^{1,2}$ composition is
classical psh, or a bounded difference of classical psh.
\end{theorem}

\noindent\textbf{Status: \texttt{NEW\_PROVED}}.

\begin{proof}
On a bounded interval the continuous representative satisfies
the elementary estimate
\[
 \|H\|_{\infty,I}
 \le |I|^{-1/2}\|H\|_{L^2(I)}
                         +|I|^{1/2}\|H'\|_{L^2(I)}.
\]
Thus the stated Sobolev convergence implies uniform
convergence on $I$, and hence after composition on the
chosen core.  For smooth $F_j,F_l$, the chain rule and
(EC6) give exactly (EC10).

The derivative products converge in $L^1(I)$, since
\[
 \|F'_jG'_j-F'G'\|_{L^1(I)}
 \le\|F'_j-F'\|_2\|G'_j\|_2
                       +\|F'\|_2\|G'_j-G'\|_2.
\]
The uniformly bounded pushforward density from (EC5) turns
this into convergence in total variation of the weighted
energy measures, uniformly over $\mathbf v$.  This proves
(EC11) and independence.  The absolute continuity in (LC12)
shows that the choice of an almost-everywhere representative
of $F',G'$ is immaterial.  Localizations agree because the
limits agree on every original compact test, so the
construction respects ordinary restriction and all actual
moment refinements.
\end{proof}

\paragraph{Scope in the final BT program.}
Theorem~\ref{thm:r11-outer-coarea-capacity} gives a way, valid for every allowed background, to remove bad
scalar values in uniform mixed energy estimates.  It also supplies a quantitative replacement
for informal claims that a gradient cannot charge its own
level set.  The family in (EC3), however, is formed from
energy measures of the fixed function $f$.  Curvature measures
$D(fQ)$ may charge those same levels.  Consequently (EC8)
does not identify arbitrary BT-capacity-zero sets, does not
prove unrestricted Borel regularity of $C_S$, and does not
declare every level set pluripolar.  Those stronger assertions
remain outside the theorem.

% END REVIEWED MODULE research_r11/capacity/uniform_coarea_capacity.tex


% BEGIN REVIEWED MODULE research_r11/energy/relative_nonpolar_truncation.tex
\section{Closed relative truncation products of unbounded classical potentials}
\label{sec:r11-relative-truncation}

Throughout this section $K/\mathbf Q_p$ is the fixed original finite
extension with its original absolute value.  The space $X$ is good,
Hausdorff, boundaryless, quasi-smooth and purely $n$-dimensional; $U$ is
ordinary open.  The background $S$ is an arbitrary original symmetric
CLD middle-positive, separately $d'$- and $d''$-closed current of
degree $(p,p)$.  Write
\[
 T_{S,I}(\mathbf u)=\bigwedge_{i\in I}D u_i\wedge S
\]
for the frozen bounded Bedford--Taylor product; $T_{S,\varnothing}=S$.
Products are always the original current products.  This notation
asserts neither a quotient for all actual delta-form presentations nor
equality against additional GK delta tests.

The following frozen inputs are used: bounded mixed products, their
weighted decreasing convergence, relative CLN, the bounded energy and
chain identities, and strict Borel maximum locality.  These are
\texttt{INTERNAL\_INPUT}, with sources R4 and the R2--R3 modules
preserved in the R10 archive.  In particular R3,
\path{pluripolar_zero_extension.tex}, is not reproved below: its
closed exceptional-set hypothesis does not apply directly to the
possibly nonclosed pole sets in the present construction.

\subsection{First derivatives and a plateau estimate}

\begin{lemma}[First-order coefficient measures for bounded potentials]
\label{lem:r11-bounded-first-order}
Let $T$ be any allowed original positive closed current of degree
$(P,P)$, $d=n-P\geq1$, and let $v$ be a locally bounded classical
psh function.  The currents $d'(vT)$ and $d''(vT)$ are coefficient-Radon.
For an original $(d-1,d)$-test $\alpha$, choose on its compact stage a
nonnegative smooth cutoff $\chi=1$ near its support and a positive
moment control $\beta=\sum_a d't_a\wedge d''t_a$.  There is a finite
constant $C_\alpha$ such that, for every bounded Borel $b$,
\[
 \left|\int b\,d'(vT)\wedge\alpha\right|^2
 \leq C_\alpha
 \left(\int\chi\,\Gamma_T(v)\wedge\beta^{d-1}\right)
 \left(\int\chi b^2\,\beta^d\wedge T\right).
 \tag{NP1}
\]
There is the analogous double-prime inequality.  Here
\[
 \Gamma_T(v)=\tfrac12D(v^2T)-vD(vT)
\]
is the frozen positive bounded-energy current.  For every $C^1$
function $F$ on a neighborhood of the bounded range under consideration,
\[
 d'(F(v)T)=F'(v)d'(vT),\qquad
 d''(F(v)T)=F'(v)d''(vT).
 \tag{NP2}
\]
These statements are local on ordinary opens and include arbitrary
additional actual moment coordinates used in $v$ or the tests.
\end{lemma}

\noindent\textbf{Status: \texttt{DIRECT\_COROLLARY}} of the
frozen weighted bounded calculus, with the first-order passage detailed
here to make the later Borel restrictions legitimate.

\begin{proof}
Choose bounded smooth decreasing regularizations $v_m$ on a fixed
neighborhood of the compact stage.  The bounded energy identity and
weighted joint decreasing convergence give
\[
 \int\chi\,\Gamma_T(v_m-v_l)\wedge\beta^{d-1}\longrightarrow0
                    \quad(\min(m,l)\longrightarrow\infty).
 \tag{NP3}
\]
This is the calculation in R4,
\path{prop:r4-classical-in-energy}: expand the energy of the
difference; the four weighted curvature terms have the same limit
and cancel, and its squared-function term tends to zero against fixed
coefficient measures.  The calculation uses bounded local calculus,
not the later global faithful energy construction.

For smooth $v_m-v_l$ the positive coefficient matrix of $T$ and
measure Cauchy--Schwarz bound the total variation of every fixed
coefficient of $d'((v_m-v_l)T)$ by a constant times the square root
of (NP3).  It follows that the first derivatives converge in local
coefficient total variation.  Since $v_mT\to vT$ on all original
tests by bounded convergence against the fixed coefficients of $T$,
the limit is $d'(vT)$.  The smooth coefficient inequality passes to
this total-variation limit and gives (NP1); it also gives the inequality
with a Borel restriction inserted into both positive factors.  Borel
weights are permitted by Radon measure Cauchy--Schwarz, not by an
assertion that a Borel-restricted background is closed.

The smooth identity for $F(v_m)$ passes to distributions.  On its
right-hand side the derivative measures converge in total variation,
the multipliers are uniformly bounded, and $F'(v_m)\to F'(v)$
pointwise.  Dominated convergence for the fixed limiting derivative
measure proves (NP2).  The double-prime argument is identical.
Finite partitions handle tests requiring several presentations.
The control $\beta$ only uses the fixed test coordinates; adjoining
the units used by each individual approximation preserves its positive
coefficient inequality.  No common finite moment for the entire
infinite sequence is assumed.
\end{proof}

\begin{lemma}[No first-order energy on the zero plateau]
\label{lem:r11-zero-plateau}
With $T$ as above, let $v\geq0$ be bounded classical psh.  Then
\[
 \mathbf1_{\{v=0\}}\Gamma_T(v)=0,\qquad
 \mathbf1_{\{v=0\}}d'(vT)=
 \mathbf1_{\{v=0\}}d''(vT)=0.
 \tag{NP4}
\]
All restrictions concern coefficient measures on the actual ordinary
space.
\end{lemma}

\noindent\textbf{Status: \texttt{NEW\_PROVED}} by the following
short bounded-energy argument; no general assertion about an arbitrary
Borel level set of a current is assumed.

\begin{proof}
For $\varepsilon>0$ set
\[
 g_\varepsilon=\max(v-\varepsilon,0),\qquad
 h_\varepsilon=v-g_\varepsilon=\min(v,\varepsilon).
\]
The functions $v$ and $g_\varepsilon$ are bounded classical, and $h_\varepsilon$
is their bounded difference.  For any fixed nonnegative compact
smooth $\chi$ and its moment control, the bounded energy identity gives
\begin{align*}
 0\leq\int\chi\Gamma_T(h_\varepsilon)\wedge\beta^{d-1}
 &\leq \tfrac12\varepsilon^2
     |\mu_{T,D\chi\wedge\beta^{d-1}}|(\operatorname{supp}\chi)\\
 &\quad+\varepsilon\int\chi
       \bigl(D(vT)+D(g_\varepsilon T)\bigr)\wedge\beta^{d-1}.
\end{align*}
The last masses are uniformly bounded as $\varepsilon\downarrow0$
by relative CLN, since $0\leq g_\varepsilon\leq v$.  The displayed
quantity therefore tends to zero.  Upper semicontinuity of $v$
makes $\{v<\varepsilon\}$ ordinarily open.  On that open
$h_\varepsilon=v$, so ordinary locality identifies their energy
currents there.  It follows that the restriction of the positive
energy current to $\{v=0\}$ is zero, for every original positive
control.  The restricted form of (NP1) now makes every first-order
coefficient zero on that same set.  Refining the moment presentation
does not change any of these identities.
\end{proof}

\subsection{Consistent truncations and closedness}

Let $u_1,\ldots,u_k\leq0$ be extended classical psh functions,
$1\leq k\leq n-p$.  Put
\[
 u_{i,j}=\max(u_i,-j),\qquad
 E_j=\bigcap_{i=1}^k\{u_i>-j\},\qquad
 T_j=T_{S,\{1,\ldots,k\}}(\mathbf u_j),\qquad
 R_j=\mathbf1_{E_j}T_j.
 \tag{NP5}
\]
The sets $E_j$ are Borel and need not be ordinarily open.  The
restriction $R_j$ need not be closed.

\begin{definition}[Finite relative truncation mass]
\label{def:r11-finite-truncation-mass}
The tuple has locally finite relative truncation mass if for every
original compact positive coefficient control $\Theta$ of degree
$(n-p-k,n-p-k)$ one has
\[
                         \sup_j\langle R_j,\Theta\rangle<\infty.
 \tag{NP6}
\]
It suffices to impose (NP6) on the finite positive controls dominating
each original compact test stage; the quantifier ranges over every
stage and its actual moment refinements.  In particular it is not a
single selected tropical trace bound.
\end{definition}

\begin{theorem}[Closed relative truncation product]
\label{thm:r11-closed-nonpolar-product}
Under (NP6), the increasing coefficient limit
\[
 \mathcal N_S(u_1,\ldots,u_k)=\lim_{j\to\infty}R_j
 \tag{NP7}
\]
is a well-defined original symmetric middle-positive,
separately $d'$- and $d''$-closed coefficient-Radon current of degree
$(p+k,p+k)$.  If
$P=\bigcup_i\{u_i=-\infty\}$, then
\[
 \mathbf1_P\mathcal N_S(\mathbf u)=0,
 \qquad
 \mathbf1_{E_j}\mathcal N_S(\mathbf u)=R_j.
 \tag{NP8}
\]
The construction is ordinary-local, symmetric in the potentials,
compatible with all original tests and actual moment refinements,
and independent of any cofinal choice of separate truncation levels.
For locally bounded potentials it is the frozen bounded BT product.
The background $S$ is arbitrary throughout.
\end{theorem}

\noindent\textbf{Status: \texttt{NEW\_PROVED}}.  The finite-mass
condition is an explicit necessary condition for the coefficient limit
to be a current; the new conclusion is its separate closedness across
all truncation interfaces, including a possibly nonclosed pole set.

\begin{proof}
For $l\geq j$, strict Borel maximum locality applied successively in
each slot gives
\[
                       \mathbf1_{E_j}T_l=\mathbf1_{E_j}T_j.
 \tag{NP9}
\]
Consequently $R_l-R_j$ is positive.  The local bounds (NP6) and the
finite positive controls yield Radon limits of all diagonal coefficient
measures; polarization gives the remaining coefficients.  Their linear
and refinement identities pass to the limits.  These same controls
give continuity on every original compact test stage.  Thus (NP7)
defines a current $N$, and (NP8) follows from (NP9) and monotone
measure convergence.  Its coefficients are carried by the union
of the $E_j$, which is $U\setminus P$.

If $p+k=n$, closedness is automatic by degree.  Suppose
$d=n-p-k\geq1$.  Fix an original $(d-1,d)$-test $\alpha$,
and choose $\chi,\beta$ on a compact stage as in (NP1).  Set
\[
 \psi=\sum_{i=1}^k u_i,\qquad
 v_j=\max(0,1+\psi/j),\qquad a_j=2v_j-v_j^2.
 \tag{NP10}
\]
The functions $v_j$ are actual bounded classical psh, between zero
and one.  This follows either by bounded decreasing approximation of
the summands or the extended classical maximum property.  They
increase to one on $U\setminus P$ and vanish on $P$.
Since all $u_i\leq0$, the set $F_j=\{v_j>0\}=\{\psi>-j\}$
is contained in $E_j$.  In particular, for $r=1,2$,
\[
 v_j^rT_j=v_j^rN,\qquad a_jT_j=a_jN.
 \tag{NP11}
\]
These are coefficient identities: on $F_j$ use (NP9), and on its
complement the multipliers vanish.  They do not assert closedness of
either Borel restriction.

The bounded energy identity, compact Stokes and positivity give
\begin{align*}
 B_j&:=\int\chi\Gamma_{T_j}(v_j)\wedge\beta^{d-1}\\
 &=\tfrac12\langle v_j^2T_j,D\chi\wedge\beta^{d-1}\rangle
      -\int\chi v_jD(v_jT_j)\wedge\beta^{d-1}\\
 &\leq\tfrac12|\mu_{N,D\chi\wedge\beta^{d-1}}|
                      (\operatorname{supp}\chi)=:B<\infty.
 \tag{NP12}
\end{align*}
In the last line we used (NP11).  Thus the bound does not require
uniform mass of the uncut currents $T_j$ at their moving interfaces.

By (NP2), $d'(a_jT_j)=2(1-v_j)d'(v_jT_j)$.  Lemma
\ref{lem:r11-zero-plateau} allows the first-order factor to be
restricted to $F_j$.  Apply (NP1) with the Borel multiplier
$\mathbf1_{F_j}(1-v_j)$.  Equations (NP9) and (NP12) yield
\[
 |\langle d'(a_jT_j),\alpha\rangle|^2
 \leq4C_\alpha B
       \int\chi(1-v_j)^2\,\beta^d\wedge N
                       \longrightarrow0.
 \tag{NP13}
\]
Dominated convergence applies to the fixed finite coefficient measure
on the right, because $N$ has no coefficient mass on $P$.
At the same time (NP11) and dominated convergence give $a_jT_j\to N$
on all original tests.  Distributional continuity now implies
$d'N=0$.  The double-prime estimate proves $d''N=0$.

The construction and proof commute with ordinary restriction and
common actual refinement.  For separate level vectors, the same
strict Borel locality gives a directed consistent family; diagonal
integer levels are cofinal.  This proves independence of that choice
and symmetry.  If all potentials are locally bounded, their
truncations agree with them on any fixed compact neighborhood once
the levels are large enough, proving compatibility with bounded BT.
\end{proof}

\begin{corollary}[Background linearity and pole restrictions]
\label{cor:r11-nonpolar-background-restriction}
On the finite-mass domain of the preceding theorem, $\mathcal N_S$
is additive and nonnegative-real homogeneous in the original positive
closed background.  If $A$ is a complete classical pole set, then
\[
 \mathbf1_A\mathcal N_S(\mathbf u)
                    =\mathcal N_{\mathbf1_A S}(\mathbf u).
 \tag{NP14}
\]
The right side is defined by the same truncations and automatically
has locally finite coefficient mass.  If $S$ has no coefficient mass
on any local pluripolar set, the same is true of $\mathcal N_S$.
For arbitrary $S$, (NP8) only removes the poles of the stated tuple.
\end{corollary}

\noindent\textbf{Status: \texttt{DIRECT\_COROLLARY}} using the
frozen complete-pole restriction theorem for bounded products.

\begin{proof}
Background linearity holds at every truncated level and passes through
the increasing positive limits.  The frozen pole restriction identity
gives $\mathbf1_A T_j=T_{\mathbf1_A S,k}(\mathbf u_j)$;
the restricted background is separately closed by that theorem.
Restrict also to $E_j$ and take the increasing coefficient limits.
The required mass bounds are inherited from $N$.  Localizing the
result to a complete pole witness proves the final assertion.
\end{proof}

\begin{proposition}[Potential additivity can fail over a polar background]
\label{prop:r11-nonpolar-additivity-failure}
For every $n\geq2$ there are an allowed original $n$-dimensional
ordinary domain, a positive closed background $S$ of degree $(1,1)$,
a bounded classical $u$, and an extended classical $v$ for which all
three relative truncation products exist and
\[
 \mathcal N_S(u+v)=0=\mathcal N_S(v),
                   \qquad \mathcal N_S(u)\ne0.
 \tag{NP15}
\]
Thus unrestricted additivity in unbounded potentials, while retaining
every polar background component, is false for (NP7).
\end{proposition}

\noindent\textbf{Status: \texttt{NEW\_PROVED}} for the explicit
original example; the unrestricted additivity assertion is
\texttt{FALSE}.  This is not a counterexample to the frozen bounded
multilinearity theorem.

\begin{proof}
Take an ordinary product of the open unit disc in coordinate $z$
and $n-1$ bounded open annuli in unit coordinates $t_1,\ldots,t_{n-1}$.
Let $Z=\{z=0\}$ and $S=[Z]$.  Put
\[
 v=\log|z|,\qquad
 u=\tfrac12(\log|t_1|)^2-C,
\]
where $C$ makes $u\leq0$ on the domain.  The logarithm is an extended
classical psh function, with value $-\infty$ on $Z$; $u$ is bounded
classical and its curvature on $Z$ is the nonzero original positive
form $d'\log|t_1|\wedge d''\log|t_1|$.  Hence $D(uS)\ne0$;
one can test it with a nonnegative compact product test in the
remaining annular directions.  The ordinary original product
calibration has positive degree, so the pairing is strictly positive.

Each truncated product over $S$ is supported on $Z$, by ordinary
locality of the frozen calculus.  The finite loci of $v$ and $u+v$
are disjoint from $Z$.  Their cut currents $R_j$ are therefore zero
at every level.  This gives (NP15) and all its mass assertions.
The example uses a special background only to disprove the universal
additivity assertion; no restriction on $S$ was used in the theorem.
\end{proof}

\paragraph{The remaining unrestricted issue.}
Theorem~\ref{thm:r11-closed-nonpolar-product} does not make (NP6)
automatic for every extended classical tuple.  Nor does it identify
relative polar null sets with ambient pluripolarity for every background.
The finite-energy criterion in the next section gives a checkable large
domain on which the uncut products also converge and the usual
recursive differential identity survives.

% END REVIEWED MODULE research_r11/energy/relative_nonpolar_truncation.tex


% BEGIN REVIEWED MODULE research_r11/energy/mixed_superlinear_energy.tex
\section{A mixed superlinear trace class and unbounded BT continuity}
\label{sec:r11-superlinear-energy}

Keep the original analytic category and arbitrary background $S$ from
Section~\ref{sec:r11-relative-truncation}.  Fix $1\leq m\leq n-p$
and an ordered tuple of nonpositive extended classical psh functions
$\mathbf u=(u_1,\ldots,u_m)$.  Repetition of a function is allowed.
For $I\subseteq\{1,\ldots,m\}$ let
\[
 T_{I,j}=T_{S,I}(u_{1,j},\ldots,u_{m,j}),\qquad
 u_{i,j}=\max(u_i,-j),\qquad
 W_j=1-\sum_{i=1}^m u_{i,j},\qquad W=1-\sum_i u_i.
 \tag{ME1}
\]
The empty product is $T_{\varnothing,j}=S$.  Each trace below uses
an original nonnegative compact positive control $\Theta$ of the
complementary degree.  This includes every finite control family
needed to dominate an arbitrary original compact smooth test, and
all additional actual moment coordinates.  No one fixed background
trace measure is presumed to dominate all singular products.

\begin{definition}[The local mixed superlinear trace class]
\label{def:r11-mixed-superlinear-class}
For $s>1$, write $\mathbf u\in\mathcal E^{s,\mathrm{mix}}_{S,\mathrm{loc}}(U)$
if for every subset $I$, every such control $\Theta$, and every
compact stage on which it is defined, there is a finite constant
$C_{I,\Theta}$ such that
\[
               \sup_{j\geq1}\int W_j^s\,\Theta\wedge T_{I,j}
                                  \leq C_{I,\Theta}.
 \tag{ME2}
\]
This is a condition on bounded products already constructed by the
frozen theory.  In particular it neither presupposes an unbounded
product nor its continuity or closedness.  The background itself is
arbitrary; integrability is a condition on the tuple relative to it.
Every locally bounded classical tuple belongs to this local class.
\end{definition}

\begin{lemma}[Mixed tail control]
\label{lem:r11-mixed-tail-control}
Under (ME2), each nonempty partial tuple has locally finite truncation
mass, and hence a closed relative truncation product $N_I$ by
Theorem~\ref{thm:r11-closed-nonpolar-product}.  Set $N_\varnothing=S$.
Then
\[
 \int W^s\,\Theta\wedge N_I\leq C_{I,\Theta},\qquad
 \mathbf1_{\bigcup_{i=1}^m\{u_i=-\infty\}}N_I=0.
 \tag{ME3}
\]
With $E_j=\bigcap_{i=1}^m\{u_i>-j\}$, one has
\[
 \left\|\Theta\wedge(T_{I,j}-N_I)\right\|_{\mathrm{TV}}
          \leq\frac{2C_{I,\Theta}}{(1+j)^s}.
 \tag{ME4}
\]
For every $i$, the product $u_iN_I$ is coefficient-defined on all
original tests and
\[
 \left\|\Theta\wedge(u_{i,j}T_{I,j}-u_iN_I)\right\|_{\mathrm{TV}}
          \leq\frac{2C_{I,\Theta}}{(1+j)^{s-1}}.
 \tag{ME5}
\]
On arbitrary signed coefficient tests the same assertions hold with
the corresponding finite sums of positive-control constants.
\end{lemma}

\noindent\textbf{Status: \texttt{NEW\_PROVED}}.  The estimates
also control poles of potentials that are not differentiated in $I$.
This is why all partial products and the common weight occur in the
definition.

\begin{proof}
Since $W_j\geq1$, (ME2) bounds all uncut traces, hence the cut traces.
For $I\ne\varnothing$ write
$R_{I,l}=\mathbf1_{\bigcap_{i\in I}\{u_i>-l\}}T_{I,l}$.
For fixed $j$ and $l\geq j$, monotonicity $W_j\leq W_l$ gives
\[
 \int W_j^s\Theta\wedge R_{I,l}
       \leq\int W_l^s\Theta\wedge T_{I,l}\leq C_{I,\Theta}.
\]
First let $l\to\infty$ using the increasing scalar coefficient
measures of $R_{I,l}$, and then let $j\to\infty$ by monotone
convergence for the fixed measure $\Theta\wedge N_I$.
This proves the first assertion of (ME3).  For $I=\varnothing$
it follows directly by monotone convergence against $\Theta\wedge S$.
The common pole union is the set where $W=+\infty$, so the second
assertion follows from the first and all original positive controls.

Strict Borel locality gives
$\mathbf1_{E_j}T_{I,j}=\mathbf1_{E_j}N_I$.  On $E_j^c$ at least
one $u_{i,j}$ equals $-j$, and consequently $W_j\geq1+j$.
Also $W\geq1+j$ there.  Therefore
\[
 (\Theta\wedge T_{I,j})(E_j^c)\leq C_{I,\Theta}(1+j)^{-s},
 \qquad
 (\Theta\wedge N_I)(E_j^c)\leq C_{I,\Theta}(1+j)^{-s}.
\]
The two measures agree on $E_j$, which proves (ME4).
On that same set the two weighted measures in (ME5) agree.
On its complement use $|u_{i,j}|\leq W_j$ and $|u_i|\leq W$;
the two weighted tails are each at most
$C_{I,\Theta}(1+j)^{1-s}$.  This proves (ME5) and local absolute
integrability.  Finite positive coefficient domination proves current
continuity of the multipliers and the statement about general tests.
No varying-measure weak limit has been restricted to a Borel set.
\end{proof}

\begin{theorem}[Recursive unbounded products in every permitted degree]
\label{thm:r11-superlinear-recursion}
For $\mathbf u\in\mathcal E^{s,\mathrm{mix}}_{S,\mathrm{loc}}(U)$,
the partial products satisfy
\[
                       D(u_iN_I)=N_{I\cup\{i\}}
              \qquad (i\notin I).
 \tag{ME6}
\]
Thus they give a symmetric, separately closed, middle-positive
unbounded Bedford--Taylor calculus on this class over the original
arbitrary $S$.  Every permutation of the recursive construction gives
the same current, and every identity is valid on all original compact
smooth complementary tests and actual moment refinements.  The
empty-slot case is $D(u_iS)=N_{\{i\}}$.
\end{theorem}

\noindent\textbf{Status: \texttt{NEW\_PROVED}}.

\begin{proof}
The frozen bounded recursion says
$D(u_{i,j}T_{I,j})=T_{I\cup\{i\},j}$.
The left-hand multiplier converges on every original compact test
by (ME5), so its distributional double derivative converges to
$D(u_iN_I)$.  The right side converges on every original test by
(ME4), to $N_{I\cup\{i\}}$.  This proves (ME6).
The same argument applies in every order.  Symmetry and all
refinement identities are inherited from the bounded products before
the limits are taken.  Separate closedness and positivity follow
either from the preceding section or from (ME4) and the uncut
positive closed approximants.
\end{proof}

\begin{lemma}[Continuous bounded weights in the joint bounded limit]
\label{lem:r11-continuous-bounded-joint-weight}
Suppose $a_i^{(l)}\downarrow a_i$ are bounded classical psh functions
with one common finite bound on a compact regularization stage.
For every continuous real function $F$ on the resulting compact
range box and each bounded mixed product in these potentials,
\[
 F(\mathbf a^{(l)})T_{S,I}(\mathbf a^{(l)})
             \longrightarrow F(\mathbf a)T_{S,I}(\mathbf a)
 \tag{ME7}
\]
on all original compact tests on that stage.
\end{lemma}

\noindent\textbf{Status: \texttt{DIRECT\_COROLLARY}} of the
frozen weighted joint decreasing theorem and elementary polynomial
approximation.  This extends the allowed scalar weights only in the
specified joint bounded limit; it is not a statement for arbitrary
Borel functions of a moving tuple.

\begin{proof}
First let $F$ be a polynomial.  On a slightly larger bounded box
write $F=G-H$, where both $G,H$ are convex and coordinatewise
nondecreasing: add a sufficiently large positive quadratic form and
then a sufficiently large positive linear function to $F$, and use
that same quadratic-plus-linear function for $H$.  Their bounded
compositions with the tuple are classical psh and decrease jointly.
The frozen theorem with one bounded undifferentiated classical weight
applies to both, proving (ME7) for $F$.  Convex monotone composition
is checked on the same bounded range and after every complete
extension, as in the classical definition; no current positivity to
classicality implication is used.  Continuous $F$ are uniformly
approximated by polynomials on the box.  The uniform local CLN mass
bound controls the resulting error for each original test, and proves
the assertion.
\end{proof}

\begin{theorem}[Joint decreasing convergence under uniform mixed trace energy]
\label{thm:r11-unbounded-joint-continuity}
Let $\mathbf u^{(l)}=(u_1^{(l)},\ldots,u_m^{(l)})$ be nonpositive
extended classical tuples with $u_i^{(l)}\downarrow u_i$ at every
actual point.  Suppose that, for one $s>1$, (ME2) holds with constants
$C_{I,\Theta}$ independent of both $l$ and the truncation level $j$.
Then the limiting tuple also satisfies (ME2), and
\[
 \mathcal N_S(\mathbf u_I^{(l)})
                  \longrightarrow\mathcal N_S(\mathbf u_I)
 \tag{ME8}
\]
on every original test, for every partial tuple $I$.

More generally, if $F:\mathbf R^m\to\mathbf R$ is continuous on the
nonpositive orthant and for some $0\leq a<s$ satisfies
\[
 |F(x_1,\ldots,x_m)|\leq A(1-\textstyle\sum_i x_i)^a,
 \tag{ME9}
\]
then the coefficient-defined currents satisfy
\[
 F(\mathbf u^{(l)})\mathcal N_S(\mathbf u_I^{(l)})
             \longrightarrow F(\mathbf u)\mathcal N_S(\mathbf u_I).
 \tag{ME10}
\]
Values of the multipliers on their tuple pole unions are immaterial,
because (ME3) makes those sets coefficient-null.  All conclusions are
ordinary-local, include every original compact smooth test, and are
unchanged by every actual moment refinement of those tests.
\end{theorem}

\noindent\textbf{Status: \texttt{NEW\_PROVED}}.  The uniform
superlinear hypothesis is retained; arbitrary decreasing unbounded
sequences are not asserted to have continuous products.

\begin{proof}
Fix $j$.  The truncated rows $u_i^{(l)}\vee(-j)$ decrease to
$u_i\vee(-j)$ and lie in $[-j,0]$.  Lemma
\ref{lem:r11-continuous-bounded-joint-weight}, with
$F(\mathbf x)=(1-\sum_i x_i)^s$, passes the uniform (ME2) bound
to the limiting row.  This works for every $j,I,\Theta$, so the
limit belongs to the stated class.

For a fixed original test $\eta$, dominate its coefficients by
finitely many controls $\Theta$.  Equation (ME4) bounds, uniformly
in $l$, the difference between its unbounded product pairing and
the pairing with the $j$-truncated product by a constant times
$(1+j)^{-s}$.  The analogous estimate holds for the limiting tuple.
Between these two errors use the bounded decreasing theorem at fixed
$j$.  First let $l\to\infty$ and then $j\to\infty$ to prove (ME8).

For (ME10), let $W_j^{(l)},W^{(l)}$ denote the corresponding common
weights, and $E_j^{(l)}$ the common finite-level set.  On that set
the weighted truncated product and the weighted unbounded product
agree.  On the complement, (ME3), (ME9) and the lower bound
$W_j^{(l)},W^{(l)}\geq1+j$ give the uniform tail estimate
\[
 \left\|\Theta\wedge\left(
 F(\mathbf u_j^{(l)})T_{I,j}^{(l)}
       -F(\mathbf u^{(l)})N_I^{(l)}\right)\right\|_{\mathrm{TV}}
       \leq2A C_{I,\Theta}(1+j)^{a-s}.
 \tag{ME11}
\]
Apply the same estimate to the limiting tuple.  At fixed $j$, (ME7)
gives the middle limit with the continuous bounded weight $F$ on
$[-j,0]^m$.  The two tail bounds tend to zero because $a<s$.
This proves (ME10), including the ordinary coefficient-current
continuity required for the stated distributions.
\end{proof}

\paragraph{Scope relative to the final theory.}
These results provide an unbounded mixed calculus for a quantified
class over every original background.  They do not identify this class
with every finite-capacity or finite-Dirichlet-energy class, prove an
unrestricted polar-null equivalence, or solve the nonlinear Dirichlet
problem.  In particular the counterexample (NP15) explains why
unbounded potential additivity cannot simply be asserted after
discarding moving pole sets over a polar background.

% END REVIEWED MODULE research_r11/energy/mixed_superlinear_energy.tex


% BEGIN REVIEWED MODULE research_r11/energy/unbounded_energy_chain_calculus.tex
\section{Unbounded energy currents and a full multivariable chain rule}
\label{sec:r11-unbounded-energy-chain}

This section strengthens the finite-energy domain in order to control
energy measures as well as curvature currents.  It works in every
ambient dimension and every permitted mixed degree, with the original
arbitrary middle-positive closed background $S$.  Its stronger moment
condition is stated before any unbounded energy is defined.

Put $q=n-p\geq1$, and let $u_1,\ldots,u_m\leq0$ be a finite family
of extended classical psh functions; here $m$ need not be at most $q$.
For a multiindex $\alpha\in\mathbf N^m$, $|\alpha|\leq q$, write
\[
 Q_{\alpha,j}=
       (Du_{1,j})^{\alpha_1}\wedge\cdots\wedge
       (Du_{m,j})^{\alpha_m}\wedge S,
 \qquad u_{i,j}=\max(u_i,-j),\qquad
 W_j=1-\sum_i u_{i,j}.
 \tag{UE1}
\]
Every product on the right is a frozen bounded BT product, including
repeated slots.  Set $W=1-\sum_i u_i$ and
$E_j=\bigcap_i\{u_i>-j\}$.

\begin{definition}[Complete mixed trace control]
\label{def:r11-complete-mixed-trace}
The family has complete mixed trace control of exponent $s>2$ if for
every $|\alpha|\leq q$ and every original compact positive coefficient
control $\Theta$ of complementary degree,
\[
                 \sup_j\int W_j^s\,\Theta\wedge Q_{\alpha,j}
                              \leq C_{\alpha,\Theta}<\infty.
 \tag{UE2}
\]
The quantifier includes $\alpha=0$, all original test stages and
their actual moment refinements.  It is a finite mixed-product
condition for each fixed stage and finite family.  The requirement
on every stage ensures the full original-current meaning.
\end{definition}

\begin{lemma}[Truncation locality for bounded energy currents]
\label{lem:r11-energy-truncation-locality}
For $l\geq j$ and $|\alpha|\leq q-1$ one has
\[
 \mathbf1_{E_j}\Gamma_{Q_{\alpha,l}}(u_{i,l},u_{a,l})
       =\mathbf1_{E_j}\Gamma_{Q_{\alpha,j}}(u_{i,j},u_{a,j}).
 \tag{UE3}
\]
These are identities of original coefficient measures.  In particular
the positive diagonal cut energy currents are increasing with $j$.
\end{lemma}

\noindent\textbf{Status: \texttt{DIRECT\_COROLLARY}} of strict
Borel bounded maximum locality, with an explicit reduction to bounded
classical slots.

\begin{proof}
First keep a fixed positive closed background $Q$ and a bounded
classical $v\geq-l$.  Put $v_j=\max(v,-j)$ and choose $c>l$.
The functions $v+c$ and $v_j+c$ are positive classical, and
\[
 (v_j+c)^2=\max\{(v+c)^2,(c-j)^2\}.
\]
All functions inside the maximum are bounded classical.  Strict
Borel maximum locality, both for this identity and for $v_j$, and
the energy formula
\[
 \Gamma_Q(v)=\tfrac12D((v+c)^2Q)-(v+c)D(vQ)
\]
show that $\Gamma_Q(v_j)=\Gamma_Q(v)$ after restriction to
$\{v>-j\}$.  Next switch the bounded background tails from
$Q_{\alpha,l}$ to $Q_{\alpha,j}$ on $E_j$.  Each of the two terms
in this energy formula is a bounded mixed product with a fixed
bounded classical first slot, followed by an actual Borel multiplier.
Applying strict maximum locality in the background slots therefore
gives the diagonal case of (UE3).

For cross energies use polarization.  The sum
$u_{i,j}+u_{a,j}$ is the maximum of the four classical branches
$u_{i,l}+u_{a,l}$, $u_{i,l}-j$, $u_{a,l}-j$, and $-2j$.
On $E_j$ the first branch is strictly largest.  Add a constant
making every branch positive and square; the same maximum locality
argument applies to its energy.  Polarization now proves (UE3).
No ordinary openness of $E_j$ is used.
\end{proof}

\begin{lemma}[A uniform superquadratic energy moment]
\label{lem:r11-superquadratic-energy-moment}
Assume (UE2).  For $|\alpha|\leq q-1$, every fixed compact moment
stage, and a nonnegative smooth cutoff $\chi$ on that stage, put
$d=q-|\alpha|\geq1$ and let $\beta$ be its original positive
moment control.  There is a finite constant $A_{\alpha,\chi,\beta}$
such that
\[
 \sup_j\sum_{i=1}^m
  \int\chi W_j^{s-2}\,
       \Gamma_{Q_{\alpha,j}}(u_{i,j})\wedge\beta^{d-1}
                         \leq A_{\alpha,\chi,\beta}.
 \tag{UE4}
\]
The constant uses only finitely many of the original trace bounds
in (UE2) on that compact stage.
\end{lemma}

\noindent\textbf{Status: \texttt{NEW\_PROVED}}.

\begin{proof}
On the nonpositive orthant set
\[
 F(x)=(1+\sum_i x_i^2)^{s/2},\quad
 r(x)=(1+\sum_i x_i^2)^{1/2},\quad
 w(x)=1-\sum_i x_i.
\]
Then $r\leq w\leq\sqrt{m+1}\,r$, and
\[
 0\leq F\leq w^s,\quad |F_i|\leq s w^{s-1},\quad
 (F_{ia})\geq c_{s,m}w^{s-2}I_m,
 \quad c_{s,m}=s(m+1)^{-(s-2)/2}>0.
 \tag{UE5}
\]
Indeed the Hessian is
$s r^{s-2}I_m+s(s-2)r^{s-4}xx^{\mathsf T}$.

Apply the bounded multivariable chain rule to $F(\mathbf u_j)$
and the closed background $Q_{\alpha,j}$.  This rule follows from
the frozen smooth chain formula by the same bounded convex-increasing
decomposition and weighted limit as
Lemma~\ref{lem:r11-bounded-difference-chain}; no unbounded chain
rule is being used here.  Compact Stokes gives
\begin{align*}
 &\int\chi\sum_{i,a}F_{ia}(\mathbf u_j)
       \Gamma_{Q_{\alpha,j}}(u_{i,j},u_{a,j})\wedge\beta^{d-1}\\
 &\quad=\langle F(\mathbf u_j)Q_{\alpha,j},
                                  D\chi\wedge\beta^{d-1}\rangle
     -\sum_i\int\chi F_i(\mathbf u_j)
                              Q_{\alpha+e_i,j}\wedge\beta^{d-1}.
 \tag{UE6}
\end{align*}
The energy matrix is positive semidefinite after pairing with any
original positive coefficient control.  The Hessian bound in (UE5)
therefore bounds the left side below by $c_{s,m}$ times the quantity
in (UE4).  The absolute value of the first term on the right is
controlled by finitely many (UE2) bounds for $Q_{\alpha,j}$; include
the coordinates of $D\chi$ in the compact controls.  The remaining
terms are controlled by (UE2) for $\alpha+e_i$, since
$|F_i|\leq sW_j^{s-1}\leq sW_j^s$.  These indices have degree at
most $q$.  All bounds are independent of $j$, proving (UE4).
\end{proof}

\begin{theorem}[Unbounded energy currents and strong truncation limits]
\label{thm:r11-unbounded-energy-current}
Under (UE2), let $N_\alpha$ be the closed relative truncation product
of the corresponding repeated tuple, with $N_0=S$.
For $|\alpha|\leq q-1$ the distribution
\[
 \mathcal G^\alpha_{ia}
 =\tfrac12\left\{D(u_i u_aN_\alpha)
             -u_iN_{\alpha+e_a}-u_aN_{\alpha+e_i}\right\}
 \tag{UE7}
\]
is a coefficient-Radon current.  The matrix of these currents is
positive semidefinite in the original middle cone, so in particular
$\mathcal G^\alpha_{ii}\geq0$.  Every entry gives zero coefficient
mass to the common pole union $\bigcup_i\{u_i=-\infty\}$.
For each original compact positive control $\Theta$ there is a
finite constant $A_\Theta$ such that
\[
 \left\|\Theta\wedge\left(
       \Gamma_{Q_{\alpha,j}}(u_{i,j},u_{a,j})
                          -\mathcal G^\alpha_{ia}\right)\right\|_{\mathrm{TV}}
                         \leq A_\Theta(1+j)^{2-s}.
 \tag{UE8}
\]
There are also the weighted bounds
\[
 \int W^{s-2}\,\Theta\wedge\mathcal G^\alpha_{ii}<\infty.
 \tag{UE9}
\]
The definitions and all limits are intrinsic to the actual potentials
and original background, on every original test and actual refinement.
These energy currents need not be closed.
\end{theorem}

\noindent\textbf{Status: \texttt{NEW\_PROVED}}.

\begin{proof}
The proof of Lemma~\ref{lem:r11-mixed-tail-control} applies verbatim
to the multiindices in (UE1), with the common weight $W_j$ and
finite locus $E_j$.  Repeated slots cause no change to strict maximum
locality.  It gives closed $N_\alpha$, moment bounds
$\int W^s\Theta\wedge N_\alpha\leq C_{\alpha,\Theta}$,
and local total-variation convergence of the uncut curvatures.
For the empty product it is simply monotone integrability against $S$.
It also gives the recursive identity
$D(u_iN_\alpha)=N_{\alpha+e_i}$, by the same weighted tail
argument and bounded recursion as (ME6).

For each diagonal entry form the positive cut energy current
\[
 B_{i,j}=\mathbf1_{E_j}
                   \Gamma_{Q_{\alpha,j}}(u_{i,j}).
\]
It is increasing by (UE3).  Estimate (UE4), with finite positive
controls, makes its coefficient limit a locally finite positive
Radon current $B_i$.  For fixed $j$ pass $l\to\infty$ in
$\int W_j^{s-2}\Theta\wedge B_{i,l}$, using $W_j\leq W_l$,
then let $j\to\infty$.  This proves the weighted bound (UE9)
for $B_i$ and its zero restriction to the common pole union.
On $E_j$ the full energy and $B_i$ agree.  On $E_j^c$ both
weights $W_j,W$ are at least $1+j$, so (UE4) and (UE9) bound
both energy tails by a constant times $(1+j)^{2-s}$.  Hence the
uncut diagonal energies converge in local coefficient total variation.
Polarization, or the measure Cauchy--Schwarz bound for the energy
matrix, proves the same assertion for its cross entries.

It remains to identify these limits intrinsically.  The functions
$u_i u_a$ have growth at most $W^2$.  Since $s>2$, the weighted
tail proof (ME11), now for a fixed tuple, gives
$u_{i,j}u_{a,j}Q_{\alpha,j}\to u_i u_aN_\alpha$ on all original
tests, indeed in local coefficient total variation.  The same is true
for $u_{i,j}Q_{\alpha+e_a,j}$ and $u_{a,j}Q_{\alpha+e_i,j}$.
Pass through the bounded energy identity in distributions.  It
identifies the limits with (UE7), proving all assertions.
\end{proof}

\begin{theorem}[The unbounded multivariable chain rule]
\label{thm:r11-unbounded-multivariable-chain}
Assume (UE2), $|\alpha|\leq q-1$, and let $F$ be $C^2$ on a
neighborhood of the nonpositive orthant.  Suppose that for some
$2\leq r<s$ and constant $A$,
\[
 |F(x)|\leq A w(x)^r,\quad
 |F_i(x)|\leq A w(x)^{r-1},\quad
 |F_{ia}(x)|\leq A w(x)^{r-2}.
 \tag{UE10}
\]
Then every term in the following equality is coefficient-defined,
and the equality holds on all original compact smooth tests:
\[
 D(F(\mathbf u)N_\alpha)
   =\sum_i F_i(\mathbf u)N_{\alpha+e_i}
      +\sum_{i,a}F_{ia}(\mathbf u)\mathcal G^\alpha_{ia}.
 \tag{UE11}
\]
If $F$ is convex and coordinatewise nondecreasing, its right side is
middle-positive.  The formula includes the usual energy identity and
compactly localized Green identities for unbounded potentials in this
class.  It does not assert an algebra of arbitrary Borel functions.
\end{theorem}

\noindent\textbf{Status: \texttt{NEW\_PROVED}}.

\begin{proof}
Apply the bounded multivariable chain rule at each level $j$.
The weighted curvature tail bounds give convergence of
$F(\mathbf u_j)Q_{\alpha,j}$ and
$F_i(\mathbf u_j)Q_{\alpha+e_i,j}$ to the corresponding terms
in (UE11).  For the energy term use (UE3), (UE4), and (UE9).
On $E_j$ the bounded and limiting energy measures and multipliers
agree.  On the complement, (UE10) bounds the two tails by a constant
times $(1+j)^{r-s}$, which tends to zero.  All these convergences
are in the local total variation of the required coefficient
measures.  Distributional continuity proves (UE11).
Finally convexity makes the Hessian positive semidefinite and
coordinatewise monotonicity gives $F_i\geq0$.  Positivity of the
curvatures and the energy matrix proves the last claim.  None of
these steps specializes the background or changes its absolute value.
\end{proof}

\begin{corollary}[Joint decreasing continuity of unbounded energies]
\label{cor:r11-unbounded-energy-decreasing}
Let $\mathbf u^{(l)}\downarrow\mathbf u$ pointwise be extended
classical families satisfying (UE2) with one exponent $s>2$ and
constants independent of $l,j$.  Then their energy currents in
(UE7) converge on every original compact test.  The same remains
true after multiplication by any continuous function $H(\mathbf u)$
with $|H(x)|\leq A w(x)^b$ for $0\leq b<s-2$.
\end{corollary}

\noindent\textbf{Status: \texttt{DIRECT\_COROLLARY}} of the
uniform tail bounds and the frozen bounded weighted energy theorem.

\begin{proof}
For fixed truncation level, first write each polynomial
$x_i x_a=G(x)-H(x)$ on the bounded range box, where $G,H$ are
convex and coordinatewise nondecreasing, by adding a sufficiently
large positive quadratic and a sufficiently large positive linear
function.  The three-term energy formula is then a finite sum of
bounded BT currents $D(G(\mathbf u)Q)$ and $D(H(\mathbf u)Q)$,
and bounded actual multipliers of such currents.  Each composition
decreases jointly with the bounded tuple.  A continuous additional
weight is handled by the bounded polynomial approximation and
convex-increasing decompositions in (ME7), using the uniform positive
coefficient masses of these finitely many currents.  Frozen weighted
bounded BT convergence therefore proves weighted convergence of the
energy currents; no nonsmooth weight is applied to an unidentified
distributional derivative.  Estimates (UE4)--(UE9) are uniform
in $l$.  They bound the two truncation errors, including the weight
$H$, by a constant times $(1+j)^{b+2-s}$.  First take the bounded
limit at fixed $j$, then let $j\to\infty$.  Since $b+2<s$,
the errors vanish.  This also proves compatibility with all original
tests and actual refinements.
\end{proof}

% END REVIEWED MODULE research_r11/energy/unbounded_energy_chain_calculus.tex


% BEGIN REVIEWED MODULE research_r11/comparison/unbounded_comparison_and_defect.tex
\section{Comparison beyond bounded potentials by controlling truncation loss}
\label{sec:r11-unbounded-comparison}

Fix the original finite extension $K/\mathbf Q_p$ and its original
absolute value. Let $X$ be good, Hausdorff, boundaryless, quasi-smooth
and pure of dimension $n$, let $V\subset X$ be ordinary open, and
let $S$ be an arbitrary original symmetric CLD middle-positive
$(p,p)$-current with $d'S=d''S=0$. Put $q=n-p\geq1$.
The background is never replaced by a smooth or geometric current.
All bounded products below are the frozen R4 products, on every
original compact smooth CLD test and every actual moment refinement.

We use three precise \texttt{INTERNAL\_INPUT}s: the fixed-background
bounded mixed calculus; its strict Borel maximum locality
(R1/R4, \path{thm:maximum-locality-relative}); and its compact-gap
comparison theorem (R1/R4, \texttt{thm:relative-comparison}).
The quantitative Dini theory already proved in R3 is not counted as
new here. The new assertion is a one-sided extension of comparison
to unbounded actual potentials, with an explicit and separately
checkable truncation-loss hypothesis.

\subsection{Top measures, their loss, and a noncircular sufficient class}

By a negative extended classical psh function we mean a Borel function
$u:V\to[-\infty,0]$ in the all-extensions classical psh class;
its actual point values are retained. Its canonical truncations
$u_j=\max\{u,-j\}$ belong to the frozen bounded class, for every
integer $j\geq1$. In particular $u$ is upper semicontinuous:
each $u_j$ is upper semicontinuous by the definition of the frozen
bounded class, and $u=\inf_j u_j$. Write
\[
 \mu_{u,j}=T_{S,q}(u_j,\ldots,u_j),\qquad
 E_{u,j}=\{u>-j\},\qquad
 \nu_{u,j}=\mathbf1_{\{u\leq-j\}}\mu_{u,j}.
 \tag{UC1}
\]
These are original positive Radon measures before restriction and
positive Borel measures after restriction.

\begin{lemma}[Consistent top-degree truncations]
\label{lem:r11-consistent-top-truncations}
For $k\geq j$, one has
\[
 \mathbf1_{E_{u,j}}\mu_{u,k}
       =\mathbf1_{E_{u,j}}\mu_{u,j}.
 \tag{UC2}
\]
Consequently the increasing sequence
$\mathbf1_{E_{u,j}}\mu_{u,j}$ defines a positive Borel measure
\[
                \mu_u^{\mathrm{np}}
       :=\lim_{j\to\infty}\mathbf1_{E_{u,j}}\mu_{u,j},
 \tag{UC3}
\]
which gives zero mass to $\{u=-\infty\}$. Local finiteness is
not asserted without an additional hypothesis. On every $E_{u,j}$,
$\mu_u^{\mathrm{np}}=\mu_{u,j}$.
\end{lemma}

\noindent\textbf{Status: \texttt{DIRECT\_COROLLARY}} of frozen
strict Borel maximum locality. This top-degree measure construction
does not require a closedness theorem for lower-degree truncations.

\begin{proof}
Since $u_j=\max\{u_k,-j\}$, strict maximum locality replaces
each of the $q$ differentiated slots on $\{u_k>-j\}=E_{u,j}$.
This gives (UC2). The carriers $E_{u,j}$ increase, so consistency
gives monotonicity of their restricted measures. The increasing
limit of positive measures is a measure, and every restricted
measure vanishes on the displayed pole set. No assertion that
$E_{u,j}$ is ordinary open is used.
\end{proof}

\begin{definition}[Local absence of truncation loss]
\label{def:r11-no-truncation-loss}
Write $u\in\mathcal F_{S,\mathrm{loc}}(V)$ if
$\mu_u^{\mathrm{np}}$ is locally finite and, for every compact
$H\Subset V$,
\[
                       \nu_{u,j}(H)\longrightarrow0.
 \tag{UC4}
\]
This is a condition on the original bounded products (UC1), not
an assumed comparison principle. It will be verified by a weighted
estimate below. We abbreviate $\mu_u^{\mathrm{np}}$ to $\mu_u$
when no confusion with a bounded product is possible.
\end{definition}

\begin{proposition}[Local total variation and a gauge criterion]
\label{prop:r11-loss-tv-and-gauge}
If $u\in\mathcal F_{S,\mathrm{loc}}(V)$, then
$\mu_{u,j}\to\mu_u$ in total variation on every compact
$H\Subset V$, with the estimate
\[
 \|\mu_{u,j}-\mu_u\|_{\mathrm{TV}(H)}
 \leq \nu_{u,j}(H)+\mu_u(H\cap\{u\leq-j\}).
 \tag{UC5}
\]
Conversely, local total-variation convergence to (UC3) implies
(UC4).

More concretely, let $\Phi:[0,\infty)\to[1,\infty)$ be
nondecreasing and Borel, with $\Phi(t)\to\infty$. Suppose
that for every compact $H\Subset V$ there is $C_H<\infty$ such
that
\[
       \sup_{j\geq1}\int_H\Phi(-u_j)\,d\mu_{u,j}\leq C_H.
 \tag{UC6}
\]
Then $u\in\mathcal F_{S,\mathrm{loc}}(V)$,
\[
 \int_H\Phi(-u)\,d\mu_u\leq C_H,
 \qquad
 \nu_{u,j}(H)\leq\frac{C_H}{\Phi(j)},
 \qquad
 \|\mu_{u,j}-\mu_u\|_{\mathrm{TV}(H)}
                         \leq\frac{2C_H}{\Phi(j)}.
 \tag{UC7}
\]
The integral at a pole is interpreted on $\{u> -\infty\}$;
$\mu_u$ is zero on the omitted set. Every locally bounded
negative classical psh function satisfies (UC6) for every fixed
finite-valued $\Phi$.
\end{proposition}

\noindent\textbf{Status: \texttt{NEW\_PROVED}} in this original
arbitrary-background truncation setting.

\begin{proof}
On $E_{u,j}$ the measures agree by (UC2). On its complement
they are positive measures of masses $\nu_{u,j}(H)$ and
$\mu_u(H\cap\{u\leq-j\})$. Their difference has total
variation at most the sum, proving (UC5). The second mass tends
to zero by local finiteness and continuity from above, since
the intersection is the $\mu_u$-null pole set. Conversely,
\[
 \nu_{u,j}(H)
 \leq\|\mu_{u,j}-\mu_u\|_{\mathrm{TV}(H)}
       +\mu_u(H\cap\{u\leq-j\}),
\]
which proves the converse.

Under (UC6), the measures in (UC3) have mass at most $C_H$ on
$H$, so their limit is locally finite. On $E_{u,l}$ and for
$j\geq l$ one has $u_j=u$ and
$\mu_{u,j}=\mu_u$. Hence
\[
       \int_{H\cap E_{u,l}}\Phi(-u)\,d\mu_u\leq C_H.
\]
Letting $l\to\infty$ proves the first inequality in (UC7).
On $\{u\leq-j\}$ the weight $\Phi(-u_j)$ is $\Phi(j)$,
giving the second. The first inequality likewise bounds
$\mu_u(H\cap\{u\leq-j\})$ by $C_H/\Phi(j)$; insert
both bounds into (UC5).

If $u$ is locally bounded, it is bounded on an ordinary
neighborhood of $H$. For all sufficiently large $j$, the
truncation equals $u$ on that neighborhood, with a uniformly
bounded weight. The remaining finitely many truncated measures
have finite mass on $H$. This verifies the last assertion.
\end{proof}

\subsection{One-sided comparison and its explicit defect}

For negative extended potentials $u,v$, write
\[
 \mathrm{BG}(u,v):\qquad
     \overline{\{u<v-\varepsilon\}}\Subset V
                   \quad\text{for every }\varepsilon>0.
 \tag{UC8}
\]
The inequality is an inequality of extended real values. In
particular $\{u<v-\varepsilon\}$ contains no point at which
$v=-\infty$, even when $u=-\infty$ there. Its closure is
the ordinary topological closure.

\begin{theorem}[One-sided comparison with truncation loss]
\label{thm:r11-unbounded-defect-comparison}
Let $u,v$ be negative extended classical psh functions satisfying
$\mathrm{BG}(u,v)$. Assume only that $\mu_u$ in (UC3) is
locally finite; $\mu_v$ may initially have infinite mass on
compact sets. For $\varepsilon>0$, put
\[
 \begin{split}
 A_\varepsilon&=\{u<v-\varepsilon\},\\
 B_\varepsilon&=\{v> -\infty,\ u\leq v-\varepsilon\},\\
 H_\varepsilon&=\overline{\{u<v-\varepsilon/2\}}.
 \end{split}
\]
Then
\[
 \mu_v(A_\varepsilon)
 \leq \mu_u(B_\varepsilon)
             +\liminf_{j\to\infty}\nu_{u,j}(H_\varepsilon).
 \tag{UC9}
\]
In particular, if $u\in\mathcal F_{S,\mathrm{loc}}(V)$,
then the full unbounded comparison principle holds:
\[
                   \boxed{\quad
       \mu_v(\{u<v\})\leq\mu_u(\{u<v\}).\quad}
 \tag{UC10}
\]
Either side of (UC10) may be infinite. Only the potential $u$
on the right is required to have no local truncation loss.
Every assertion keeps the original $S$, domain and absolute value.
\end{theorem}

\noindent\textbf{Status: \texttt{NEW\_PROVED}}. The bounded
comparison input alone does not suppress the extra term in (UC9).

\begin{proof}
Write $u_j=\max(u,-j)$ and $v_j=\max(v,-j)$. If
$u_j\leq v_j-\varepsilon$, then $v_j>-j$, so $v_j=v$,
$v$ is finite, and $u\leq v-\varepsilon$. Thus
\[
 B_{j,\varepsilon}:=\{u_j\leq v_j-\varepsilon\}
       \subset B_\varepsilon\subset H_\varepsilon.
 \tag{UC11}
\]
The same inclusion with strict inequalities and any positive
margin proves $\mathrm{BG}(u_j,v_j)$. The level-set form of
frozen bounded comparison, obtained in its proof before letting
the level tend to zero, therefore gives
\[
 \mu_{v,j}(A_{j,\varepsilon})
                  \leq\mu_{u,j}(B_{j,\varepsilon}),
 \qquad A_{j,\varepsilon}=\{u_j<v_j-\varepsilon\}.
 \tag{UC12}
\]
Split the right-hand measure into its consistent part and
discarded part. From (UC11),
\[
 \mu_{u,j}(B_{j,\varepsilon})
          \leq\mu_u(B_\varepsilon)
                           +\nu_{u,j}(H_\varepsilon).
 \tag{UC13}
\]

Fix $k$. For $j\geq k$, on $E_{v,k}=\{v>-k\}$,
consistency gives $\mu_{v,j}=\mu_v$, while $v_j=v$.
There $A_{j,\varepsilon}$ increases to $A_\varepsilon$ as
$j\to\infty$. Restrict the left side of (UC12) to this set
and use continuity from below of the fixed positive measure:
\[
 \mu_v(A_\varepsilon\cap E_{v,k})
 \leq\mu_u(B_\varepsilon)
                        +\liminf_j\nu_{u,j}(H_\varepsilon).
\]
Now let $k\to\infty$. The union contains every point of
$A_\varepsilon$, since $v$ is finite on that set. This proves
(UC9), even if $\mu_v$ was not known locally finite.

Under (UC4) the error is zero. Along any positive sequence
$\varepsilon\downarrow0$, both $A_\varepsilon$ and
$B_\varepsilon$ increase to $\{u<v\}$. Continuity from
below gives (UC10). No limit over a changing background is used,
and no mass at a pole of the other potential has been discarded:
the condition $v> -\infty$ in $B_\varepsilon$ records exactly
which extended-value gaps occur.
\end{proof}

\subsection{Bounded perturbations, domination and uniqueness}

\begin{lemma}[A bounded perturbation on the finite-value locus]
\label{lem:r11-unbounded-positive-perturbation}
Let $v$ be negative extended classical psh, let $h$ be original
bounded classical psh with $0\leq h\leq1$, and let
$a>0$, $b\geq0$. Set $w=v+a(h-1)-b$. Then $w$ has the
same pole set as $v$, and its measure (UC3) satisfies
\[
       \mu_w\geq\mu_v+
                 a^q\mathbf1_{\{v> -\infty\}}\mu_h.
 \tag{UC14}
\]
No no-loss hypothesis on $v$ or $w$ is needed here.
\end{lemma}

\noindent\textbf{Status: \texttt{NEW\_PROVED}} as an unbounded
finite-locus consequence of the original bounded mixed expansion.

\begin{proof}
Fix $k$ and choose $j>k+a+b$. The bounded function
$b_k=-k+a(h-1)-b$ satisfies $b_k>-j$. Consequently
\[
 v_k+a(h-1)-b
       =\max\{w,b_k\}=\max\{w_j,b_k\}.
\]
On $\{v>-k\}$ its first branch is strictly larger. Strict
Borel maximum locality, in every differentiated slot, therefore
identifies its bounded top measure with $\mu_{w,j}$ on this
carrier. The bounded positive mixed expansion gives
\[
 \mathbf1_{\{v>-k\}}\mu_{w,j}
 \geq\mathbf1_{\{v>-k\}}
                    (\mu_{v,k}+a^q\mu_h).
\]
On that same carrier $w>-k-a-b>-j$, so (UC2) identifies
the left side with $\mu_w$; it identifies the first term
on the right with $\mu_v$. Increase $k$ to infinity and
use that both unbounded measures vanish on their common pole
set. This proves (UC14). In particular mere Borel equality of
potentials was never used as a substitute for maximum locality.
\end{proof}

Use the raw fixed-domain relative capacity
\[
 \operatorname{Cap}_S(E;V)
    =\sup_{\substack{h\in\mathcal C_b(V)\\0\leq h\leq1}}
                              \mu_h(E)
 \qquad(E\subset V\text{ Borel}).
 \tag{UC15}
\]

\begin{theorem}[Unbounded domination and order modulo capacity]
\label{thm:r11-unbounded-capacity-domination}
Let $u\in\mathcal F_{S,\mathrm{loc}}(V)$ and let $v$ be
negative extended classical psh, with $\mathrm{BG}(u,v)$.
Each of the following separate hypotheses implies
\[
                \operatorname{Cap}_S(\{u<v\};V)=0:
 \tag{UC16}
\]
\begin{enumerate}
\item $\mu_u(\{u<v\})=0$;
\item $\mathbf1_{\{u<v\}}\mu_u\leq
       \mathbf1_{\{u<v\}}\mu_v$ as positive Borel measures.
\end{enumerate}
In particular, if $u,v\in\mathcal F_{S,\mathrm{loc}}(V)$,
the compact-gap condition holds in both orders, and
$\mu_u=\mu_v$, then
\[
           \operatorname{Cap}_S(\{u\ne v\};V)=0.
 \tag{UC17}
\]
Equality of extended values includes equality at common poles.
The statements do not assert ordinary pointwise order, equality
of raw and outer capacities, or an unrestricted existence theorem.
\end{theorem}

\noindent\textbf{Status: \texttt{NEW\_PROVED}}. This extends the
frozen bounded domination/uniqueness statements without assuming
that bounded perturbations preserve the no-loss class.

\begin{proof}
Fix $\delta>0$, choose $0<a<\delta$, and choose any tester
$h$ in (UC15). Put $w=v+a(h-1)-\delta$ and $E=\{u<w\}$.
The inclusions
\[
         \{u<v-2\delta\}\subset E\subset\{u<v-\delta\}
 \tag{UC18}
\]
hold with the literal extended values. They also show
$\mathrm{BG}(u,w)$ and $\overline E\Subset V$.
One-sided comparison is applicable because only $u$, not $w$,
must satisfy (UC4). On $E$ the potential $v$ is finite.
Thus (UC14) and (UC10) give
\[
 \mu_v(E)+a^q\mu_h(E)
            \leq\mu_w(E)\leq\mu_u(E)<\infty.
 \tag{UC19}
\]
Under hypothesis (1) the last quantity is zero. Under (2),
$\mu_u(E)\leq\mu_v(E)$, so the finite quantity
$\mu_v(E)$ cancels and again $\mu_h(E)=0$.
In either case (UC18) makes
$\mu_h(\{u<v-2\delta\})=0$. Take the supremum over $h$,
then a countable union in $\delta\downarrow0$ and use
countable subadditivity of the raw capacity. This proves (UC16).
For (UC17), apply hypothesis (2) in each order and take the
union of the two strict-gap sets.
\end{proof}

\begin{corollary}[An unbounded sublevel estimate]
\label{cor:r11-unbounded-sublevel-capacity}
Under the hypotheses of Theorem~\ref{thm:r11-unbounded-capacity-domination}
before either additional measure condition is imposed, for every
$\delta>0$ one has
\[
 \delta^q\operatorname{Cap}_S(\{u<v-2\delta\};V)
                 \leq\mu_u(\{u<v-\delta\}).
 \tag{UC20}
\]
In particular, suppose $\mathrm{BG}(u,0)$ and (UC6) hold.
Fix $\delta_0>0$ and let
$H=\overline{\{u<-\delta_0/2\}}\Subset V$.
Then for every $\delta\geq\delta_0$,
\[
 \operatorname{Cap}_S(\{u<-2\delta\};V)
                      \leq\frac{C_H}{\delta^q\Phi(\delta)}.
 \tag{UC21}
\]
Consequently $\operatorname{Cap}_S(\{u=-\infty\};V)=0$.
\end{corollary}

\noindent\textbf{Status: \texttt{DIRECT\_COROLLARY}} of the
new one-sided comparison theorem and gauge estimate. No general
raw-capacity outer regularity is used or inferred.

\begin{proof}
In (UC19), discard its first nonnegative term and bound
$\mu_u(E)$ by $\mu_u(\{u<v-\delta\})$. Take the
supremum over $h$ and then let $a\uparrow\delta$ to obtain
(UC20). With $v=0$, (UC7) and monotonicity of $\Phi$ bound
the right side by $C_H/\Phi(\delta)$; this proves (UC21).
The pole set is contained in every such sublevel, whose bounds
tend to zero.
\end{proof}

\begin{corollary}[Outer-polarity and quasicontinuity in the no-loss boundary class]
\label{cor:r11-unbounded-outer-quasicontinuity}
Let $u\in\mathcal F_{S,\mathrm{loc}}(V)$ satisfy
$\mathrm{BG}(u,0)$. Define the relative outer capacity by
\[
 \operatorname{Cap}_S^*(A;V)
     =\inf_{\substack{G\subset V\text{ ordinary open}\\A\subset G}}
                           \operatorname{Cap}_S(G;V).
 \tag{UC22}
\]
Fix $\delta_0>0$ and put
$H=\overline{\{u<-\delta_0/2\}}\Subset V$.
For $\delta\geq\delta_0$ one has
\[
 \operatorname{Cap}_S^*(\{u<-2\delta\};V)
    =\operatorname{Cap}_S(\{u<-2\delta\};V)
    \leq \delta^{-q}\mu_u(H\cap\{u<-\delta\})
                        \longrightarrow0.
 \tag{UC23}
\]
Consequently the pole set is outer null,
\[
                 \operatorname{Cap}_S^*(\{u=-\infty\};V)=0,
 \tag{UC24}
\]
and the actual canonical truncations converge in relative outer
capacity: for every $\varepsilon>0$,
\[
       \operatorname{Cap}_S^*(\{u_j-u>\varepsilon\};V)
                                      \longrightarrow0.
 \tag{UC25}
\]
The difference in (UC25) is $+\infty$ at a pole.
For every compact $K_0\Subset V$ and every $\eta>0$, there is
an ordinary open $G\subset V$ with
$\operatorname{Cap}_S^*(G;V)<\eta$ such that
$u|_{K_0\setminus G}$ is finite and continuous in the ordinary
subspace topology. Under (UC6), the right side of (UC23) is further
bounded by $C_H/(\delta^q\Phi(\delta))$.
\end{corollary}

\noindent\textbf{Status: \texttt{NEW\_PROVED}} for the unbounded
no-loss boundary class. The bounded relative quasicontinuity input
is the analytic Laurent-domain extension in
Corollary~\ref{cor:r11-analytic-bounded-relative-QC}; its bounded
regularization and energy ingredients are frozen
\texttt{INTERNAL\_INPUT}s.
This does not assert outer regularity for arbitrary Borel sets.

\begin{proof}
Upper semicontinuity makes $\{u<-2\delta\}$ ordinary open.
Raw and outer capacities agree on an ordinary open set by (UC22)
and monotonicity. Apply (UC20) with $v=0$. Its right side is
contained in $H$ for the indicated range of $\delta$.
The finite measure $\mu_u|_H$ gives zero mass to the pole set,
and its masses on $H\cap\{u<-\delta\}$ tend to zero by
continuity from above. This proves (UC23). The pole set is
contained in every one of these open sublevels, proving (UC24).
For finite $u$ the exact identity
\[
                   \{u_j-u>\varepsilon\}
                         =\{u<-j-\varepsilon\}
\]
holds, and the same identity holds with the stated convention at
a pole. Equation (UC23) therefore proves (UC25).

For quasicontinuity, choose $\delta$ so large that
$G_0=\{u<-2\delta\}$ has capacity less than $\eta/2$, and
choose an integer $j>2\delta$. Outside $G_0$ one has $u=u_j$
with finite values. Cover $K_0$ by finitely many smaller ordinary
analytic Laurent-box cores $V_{0,i}\Subset\Omega_i\Subset V$ to
which Corollary~\ref{cor:r11-analytic-bounded-relative-QC} applies.
For each $i$,
choose an ordinary open $G_i\subset\Omega_i$ of
$\operatorname{Cap}_S(\,\cdot\,;\Omega_i)$ less than
$\eta/(2N)$ such that $u_j$ is continuous on
$V_{0,i}\setminus G_i$, where $N$ is the number of cores.
Restriction of a normalized tester from $V$ to $\Omega_i$ gives
\[
 \operatorname{Cap}_S(G_i;V)
                  \leq\operatorname{Cap}_S(G_i;\Omega_i).
\]
Thus $G=G_0\cup\bigcup_iG_i$ is ordinary open and has raw,
hence outer, capacity less than $\eta$. On $K_0\setminus G$
the function $u$ equals $u_j$, whose continuity holds locally
on the stated finite cover. This proves the required continuity
without changing any actual values. Finally (UC7) supplies the
gauge bound. No measure-theoretic replacement of the pole set,
reference volume, or background regularization enters the proof.
\end{proof}

\paragraph{Precise remaining gates.}
The hypotheses (UC4) and (UC6) are not claimed for every
unbounded classical psh function. Their failure leaves an explicit
term in (UC9). A capacity-zero strict gap need not be empty for
arbitrary $S$. The upgrade to outer nullity above uses the specific
ordinary open sublevels and the printed zero-boundary gap condition;
no such upgrade is claimed for a general raw-capacity null strict gap.
This section gives no new
ambient positivity converse, no general nonlinear Dirichlet
existence, and no actual PS/PL relation quotient. It applies to
every original background satisfying the initial hypotheses,
including degenerate or pluripolar backgrounds, with its printed
boundary and truncation assumptions intact.

% END REVIEWED MODULE research_r11/comparison/unbounded_comparison_and_defect.tex


% BEGIN REVIEWED MODULE research_r11/comparison/capacity_tail_boundary_class.tex
\section{An intrinsic no-loss boundary class and its mixed top calculus}
\label{sec:r11-capacity-tail-boundary-class}

Retain the original field, analytic space, arbitrary background $S$,
ordinary domain $V$, and $q=n-p\geq1$ from
Section~\ref{sec:r11-unbounded-comparison}. Every potential below is
negative extended classical psh, with its actual values retained.
The capacity always has the fixed tester domain $V$. The results
of this section concern top-degree measures. They do not construct
all lower-degree products without the separate energy hypotheses.

\begin{lemma}[Compact conservation of truncation mass]
\label{lem:r11-compact-truncation-conservation}
If $\mathrm{BG}(u,0)$ holds, then $\mu_u$ in (UC3) is locally
finite, without assuming absence of truncation loss. More precisely,
fix $M>0$ and $H\Subset V$, set
$K_M=\overline{\{u<-M\}}\Subset V$, and choose an original smooth
compact function $0\leq\chi\leq1$ which is one on an ordinary
neighborhood of $H\cup K_M$. Then for integers $j,l>M$,
\[
                       \int\chi\,d\mu_{u,j}
                           =\int\chi\,d\mu_{u,l}.
 \tag{BC1}
\]
In particular $\sup_j\mu_{u,j}(H)<\infty$.
\end{lemma}

\noindent\textbf{Status: \texttt{NEW\_PROVED}} as a compact
conservation statement for unbounded actual potentials, using only
the frozen bounded telescoping identity and original compact Stokes.

\begin{proof}
Outside $K_M$ one has $u\geq-M$, hence $u_j=u_l=u$ there.
The bounded mixed calculus gives
\[
 \mu_{u,j}-\mu_{u,l}
   =D\left((u_j-u_l)
       \sum_{r=0}^{q-1}(Du_j)^r(Du_l)^{q-1-r}\wedge S\right).
 \tag{BC2}
\]
Every product in the parentheses is a frozen original positive
closed product. Its bounded Borel multiplier is zero outside
$K_M$, so the entire current in parentheses has support in $K_M$.
Pair (BC2) with $\chi$. Since $D\chi$ vanishes near $K_M$,
the pairing is zero, proving (BC1). Fixing one $l>M$ bounds
the nonnegative masses on $H$ for every $j>M$; the remaining
finitely many $j$ have finite mass. The increasing restricted
measures in (UC3) have the same bound on $H$, proving local
finiteness of their limit. No bound for a lower-degree truncation
product is inferred from this top-degree Stokes argument.
\end{proof}

\begin{definition}[The no-loss zero-boundary class]
\label{def:r11-no-loss-boundary-class}
Set
\[
 \mathcal F_{S,\partial}(V)
   =\{u\in\mathcal F_{S,\mathrm{loc}}(V):
                                      \mathrm{BG}(u,0)\}.
 \tag{BC3}
\]
The notation records compact strict negative sublevels, not an
unstated ordinary boundary regularity hypothesis.
\end{definition}

\begin{theorem}[An exact capacity-tail criterion]
\label{thm:r11-capacity-tail-criterion}
Let $u$ satisfy $\mathrm{BG}(u,0)$. The following are equivalent:
\begin{enumerate}
\item $u\in\mathcal F_{S,\partial}(V)$;
\item $\displaystyle\lim_{t\to\infty}
              t^q\operatorname{Cap}_S(\{u<-t\};V)=0$;
\item $\displaystyle\lim_{j\to\infty}\nu_{u,j}(V)=0$.
\end{enumerate}
The capacity in (2) can equivalently be replaced by its ordinary-open
outer regularization. If these conditions hold, then
$\mu_{u,j}-\mu_u$ is a finite signed measure on $V$ for every
sufficiently large $j$, and
\[
 \begin{split}
 t^q\operatorname{Cap}_S(\{u<-t\};V)
       &\leq2^q\mu_u(\{u<-t/2\}),\\
 \nu_{u,j}(V)&\leq4^q\mu_u(\{u<-j/4\}),\\
 \|\mu_{u,j}-\mu_u\|_{\mathrm{TV}(V)}
       &\leq(4^q+1)\mu_u(\{u<-j/4\})\longrightarrow0.
 \end{split}
 \tag{BC4}
\]
The measures themselves need not have finite total mass near the
boundary; their difference is supported in one fixed compact set
for all sufficiently large $j$.
\end{theorem}

\noindent\textbf{Status: \texttt{NEW\_PROVED}}. The condition in
(2) is expressed solely through the frozen bounded capacity.

\begin{proof}
Lemma~\ref{lem:r11-compact-truncation-conservation} supplies local
finiteness automatically. For $j>M$, the discarded measure is
carried by $\{u\leq-j\}\subset K_M$. Thus (1) and (3) are
equivalent. Under (1), (UC20) with $v=0$ and $\delta=t/2$
gives the first inequality of (BC4). Its right side tends to zero:
for large $t$ the carrier lies in the fixed compact $K_M$, on
which $\mu_u$ is finite, and the decreasing carriers intersect
in the $\mu_u$-null pole set. Upper semicontinuity makes
$\{u<-t\}$ ordinary open, so its raw and outer capacities agree.
This proves (2).

Conversely, the bounded function $h_j=1+u_j/j$ is a normalized
capacity tester and $\mu_{u,j}=j^q\mu_{h_j}$. Therefore
\[
 \nu_{u,j}(V)
    \leq j^q\operatorname{Cap}_S(\{u<-j/2\};V).
 \tag{BC5}
\]
Condition (2) makes this tend to zero and proves (3).
Combining (BC5) with the first inequality of (BC4), at $t=j/2$,
gives the second. The measures agree on $\{u>-j\}$; on its
complement their variation difference is at most the sum of their
masses. The second inequality and
$\{u\leq-j\}\subset\{u<-j/4\}$ give the last bound in
(BC4). All these tail sets lie in $K_M$ for large $j$, as claimed.
\end{proof}

\begin{theorem}[A solid convex cone of actual potentials]
\label{thm:r11-solid-boundary-cone}
The class $\mathcal F_{S,\partial}(V)$ contains $0$ and has the
following closure properties.
\begin{enumerate}
\item If $u$ belongs to the class and $v$ is any negative extended
classical psh function with $u\leq v\leq0$, then $v$ belongs
to the class.
\item If $u_1,\ldots,u_m$ belong to the class and $a_i\geq0$,
then $\sum_i a_i u_i$ belongs to the class, with zero-coefficient
terms omitted.
\item Finite maxima of its members belong to the class.
\end{enumerate}
No continuity, smoothness, geometric origin or full-support
assumption is imposed on $S$.
\end{theorem}

\noindent\textbf{Status: \texttt{NEW\_PROVED}}. These are
closure assertions for the explicit no-loss class, not for all
unbounded classical psh functions.

\begin{proof}
For (1), $\{v<-t\}\subset\{u<-t\}$ gives both the compact
boundary-gap condition and the capacity-tail criterion.
For (2), assume $A=\sum_i a_i>0$. The elementary inclusion
\[
 \left\{\sum_i a_i u_i<-t\right\}
                  \subset\bigcup_{i:a_i>0}\{u_i<-t/A\}
 \tag{BC6}
\]
gives compact sublevels and
\[
 t^q\operatorname{Cap}_S
          \left(\left\{\sum_i a_i u_i<-t\right\};V\right)
 \leq A^q\sum_{i:a_i>0}
       (t/A)^q\operatorname{Cap}_S(\{u_i<-t/A\};V)
                       \longrightarrow0.
\]
The sum is an extended classical psh function by the frozen
curve definition and positivity. The zero sum is immediate.
Finally a finite maximum is classical psh and dominates any one
of its members, so (1) proves (3).
\end{proof}

\subsection{Top mixed products without a superlinear energy assumption}

\begin{theorem}[A positive mixed top calculus in the boundary class]
\label{thm:r11-boundary-mixed-top-calculus}
For $u_1,\ldots,u_q\in\mathcal F_{S,\partial}(V)$ put
\[
 M_j=T_{S,q}(u_{1,j},\ldots,u_{q,j}),\qquad
 E_j=\bigcap_{i=1}^q\{u_i>-j\}.
\]
The increasing consistent restrictions define a locally finite
positive measure
\[
            M(u_1,\ldots,u_q)=\lim_j\mathbf1_{E_j}M_j.
 \tag{BC7}
\]
It gives zero mass to the union of the $q$ pole sets. The uncut
measures converge to it in total variation on the whole domain,
in the sense of a finite signed difference:
\[
               \|M_j-M(u_1,\ldots,u_q)\|_{\mathrm{TV}(V)}
                                           \longrightarrow0.
 \tag{BC8}
\]
Independent canonical truncation levels have the same limit as
their minimum tends to infinity. The resulting product is symmetric,
additive and positively homogeneous in every slot on this cone, and
\[
 M(u_1,\ldots,u_q)
   ={1\over q!}\sum_{I\subset\{1,\ldots,q\}}
        (-1)^{q-|I|}\mu_{\sum_{i\in I}u_i}.
 \tag{BC9}
\]
Here the empty sum has zero top measure. Every equality is an
equality of actual Borel measures, and hence holds against every
original compact smooth scalar test and all its actual refinements.
\end{theorem}

\noindent\textbf{Status: \texttt{NEW\_PROVED}}. This top-degree
result requires no superlinear weighted energy estimate and does
not assert the corresponding lower-degree recursive calculus.

\begin{proof}
Let $w=\sum_i u_i$ and $s_j=\sum_i u_{i,j}$. The solid-cone
theorem puts $w$ in $\mathcal F_{S,\partial}(V)$. The exact identity
\[
                      \max(s_j,-j)=\max(w,-j)=w_j
 \tag{BC10}
\]
holds: if no slot is truncated then $s_j=w$, while if a slot is
truncated then $s_j\leq-j$. Strict Borel maximum locality thus
identifies the top measure of $s_j$ with $\mu_w$ on $\{w>-j\}$.
As $-qj\leq s_j\leq0$, capacity normalization gives
\[
 \mu_{s_j}(\{w\leq-j\})
       \leq(qj)^q\operatorname{Cap}_S(\{w<-j/2\};V)
                                     \longrightarrow0.
 \tag{BC11}
\]
The mass of $\mu_w$ on the same compact tails tends to zero.
Hence $\mu_{s_j}\to\mu_w$ in total variation on $V$.

The positive bounded mixed expansion gives $q!M_j\leq\mu_{s_j}$.
The complement of $E_j$ is contained in $\{w\leq-j\}$.
Equation (BC11) therefore makes the discarded part
$\mathbf1_{V\setminus E_j}M_j$ tend to zero in mass. Frozen
strict maximum locality gives consistency on $E_j$, exactly as
in (UC2). Their increasing limit is bounded locally by the
already bounded family $\mu_{s_j}/q!$, so is locally finite.
Its mass on $V\setminus E_j$ tends to zero, since these tails
are eventually compact and decrease to the union of the pole sets.
This proves (BC7)--(BC8).

For completeness consider independent levels $l_i\geq m$. Fix
$M>0$ and let $K=\overline{\{w<-M\}}$. Whenever all levels
exceed $M$, all the truncated slots agree with their actual values
outside $K$. A slot-by-slot bounded telescoping identity and
compact Stokes, exactly as in (BC2), show that the integral of
the mixed measure against any fixed cutoff $\chi=1$ near $K$
is independent of all the levels. Passing to the common-level
limit using (BC8) shows that this integral equals $\int\chi\,dM$.
For $m>M$, strict locality gives equality with $M$ on $E_m$,
whereas $V\setminus E_m\subset K$. Subtracting these equal
cutoff integrals yields
\[
 T_{S,q}(u_{1,l_1},\ldots,u_{q,l_q})(V\setminus E_m)
                            =M(V\setminus E_m).
\]
Consequently the total variation of the difference is at most
$2M(V\setminus E_m)$, which tends to zero. This proves independence
of the truncation rates.

Bounded polarization gives (BC9) at level $j$. For every subset
$I$, the preceding argument with $s_{I,j}=\sum_{i\in I}u_{i,j}$
shows $\mu_{s_{I,j}}\to\mu_{\sum_{i\in I}u_i}$ in total
variation; the empty subset contributes zero. Passing in the finite
polarization identity proves (BC9).

Here is also a direct justification of all multilinearity claims.
For any finite family in the cone and fixed positive coefficients
$c_i$, put $A=\sum_i c_i$ and $c=\min_i c_i>0$. With
$s_j=\sum_i c_i u_{i,j}$ and $w=\sum_i c_i u_i$ one has
\[
                    \max(s_j,-cj)=\max(w,-cj).
\]
The preceding locality proof applies, and its discarded mass is
at most $(Aj)^q\operatorname{Cap}_S(\{w<-cj/2\};V)$, which
tends to zero by the tail criterion. Thus the positive bounded
polynomial expansion of $\mu_{s_j}$ passes to the limit for
every set of coefficients. Comparing coefficients in these finite
polynomial identities proves additivity and positive homogeneity
in each slot, with zero coefficients omitted. Symmetry is inherited
from the bounded products. No signed current order or PS quotient
is used.
\end{proof}

\begin{lemma}[Uniform tails above one no-loss lower barrier]
\label{lem:r11-uniform-upper-tail}
If $u,v\in\mathcal F_{S,\partial}(V)$ and $u\leq v\leq0$,
then for every $t>0$,
\[
                 \mu_v(\{v<-t\})
                      \leq2^q\mu_u(\{u<-t/2\}).
 \tag{BC12}
\]
The right side is finite, and tends to zero as $t\to\infty$.
\end{lemma}

\noindent\textbf{Status: \texttt{NEW\_PROVED}} from the
one-sided comparison theorem and the solid boundary cone.

\begin{proof}
Scaling and addition of constants have their expected action on
(UC3): $\mu_{au-b}=a^q\mu_u$ for $a>0$, $b\geq0$.
Indeed, on each $\{u>-k\}$ and for $j>ak+b$, write
$au_k-b=\max((au-b)_j,-ak-b)$ and apply strict maximum
locality and bounded scaling. Exhaust the finite-value locus;
both sides vanish on its complement.
Now $2u$ belongs to the boundary class and
$\mathrm{BG}(2u,v-t)$ holds because its positive-margin gaps lie
in compact strict negative sublevels of $u$. Moreover,
\[
       \{-\infty<v<-t\}\subset\{2u<v-t\}
                                      \subset\{u<-t/2\}.
\]
Apply (UC10) to $2u,v-t$ and use the scaling and translation
identities. The omitted set $\{v=-\infty\}$ is $\mu_v$-null
by its definition; no literal strict-gap inclusion at a common
pole is asserted. This proves (BC12). The limiting assertion follows
from local finiteness on a fixed compact tail and zero mass on
the pole set.
\end{proof}

\begin{theorem}[Joint decreasing continuity in the boundary class]
\label{thm:r11-boundary-joint-decreasing}
For $1\leq i\leq q$, suppose $v_{i,l}$ are negative extended
classical psh functions and
\[
              v_{i,l}\downarrow u_i\in\mathcal F_{S,\partial}(V)
                         \quad\hbox{pointwise on }V.
 \tag{BC13}
\]
Then every $v_{i,l}$ belongs to the same boundary class, and
\[
       M(v_{1,l},\ldots,v_{q,l})
                       \longrightarrow M(u_1,\ldots,u_q)
 \tag{BC14}
\]
against every original compact smooth scalar test. No uniform
superlinear energy bound is needed. The conclusion is ordinary
current convergence, not convergence on arbitrary delta tests.
\end{theorem}

\noindent\textbf{Status: \texttt{NEW\_PROVED}} in this top-degree
no-loss boundary class.

\begin{proof}
Membership follows from upward stability. Put $w=\sum_i u_i$
and $w_l=\sum_i v_{i,l}$, so $w\leq w_l\leq0$. Let
$M_l=M(v_{1,l},\ldots,v_{q,l})$ and let $M_{l,k}$ be the
bounded mixed product with every slot truncated at $-k$.
They agree by strict locality on
$E_{l,k}=\bigcap_i\{v_{i,l}>-k\}$. The positive mixed expansion
and (BC12) give
\[
 M_l(V\setminus E_{l,k})
 \leq {1\over q!}\mu_{w_l}(\{w_l\leq-k\})
 \leq {2^q\over q!}\mu_w(\{w<-k/4\}).
 \tag{BC15}
\]
In the last step use $\{w_l\leq-k\}\subset\{w_l<-k/2\}$
before applying (BC12).
The sum of the $q$ bounded truncated slots lies in $[-qk,0]$.
Also $V\setminus E_{l,k}\subset\{w\leq-k\}$. Capacity
normalization, the positive expansion and the first inequality
of (BC4) therefore give
\[
 M_{l,k}(V\setminus E_{l,k})
 \leq{(qk)^q\over q!}\operatorname{Cap}_S(\{w<-k/2\};V)
 \leq{q^q4^q\over q!}\mu_w(\{w<-k/4\}).
 \tag{BC16}
\]
The right sides are independent of $l$ and tend to zero.
Because the measures agree on $E_{l,k}$, their sum bounds the
total variation of $M_l-M_{l,k}$ on all of $V$.

For fixed $k$, the bounded functions
$\max(v_{i,l},-k)$ decrease jointly to $\max(u_i,-k)$.
Frozen bounded joint decreasing continuity gives convergence of
$M_{l,k}$ against every original compact smooth test as
$l\to\infty$. Apply the uniform tail bounds to the two ends
of this convergence, then let $k\to\infty$ and use (BC8) for
the limiting tuple. This proves (BC14) for the full original
test family. The two limits are separated by explicit uniform
total-variation bounds; no varying-background continuity or
unproved capacity regularity is inserted.
\end{proof}

\begin{corollary}[Outer-capacity domination in the boundary class]
\label{cor:r11-boundary-outer-domination}
Let $u\in\mathcal F_{S,\partial}(V)$ and let $v$ be any
negative extended classical psh function. The strict actual gap
$A=\{u<v\}$ is relatively quasi-open for the fixed capacity:
for each $\eta>0$ there is an ordinary open $G\subset V$ with
$\operatorname{Cap}_S(G;V)<\eta$ such that $A\cup G$ is
ordinary open. In particular,
\[
                 \operatorname{Cap}_S^*(A;V)
                              =\operatorname{Cap}_S(A;V).
 \tag{BC17}
\]
If either $\mu_u(A)=0$ or
$\mathbf1_A\mu_u\leq\mathbf1_A\mu_v$, then
\[
                        \operatorname{Cap}_S^*(A;V)=0.
 \tag{BC18}
\]
If $u,v\in\mathcal F_{S,\partial}(V)$ and $\mu_u=\mu_v$,
then $\operatorname{Cap}_S^*(\{u\ne v\};V)=0$.
These conclusions do not assert actual pointwise equality for a
degenerate background.
\end{corollary}

\noindent\textbf{Status: \texttt{NEW\_PROVED}} from the solid
cone, the new unbounded outer quasicontinuity, and one-sided
domination. Equality of raw and outer capacities is proved only
for these specified strict-gap sets.

\begin{proof}
Put $z=\max(u,v)$. Upward stability makes $z$ a member of the
boundary class even when $v$ itself is not. The exact identity
$A=\{u<z\}$ holds, including the extended values. Since $z\leq0$,
each set $\{u<z-1/m\}$ has compact closure in $V$.
Choose compact neighborhoods $H_m\Subset V$ of these closures,
with ordinary interiors, and set $W=\bigcup_m\operatorname{Int}_X H_m$.
Then $A\subset W$.

Apply Corollary~\ref{cor:r11-unbounded-outer-quasicontinuity} to
$u$ and $z$ on each $H_m$, with total error budget less than
$\eta2^{-m}$. Let $G$ be the union of the resulting ordinary
open exceptional sets. Countable subadditivity gives
$\operatorname{Cap}_S(G;V)<\eta$, and $u,z$ are finite and
continuous on $W\setminus G$: this follows locally from the
cover by the interiors of the $H_m$. Thus $A\setminus G$ is
relatively open in $W\setminus G$. There is an ordinary open
$O\subset W$ with $O\setminus G=A\setminus G$.
Consequently $A\cup G=O\cup G$ is ordinary open. By (UC22),
\[
 \operatorname{Cap}_S^*(A;V)
 \leq\operatorname{Cap}_S(A\cup G;V)
 \leq\operatorname{Cap}_S(A;V)+\operatorname{Cap}_S(G;V).
\]
Let $\eta\downarrow0$ and use the opposite inequality from
monotonicity to prove (BC17).

The condition $\mathrm{BG}(u,v)$ is automatic, since
$\{u<v-\varepsilon\}\subset\{u<-\varepsilon\}$.
Theorem~\ref{thm:r11-unbounded-capacity-domination} therefore
gives raw nullity under either printed measure hypothesis;
(BC17) upgrades it to (BC18). For equal measures and two class
members, apply this argument in both orders and take the union
of the two strict gaps.
\end{proof}

\paragraph{Scope and unsolved gates.}
The boundary class is strictly smaller than all negative extended
classical psh functions: the logarithmic example in
Proposition~\ref{prop:r11-logarithmic-comparison-failure} has
persistent discarded mass and fails the equivalent capacity-tail
criterion. The present top calculus has genuine unbounded members,
for example all radial attenuations $u_a$ with $0<a<1$ from
Section~\ref{sec:r11-radial-unbounded-energy}. It supplies no lower-degree
closedness without additional input, no arbitrary-boundary Perron
solution, and no GQ or CTO relation descent.

% END REVIEWED MODULE research_r11/comparison/capacity_tail_boundary_class.tex


% BEGIN REVIEWED MODULE research_r11/comparison/boundary_gap_quasicontinuity.tex
\section{Quasicontinuity before no-loss, and explicit power attenuation}
\label{sec:r11-boundary-gap-quasicontinuity}

The original field, space, ordinary domain $V$, arbitrary original
middle-positive closed background $S$, and $q=n-p\geq1$ retain
their standing hypotheses. This section distinguishes two already
frozen R2 facts from the new assertions. The R2 corollary
\texttt{cor:r2-extended-psh-qc} proves extended-potential
quasicontinuity for the ambient background $[1]$. Its arbitrary-$S$
criterion \texttt{cor:r2-no-pole-trace} assumes that a specified
background trace avoids the pole and uses a measure-dependent
attenuation. Neither result is counted again here. The following
theorem instead uses the explicit compact-sublevel condition
$\mathrm{BG}(u,0)$, keeps arbitrary $S$, and assumes no trace
non-concentration. Its quantitative tail follows from the new
compact mass conservation (BC1).

\begin{theorem}[Outer quasicontinuity for every compact-sublevel potential]
\label{thm:r11-BG-outer-quasicontinuity}
Let $u:V\to[-\infty,0]$ be extended classical psh and suppose
$\mathrm{BG}(u,0)$. No absence of truncation loss is assumed.
For every fixed $t_0>0$ there is a finite constant $C$ such that
\[
             \operatorname{Cap}_S^*(\{u<-2t\};V)
               =\operatorname{Cap}_S(\{u<-2t\};V)
                       \leq C t^{-q}\qquad(t\geq t_0).
 \tag{BQ1}
\]
The constant is an explicit original compact cutoff mass:
one may take $C=\int\chi\,d\mu_{u,l_0}$, where
$l_0>t_0$ is an integer and $\chi$ is one near
$\overline{\{u<-t_0\}}$ and compactly supported in $V$.
Consequently
\[
 \operatorname{Cap}_S^*(\{u=-\infty\};V)=0,
 \qquad
 \operatorname{Cap}_S^*(\{u_j-u>\varepsilon\};V)\to0
                           \quad(\varepsilon>0).
 \tag{BQ2}
\]
For every compact $H\Subset V$ and every $\eta>0$, an ordinary
open set $G\subset V$ can be chosen with
$\operatorname{Cap}_S^*(G;V)<\eta$ such that
$u|_{H\setminus G}$ is finite and continuous.
The conclusions retain the actual values of $u$.
\end{theorem}

\noindent\textbf{Status: \texttt{NEW\_PROVED}} for arbitrary
original $S$ under the printed compact-sublevel hypothesis.
This is a statement about capacity and potentials, not about
continuity of their uncut Monge--Ampere measures.

\begin{proof}
Let $K=\overline{\{u<-t_0\}}\Subset V$ and choose the stated
cutoff $0\leq\chi\leq1$. Compact conservation (BC1), with
$M=t_0$, gives
\[
                        \int\chi\,d\mu_{u,j}=C
                                  \qquad(j>t_0).
\]
Fix $t\geq t_0$ and an integer $j>2t$. The bounded potential
$u_j$ satisfies $\mathrm{BG}(u_j,0)$ and belongs to the no-loss
class. The bounded case of (UC20), with $v=0$, gives
\[
 t^q\operatorname{Cap}_S(\{u_j<-2t\};V)
                   \leq\mu_{u,j}(\{u_j<-t\}).
\]
Both strict sublevels agree with those of $u$ for this choice
of $j$. The right-hand carrier lies in $K$, so its mass is at
most $C$. This proves the raw-capacity bound in (BQ1).
Upper semicontinuity makes each sublevel ordinary open; its
raw and outer capacities therefore agree.

The pole is contained in every such open sublevel, proving its
outer nullity. The exact identity
$\{u_j-u>\varepsilon\}=\{u<-j-\varepsilon\}$, with infinite
difference at a pole, proves the truncation convergence.
To obtain quasicontinuity on $H$, first remove an open deep
sublevel of capacity less than $\eta/2$ and choose $j$ such
that $u=u_j$ outside it. Apply the analytic bounded relative
quasicontinuity of Corollary
\ref{cor:r11-analytic-bounded-relative-QC} to $u_j$ on $H$,
with error less than $\eta/2$. The union of the two ordinary
open exceptions has the required capacity, and the unchanged
function is finite and continuous on its complement in $H$.
The logarithmic example (RE10) shows why these conclusions
alone do not remove the truncation-loss term in comparison.
\end{proof}

\begin{theorem}[A universal power attenuation into the nonlinear no-loss class]
\label{thm:r11-power-attenuation-no-loss}
Let $u$ satisfy the preceding theorem and let $0<\theta<1$.
Define, with the actual value $-\infty$ at every pole,
\[
                  v_\theta=-\bigl((1-u)^\theta-1\bigr).
 \tag{BQ3}
\]
Then $v_\theta$ is negative extended classical psh,
$\mathrm{BG}(v_\theta,0)$ holds, and
\[
 \{v_\theta=-\infty\}=\{u=-\infty\},\qquad
            v_\theta\in\mathcal F_{S,\partial}(V).
 \tag{BQ4}
\]
Writing $\beta=q(1-\theta)/\theta>0$, one has, for all
sufficiently large $t$ and finite constants independent of $t$,
\[
 \operatorname{Cap}_S(\{v_\theta<-t\};V)
               \leq C_1t^{-q/\theta},\qquad
 \mu_{v_\theta}(\{v_\theta\leq-t\})\leq C_2t^{-\beta}.
 \tag{BQ5}
\]
For every $0<s<\beta$ and every compact $H\Subset V$, its
canonical truncations satisfy the explicit local gauge condition
\[
 \sup_{j\geq1}\int_H(1-v_{\theta,j})^s\,d\mu_{v_\theta,j}
                                                   <\infty.
 \tag{BQ6}
\]
Thus any prescribed local power gauge $s>0$ is obtained by
choosing $\theta<q/(q+s)$. This universal threshold is sharp:
the original radial logarithmic witness in dimension $q$ fails
(BQ6) at and above that threshold after the same attenuation.
No claim of finite \emph{global} power energy is made from these
local estimates alone; its behavior near the ordinary boundary
requires the full integral $J_s$.
\end{theorem}

\noindent\textbf{Status: \texttt{NEW\_PROVED}}. The explicit
power composition and the nonlinear no-loss conclusion are
different from the frozen measure-dependent $L^1$ attenuation.

\begin{proof}
The scalar function
$f(r)=-((1-r)^\theta-1)$ on $r\leq0$ is increasing and
convex: $f'(r)=\theta(1-r)^{\theta-1}>0$ and
$f''(r)=\theta(1-\theta)(1-r)^{\theta-2}>0$.
Convex increasing composition of the bounded classical $u_j$
is classical, and their decreasing limit is (BQ3), with its
printed actual values. This uses the frozen complete-extension
curve composition and decreasing-closure properties.

The exact sublevel identity is
\[
 \{v_\theta<-t\}
       =\{u<-R_\theta(t)\},\qquad
             R_\theta(t)=(1+t)^{1/\theta}-1.
 \tag{BQ7}
\]
It proves the compact-sublevel condition and equality of pole
sets. Equation (BQ1) gives a bound by
$2^qC R_\theta(t)^{-q}=O(t^{-q/\theta})$ for large $t$.
Since $q/\theta>q$, the capacity-tail criterion proves
$v_\theta\in\mathcal F_{S,\partial}$ and the first part
of (BQ5).

Put $v=v_\theta$. On the annulus
$A_R=\{-2R<v\leq-R\}$, the measure $\mu_v$ equals the
bounded measure of $\max(v,-2R)$. Capacity normalization and
the strict carrier enlargement give, for large $R$,
\[
 \mu_v(A_R)
     \leq(2R)^q\operatorname{Cap}_S(\{v<-R/2\};V)
     \leq C_3 R^{-\beta}.
 \tag{BQ8}
\]
The dyadic annuli $A_{2^lt}$, $l\geq0$, partition the finite
part of $\{v\leq-t\}$. Its pole has zero $\mu_v$-mass.
Summing the geometric series in (BQ8) proves the second part
of (BQ5).

Every bounded truncation $v_j$ lies above $v$ and belongs to
the same boundary class. The uniform upper-tail estimate (BC12)
therefore gives
\[
          \mu_{v,j}(\{v_j<-t\})
                   \leq2^q\mu_v(\{v<-t/2\})
                   \leq C_4t^{-\beta}
 \tag{BQ9}
\]
for large $t$, uniformly in $j$. On any fixed compact $H$,
compact conservation bounds $\sup_j\mu_{v,j}(H)$.
Layer cake for the weight $(1-v_j)^s$ now bounds (BQ6) by
a finite low-value contribution plus a constant times
\[
                        \int_1^\infty t^{s-\beta-1}\,dt,
\]
which is finite for $s<\beta$. Finally take $u=-t$ on the
original unit polydisc with $S=[V]$. Its attenuation is exactly
$u_\theta$ in (RE3), and the sharp local gauge calculation
(RE6) fails when $\theta\geq q/(q+s)$. This proves the final
sharpness assertion with the same original field and all
original tests.
\end{proof}

% END REVIEWED MODULE research_r11/comparison/boundary_gap_quasicontinuity.tex


% BEGIN REVIEWED MODULE research_r11/capacity/capacity_dimension_gain.tex
\section{A dimension gain from boundary capacity to all original degrees}
\label{sec:r11-capacity-dimension-gain}

This section imposes an explicit geometric domain hypothesis.
The field and analytic space retain the standing original assumptions,
and there is one compact $K$-affinoid domain $B\subset X$ such that
\[
             \overline V\Subset\operatorname{Int}_X(B).
 \tag{DG1}
\]
Here $V$ is ordinary open, and the displayed closure is compact.
The background $S$ is any original symmetric middle-positive,
separately $d'$- and $d''$-closed current of degree $(p,p)$ on $V$;
put $q=n-p\geq1$. Its origin, support and regularity are unrestricted.
Every capacity below uses all normalized bounded classical testers
on this same $V$.

The new input from Section~\ref{sec:r11-affinoid-energy} is the
exact replacement of local analytic units by quotients of functions
on $B$, followed by its explicit global smooth psh compensator.
The new boundary class
$\mathcal F_{S,\partial}(V)$ and its equivalent capacity-tail
criterion are those of Section~\ref{sec:r11-capacity-tail-boundary-class}.
The frozen bounded mixed calculus and coefficient domination remain
\texttt{INTERNAL\_INPUT}. No energy assumption on the background
is introduced.

\subsection{Every original coefficient has a global tester control}

\begin{lemma}[Compensating arbitrary original complementary tests]
\label{lem:r11-global-capacity-coefficient-control}
Let $0\leq k\leq q$ and let $\eta$ be any original compact smooth
test of complementary degree $(q-k,q-k)$ on $V$. There are a
normalized original smooth classical psh function $0\leq a\leq1$
on $V$ and a finite constant $A_\eta$ such that every original
middle-positive current $T$ of degree $(p+k,p+k)$ satisfies
\[
       |\mu_{T,\eta}|
       \leq A_\eta\,\mathbf1_{\operatorname{supp}\eta}
                         (Da)^{q-k}\wedge T .
 \tag{DG2}
\]
A finite sum of such controls may equivalently be used.
For $k=q$ the assertion is the scalar estimate with $(Da)^0=1$.

In particular, let $b_1,\ldots,b_k$ be bounded classical psh
functions on $V$ with $-L\leq b_i\leq0$, where $L>0$.
For every Borel set $E\subset V$ one has
\[
 |\mu_{T_{S,k}(\mathbf b),\eta}|(E)
       \leq C_\eta L^k C_S(E;V),\qquad
       C_\eta=A_\eta\,\frac{q^q}{q!}.
 \tag{DG3}
\]
At $k=0$ this is $|\mu_{S,\eta}|(E)\leq C_\eta C_S(E;V)$,
with no potential or $L$ involved. Consequently every raw
relative-capacity-null Borel set is coefficient-null for $S$
under (DG1). No converse is asserted.
\end{lemma}

\noindent\textbf{Status: \texttt{DIRECT\_COROLLARY}} of the new
affinoid compensator and the frozen original coefficient theorem.

\begin{proof}
Cover the compact test support by finitely many original moment
stages. On each stage exact affinoid-source unit replacement from
Lemma~\ref{lem:r11-affinoid-rational-unit-replacement} realizes
the finitely many moment functions as logarithmic norms of
quotients of functions on $B$. The original finite-dimensional
coefficient theorem bounds the variation of the localized test
pairing by a constant times
$\beta^{q-k}\wedge T$, where
$\beta=\sum_i d't_i\wedge d''t_i$ for those actual coordinates.

Lemma~\ref{lem:r11-affinoid-compensator} gives a global normalized
smooth classical psh function $a_l$ on $V$ with
$Da_l\geq c_l\beta$ near the relevant compact support.
Average the finitely many $a_l$. The resulting normalized $a$
has $Da\geq c'_l\beta$ on each corresponding stage. Wedge powers
preserve this inequality in the required direction: the difference
between the two powers is a sum of products of positive $(1,1)$
forms and hence is strongly positive. Pairing with the original
middle-positive $T$ is therefore legitimate. A finite partition
and a larger constant prove (DG2), including arbitrary actual
moment refinements of the starting test.

For (DG3), put $h_i=1+b_i/L\in[0,1]$. Bounded smooth-slot
compatibility gives
\[
 (Da)^{q-k}\wedge T_{S,k}(\mathbf b)
     =L^kT_{S,q}(h_1,\ldots,h_k,a,\ldots,a).
\]
The positive expansion of the $q$-th power of the average of
these $q$ normalized testers bounds this mixed measure by
$(q^q/q!)L^k C_S$ on every Borel set. This proves (DG3).
For $k=0$ there are only $q$ copies of $a$, so the same bound
applies directly to every original coefficient of $S$.
\end{proof}

\subsection{The tail modulus and partial products}

Fix a finite family $u_1,\ldots,u_m\in\mathcal F_{S,\partial}(V)$.
Its members may coincide. Set
\[
 w=\sum_{i=1}^m u_i,\quad W=1-w,\quad
 P=\bigcup_i\{u_i=-\infty\}=\{w=-\infty\},
 \quad
 F_w(t)=C_S(\{w<-t\};V).
 \tag{DG4}
\]
The solid-cone theorem makes $w$ a member of the same boundary
class. In particular its strict negative sublevels have compact
closure, $C_S^*(P;V)=0$, and $t^qF_w(t)\to0$.
For sufficiently large $t$ define the finite decreasing modulus
\[
                 \delta_w(t)=\sup_{s\geq t}s^qF_w(s).
 \tag{DG5}
\]
It tends to zero as $t\to\infty$.

For a multi-index $\alpha\in\mathbf N^m$, $|\alpha|=k\leq q$,
write $T_{\alpha,j}$ for the bounded mixed product containing
$\alpha_i$ copies of $u_{i,j}=\max(u_i,-j)$ over $S$.
The convention is $T_{0,j}=S$. Let
\[
                 E_j=\{w>-j\}.
 \tag{DG6}
\]
On $E_j$ every $u_i$ exceeds $-j$.

\begin{theorem}[Closed products in every degree from the capacity tail]
\label{thm:r11-boundary-all-degree-products}
Under (DG1), for every $0\leq|\alpha|\leq q$, the consistent
restrictions $\mathbf1_{E_j}T_{\alpha,j}$ increase to a locally
finite original middle-positive coefficient-Radon current
$N_\alpha$. For $\alpha=0$, $N_0=S$.
For $\alpha\ne0$ this is exactly the relative truncation product
of the differentiated tuple in
Theorem~\ref{thm:r11-closed-nonpolar-product}; it is separately
$d'$- and $d''$-closed. Every coefficient of $N_\alpha$ gives
zero variation to the entire $P$, including poles of members
not differentiated in $\alpha$.

The uncut currents converge in local coefficient total variation:
\[
                          T_{\alpha,j}\longrightarrow N_\alpha.
 \tag{DG7}
\]
Independent canonical truncation levels have the same limit as
their minimum tends to infinity. These assertions include every
original compact smooth test and every actual moment refinement.
For $k=q$ they recover the previously proved top-degree boundary
calculus; the new dimension gain is $k<q$.
\end{theorem}

\noindent\textbf{Status: \texttt{NEW\_PROVED}} below top degree,
relative to the printed one-affinoid domain condition.

\begin{proof}
Strict Borel maximum locality gives consistency on $E_j$,
and these carriers increase to $V\setminus P$. By (DG3) and
$C_S(P;V)=0$, every coefficient of $T_{\alpha,j}$ is zero on
$P$, for each fixed $j$; the case $\alpha=0$ gives
$\mathbf1_P S=0$. Thus the increasing zero-degree limit is $S$.

Suppose $k<q$, and fix one positive compact original control
$\Theta$ of degree $(q-k,q-k)$. On a value shell
\[
                  A_R=\{-2R<w\leq-R\},\qquad R\geq1,
\]
all members $u_i$ exceed $-2R$. Its restriction of any sufficiently
high truncation current therefore equals that of $T_{\alpha,2R}$
(use real levels or first take integral dyadic levels).
The original measure estimate (DG3) gives
\[
 \begin{split}
 (\Theta\wedge T_{\alpha,2R})(A_R)
 &\leq C_\Theta(2R)^k F_w(R/2)\\
 &\leq C_\Theta 2^{k+q}R^{k-q}\delta_w(R/2).
 \end{split}
 \tag{DG8}
\]
The harmless replacement of the weak bound $w\leq-R$ by the
strict bound $w<-R/2$ retains all equality-level masses.
The sum of (DG8) over dyadic $R$ converges because $k<q$.
On $\{w>-1\}$ consistency identifies every restricted current
with a restriction of the fixed bounded product $T_{\alpha,1}$.
Consequently all increasing cut traces have a common finite
local bound, and define the asserted coefficient-Radon limit.
Positivity and the coefficient theorem identify a single original
current on every compact stage.

To identify this with the standard product, let
$F_{\alpha,j}=\bigcap_{\alpha_i>0}\{u_i>-j\}$.
One has $E_j\subset F_{\alpha,j}$, so the standard cut limit
dominates the just constructed limit. Conversely, for fixed $j$
the current $\mathbf1_{F_{\alpha,j}}T_{\alpha,j}$ is zero on
$P$. Exhaust its complement by $E_l$, $l\geq j$.
On $F_{\alpha,j}\cap E_l$ strict maximum locality identifies it
with the corresponding restriction of $T_{\alpha,l}$.
It is therefore dominated by $N_\alpha$. Taking the increasing
limit in $j$ proves equality of the two cut constructions and
finite standard truncation mass. The closedness theorem cited
in the statement now applies. No Borel-restricted finite-stage
current has been declared closed.

For common level $j$, the uncut current and $N_\alpha$ agree
on $E_j$. The uncut trace on its complement is at most
$C_\Theta j^kF_w(j/2)$ by (DG3). The limiting trace there is
bounded by summing (DG8) with $R=2^lj$, $l\geq0$; $P$
contributes zero. Both tend to zero. This proves (DG7).

There is also a uniform estimate for independent levels
$l_i\geq j$, regardless of their maximum. On $A_R$, every
potential $\max(u_i,-l_i)$ exceeds $-2R$. Strict maximum
locality permits replacing it there by
$\max(u_i,-\min(l_i,2R))$, which is globally bounded between
$-2R$ and zero. Hence the same bound (DG8) holds on that shell,
uniformly in all $l_i$. On $E_j$ the independent product equals
$N_\alpha$. Summing the uniform shell bounds proves its
local coefficient total-variation convergence as $j\to\infty$.
This avoids using compact top-mass conservation in lower degree.

For $k=q$, Theorem~\ref{thm:r11-boundary-mixed-top-calculus}
already proves local finiteness, the standard cut construction,
uncut total-variation convergence, and independent-rate convergence.
If the family contains undifferentiated members, their pole union
is still coefficient-null for each bounded product by (DG3).
Exhausting its complement as above identifies the common-$w$
cut with the same top product. Separate closedness in top degree
is automatic. These arguments prove every assertion for all $k$.
\end{proof}

\begin{theorem}[Subcritical weights and a Dini endpoint]
\label{thm:r11-capacity-dimension-weighted}
Let $k=|\alpha|<q$, let $a\geq0$ and $a+k<q$, and fix a
positive compact original control $\Theta$.
Then $W^a\Theta\wedge N_\alpha$ is locally finite and, for all
sufficiently large $j$,
\[
 \int_{\{w\leq-j\}}W^a\,\Theta\wedge N_\alpha
       \leq C_{\Theta,\alpha,a,m}\,
               j^{a+k-q}\delta_w(j/2).
 \tag{DG9}
\]
With $W_j=1-\sum_i u_{i,j}$ one also has
\[
 \left\|\Theta\wedge
       (W_j^aT_{\alpha,j}-W^aN_\alpha)\right\|_{\rm TV}
       \leq C'_{\Theta,\alpha,a,m}\,
               j^{a+k-q}\delta_w(j/2).
 \tag{DG10}
\]
The corresponding coefficient-weighted assertions hold for any
continuous $F$ on the nonpositive orthant satisfying
$|F(\mathbf x)|\leq A(1-\sum_i x_i)^a$.

At the borderline $a+k=q$, both conclusions, with a tail tending
to zero in place of the right sides, hold under the explicit
capacity Dini condition
\[
                 \int_1^\infty t^{q-1}F_w(t)\,dt<\infty.
 \tag{DG11}
\]
More precisely their right sides can then be bounded by a fixed
constant times
$\int_{j/4}^\infty t^{q-1}F_w(t)\,dt$.
At $k=q,a=0$ unweighted convergence was already proved without
(DG11); the condition is only a sufficient extra quantitative
bound at that endpoint.
\end{theorem}

\noindent\textbf{Status: \texttt{NEW\_PROVED}}.
The index $\alpha=0$ is included. In particular the hypotheses
give local coefficient integrability of every $u_i$ against $S$
when $q\geq2$; no such unconditional first-moment assertion is
made when $q=1$.

\begin{proof}
On the shell $A_R$ one has $W\leq1+2R\leq3R$.
Multiplying (DG8) by $(3R)^a$ and summing over $R=2^lj$
gives a convergent geometric series of ratio $2^{a+k-q}$.
This proves (DG9). On the complement of $E_j$, one has
$W_j\leq1+mj\leq(m+1)j$. Equation (DG3) bounds the weighted
uncut tail there by
\[
                 C_\Theta(m+1)^a j^{a+k}F_w(j/2).
\]
On $E_j$ the weighted measures agree. The sum of this last
bound and (DG9) proves (DG10), and also proves the claimed
coefficient-weighted version for $F$ by the same absolute
growth bound. On the finite-value complement the actual values
are used; all multipliers on $P$ are immaterial because the
coefficient measures are zero there.

For the borderline, use the first line of (DG8), rather than
its bound by $\delta_w$. The weighted shell sum is at most a
constant times
\[
                         \sum_{l\geq0}(2^lj)^qF_w(2^{l-1}j).
\]
Since $F_w$ is decreasing,
\[
 R^qF_w(R)
    \leq \frac{q}{1-2^{-q}}
                     \int_{R/2}^{R}t^{q-1}F_w(t)\,dt .
\]
The dyadic intervals are disjoint. Their sum bounds the preceding
series by a fixed multiple of the tail integral in the statement.
The weighted uncut tail is bounded by the same integral over
$[j/4,j/2]$. This proves the borderline assertion, including
local weighted finiteness after adding the fixed finite-value core.
\end{proof}

\subsection{Continuity, multilinearity and recursive products}

\begin{theorem}[Joint continuity on the boundary cone in all degrees]
\label{thm:r11-boundary-all-degree-continuity}
Under (DG1), the products in
Theorem~\ref{thm:r11-boundary-all-degree-products} are symmetric,
additive and positively homogeneous in every differentiated slot
on $\mathcal F_{S,\partial}(V)$.
Suppose $v_{i,l}\downarrow u_i$ are negative extended classical
psh functions with each $u_i\in\mathcal F_{S,\partial}(V)$.
Then, for every $|\alpha|\leq q$,
\[
                    N_\alpha(\mathbf v_l)\longrightarrow
                                      N_\alpha(\mathbf u)
 \tag{DG12}
\]
on every original compact smooth complementary test.
At top degree, this also holds with every fixed bounded continuous
composition of the tuple as an undifferentiated weight, without
the Dini condition.
For $a+|\alpha|<q$, the same holds after multiplication by a
continuous $F$ of growth at most $(1-\sum_i x_i)^a$.
For $a+|\alpha|=q$ it holds under (DG11) for the limiting
sum $w=\sum_i u_i$. No background is varied in these passages.
\end{theorem}

\noindent\textbf{Status: \texttt{NEW\_PROVED}} below top degree.

\begin{proof}
Upward stability puts all $v_{i,l}$ in the same cone.
Since $\sum_i v_{i,l}\geq w$, its capacity-tail modulus is
bounded by $\delta_w$. Thus (DG9)--(DG10) give tail bounds
independent of $l$ in every subcritical degree and weight;
at the endpoint the limiting sum's Dini tail controls all rows.
For fixed truncation level $j$, frozen joint bounded decreasing
continuity applies. Its continuous bounded weight extension is
Lemma~\ref{lem:r11-continuous-bounded-joint-weight}.
Separate the two limits by the uniform weighted total-variation
tail estimates, first letting $l\to\infty$ at fixed $j$, then
$j\to\infty$. This proves the asserted weighted convergence and
(DG12) below top degree. At top degree (DG12) is precisely
Theorem~\ref{thm:r11-boundary-joint-decreasing}; its uniform
tail proof also allows a fixed bounded continuous multiplier.
For top-degree endpoint estimates under (DG11), use the Dini
bounds just proved.

For completeness, multilinearity requires a passage for sums
of truncated slots, not just equality of their values on a
Borel set. For positive coefficients $c_i$, set
$A=\sum_i c_i$, $c=\min_i c_i>0$,
$s_j=\sum_i c_i u_{i,j}$, and $z=\sum_i c_i u_i$.
Then
\[
                       \max(s_j,-cj)=\max(z,-cj).
 \tag{DG13}
\]
If any slot is truncated its contribution is at most $-cj$;
otherwise $s_j=z$. Strict maximum locality therefore identifies
the $k$-fold bounded product of $s_j$ with the unbounded product
of $z$ on $\{z>-cj\}$.
The trace of its complementary tail is at most
$C_\Theta(Aj)^k C_S(\{z<-cj/2\};V)$ by (DG3).
For $k<q$ this tends to zero, and the unbounded tail tends to
zero by (DG9). For $k=q$ the established top argument applies.
Hence the positive bounded polynomial expansion of the
$k$-fold product of $s_j$ passes to the corresponding unbounded
diagonal product for every positive coefficient vector.
Comparing coefficients in these finite polynomial identities,
also after replacing one family member by a sum, proves
additivity and positive homogeneity. Zero coefficients are
omitted before evaluation at an infinite value. Symmetry is
already inherited from every bounded approximant.
\end{proof}

\begin{corollary}[Subcritical recursive identities and Dini completion]
\label{cor:r11-boundary-recursion-dini}
Let $|\alpha|=k\leq q-2$. Then every multiplier $u_iN_\alpha$
is coefficient-defined on all original tests and
\[
                         D(u_iN_\alpha)=N_{\alpha+e_i}.
 \tag{DG14}
\]
If the common family satisfies (DG11), (DG14) also holds for
$k=q-1$. Thus that Dini subclass has the full original recursive
mixed calculus through top degree. Repetitions of a potential
are permitted throughout.
\end{corollary}

\noindent\textbf{Status: \texttt{NEW\_PROVED}}.
At $k=q-1$ without (DG11), automatic coefficient integrability
of the multiplier remains \texttt{OPEN} in this argument.
Once that integrability is verified separately, the no-loss
endpoint argument (CR6)--(CR7) in
Section~\ref{sec:r11-critical-recursive-residue} proves the
recursive identity without the Dini hypothesis. Thus (DG11)
is an explicit sufficient integrability condition, not a
necessary condition for endpoint recursion. The closed top
current itself already exists without that extra condition.

\begin{proof}
For $k\leq q-2$, take $a=1$ in
Theorem~\ref{thm:r11-capacity-dimension-weighted}.
It proves local integrability of $u_i$ against every original
coefficient of $N_\alpha$ and total-variation convergence of
$u_{i,j}T_{\alpha,j}$ to $u_iN_\alpha$ on every compact stage.
The bounded recursion gives
\[
                  D(u_{i,j}T_{\alpha,j})=T_{\alpha+e_i,j}.
\]
Its left side converges distributionally to $D(u_iN_\alpha)$;
its right side converges on all original tests by (DG7).
This proves (DG14). When $k=q-1$, the identical argument uses
the borderline $a=1$ estimate under (DG11). No multiplication
by an undefined unbounded function is performed.
\end{proof}

\paragraph{Relation to polar and negligible-set questions.}
Under the explicit domain condition (DG1), a raw-capacity-null
Borel set cannot carry a coefficient of the original background,
by the zero-degree case of (DG3). For the pole union of a boundary
class tuple, outer nullity is already known, and all partial
unbounded products and subcritical weighted products avoid it.
This is a relative implication with quantified hypotheses.
It does not prove the converse that every background-null Borel
set is capacity-null, a characterization of arbitrary negligible
sets, or general raw-capacity outer regularity.

The coefficient compensation in (DG2) is \texttt{FALSE} on
arbitrary compact boundaryless domains. For example take
$V=(\mathbf P^1_K)^{\rm an}$ and $S=[V]$. For every normalized
bounded classical tester $a$, the positive measure $D(aS)$ has
total mass $\langle aS,D1\rangle=0$ by compact Stokes. Thus
$C_S(E;V)=0$ for every Borel $E$, while $S$ is a nonzero original
current detected by a compact smooth form in an affine annulus.
This direct counterexample uses only frozen positivity and Stokes.
Dropping (DG1) from this argument is therefore a
\texttt{FAILED\_ATTEMPT}, not a harmless localization of its
fixed-domain tester class. The results also concern original
current equality; no additional GK delta-test equality or
analytic PS/PL relation quotient is obtained.


% END REVIEWED MODULE research_r11/capacity/capacity_dimension_gain.tex


% BEGIN REVIEWED MODULE research_r11/comparison/power_energy_mixed_stability.tex
\section{Power energies, integral mixed inequalities, and capacity stability}
\label{sec:r11-power-energy-stability}

Keep the original arbitrary $S$, ordinary domain $V$, and $q=n-p\geq1$.
All top products are those of
Theorem~\ref{thm:r11-boundary-mixed-top-calculus}. The statements
below use integrated positivity, comparison, and fixed-domain
capacity. They do not assume a pointwise Alexandrov--Fenchel or
hyperbolicity property for the middle-positive background. Such a
property is not available for arbitrary $S$.

For a negative extended classical psh $u$ with $\mathrm{BG}(u,0)$
and $s>0$, define the extended nonnegative quantities
\[
 \begin{split}
 e_s(u)&=\int_V(-u)^s\,d\mu_u,\\
 J_s(u)&=\int_0^\infty t^{q+s-1}
                       \operatorname{Cap}_S(\{u<-t\};V)\,dt.
 \end{split}
 \tag{PE1}
\]
The first integral is over the finite-value locus; $\mu_u$ is zero
on the omitted pole set. The integrand defining $J_s$ is Borel
because its capacity factor is nonincreasing in $t$.

\begin{theorem}[A capacity-only criterion for finite power energy]
\label{thm:r11-power-energy-capacity-equivalence}
For $u$ as above and $s>0$, the following are equivalent:
\[
 J_s(u)<\infty
 \quad\Longleftrightarrow\quad
 u\in\mathcal F_{S,\partial}(V)\ \hbox{and}\ e_s(u)<\infty.
 \tag{PE2}
\]
On this class, more generally with extended values allowed after
membership in $\mathcal F_{S,\partial}$ is known,
\[
 {s\over2^{q+s}}J_s(u)
       \leq e_s(u)
       \leq {s2^q\over1-2^{-s}}J_s(u).
 \tag{PE3}
\]
Denote the equivalent class in (PE2) by
$\mathcal E^s_{S,\partial}(V)$. Raw or outer capacity gives the
same $J_s$, because all its sublevels are ordinary open.
\end{theorem}

\noindent\textbf{Status: \texttt{NEW\_PROVED}} for the
all-original arbitrary-background top calculus. Finiteness of
$e_s$ alone, outside the no-loss class, is not a sufficient
condition; the logarithmic counterexample has $\mu_u=0$.

\begin{proof}
For $u\in\mathcal F_{S,\partial}$, the first inequality of (BC4)
and Tonelli's theorem give
\[
 \begin{split}
 J_s(u)
 &\leq2^q\int_0^\infty t^{s-1}\mu_u(\{u<-t/2\})\,dt\\
 &= {2^{q+s}\over s}\int_V(-u)^s\,d\mu_u.
 \end{split}
 \tag{PE4}
\]
For the opposite direction no absence of loss is needed.
Real-level truncations are bounded classical psh functions and
obey the same consistency as integer truncations: one may express
them as the maximum of an integer truncation and a constant.
On $\{-2t<u<-t\}$ the measure $\mu_u$ equals the bounded
measure of $\max(u,-2t)$. Normalizing this potential to $[0,1]$
therefore gives
\[
       \mu_u(\{-2t<u<-t\})
          \leq(2t)^q\operatorname{Cap}_S(\{u<-t\};V).
 \tag{PE5}
\]
For a fixed finite value $a=-u>0$, the parameter $t$ in this
annulus ranges over $(a/2,a)$. Tonelli consequently yields
\[
 {1-2^{-s}\over s}e_s(u)
   =\int_0^\infty t^{s-1}
                 \mu_u(\{-2t<u<-t\})\,dt
   \leq2^qJ_s(u).
 \tag{PE6}
\]
Endpoint levels cause no loss in this integration, since for each
fixed value they form a Lebesgue-null set of parameters.

It remains to show that finite $J_s$ forces no loss. Write
$\kappa=q+s$. Monotonicity of the sublevel capacity implies
\[
 {1-2^{-\kappa}\over\kappa}\,
       t^\kappa\operatorname{Cap}_S(\{u<-t\};V)
 \leq\int_{t/2}^t r^{\kappa-1}
                 \operatorname{Cap}_S(\{u<-r\};V)\,dr.
 \tag{PE7}
\]
When $J_s<\infty$, this proves
$t^q\operatorname{Cap}_S(\{u<-t\};V)\to0$.
Theorem~\ref{thm:r11-capacity-tail-criterion} gives
$u\in\mathcal F_{S,\partial}$, and (PE6) gives finite $e_s$.
Conversely, (PE4) proves finiteness of $J_s$ from the right side
of (PE2). This proves every assertion without an assumption on
an unrestricted unbounded Monge--Ampere operator.
\end{proof}

Put
\[
 \kappa=q+s,\qquad
 A_{q,s}={2^{2q+s}\over1-2^{-s}},\qquad
 C_{q,s}={A_{q,s}(q+1)^{\kappa+1}\over q!}.
 \tag{PE8}
\]
These constants are explicit and are not claimed optimal.

\begin{theorem}[Power-energy cone and an integral mixed Holder inequality]
\label{thm:r11-power-mixed-holder}
For every $s>0$, the class $\mathcal E^s_{S,\partial}(V)$ is
upward closed among negative extended classical psh functions,
closed under finite positive sums, and closed under finite maxima.
If $u\leq v\leq0$ with $u$ in the class, then
\[
                         e_s(v)\leq2^{q+s}e_s(u).
 \tag{PE9}
\]
For $u_1,\ldots,u_m$ in the class,
\[
                  e_s\left(\sum_{i=1}^m u_i\right)
                    \leq A_{q,s}m^\kappa\sum_{i=1}^m e_s(u_i).
 \tag{PE10}
\]
For any $u_0,u_1,\ldots,u_q$ in the class, the full mixed
weighted integral is finite and satisfies
\[
 \boxed{\quad
 \int_V(-u_0)^s\,dM(u_1,\ldots,u_q)
 \leq C_{q,s}\,e_s(u_0)^{s/\kappa}
                         \prod_{i=1}^q e_s(u_i)^{1/\kappa}.
 \quad}
 \tag{PE11}
\]
A zero factor makes the right side zero. The mixed measure gives
zero mass to the pole of $u_0$, so the integral retains the
actual representative without any undefined product at a pole.
\end{theorem}

\noindent\textbf{Status: \texttt{NEW\_PROVED}}. This is an
integral Cegrell-type estimate with a nonsharp constant. It is
not a pointwise determinant inequality for $S$.

\begin{proof}
Upward membership follows from the solid boundary cone and
$J_s(v)\leq J_s(u)$. Integrating (BC12) by layer cake gives
the sharper numerical estimate (PE9). For a sum, (BC6) with
all coefficients one yields
\[
 J_s\left(\sum_i u_i\right)\leq m^\kappa\sum_iJ_s(u_i).
\]
Combine both inequalities of (PE3) to obtain (PE10).
Positive homogeneity $e_s(cu)=c^\kappa e_s(u)$ follows from
the exact scaling of the top measure. These facts prove the
stated cone properties; a maximum follows from upward closure.

Let $w=u_0+\cdots+u_q$. Its positive mixed expansion contains
$q!M(u_1,\ldots,u_q)$, hence
\[
                 q!M(u_1,\ldots,u_q)\leq\mu_w.
 \tag{PE12}
\]
Since $w=-\infty$ at every pole of $u_0$, the right side gives
zero mass to that pole, and so does the left side. On the
remaining locus, $(-u_0)^s\leq(-w)^s$. Thus the left side
of (PE11) is at most $e_s(w)/q!$ and is finite.

For $\varepsilon>0$ set
$v_i=u_i/(e_s(u_i)+\varepsilon)^{1/\kappa}$.
Then $e_s(v_i)\leq1$. By (PE10) with $m=q+1$,
$e_s(\sum_{i=0}^qv_i)/q!\leq C_{q,s}$.
Apply (PE12) to this normalized tuple and then use positive
homogeneity in all $q$ differentiated slots and the degree-$s$
homogeneity of the weight. This gives
\[
 \int(-u_0)^s\,dM(u_1,\ldots,u_q)
 \leq C_{q,s}(e_s(u_0)+\varepsilon)^{s/\kappa}
                       \prod_{i=1}^q(e_s(u_i)+\varepsilon)^{1/\kappa}.
\]
Letting $\varepsilon\downarrow0$ proves (PE11), including every
degenerate zero-energy case. No nondegeneracy of $S$ is used.
\end{proof}

\begin{proposition}[Capacity domination and quantitative top truncation tails]
\label{prop:r11-power-capacity-domination}
Let $u_1,\ldots,u_q\in\mathcal E^s_{S,\partial}(V)$, put
$w=\sum_i u_i$, and write $M=M(u_1,\ldots,u_q)$.
For every Borel $E\subset V$ and every $t>0$,
\[
 M(E)\leq{1\over q!}
       \left(q^qt^q\operatorname{Cap}_S(E;V)+t^{-s}e_s(w)\right).
 \tag{PE13}
\]
In particular, for finite-capacity sets,
\[
 M(E)\leq B_{q,s}
       \left(\prod_{i=1}^q e_s(u_i)^{1/\kappa}\right)
                    \operatorname{Cap}_S(E;V)^{s/\kappa},
 \quad
 B_{q,s}={q^q+1\over q!}
                  \left(A_{q,s}q^{\kappa+1}\right)^{q/\kappa}.
 \tag{PE14}
\]
For infinite capacity the estimate is vacuous. A zero energy
factor makes $M$ the zero measure.
If $M_k=T_{S,q}(\max(u_1,-k),\ldots,\max(u_q,-k))$, then
\[
 \|M_k-M\|_{\mathrm{TV}(V)}
       \leq {q^q4^{q+s}+1\over q!}\,k^{-s}e_s(w).
 \tag{PE15}
\]
All inequalities hold on actual Borel sets, not only on a chosen
moment image or a restricted family of test functions.
\end{proposition}

\noindent\textbf{Status: \texttt{NEW\_PROVED}}.

\begin{proof}
On $E_t=\bigcap_i\{u_i>-t\}$, strict locality identifies $M$
with the bounded mixed product of the real-level truncations.
Their sum lies in $[-qt,0]$. The positive expansion and capacity
normalization give
\[
                         M(E\cap E_t)
                        \leq{(qt)^q\over q!}\operatorname{Cap}_S(E;V).
\]
The complement of $E_t$ is contained in $\{w\leq-t\}$.
The positive expansion $q!M\leq\mu_w$ and the power moment
bound give $M(V\setminus E_t)\leq e_s(w)/(q!t^s)$.
This proves (PE13).

If $c=\operatorname{Cap}_S(E;V)$ and $e=e_s(w)$ are positive
and finite, choose $t=(e/c)^{1/\kappa}$ in (PE13), obtaining
\[
                   M(E)\leq{q^q+1\over q!}
                                   e^{q/\kappa}c^{s/\kappa}.
\]
The cases $c=0$ or $e=0$ follow by letting $t\to\infty$ or
$t\downarrow0$, respectively. Normalize the $q$ slots by
$(e_s(u_i)+\varepsilon)^{1/\kappa}$ as in the preceding proof;
(PE10) bounds the energy of their sum by
$A_{q,s}q^{\kappa+1}$. Scaling back and letting
$\varepsilon\downarrow0$ gives (PE14). If an energy is zero,
this proves vanishing on every compact set, whose relative
capacity is finite by frozen CLN; the Radon measure is therefore
zero.

For (PE15), use the exact agreement on $E_k$. On its complement,
the bounded positive expansion and (BC4) give
\[
 \begin{split}
 M_k(V\setminus E_k)
  &\leq{(qk)^q\over q!}\operatorname{Cap}_S(\{w<-k/2\};V)\\
  &\leq{q^q4^q\over q!}\mu_w(\{w<-k/4\})
   \leq{q^q4^{q+s}\over q!}k^{-s}e_s(w).
 \end{split}
\]
The limiting mass on the complement is at most
$k^{-s}e_s(w)/q!$. Their sum bounds total variation and proves
(PE15).
\end{proof}

\subsection{The exact capacity quantifier for nonmonotone convergence}

\begin{definition}[Truncated local convergence in original capacity]
\label{def:r11-truncated-local-capacity}
For negative extended classical psh functions $v_l,u$ on $V$,
write $v_l\to u$ in truncated local original capacity if the
following holds. For every ordinary analytic Laurent-box domain
$\Omega\Subset V$ whose compact closure is contained in an
original unit-moment chart, every compact $L\Subset\Omega$,
every $k>0$, and every $\delta>0$,
\[
 \operatorname{Cap}_{S|\Omega}
 \left(L\cap\{|\max(v_l,-k)-\max(u,-k)|>\delta\};\Omega\right)
                                               \longrightarrow0.
 \tag{PE16}
\]
The quantifier includes every actual finite refinement of that
moment chart and every further Laurent-box shrink. Under the
additional one-affinoid condition (DG1), Corollary
\ref{cor:r11-affinoid-local-capacity-topology} proves its equivalence
with fixed-domain capacity convergence on every compact core of $V$.
Without that domain hypothesis the printed all-local definition
remains the hypothesis used here.
\end{definition}

\begin{lemma}[The bounded finite-stage stability input on analytic boxes]
\label{lem:r11-bounded-analytic-capacity-stability}
For uniformly bounded classical psh families, (PE16) implies the
frozen weighted BT stability on every original compact test, in
every degree $0\leq r\leq q$. The same multiplier conclusion
holds for a bounded continuous scalar composition of one of
the potentials.
\end{lemma}

\noindent\textbf{Status: \texttt{DIRECT\_COROLLARY}} of the
frozen R1/R4 theorem \texttt{thm:capacity-stability}, with the
analytic bounded quasicontinuity of
Corollary~\ref{cor:r11-analytic-bounded-relative-QC} supplying
its domain extension. The bounded calculus itself is not new.

\begin{proof}
We record the precise controls, so no stronger topology is
silently substituted. Localize an original test to finitely many
of the stated boxes. Its finitely many moment coordinates $t_a$
are bounded on each compact closure. The bounded smooth classical
psh function $\rho=\frac12\sum_a t_a^2$ has
$D\rho=\sum_a d't_a\wedge d''t_a$. Finite original positive
matrix domination bounds every coefficient measure of a degree-$r$
bounded product on the compact stage by a constant times
\[
          T_{S,q}(a_1,\ldots,a_r,\rho,\ldots,\rho).
\]
Mixed capacity domination therefore bounds it by
$C\operatorname{Cap}_{S|\Omega}(\,\cdot\,;\Omega)$, uniformly
in the family. Only wedge with positive $(1,1)$ forms is used,
so the arbitrary middle-positive cone is retained.

The analytic Laurent collar theorem gives the exact bounded
quasicontinuity needed for the frozen capacity-controlled
multiplier lemma. Its proof separates a changing bounded weight
on its capacity-small gap set and replaces the fixed limit weight
by a continuous function off a small ordinary open exception.
The coefficient domination just given controls both errors.
Induction through $D(uP)$ then proves the bounded stability in
every degree, exactly as in the frozen proof. A bounded continuous
composition has capacity convergence by uniform continuity on
the bounded range and is quasicontinuous by composition, so the
same multiplier lemma applies to it. All controls in this proof use
the local capacity in (PE16). Under (DG1), the separate analytic
extension theorem (AL8)--(AL9) gives the equivalent fixed-domain
criterion; that comparison is not assumed in the multiplier proof.
\end{proof}

\begin{theorem}[Nonmonotone stability and lower semicontinuity of power energy]
\label{thm:r11-power-capacity-stability}
Fix $s>0$. For $1\leq i\leq q$, let
$v_{i,l}\in\mathcal E^s_{S,\partial}(V)$ and let $u_i$ be
negative extended classical psh with $\mathrm{BG}(u_i,0)$.
Assume (PE16) in every slot and
\[
                         \sup_{i,l}e_s(v_{i,l})<\infty.
 \tag{PE17}
\]
Then $u_i\in\mathcal E^s_{S,\partial}(V)$,
\[
 e_s(u_i)\leq\liminf_l e_s(v_{i,l}),\qquad
 M(v_{1,l},\ldots,v_{q,l})\longrightarrow M(u_1,\ldots,u_q)
 \tag{PE18}
\]
on every original compact smooth scalar test. No monotonicity of
the approximating potentials is assumed.

Suppose additionally that $v_{0,l}$ and $u_0$ satisfy the same
hypotheses. If $\Psi:[0,\infty)\to\mathbf R$ is continuous and
\[
                         |\Psi(t)|/t^s\longrightarrow0
                                      \quad(t\to\infty),
 \tag{PE19}
\]
then the locally finite weighted measures satisfy
\[
 \Psi(-v_{0,l})M(v_{1,l},\ldots,v_{q,l})
        \longrightarrow\Psi(-u_0)M(u_1,\ldots,u_q)
 \tag{PE20}
\]
on the same full original test family. Their weights at poles
are immaterial because the measures give those poles zero mass.
In particular powers $t^a$ with $0<a<s$ are allowed.
\end{theorem}

\noindent\textbf{Status: \texttt{NEW\_PROVED}}. The exact
localized capacity hypothesis and the global power-energy bound
are both retained.

\begin{proof}
First prove membership of a limit $u$, before applying any
unbounded continuity assertion. Fix $t,\delta>0$. The set
$\{u<-t-\delta\}$ has compact closure by its boundary-gap
condition. Choose a finite cover of that closure by compact
cores $L_a\Subset\Omega_a$ of admissible boxes. Choose
$k>t+\delta$. On the strict sublevel one has
\[
 \{u<-t-\delta\}
 \subset\{v_l<-t\}\cup
       \{|\max(v_l,-k)-\max(u,-k)|>\delta\}.
\]
The second set need only be considered on the fixed compact
cover. Domain restriction and (PE16) make its contribution to
the capacity measured on $V$ tend to zero. Consequently
\[
 \operatorname{Cap}_S(\{u<-t-\delta\};V)
       \leq\liminf_l\operatorname{Cap}_S(\{v_l<-t\};V).
\]
Let $\delta\downarrow0$. The raw capacity is continuous from
below: this follows by commuting the two increasing suprema
over capacity testers and over the increasing carriers. Hence
the same inequality holds with $\{u<-t\}$ on the left.
Fatou and (PE4) give
\[
 J_s(u)\leq\liminf_lJ_s(v_l)
       \leq {2^{q+s}\over s}\liminf_l e_s(v_l)<\infty.
 \tag{PE21}
\]
Theorem~\ref{thm:r11-power-energy-capacity-equivalence} proves
that every $u_i$ lies in the energy class. This step used only
bounded capacity and the printed boundary condition.

Put $w_l=\sum_i v_{i,l}$ and $w=\sum_i u_i$.
Equation (PE10) gives one finite uniform bound for $e_s(w_l)$
and $e_s(w)$. By (PE15), replacing every differentiated slot
by its truncation at $-k$ changes its mixed top measure by at
most a constant times $k^{-s}$ in total variation, uniformly
in $l$. For fixed $k$, Lemma
\ref{lem:r11-bounded-analytic-capacity-stability} gives convergence
of the bounded products on every original compact test.
Pass first in $l$ and then in $k$; the explicit uniform
variation bound proves the unweighted convergence in (PE18).

We next prove (PE20). The integral inequality (PE11) and the
uniform energy bound give
\[
 \sup_l\int(-v_{0,l})^s\,dM(v_{1,l},\ldots,v_{q,l})<\infty,
 \qquad
 \int(-u_0)^s\,dM(u_1,\ldots,u_q)<\infty.
 \tag{PE22}
\]
For $R>0$ let $\Psi_R(t)=\Psi(\min(t,R))$ and put
$\epsilon_R=\sup_{t\geq R}|\Psi(t)|/t^s$ for large $R$.
Then $\epsilon_R\to0$ and
\[
          |\Psi(t)-\Psi_R(t)|
                \leq2\epsilon_R t^s\mathbf1_{\{t>R\}}.
\]
By (PE22), the errors from replacing the weight by $\Psi_R$
are uniformly small on every compact test as $R\to\infty$.
For fixed $R$, the weights are bounded continuous compositions
of truncations of the zeroth potential. Replacing the other
slots by their truncations at $-k$ has a uniform error bounded
by $\|\Psi_R\|_\infty$ times (PE15). For fixed $R,k$, the
bounded multiplier statement of the preceding lemma passes
the limit in $l$. Let $k\to\infty$ and finally $R\to\infty$.
This proves (PE20). Local finiteness follows from (PE22) on
the large-value region and the uniform finite compact masses
on the remaining bounded-value region. The positive expansion
with the zeroth potential proves that its pole has zero mixed
mass, as in (PE12).

Finally apply the bounded-weight case just proved to a single
potential repeated in every differentiated slot, with weight
$\min(t^s,R^s)$ and any nonnegative original compact smooth
$0\leq\chi\leq1$. It gives
\[
 \int\chi\min((-u)^s,R^s)\,d\mu_u
  =\lim_l\int\chi\min((-v_l)^s,R^s)\,d\mu_{v_l}
                         \leq\liminf_l e_s(v_l).
\]
Increase $R$ and then take the supremum over such compact
cutoffs. Radon inner regularity gives $e_s(u)$ on the left.
This proves the sharp lower-semicontinuity assertion in (PE18),
improving the preliminary bound from (PE21).
\end{proof}

\begin{corollary}[A precise closure principle for class-valued nonlinear Dirichlet data]
\label{cor:r11-power-Dirichlet-closure}
Fix $s>0$ and a positive Radon datum $\mu$ on $V$. Suppose there
are bounded negative classical psh $u_l$ satisfying
$\mathrm{BG}(u_l,0)$,
\[
 (Du_l)^q\wedge S=\mu_l,\qquad
 \mu_l\longrightarrow\mu\ \hbox{on all original compact tests},
 \qquad\sup_l e_s(u_l)<\infty,
 \tag{PE23}
\]
and suppose $u_l$ has a limit $u$ in the exact sense (PE16),
where $u$ is actual negative extended classical psh and
$\mathrm{BG}(u,0)$. Then
\[
                    u\in\mathcal E^s_{S,\partial}(V),
                    \qquad\mu_u=\mu.
 \tag{PE24}
\]
Any two solutions in this class with datum $\mu$ agree outside
a relative outer-capacity-null set.
\end{corollary}

\noindent\textbf{Status: \texttt{CONDITIONAL}} as a nonlinear
Dirichlet existence mechanism: bounded solvability, the uniform
power-energy estimate, an actual classical capacity limit, and
its zero-boundary gap condition are the explicit premises.
The passage to the limit and the stated uniqueness conclusion
are fully proved. No existence or compactness premise is declared
automatic for arbitrary original $S$.

\begin{proof}
Apply Theorem~\ref{thm:r11-power-capacity-stability} with the
same potential in every slot. The original current limit is
both $\mu_u$ and $\mu$, proving (PE24). Outer-capacity uniqueness
is Corollary~\ref{cor:r11-boundary-outer-domination}.
\end{proof}

\paragraph{Precise remaining gates.}
The theory above supplies full-dimensional arbitrary-background
energy inequalities and stability in a specified unbounded
boundary class. It does not prove capacity compactness of an
arbitrary energy-bounded family, regularization of a proposed
unbounded envelope, pointwise uniqueness for degenerate $S$,
or unrestricted nonlinear Dirichlet existence. The original
all-refinement PS/PL quotient and CTO remain separate open gates.

% END REVIEWED MODULE research_r11/comparison/power_energy_mixed_stability.tex


% BEGIN REVIEWED MODULE research_r11/energy/power_energy_all_degree_calculus.tex
\section{Power energy controls all original degrees and unbounded energies}
\label{sec:r11-power-all-degree}

Keep the explicit domain condition (DG1): the compact closure of
the ordinary open set $V$ lies in the ordinary interior of one
original $K$-affinoid domain. The background $S$ is an arbitrary
original symmetric middle-positive separately closed $(p,p)$-current
on $V$, and $q=n-p\geq1$. No smoothness, geometric origin or prior
BT representation of $S$ is assumed.

Fix $\sigma>0$ and a finite family
$u_1,\ldots,u_m\in\mathcal E^\sigma_{S,\partial}(V)$. Put
\[
 \kappa=q+\sigma,\quad w=\sum_i u_i,\quad W=1-w,\quad
 W_j=1-\sum_i\max(u_i,-j),\quad E_j=\{w>-j\}.
 \tag{PG1}
\]
Let $T_{\alpha,j}$ and $N_\alpha$ be the bounded and unbounded
mixed currents from Section~\ref{sec:r11-capacity-dimension-gain};
$|\alpha|=k\leq q$ includes repeated slots and $\alpha=0$.
The power-energy cone theorem and (PE2) give
\[
 {\cal J}_w(R):=\int_R^\infty
       t^{\kappa-1}C_S(\{w<-t\};V)\,dt<\infty
       \quad(R>0),\qquad {\cal J}_w(R)\longrightarrow0.
 \tag{PG2}
\]
This is an actual capacity integral in the fixed domain $V$.

\begin{lemma}[Graded coefficient moments including the endpoint]
\label{lem:r11-graded-curvature-moments}
Let $a\geq0$ and $a+k\leq\kappa$. For every original compact
positive coefficient control $\Theta$ of degree $(q-k,q-k)$,
the measures $W^a\Theta\wedge N_\alpha$ are locally finite.
There is a finite constant $C_\Theta$, independent of $j,R$, such
that for $R\geq1$,
\[
 \begin{split}
 \sup_{j\geq1}\int_{\{w\leq-R\}}W_j^a\,
                           \Theta\wedge T_{\alpha,j}
     &\leq C_\Theta R^{a+k-\kappa}{\cal J}_w(R/4),\\
 \int_{\{w\leq-R\}}W^a\,\Theta\wedge N_\alpha
     &\leq C_\Theta R^{a+k-\kappa}{\cal J}_w(R/4).
 \end{split}
 \tag{PG3}
\]
The constant may depend on the family size, degree, exponent,
and coefficient control. Each full local moment in the first
line, with the tail indicator removed, is uniformly bounded in $j$.

For every continuous $F$ on the nonpositive orthant with
$|F(x)|\leq A(1-\sum_i x_i)^a$, there is local coefficient
total-variation convergence
\[
               F(\mathbf u_j)T_{\alpha,j}
                          \longrightarrow F(\mathbf u)N_\alpha.
 \tag{PG4}
\]
The endpoint $a+k=\kappa$ is included for this fixed family.
Multipliers on its pole union are immaterial because all these
currents have zero coefficient mass there.
\end{lemma}

\noindent\textbf{Status: \texttt{NEW\_PROVED}} as a
degree-dependent strengthening of (DG9)--(DG11), using the
proved power-energy capacity criterion.

\begin{proof}
Write $F_w(t)=C_S(\{w<-t\};V)$. On
$A_T=\{-2T<w\leq-T\}$ all $u_i>-2T$. For any $j$, strict
bounded maximum locality permits replacing each $u_{i,j}$
on this shell by $\max(u_i,-\min(j,2T))$, whose global range
is $[-2T,0]$. Estimate (DG3) and $W_j\leq W\leq1+2T$ give
\[
 \int_{A_T}W_j^a\Theta\wedge T_{\alpha,j}
       \leq C T^{a+k}F_w(T/2).
 \tag{PG5}
\]
This is uniform even when $j<T$. Each finite-stage current is
zero on the pole union by (DG3). The same shell estimate holds
for $N_\alpha$ by consistent restriction. Monotonicity gives
\[
 T^\kappa F_w(T/2)
       \leq C_\kappa\int_{T/4}^{T/2}
                         t^{\kappa-1}F_w(t)\,dt.
 \tag{PG6}
\]
Sum (PG5) over $T=2^\ell R$, $\ell\geq0$. The factors
$T^{a+k-\kappa}$ are bounded by $R^{a+k-\kappa}$, and the
integration intervals in (PG6) are disjoint. This proves (PG3).
On $\{w>-1\}$ consistency with the fixed level-one product
controls the remaining compact coefficient masses.

The bounded and limiting weighted currents agree on $E_j$.
Their variations on its complement are bounded by (PG3) at
$R=j$ and the growth bound for $F$. The bounds tend to zero,
also at equality, because $ {\cal J}_w(j/4)\to0$.
Finite coefficient domination proves (PG4) for every original
test and actual refinement.
\end{proof}

\begin{theorem}[Full recursive calculus in every positive power class]
\label{thm:r11-power-full-recursion}
For every $\sigma>0$, every finite family in
$\mathcal E^\sigma_{S,\partial}(V)$, and every
$|\alpha|\leq q-1$, the multiplier $u_iN_\alpha$ is defined
on all original coefficients and
\[
                         D(u_iN_\alpha)=N_{\alpha+e_i}.
 \tag{PG7}
\]
All repeated and mixed products in this class therefore have
the full recursive differential meaning through top degree,
in addition to separate closedness, symmetry, multilinearity,
and truncation independence.
\end{theorem}

\noindent\textbf{Status: \texttt{NEW\_PROVED}} as an
all-degree consequence of one top power-energy condition per
potential. The affinoid domain hypothesis remains in force.

\begin{proof}
For $k=|\alpha|\leq q-1$, one has $1+k\leq q<q+\sigma$.
Apply (PG4) with $F(x)=x_i$ before differentiating. The bounded
identity $D(u_{i,j}T_{\alpha,j})=T_{\alpha+e_i,j}$ then passes
distributionally to (PG7). The other properties were proved
without this recursion in the dimension-gain theorem; there
is no circular definition of the top product.
\end{proof}

\subsection{A direct capacity bound for bounded energy currents}

\begin{lemma}[The extra two powers in an energy trace]
\label{lem:r11-bounded-energy-capacity-power}
Let $-L\leq b_i\leq0$ be bounded classical psh functions on $V$,
and let $Q_\alpha$ be their bounded mixed product over $S$,
$|\alpha|=k\leq q-1$. For every original compact positive
control $\Theta$ of degree $(q-k-1,q-k-1)$ and Borel $A\subset V$,
\[
       \int_A\Theta\wedge\Gamma_{Q_\alpha}(b_i)
                 \leq C_\Theta L^{k+2}C_S(A;V).
 \tag{PG8}
\]
The constant is independent of the bounded potentials and $L>0$.
Mixed energy coefficient variations have the corresponding
bound by the sum of the diagonal traces.
\end{lemma}

\noindent\textbf{Status: \texttt{DIRECT\_COROLLARY}} of
the bounded energy identity and the global original-coefficient
compensator, retaining the arbitrary background.

\begin{proof}
Put $v=b_i+L$, so $0\leq v\leq L$ and $v^2$ is bounded
classical psh. The energy identity gives
\[
 0\leq\Gamma_{Q_\alpha}(b_i)
       =\tfrac12D(v^2Q_\alpha)-vD(b_iQ_\alpha)
       \leq\tfrac12D(v^2Q_\alpha).
 \tag{PG9}
\]
These are middle-cone inequalities; the nonnegative Borel
multiplier in the second term preserves positivity.
Use (DG2) for the energy current and its original test.
Normalize the $k$ slots in $Q_\alpha$ by $L$, the first
slot $v^2$ by $L^2$, and fill the remaining $q-k-1$ slots
with the global normalized smooth compensator. Positive
mixed expansion gives the factor $(q^q/q!)C_S$ and proves
(PG8). The positive energy matrix and its measure
Cauchy--Schwarz inequality control cross variations.
No energy current has been assumed closed.
\end{proof}

\begin{theorem}[Unbounded energy currents in the graded power domain]
\label{thm:r11-power-graded-energy}
Assume $|\alpha|=k\leq q-1$ and $k+2\leq\kappa=q+\sigma$.
There is an intrinsic matrix of coefficient-Radon currents
\[
 G^\alpha_{ia}
  =\tfrac12\left[
     D(u_i u_aN_\alpha)-u_iN_{\alpha+e_a}
                           -u_aN_{\alpha+e_i}\right].
 \tag{PG10}
\]
It is positive semidefinite in the original middle cone,
has zero coefficient mass on the pole union, and satisfies
\[
       \Gamma_{T_{\alpha,j}}(u_{i,j},u_{a,j})
                                \longrightarrow G^\alpha_{ia}
 \tag{PG11}
\]
in local coefficient total variation. If $b\geq0$ and
$b+k+2\leq\kappa$, then for every original compact positive
control $\Theta$,
\[
 \begin{split}
 \sup_j\int_{\{w\leq-R\}} W_j^b
           \Theta\wedge\Gamma_{T_{\alpha,j}}(u_{i,j})
   &\leq C_\Theta R^{b+k+2-\kappa}{\cal J}_w(R/4),\\
 \int_{\{w\leq-R\}} W^b\Theta\wedge G^\alpha_{ii}
   &\leq C_\Theta R^{b+k+2-\kappa}{\cal J}_w(R/4).
 \end{split}
 \tag{PG12}
\]
The corresponding full local moments are finite. A continuous
weight of growth at most $W^b$ may be inserted in (PG11),
with the same total-variation conclusion for this fixed tuple.

In particular $\sigma\geq1$ gives these energy currents in
every permitted energy degree. If $0<\sigma<1$, the conclusion
applies to $k\leq q-2$. The energy currents need not be closed.
\end{theorem}

\noindent\textbf{Status: \texttt{NEW\_PROVED}}. This improves
the sufficient superquadratic all-trace hypothesis (UE2) on the
printed domain, and includes the finite first-power energy endpoint.

\begin{proof}
On $A_T=\{-2T<w\leq-T\}$ replace every truncated slot by its
further truncation at $-2T$. Lemma
\ref{lem:r11-energy-truncation-locality} permits this in the
background slots and the differentiated energy slot. Its proof
uses shifted squares and four-branch maxima, so also applies
with unequal initial truncation levels. Estimate (PG8) gives
\[
 \int_{A_T}W_j^b
        \Theta\wedge\Gamma_{T_{\alpha,j}}(u_{i,j})
                  \leq C T^{b+k+2}F_w(T/2).
 \tag{PG13}
\]
Each finite-stage energy is zero on the pole union by (PG8)
and $C_S(P;V)=0$. Sum (PG13) with (PG6) to obtain the first
bound of (PG12).

Bounded energy locality gives consistency of the restrictions
on $E_j$. The positive cut diagonal energies increase, and
the shell estimates and level-one core give locally finite
limits. Their tails and those of the uncut energies satisfy
(PG12), first with $b=0$, then with the stated weights.
Both tails tend to zero, proving total-variation convergence.
Polarization gives the cross limits, the positive matrix
assertion, and all coefficient variation bounds.

The three terms in (PG10) are defined before identification:
use (PG3) with $(a,k)=(2,k)$ and $(a,k)=(1,k+1)$.
Their total exponent is $k+2\leq\kappa$. Pass the bounded
three-term energy formula through (PG4) and distributional
continuity. This identifies the limits with (PG10), including
every original compact test and actual refinement.
No Borel-restricted current has been differentiated.
\end{proof}

\begin{theorem}[A degree-dependent unbounded chain rule]
\label{thm:r11-power-graded-chain}
Assume the preceding hypotheses. Let $F$ be $C^2$ near
the nonpositive orthant and suppose, for some
$2\leq r\leq\kappa-k$ and finite $A$,
\[
 |F(x)|\leq A W(x)^r,\quad
 |F_i(x)|\leq A W(x)^{r-1},\quad
 |F_{ia}(x)|\leq A W(x)^{r-2},
 \qquad W(x)=1-\sum_i x_i.
 \tag{PG14}
\]
Then the following is an equality of original
coefficient-defined currents on ordinary opens:
\[
 D(F(\mathbf u)N_\alpha)
       =\sum_i F_i(\mathbf u)N_{\alpha+e_i}
                         +\sum_{i,a}F_{ia}(\mathbf u)G^\alpha_{ia}.
 \tag{PG15}
\]
Convex coordinatewise nondecreasing $F$ gives a middle-positive
right side. The endpoint $r=\kappa-k$ is permitted for each
fixed finite-power-energy family.
\end{theorem}

\noindent\textbf{Status: \texttt{NEW\_PROVED}} with the
printed growth, background and domain hypotheses.

\begin{proof}
Apply the bounded multivariable chain identity to common
truncations. The term to be differentiated has growth $W^r$
and degree $k$; each curvature multiplier has growth
$W^{r-1}$ and degree $k+1$; each energy multiplier has growth
$W^{r-2}$ and energy budget $k+2$. Each total exponent is
$r+k\leq\kappa$. Equations (PG4) and (PG12) therefore pass
every term in local coefficient total variation.
Distributional continuity proves (PG15). The sign assertion
uses $F_i\geq0$, the positive Hessian of $F$, and the positive
energy matrix, without specializing $S$.
\end{proof}

\paragraph{Uniform families and the critical boundary.}
For a uniform $\sigma$-power energy bound, (PE10) and (PE4)
bound $ {\cal J}_{w_l}(0)$ uniformly. The tail estimates above
are then uniform and vanish whenever the relevant total
exponent is strictly below $q+\sigma$. At equality, a bound
for the full integral does not force a uniformly vanishing
tail. This distinction is essential for varying tuples;
it does not affect the fixed-tuple endpoint constructions.

% END REVIEWED MODULE research_r11/energy/power_energy_all_degree_calculus.tex


% BEGIN REVIEWED MODULE research_r11/energy/first_order_capacity_currents.tex
\section{First-order coefficient currents below the quadratic energy threshold}
\label{sec:r11-first-order-power-currents}

Retain the original arbitrary background, affinoid domain,
finite power-energy family, and notation (PG1)--(PG2).
The construction requires only $\sigma>0$; it does not require
that every quadratic energy current exist.
Set $k=|\alpha|\leq q-1$ and $\kappa=q+\sigma$.
A first-order \emph{coefficient measure} is its pairing with
each fixed original complementary smooth form, with actual
Borel scalar multipliers inserted. All variation bounds
are on compact original test stages.

\begin{lemma}[Bounded first-order chain and Borel locality]
\label{lem:r11-first-order-borel-locality}
Let $Q$ be an arbitrary original positive separately closed
current and $v_1,\ldots,v_m$ bounded classical psh functions.
For every $C^1$ function $F$ near their bounded range box,
\[
 d'(F(\mathbf v)Q)=\sum_i F_i(\mathbf v)d'(v_iQ),\qquad
 d''(F(\mathbf v)Q)=\sum_i F_i(\mathbf v)d''(v_iQ).
 \tag{FO1}
\]
These are coefficient-Radon identities. For any real $t$,
\[
 d'(\max(v_i,t)Q)=\mathbf1_{\{v_i>t\}}d'(v_iQ),\qquad
 d''(\max(v_i,t)Q)=\mathbf1_{\{v_i>t\}}d''(v_iQ).
 \tag{FO2}
\]
If positive separately closed backgrounds $Q,Q'$ of the
same degree have equal coefficient restrictions on a Borel
set $A$ in the full original test sense, then for every
fixed bounded classical $v$,
\[
 \mathbf1_A d'(vQ)=\mathbf1_A d'(vQ'),\qquad
 \mathbf1_A d''(vQ)=\mathbf1_A d''(vQ').
 \tag{FO3}
\]
Consequently $d'(u_{i,j}T_{\alpha,j})$ and
$d''(u_{i,j}T_{\alpha,j})$ are consistent on $E_j$
as $j$ increases.
\end{lemma}

\noindent\textbf{Status: \texttt{DIRECT\_COROLLARY}} of
the bounded first-order measure theorem and strict maximum
locality, without differentiating a Borel restriction.

\begin{proof}
Take a common local bounded smooth decreasing approximation
of the tuple. Lemma~\ref{lem:r11-bounded-first-order}
gives local coefficient total-variation convergence of each
first derivative. The bounded continuous multipliers
$F_i(\mathbf v_l)$ converge pointwise. Dominated convergence
against the fixed limiting derivative measures, together
with the variation errors, passes the smooth chain formula.
Bounded convergence against the coefficients of $Q$
identifies its left side. This proves (FO1).

Choose smooth convex nondecreasing hinge approximations
$H_\epsilon$, with uniform error at most $\epsilon$,
$0\leq H_\epsilon'\leq1$, and
$H_\epsilon'(s)\to\mathbf1_{\{s>t\}}$ at every $s$,
including value zero at $t$. Formula (FO1) and dominated
convergence against the fixed derivative measure prove (FO2).

For (FO3), take one common local smooth approximation $v_l$
independent of the backgrounds. Closedness gives
$d'(v_lQ)=d'v_l\wedge Q$, whose restriction on $A$
equals that for $Q'$ in every actual refinement containing
the approximation's coordinates. Pass this equality through
coefficient total-variation convergence.
For consistency on $E_j$, first change $T_{\alpha,l}$
to $T_{\alpha,j}$ by bounded maximum locality and (FO3),
holding $u_{i,l}$ fixed. Then apply (FO2) to
$u_{i,j}=\max(u_{i,l},-j)$. The double-prime proofs are identical.
\end{proof}

\begin{lemma}[A first-order capacity bound]
\label{lem:r11-first-order-capacity-bound}
Let $-L\leq b_i\leq0$ be bounded classical psh functions,
and let $Q_\alpha$ be their bounded mixed product over $S$,
$|\alpha|=k\leq q-1$. For every original compact
complementary first-order test $\eta$ and Borel $A\subset V$,
\[
 \bigl|\mu_{d'(b_iQ_\alpha),\eta}\bigr|(A)
                  \leq C_\eta L^{k+1}C_S(A;V).
 \tag{FO4}
\]
There is the analogous double-prime estimate. The constant
is independent of the bounded potentials and $L>0$.
\end{lemma}

\noindent\textbf{Status: \texttt{NEW\_PROVED}}.

\begin{proof}
Put $d=q-k$ and choose a compact cutoff and fixed original
moment control $\beta$ for $\eta$. The restricted (NP1),
after taking a supremum over Borel multipliers of absolute
value at most one, gives
\[
 \bigl|\mu_{d'(b_iQ_\alpha),\eta}\bigr|(A)^2
 \leq C_\eta
       \left(\int_A\chi\,\Gamma_{Q_\alpha}(b_i)
                                  \wedge\beta^{d-1}\right)
       \left(\int_A\chi\,\beta^d\wedge Q_\alpha\right).
 \tag{FO5}
\]
The factors are bounded by $CL^{k+2}C_S(A;V)$ from
(PG8) and by $C'L^kC_S(A;V)$ from (DG3).
Their square root proves (FO4). The coefficient controls
remain valid on every actual refinement.
\end{proof}

\begin{theorem}[Intrinsic unbounded first-order coefficient currents]
\label{thm:r11-unbounded-first-order}
For $\sigma>0$, $|\alpha|=k\leq q-1$, and every index $i$,
the intrinsic distributional currents
\[
 {\cal J}'_{\alpha,i}=d'(u_iN_\alpha),\qquad
 {\cal J}''_{\alpha,i}=d''(u_iN_\alpha)
 \tag{FO6}
\]
are coefficient-Radon. They have zero coefficient variation
on the pole union, and
\[
 d'(u_{i,j}T_{\alpha,j})\longrightarrow{\cal J}'_{\alpha,i},
 \qquad
 d''(u_{i,j}T_{\alpha,j})\longrightarrow{\cal J}''_{\alpha,i}
 \tag{FO7}
\]
in local coefficient total variation.
For every original compact complementary test $\eta$,
$a\geq0$ with $a+k+1\leq\kappa$, and $R\geq1$,
\[
 \begin{split}
 \sup_j\int_{\{w\leq-R\}}W_j^a
       \bigl|d\mu_{d'(u_{i,j}T_{\alpha,j}),\eta}\bigr|
     &\leq C_\eta R^{a+k+1-\kappa}{\cal J}_w(R/4),\\
 \int_{\{w\leq-R\}}W^a
       \bigl|d\mu_{{\cal J}'_{\alpha,i},\eta}\bigr|
     &\leq C_\eta R^{a+k+1-\kappa}{\cal J}_w(R/4).
 \end{split}
 \tag{FO8}
\]
Full local weighted variations are finite, the double-prime
bounds are identical, and a continuous multiplier of growth
at most $W^a$ may be inserted in (FO7).
These assertions apply even when $k+2>q+\sigma$, where
the graded theorem constructs no quadratic energy current.
\end{theorem}

\noindent\textbf{Status: \texttt{NEW\_PROVED}}.
These are first-order coefficient measures, not closed
positive backgrounds for a subsequent BT iteration.

\begin{proof}
On $A_T=\{-2T<w\leq-T\}$ further truncate every current
slot at $-2T$. Lemma
\ref{lem:r11-first-order-borel-locality} permits this in
both the derivative slot and its closed background.
Apply (FO4) with $L=2T$. Since $W_j\leq1+2T$ on the shell,
\[
 \int_{A_T}W_j^a
       \bigl|d\mu_{d'(u_{i,j}T_{\alpha,j}),\eta}\bigr|
                    \leq C T^{a+k+1}F_w(T/2).
 \tag{FO9}
\]
Every finite-stage derivative is zero on the pole union
by (FO4). Summing (FO9) with (PG6) proves the first (FO8).

The cut coefficient measures are consistent on $E_j$.
Their variations increase and are locally bounded by
the preceding estimates. Hence they define Radon limits
on the finite-value locus, extended by zero on the
pole union. The second (FO8) holds for these limits.
Agreement on $E_j$ and the two tails at $R=j$ prove
coefficient total-variation convergence of the uncut
derivatives and all the stated weighted versions.
Finite linear coefficient identities and actual
refinement identities pass to these limits.

Equation (PG4) with $F(x)=x_i$ gives
$u_{i,j}T_{\alpha,j}\to u_iN_\alpha$ on every original
test. Distributional continuity identifies the limits
with (FO6), including their current continuity on
each original compact stage. The double-prime proof
is identical.
\end{proof}

\begin{theorem}[Unbounded first-order chain and conormal identities]
\label{thm:r11-unbounded-first-order-chain}
Let $F$ be $C^1$ near the nonpositive orthant and suppose
that for $1\leq r\leq\kappa-k$,
\[
 |F(x)|\leq A W(x)^r,\qquad
 |F_i(x)|\leq A W(x)^{r-1}.
 \tag{FO10}
\]
Then, on all original complementary tests,
\[
 d'(F(\mathbf u)N_\alpha)
                  =\sum_iF_i(\mathbf u){\cal J}'_{\alpha,i},
 \qquad
 d''(F(\mathbf u)N_\alpha)
                  =\sum_iF_i(\mathbf u){\cal J}''_{\alpha,i}.
 \tag{FO11}
\]
If $F(\mathbf u)$ equals a constant on the actual
finite-value locus of an ordinary germ, then
\[
       \sum_iF_i(\mathbf u){\cal J}'_{\alpha,i}
       =\sum_iF_i(\mathbf u){\cal J}''_{\alpha,i}=0
 \tag{FO12}
\]
on that germ and every actual moment refinement.
\end{theorem}

\noindent\textbf{Status: \texttt{NEW\_PROVED}} for
(FO11); (FO12) is a \texttt{DIRECT\_COROLLARY}.
The conormal equation uses an actual $C^1$ value
relation, not an analytic delta-form relation quotient.

\begin{proof}
Apply (FO1) to common bounded truncations. The multiplier
on the left has total exponent $r+k\leq\kappa$, so
(PG4) passes it to its coefficient-defined limit.
On the right the weight has growth $W^{r-1}$ and the
derivative budget is $k+1$, giving the same exponent.
Thus (FO8) passes each term in coefficient variation;
distributional continuity gives (FO11).
If the actual value relation is $F(\mathbf u)=c$,
then $F(\mathbf u)N_\alpha=cN_\alpha$, because
$N_\alpha$ is zero on the omitted poles.
Separate closedness now gives (FO12).
\end{proof}

\paragraph{Relation to the analytic relation quotient.}
This supplies actual coefficient derivatives and conormal
identities for every original test and refinement with
arbitrary original $S$. It does not supply an arbitrary
wedge of these measures, closedness of level restrictions,
or equality on additional GK delta tests. The complete
signed PS/PL quotient and the low-real-degree constant-tensor
order gate remain \texttt{OPEN}.

% END REVIEWED MODULE research_r11/energy/first_order_capacity_currents.tex


% BEGIN REVIEWED MODULE research_r11/energy/graded_capacity_stability.tex
\section{Capacity stability of the graded current and energy calculus}
\label{sec:r11-graded-capacity-stability}

Use the domain, background and power-energy notation of
Sections~\ref{sec:r11-power-all-degree} and
\ref{sec:r11-first-order-power-currents}.
The limits below are ordinary original-current limits,
including every compact smooth complementary test and
every actual moment refinement. No assertion about
additional GK delta tests is implicit.

\begin{lemma}[Capacity-controlled coefficient multipliers]
\label{lem:r11-general-coefficient-capacity-multiplier}
On a fixed ordinary compact test stage, suppose
coefficient-Radon currents $P_l$ converge to $P$ on
original smooth tests. Suppose that for each fixed
complementary test $\eta$, all its coefficient variations,
including that of $P$, are bounded by $C_\eta c(A)$ on
Borel sets $A$, where $c$ is the same specified local
relative capacity. Let $a_l,a$ be actual bounded Borel
weights, with one common bound, such that $a_l\to a$ in
that local capacity and $a$ is quasicontinuous for its
ordinary-open outer version. Then
\[
                         a_lP_l\longrightarrow aP
 \tag{GS1}
\]
on the original compact tests of that stage.
Positivity and closedness of the $P_l$ are not required.
\end{lemma}

\noindent\textbf{Status: \texttt{DIRECT\_COROLLARY}} of
the frozen capacity multiplier argument, stated for
coefficient variation rather than only positive products.

\begin{proof}
For a fixed $\eta$, split the changing-weight error over
$\{|a_l-a|>\delta\}$ and its complement. Its variation
is at most
$2M C_\eta c(\{|a_l-a|>\delta\})$ plus $\delta$ times
the uniformly bounded compact coefficient mass.
Next remove an ordinary open set $G$ of arbitrarily small
capacity on whose complement $a$ is continuous.
On the compact remainder extend this bounded continuous
function to a bounded continuous function on the compact
stage. The latter extension is available in the original
metrizable affinoid neighborhood. Both errors over $G$
are bounded by a fixed multiple of $c(G)$, for the
approximating and limiting coefficient measures separately.
Original smooth scalar functions are uniformly dense in
continuous functions on the compact stage, by the frozen
original separation and cutoff argument. Their test
convergence and the uniform coefficient mass bound give
convergence with this continuous extension. Let the
capacity error and then $\delta$ tend to zero.
The assumed bound for the limiting current has been used
explicitly; it is not inferred by restricting a weak
limit to an arbitrary Borel set.
\end{proof}

\begin{lemma}[Bounded weighted first-order and energy stability]
\label{lem:r11-bounded-first-energy-capacity-stability}
Let bounded classical tuples converge in the exact
localized capacity sense of (PE16), with a common actual
bound. Over the fixed arbitrary original background,
their bounded mixed currents, first-order currents
$d'(v_iQ_\alpha)$ and $d''(v_iQ_\alpha)$, and energy
currents $\Gamma_{Q_\alpha}(v_i,v_a)$ converge on all
original compact tests. Each conclusion remains valid
after multiplication by a continuous function of the
bounded tuple.
\end{lemma}

\noindent\textbf{Status: \texttt{NEW\_PROVED}} for the
first-order and energy extension; the bounded mixed
current assertion is the frozen input already identified
in Lemma~\ref{lem:r11-bounded-analytic-capacity-stability}.

\begin{proof}
The bounded mixed stability theorem first gives
$Q_{\alpha,l}\to Q_\alpha$ and all continuous tuple
weights needed here. For clarity, a continuous function
on the common compact range box is uniformly continuous,
so its changing-weight gap is contained in a finite
union of potential gap sets. It is quasicontinuous off
a union of the finitely many bounded-potential exceptions.
Lemma~\ref{lem:r11-general-coefficient-capacity-multiplier}
therefore proves its weighted convergence, using the
bounded mixed capacity domination, without restricting
the scalar composition to one variable.

In particular $v_{i,l}Q_{\alpha,l}\to v_iQ_\alpha$.
Differentiating distributionally proves unweighted
first-order convergence. Estimate (FO4), applied to
the common bound on each local stage, controls the
coefficient variations of these derivatives and their
limit by the required local capacity.
Apply the preceding multiplier lemma to prove the
weighted first-order assertion.

For energies use the identity
\[
 \Gamma_{Q_\alpha}(v_i,v_a)
 =\tfrac12D(v_iv_aQ_\alpha)
       -\tfrac12v_iD(v_aQ_\alpha)
       -\tfrac12v_aD(v_iQ_\alpha).
 \tag{GS2}
\]
The weighted bounded mixed theorem passes all three
terms in distributions, hence gives unweighted current
convergence of the energies. Estimate (PG8), and
polarization of the positive energy matrix, bound
their coefficient variations by the same local
capacity. The multiplier lemma proves the last
weighted assertion. Derivatives of nonsmooth weights
have not been used anywhere in this proof.
\end{proof}

\begin{theorem}[Nonmonotone graded stability]
\label{thm:r11-nonmonotone-graded-stability}
Let $v_{i,l}\in\mathcal E^\sigma_{S,\partial}(V)$,
$1\leq i\leq m$, have a uniform bound on their
$\sigma$-power energies. Let $u_i$ be actual negative
extended classical psh functions satisfying
$\mathrm{BG}(u_i,0)$, and suppose
$v_{i,l}\to u_i$ in (PE16) in every slot.
Then all $u_i$ belong to the same power class.

Let $F$ be continuous on the nonpositive orthant with
$|F(x)|\leq A(1-\sum_i x_i)^a$, $a\geq0$.
With $k=|\alpha|$, the following convergences hold
under their respective strict inequalities:
\[
\begin{array}{ll}
 F(\mathbf v_l)N_\alpha(\mathbf v_l)
       \longrightarrow F(\mathbf u)N_\alpha(\mathbf u),
       &a+k<q+\sigma,\quad k\leq q,\\[2pt]
 F(\mathbf v_l){\cal J}'_{\alpha,i}(\mathbf v_l)
       \longrightarrow F(\mathbf u){\cal J}'_{\alpha,i}(\mathbf u),
       &a+k+1<q+\sigma,\quad k\leq q-1,\\[2pt]
 F(\mathbf v_l)G^\alpha_{ih}(\mathbf v_l)
       \longrightarrow F(\mathbf u)G^\alpha_{ih}(\mathbf u),
       &a+k+2<q+\sigma,\quad k\leq q-1.
\end{array}
 \tag{GS3}
\]
There is the corresponding double-prime statement.
Only terms whose indices meet the printed inequalities
are asserted.

Equality in any of the three exponent conditions is
also permitted if the following additional uniform
tail condition holds for $w_l=\sum_i v_{i,l}$:
\[
 \lim_{R\to\infty}\sup_l
       \int_R^\infty t^{q+\sigma-1}
              C_S(\{w_l<-t\};V)\,dt=0.
 \tag{GS4}
\]
This endpoint condition is not inferred from a
uniform bound on the full power energies alone.
\end{theorem}

\noindent\textbf{Status: \texttt{NEW\_PROVED}}.
The topology is the exact all-stage capacity topology
(PE16). Under the standing one-affinoid condition (DG1),
Corollary~\ref{cor:r11-affinoid-local-capacity-topology}
now proves its equivalence with fixed-domain capacity
convergence on compact cores. This is the explicit analytic
tester-extension theorem, not an assumed domain identification.

\begin{proof}
Apply Theorem~\ref{thm:r11-power-capacity-stability} to
each slot, repeated in the $q$ differentiated positions,
to obtain membership of the limits before constructing
their unbounded currents. The energy cone estimate
(PE10) bounds $e_\sigma(w_l)$ uniformly; (PE4) therefore
bounds the full integrals $ {\cal J}_{w_l}(0)$ uniformly.

For each fixed truncation level $j$, the preceding
bounded stability lemma applies to every original
compact test, including the bounded tuple weight.
Equations (PG3), (FO8), and (PG12) bound the errors
from removing the truncations, respectively, by
constants times
\[
 j^{a+k-q-\sigma},\qquad
 j^{a+k+1-q-\sigma},\qquad
 j^{a+k+2-q-\sigma},
 \tag{GS5}
\]
uniformly in $l$. Their additional factor
$ {\cal J}_{w_l}(j/4)$ may be bounded by its full
integral. The finitely many compact core controls
are uniform by bounded CLN at level one.
The corresponding limiting errors tend to zero
by the fixed-tuple estimates already proved.
First pass $l\to\infty$ at fixed $j$, then let
$j\to\infty$. This proves (GS3) whenever the
displayed exponent is negative.

At exponent zero retain the tail factor instead
of bounding it by the full integral. Assumption
(GS4) then supplies the same two-limit argument.
This proof does not pass a moving Borel restriction
through weak convergence or presume equality of
raw and outer capacities.
\end{proof}

\begin{corollary}[Full endpoint continuity for decreasing power families]
\label{cor:r11-decreasing-graded-endpoints}
If $v_{i,l}\downarrow u_i\in
\mathcal E^\sigma_{S,\partial}(V)$ are actual negative
extended classical psh functions, then (GS3) holds
with weak inequalities at every endpoint.
In particular, for $\sigma=1$, the top energy
currents $G^\alpha_{ih}$ with $|\alpha|=q-1$
are jointly decreasing-continuous on all original tests.
\end{corollary}

\noindent\textbf{Status: \texttt{DIRECT\_COROLLARY}}
of the proved graded estimates and bounded local
decreasing capacity convergence.

\begin{proof}
Upward stability and (PE9) put all rows in the
same power class with uniformly bounded energies.
The analytic common-collar theorem gives (PE16)
for each fixed bounded truncation. Since
$w_l\geq w=\sum_i u_i$, their capacity tails
are bounded pointwise by the tail of $w$.
Thus (GS4) follows from the finite integral
$J_\sigma(w)$. Apply the endpoint assertion.
\end{proof}

\paragraph{Why an endpoint needs a hypothesis.}
The bounded shell family in
Proposition~\ref{prop:r11-power-energy-continuity-failure}
has uniformly unit $\sigma$-power energy and converges
in (PE16) to zero, while its critical weighted top
measures converge to a nonzero mass at the rigid
origin. This refutes removal of the strict exponent
or uniform-tail qualification in the top weighted
case of (GS3). It does not contradict decreasing
endpoint continuity with an actual power-class limit.

% END REVIEWED MODULE research_r11/energy/graded_capacity_stability.tex


% BEGIN REVIEWED MODULE research_r11/capacity/finite_mass_tester_core.tex
\section{A finite-mass boundary core for the full raw capacity}
\label{sec:r11-finite-mass-tester-core}

This section keeps the arbitrary original background and strengthens
only the geometry of the fixed tester domain. Let $B$ be one compact
original $K$-affinoid domain in the allowed analytic space, and choose
finitely many $f_a,g_b\in\mathcal O(B)$ and positive real thresholds
$R_a,r_b$. Assume that the closed Laurent box
\[
 C=\{|f_a|\leq R_a\ \hbox{for all }a,\quad
                       |g_b|\geq r_b\ \hbox{for all }b\}
                 \Subset\operatorname{Int}_X(B)
 \tag{TC1}
\]
is compact, and put
\[
 V=\{|f_a|<R_a\ \hbox{for all }a,\quad
                       |g_b|>r_b\ \hbox{for all }b\}.
 \tag{TC2}
\]
We assume $V$ is nonempty. Its ordinary closure lies in $C$.
The original symmetric middle-positive separately closed
$(p,p)$-current $S$ is defined on an ordinary neighborhood of $C$;
the background on $V$ is its actual restriction. Set $q=n-p\geq1$.
No change of the ground field, absolute value, test category,
or current class is made.

The bounded classical collar construction is frozen in its local
form; the new assertion below is that the \emph{entire} raw normalized
capacity can be computed with bounded finite-mass zero-boundary
testers, on this fixed domain. It is not a capacity outer-regularity
statement.

\begin{lemma}[An explicit bounded boundary barrier]
\label{lem:r11-strict-Laurent-barrier}
On a neighborhood of $C$ the function
\[
 \rho=\max\left(-1,\,
       \max_a\log\frac{|f_a|}{R_a},\,
       \max_b\log\frac{r_b}{|g_b|}\right)
 \tag{TC3}
\]
is bounded, continuous and classical psh. Terms with a zero numerator
are interpreted as $-\infty$ and are inactive near their zero set.
On $V$ one has $-1\leq\rho<0$; on the ordinary boundary
$\partial_XV$ one has $\rho=0$. Every negative sublevel of $\rho$
has compact closure in $V$, and
\[
                          \int_V(D\rho)^q\wedge S<\infty.
 \tag{TC4}
\]
Empty lists of upper or lower inequalities are simply omitted.
\end{lemma}

\noindent\textbf{Status: \texttt{DIRECT\_COROLLARY}} of the
frozen classical maximum/logarithm calculus and original local CLN.

\begin{proof}
All lower-bound functions are nonzero on $C$, hence on a
neighborhood of that compact set. On an active branch every
logarithm is the logarithmic norm of an analytic unit or its
inverse, plus a constant. After every complete extension and
curve pullback it is harmonic. The finite maximum with the
constant branch is continuous and classical psh; zero upper
functions are inactive on a whole neighborhood of their zero
loci. Compactness gives a finite bound near $C$.

Every point satisfying all inequalities strictly is an ordinary
interior point of $V$, since $C$ lies in the interior of $B$.
Thus at a boundary point at least one inequality is an equality,
and (TC3) is zero. The closed sublevel $\{\rho\leq-\epsilon\}$
inside $C$ satisfies all defining inequalities with a strict
uniform margin, so it is compact and contained in $V$.
This proves the compact-sublevel assertion.
Finally choose a compact original smooth cutoff equal to one
near $C$ in the neighborhood where both $S$ and $\rho$ are
defined. Relative CLN bounds its pairing with $(D\rho)^q\wedge S$.
This bounds the total mass on $V$ and proves (TC4).
\end{proof}

\begin{definition}[Normalized finite-mass boundary testers]
\label{def:r11-finite-mass-tester-core}
Let $\mathscr T_0(S,V)$ consist of the bounded classical psh
functions $h$ on $V$ such that
\[
 -1\leq h\leq0,\qquad \mathrm{BG}(h,0),\qquad
                         \mu_h(V)<\infty,
 \tag{TC5}
\]
where $\mu_h=T_{S,q}(h,\ldots,h)$.
Every such tester belongs to
$\mathcal E^\sigma_{S,\partial}(V)$ for every $\sigma>0$:
it has no truncation loss, and
$\int(-h)^\sigma\,d\mu_h\leq\mu_h(V)$.
The barrier $\rho$ belongs to this class.
\end{definition}

\begin{theorem}[The boundary core computes every Borel raw capacity]
\label{thm:r11-finite-mass-capacity-core}
Under (TC1)--(TC2), for every Borel $E\subset V$,
\[
 \boxed{\quad
 C_S(E;V)=\sup_{h\in\mathscr T_0(S,V)}\mu_h(E).
 \quad}
 \tag{TC6}
\]
Infinite values are permitted. No attainment by a single tester
is asserted. The same equality of the two outer regularizations
follows for arbitrary subsets by taking ordinary open covers;
it does not identify raw and outer values on a Borel set.
\end{theorem}

\noindent\textbf{Status: \texttt{NEW\_PROVED}} for the complete
normalized tester class and arbitrary original $S$ on the stated
strict analytic boxes.

\begin{proof}
Every $h\in\mathscr T_0$ gives a normalized original tester
$h+1\in[0,1]$ with the same top product. This proves the
inequality from right to left.

For the reverse inequality fix a normalized bounded classical
tester $0\leq v\leq1$ on $V$, a compact $K\subset V$, and
$\epsilon>0$. Continuity and negativity of $\rho$ give an
ordinary neighborhood of $K$ on which $\rho\leq-c_K<0$.
Choose $A>0$ so large that $A\rho<-1-\epsilon$ on that
neighborhood, and set
\[
             b=\max(v-1-\epsilon,A\rho),\qquad
             h=\frac{b}{1+\epsilon}.
 \tag{TC7}
\]
The function $h$ is bounded classical, lies in $[-1,0)$, and
equals $(v-1-\epsilon)/(1+\epsilon)$ on a neighborhood of $K$.
Its strict negative sublevels are compact because they are
contained in strict negative sublevels of $\rho$.
Thus $\mathrm{BG}(h,0)$ holds.

The branch $v-1-\epsilon$ is at most $-\epsilon$ on all of $V$.
On the ordinary boundary collar
$\{\rho>-\epsilon/(2A)\}\cap V$, the other branch is larger,
so $b=A\rho$ there. Define $b=A\rho$ across the boundary on
that collar outside $V$. The two definitions agree on their
ordinary overlap and define one bounded classical psh function
on an ordinary neighborhood of $\overline V$.
The background is already defined there by hypothesis.
A compact cutoff and relative CLN consequently give
$\mu_h(V)<\infty$. Therefore $h\in\mathscr T_0(S,V)$.

Ordinary locality and the bounded scaling identity now give
the exact original measure equality
\[
                \mathbf1_K\mu_h
                    =(1+\epsilon)^{-q}\mathbf1_K\mu_v.
 \tag{TC8}
\]
For a given Borel $E$, apply this to compact $K\subset E$.
Since $V$ lies in one compact affinoid over the original finite
$p$-adic field, it is a locally compact second-countable space;
the positive Radon measure $\mu_v$ is inner regular on its
Borel sets. Hence
\[
 \sup_{h\in\mathscr T_0}\mu_h(E)
       \geq (1+\epsilon)^{-q}
                              \sup_{\substack{K\subset E\\K\text{ compact}}}\mu_v(K)
       =(1+\epsilon)^{-q}\mu_v(E).
\]
Let $\epsilon\downarrow0$, including the case of infinite mass,
and then take the supremum over the original normalized $v$.
This proves (TC6). Finally applying the identical equality to
each ordinary open neighborhood gives the statement about the
two outer regularizations. At no point are an infimum over
neighborhoods and a supremum over testers interchanged.
\end{proof}

\begin{corollary}[The raw null ideal is detected by every positive power class]
\label{cor:r11-power-class-null-ideal}
Fix any $\sigma>0$ and a Borel set $E\subset V$.
The following conditions are equivalent:
\begin{enumerate}
\item $C_S(E;V)=0$.
\item $\mu_h(E)=0$ for every $h\in\mathscr T_0(S,V)$.
\item Every mixed top measure
$M(u_1,\ldots,u_q)$, with
$u_i\in\mathcal E^\sigma_{S,\partial}(V)$, gives $E$ zero mass.
\item For every finite family in
$\mathcal E^\sigma_{S,\partial}(V)$, every multi-index
$0\leq|\alpha|\leq q$, and every original compact smooth
complementary test $\eta$,
\[
                              |\mu_{N_\alpha,\eta}|(E)=0.
 \tag{TC9}
\]
\end{enumerate}
Condition (4) includes $\alpha=0$, hence all original
coefficients of the background itself.
Whenever the graded energy currents of (PG10) are defined,
their original coefficient variations also vanish on such an $E$.
\end{corollary}

\noindent\textbf{Status: \texttt{NEW\_PROVED}} as an exact
relative null-ideal identification on the printed fixed domain.

\begin{proof}
Equivalence of (1) and (2) is (TC6).
Every normalized core tester belongs to the fixed power class,
so (3) implies (2), and (4) implies (3) by testing the top
measure on compact original scalar cutoffs.

Assume (1). The one-affinoid coefficient control (DG3) applies
on this domain. Each bounded canonical mixed truncation has
zero variation on $E$ for every original complementary test.
The unbounded currents are increasing limits of their consistent
positive cut restrictions; their positive trace measures therefore
also give $E$ zero mass. Finite coefficient domination gives
(4). This is a limit of increasing positive cut measures, not
a Borel restriction of an uncontrolled weak limit.
For $\alpha=0$ the conclusion follows directly from (DG3).
The same argument with the bounded energy estimate (PG8) and
the increasing positive cut energy construction proves the
last assertion, followed by polarization for cross energies.
These implications prove all stated equivalences.
\end{proof}

\begin{corollary}[The power classes detect a nonzero original background]
\label{cor:r11-nonvacuous-power-testers}
If $S\ne0$ on $V$, then for every $\sigma>0$ there exists
$h\in\mathscr T_0(S,V)\subset\mathcal E^\sigma_{S,\partial}(V)$
with $\mu_h\ne0$.
\end{corollary}

\noindent\textbf{Status: \texttt{DIRECT\_COROLLARY}}.

\begin{proof}
A nonzero positive original current has a nonzero positive
coefficient trace on some compact stage. The zero-degree
case of (DG3) makes the raw capacity of some compact subset
of that stage strictly positive. Equation (TC6) then gives
a core tester with strictly positive mass on that compact.
This proves the assertion without imposing nondegeneracy
or a geometric representation on $S$.
\end{proof}

\paragraph{Exact scope and dependency audit.}
The older R2 relative capacitary envelope theorem constructs a
bounded common-collar least upper bound modulo relative capacity.
R4's close-tester pasting theorem combines close testers on separated
compact sets with an explicit cutoff cost. Neither frozen result is
being rewritten here. The new point is (TC8) with a loss tending to
one, a finite total-mass boundary extension for the same arbitrary
background, and the resulting exact full-tester equality (TC6).

The neighborhood extension of $S$ is used for finite total mass of
the pasted testers. It is not an assertion that an arbitrary current
on an arbitrary ordinary domain automatically extends across its
boundary. The raw null ideal in (TC9) is not identified with an
ambient pluripolar or negligible ideal, and no unrestricted Borel
raw/outer equality is inferred. Removing the extension or strict-box
hypotheses requires another proof; it is not part of this theorem.


% END REVIEWED MODULE research_r11/capacity/finite_mass_tester_core.tex


% BEGIN REVIEWED MODULE research_r11/comparison/barrier_energy_approximation.tex
\section{Critical energy continuity under a lower barrier and finite-mass approximation}
\label{sec:r11-barrier-energy-approximation}

Retain the arbitrary original background and ordinary domain of
Section~\ref{sec:r11-power-energy-stability}. Fix $s>0$ and
write $\kappa=q+s$. The first two results need no extra
geometric condition on $V$. The finite-mass approximation theorem
then uses the explicit strict Laurent-box domain and background
extension of (TC1)--(TC2).

\begin{lemma}[Small-value and large-value energy tails]
\label{lem:r11-two-sided-energy-tails}
For $v\in\mathcal E^s_{S,\partial}(V)$, let
$F_v(t)=\operatorname{Cap}_S(\{v<-t\};V)$ and put
$K_{q,s}=s2^q/(1-2^{-s})$. For every $\varepsilon,R>0$,
\[
 \begin{split}
 \int_{\{0<-v\leq\varepsilon\}}(-v)^s\,d\mu_v
     &\leq K_{q,s}\int_0^\varepsilon t^{\kappa-1}F_v(t)\,dt,\\
 \int_{\{-v\geq R\}}(-v)^s\,d\mu_v
     &\leq K_{q,s}\int_{R/2}^\infty t^{\kappa-1}F_v(t)\,dt.
 \end{split}
 \tag{BA1}
\]
The pole contributes zero to the second integral. If $a\leq v\leq0$
for one $a\in\mathcal E^s_{S,\partial}(V)$, the right sides
can be replaced by the same integrals with $F_a$. These bounds
then tend to zero uniformly in all such $v$ as
$\varepsilon\downarrow0$ or $R\to\infty$.
\end{lemma}

\noindent\textbf{Status: \texttt{DIRECT\_COROLLARY}} of the
annular estimate (PE5), with the parameter ranges retained.

\begin{proof}
For a fixed value $b=-v>0$, integration over $b/2<t<b$ gives
$(1-2^{-s})b^s/s$. If $b\leq\varepsilon$, this whole
parameter interval lies in $(0,\varepsilon)$; if $b\geq R$,
it lies in $(R/2,\infty)$. Tonelli and (PE5) give the two
inequalities, respectively. Endpoint values cause no change,
because the interval endpoints are null for Lebesgue measure
in $t$. Finally $a\leq v$ implies $F_v\leq F_a$, and the
integrability of $t^{\kappa-1}F_a(t)$ at both ends proves
the uniform vanishing.
\end{proof}

\begin{theorem}[Global critical mixed-energy continuity above fixed lower barriers]
\label{thm:r11-critical-energy-barrier-continuity}
For $0\leq i\leq q$, fix
$a_i\in\mathcal E^s_{S,\partial}(V)$. Suppose that
$v_{i,l},u_i$ are negative extended classical psh functions with
\[
                    a_i\leq v_{i,l}\leq0,
                    \qquad a_i\leq u_i\leq0,
 \tag{BA2}
\]
and $v_{i,l}\to u_i$ in the full truncated local original
capacity (PE16). They automatically belong to the same power
class. The finite positive weighted measures
\[
 \nu_l=(-v_{0,l})^s M(v_{1,l},\ldots,v_{q,l}),\qquad
 \nu=(-u_0)^s M(u_1,\ldots,u_q)
\]
satisfy
\[
                         \int_V\phi\,d\nu_l
                              \longrightarrow\int_V\phi\,d\nu
 \tag{BA3}
\]
for every bounded original smooth scalar function $\phi$ on $V$,
including those without compact support. In particular their
total masses converge. Taking the same potential in every slot
gives
\[
                              e_s(v_l)\longrightarrow e_s(u)
 \tag{BA4}
\]
under one common lower barrier in the power class.
Every decreasing sequence $v_l\downarrow u$ with
$u\in\mathcal E^s_{S,\partial}(V)$ satisfies (BA4).
The statements concern the actual prescribed representatives;
no pointwise changes on polar sets are made.
\end{theorem}

\noindent\textbf{Status: \texttt{NEW\_PROVED}}. The common
power-energy lower barrier is stronger than a uniform bound of
the same energies; the shell counterexample (PS6) shows why
that distinction matters.

\begin{proof}
Upward closure and (PE9) give membership and uniform power-energy
bounds. Put
$a=\sum_{i=0}^q a_i$, $w_l=\sum_{i=0}^qv_{i,l}$, and
$w=\sum_{i=0}^q u_i$. Then $a\leq w_l,w\leq0$ and
$a\in\mathcal E^s_{S,\partial}(V)$.
The positive mixed expansion gives
\[
             q!M(v_{1,l},\ldots,v_{q,l})\leq\mu_{w_l},
\]
and the analogous limiting inequality. Hence, on every Borel
carrier, the weight $(-v_{0,l})^s$ can be bounded above by
$(-w_l)^s$ against this larger measure.

For $\varepsilon>0$ choose an original compact smooth
$0\leq\chi_\varepsilon\leq1$ equal to one near
$K_\varepsilon=\overline{\{a<-\varepsilon\}}\Subset V$.
Outside this compact set, $0\leq-w_l\leq\varepsilon$.
The first estimate of (BA1) therefore gives the spatial
tightness bound
\[
 \int_V(1-\chi_\varepsilon)\,d\nu_l,
 \quad\int_V(1-\chi_\varepsilon)\,d\nu
       \leq {K_{q,s}\over q!}
                    \int_0^\varepsilon t^{\kappa-1}F_a(t)\,dt.
 \tag{BA5}
\]
Zero values of $-w_l$ contribute zero because the weight is zero.
For $R>0$, the second estimate of (BA1) likewise gives
\[
 \int_{\{-v_{0,l}>R\}}(-v_{0,l})^s\,
                        dM(v_{1,l},\ldots,v_{q,l})
       \leq {K_{q,s}\over q!}
                    \int_{R/2}^\infty t^{\kappa-1}F_a(t)\,dt,
 \tag{BA6}
\]
with the same bound for the limiting tuple. Both bounds vanish
uniformly in $l$ at their respective ends.

For fixed $\varepsilon,R$, the weight
$\min((-v_{0,l})^s,R^s)$ is a bounded continuous composition
of a truncation of the zeroth potential. The bounded-weight
case of (PE20) therefore gives convergence when paired with
the original compact test $\chi_\varepsilon\phi$.
The difference from the uncut weight is bounded by (BA6)
times $\|\phi\|_\infty$. Let $R\to\infty$.
The difference between this compact pairing and the global
one is bounded by (BA5) times the same norm. Let
$\varepsilon\downarrow0$. This proves (BA3), including
$\phi=1$, and hence (BA4).

Finally, if $v_l\downarrow u$ and $u$ belongs to the power
class, take $a=u$. The bounded truncations decrease to those
of $u$, so frozen bounded decreasing capacity convergence,
with its analytic Laurent-domain extension, gives precisely
(PE16). Thus the preceding argument applies. No monotonicity
of the numerical energies themselves is needed in this proof.
\end{proof}

\begin{theorem}[Bounded finite-mass approximation of the power class]
\label{thm:r11-finite-mass-power-approximation}
Assume the strict analytic Laurent-box hypotheses (TC1)--(TC2),
including extension of the original arbitrary background to an
ordinary neighborhood of the closed box. Let $\rho$ be the
barrier of Lemma~\ref{lem:r11-strict-Laurent-barrier}.
For every $u\in\mathcal E^s_{S,\partial}(V)$, put
\[
                            v_j=\max(u,j\rho),\qquad j\geq1.
 \tag{BA7}
\]
Then $v_j$ are bounded actual classical psh functions,
$v_j\downarrow u$ pointwise, and
\[
 \begin{gathered}
 -j\leq v_j\leq0,\qquad \mathrm{BG}(v_j,0),\qquad
 v_j/j\in\mathscr T_0(S,V),\\
 \mu_{v_j}(V)\leq j^q\mu_\rho(V)<\infty,\qquad
 e_s(v_j)\leq2^{q+s}e_s(u),\qquad
 e_s(v_j)\longrightarrow e_s(u).
 \end{gathered}
 \tag{BA8}
\]
Consequently, for an actual negative extended classical psh
$u$ with $\mathrm{BG}(u,0)$ on this domain, membership in
$\mathcal E^s_{S,\partial}(V)$ is equivalent to the existence
of a decreasing pointwise approximation $b_j\downarrow u$
by bounded negative classical psh functions satisfying
\[
 \mathrm{BG}(b_j,0),\qquad \mu_{b_j}(V)<\infty,\qquad
                         \sup_j e_s(b_j)<\infty.
 \tag{BA9}
\]
For every approximation satisfying (BA9), its energies converge
to $e_s(u)$. The boundary-gap hypothesis on the limiting $u$
is retained; it is not inferred from an unspecified boundary trace.
\end{theorem}

\noindent\textbf{Status: \texttt{NEW\_PROVED}} on the printed
strict-box and background-extension domain. This supplies an
intrinsic approximation characterization, not a general
Monge--Ampere existence theorem.

\begin{proof}
The original barrier satisfies $-1\leq\rho<0$ on $V$.
Thus the maxima in (BA7) are bounded classical and decrease
to $u$ at every actual point, including its poles. Their
negative sublevels are contained in compact negative sublevels
of $j\rho$, proving the boundary-gap condition.

Finite total mass requires a separate argument; no fixed collar
for $v_j$ is assumed. For every $c>1$, strict negativity of
$\rho$ gives
\[
                          cj\rho<j\rho\leq v_j
                                     \quad\hbox{everywhere on }V.
\]
The compact-gap comparison theorem applies to $cj\rho,v_j$:
every positive-margin gap lies in a compact negative sublevel
of $\rho$. Its strict comparison set is all of $V$, so
\[
                 \mu_{v_j}(V)\leq c^q j^q\mu_\rho(V).
\]
Let $c\downarrow1$ and use the finite mass of $\rho$.
This proves the mass estimate and the normalized core membership.
Since $u\leq v_j\leq0$, (PE9) gives the energy bound and
Theorem~\ref{thm:r11-critical-energy-barrier-continuity} gives
its convergence. This proves the forward approximation statement.

Conversely, boundedness and the boundary-gap condition make each
$b_j$ a no-loss potential; finite mass makes its $s$-energy
finite. The bounded truncations of $b_j$ decrease to those
of the actual $u$, so (PE16) holds on every admissible original
moment box. The uniform bound in (BA9), the assumed boundary
condition of $u$, and Theorem~\ref{thm:r11-power-capacity-stability}
put $u$ in $\mathcal E^s_{S,\partial}(V)$. Now $u$ itself
is a common power-energy lower barrier for all $b_j$, so
(BA4) gives convergence of their energies. This proves the
equivalence and the last assertion without any unproved
compactness or envelope-contact step.
\end{proof}

% END REVIEWED MODULE research_r11/comparison/barrier_energy_approximation.tex


% BEGIN REVIEWED MODULE research_r11/comparison/radial_unbounded_energy_and_loss.tex
\section{Unbounded examples and an unavoidable comparison defect in every dimension}
\label{sec:r11-radial-unbounded-energy}

The general theorem of Section~\ref{sec:r11-unbounded-comparison}
is stated for arbitrary original $S$. The examples here verify
that its unbounded class is nonempty in every dimension and that
the loss condition cannot simply be removed. They do not replace
that theorem by a toric-background assertion.

Let $n\geq1$ and let
\[
 V=\{x\in(\mathbf A_K^n)^{\mathrm{an}}:
                         |T_i(x)|<1\ (1\leq i\leq n)\},
 \qquad S=[V].
\]
This is an ordinary open unit polydisc in the original smooth
analytic affine space. Put
\[
 t(x)=-\log\max_i|T_i(x)|\in(0,\infty],
 \qquad \sigma(s)=\text{the diagonal radius-$e^{-s}$ Gauss point}.
 \tag{RE1}
\]
The unique point at which $t=\infty$ is the rigid origin $o$.
All real radii in (RE1) refer to seminorms over the original $K$;
no field or valuation normalization is changed.

The precise \texttt{INTERNAL\_INPUT} is the original radial-shell
calibration in R6,
\path{dirichlet_\allowbreak radial_\allowbreak measure_\allowbreak carriers.tex}:
for a bounded continuous radial classical psh profile
\[
 h(x)=-\int_0^{t(x)} A(s)\,ds,
\]
where $A$ is nonnegative and nonincreasing and is zero for all
sufficiently large $s$, the original top product is
\[
       (Dh)^n=\sigma_*\tau,
       \qquad \tau([s,\infty))=A(s)^n
 \tag{RE2}
\]
at continuity points of $A$. Finite step functions follow from
the calibrated Newton-simplex calculation, and the general bounded
profile follows by piecewise-linear approximation and frozen
locally uniform bounded BT continuity. The normalization is one
for a unit-slope shell; the factor $n!$ is already absorbed by
the simplex volume $1/n!$.

\begin{proposition}[A sharp gauge threshold for genuine unbounded potentials]
\label{prop:r11-radial-gauge-threshold}
Fix $0<a<1$ and define
\[
 u_a(x)=-\bigl((1+t(x))^a-1\bigr),\qquad u_a(o)=-\infty.
 \tag{RE3}
\]
This is an actual negative extended classical psh function with
ordinary boundary limit zero. Its non-polar truncation measure
and its discarded mass are
\[
 \begin{split}
 \mu_{u_a}
 &=\sigma_*\!\left(
     n a^n(1-a)(1+s)^{-n(1-a)-1}\,ds\right),\\
 \nu_{u_a,j}
 &=a^n(1+j)^{-n(1-a)/a}\,\delta_{\sigma(R_j)},
       \qquad R_j=(1+j)^{1/a}-1.
 \end{split}
 \tag{RE4}
\]
In particular $u_a\in\mathcal F_{[V],\mathrm{loc}}(V)$, and
indeed on the whole domain
\[
 \|\mu_{u_a,j}-\mu_{u_a}\|_{\mathrm{TV}(V)}
                       =2a^n(1+j)^{-n(1-a)/a}.
 \tag{RE5}
\]

For $\alpha>0$ and $\Phi(s)=(1+s)^\alpha$, the local
gauge criterion (UC6) holds if and only if
\[
                 a<\frac{n}{n+\alpha}.
 \tag{RE6}
\]
In this range it even holds on all of $V$, with
\[
 \sup_j\int_V(1-u_{a,j})^\alpha\,d\mu_{u_a,j}
       =\frac{n a^n(1-a)}{n(1-a)-a\alpha}.
 \tag{RE7}
\]
Thus the no-loss class contains genuinely unbounded potentials
in every dimension; the gauge condition is a useful sufficient
condition but is strictly stronger than absence of loss.
\end{proposition}

\noindent\textbf{Status: \texttt{NEW\_PROVED}} by the explicit
calculation below, using the frozen bounded radial calibration.

\begin{proof}
As a function of $r=\log\max_i|T_i|<0$, the finite profile
$-((1-r)^a-1)$ is increasing and convex. Its truncations are
continuous classical psh: the truncated convex increasing profile
is constant for sufficiently negative $r$, so the composition
extends continuously across the origin and across coordinate
zero loci. Convex piecewise-linear approximation gives classicality
after every complete extension; the decreasing limit has the
printed extended values. Since $t\to0$ at any ordinary boundary
point of the unit polydisc, the profile has boundary value zero.

The truncation $u_{a,j}=\max\{u_a,-j\}$ has radial slope
\[
 A_j(s)=
 \begin{cases}
 a(1+s)^{a-1},&0<s<R_j,\\
 0,&s>R_j.
 \end{cases}
\]
By (RE2), $\mu_{u_a,j}=\sigma_*\tau_j$, where
\[
 \begin{split}
 \tau_j={}&
   \mathbf1_{(0,R_j)}\,
        n a^n(1-a)(1+s)^{-n(1-a)-1}\,ds\\
 &+a^n(1+R_j)^{-n(1-a)}\delta_{R_j}.
 \end{split}
 \tag{RE8}
\]
One may obtain this instance of (RE2) by step-slope approximation
on compact subintervals of $(0,R_j)$, retaining the final downward
jump, and then applying bounded continuous-potential stability.
The small interval next to $s=0$ is outside any prescribed compact
analytic test core after it is made sufficiently small. At the
other endpoint the functions are constant near $o$. Thus this
passage proves the identity on all original compact smooth tests,
including refinements not represented in the coordinates $T_i$.

The finite-value carrier of the truncation is $\{t<R_j\}$.
Consequently the atom at $R_j$ is exactly the discarded part,
while the increasing restrictions retain the displayed diffuse
measure. This proves (RE4). Its diffuse tail beyond $R_j$ has
mass $a^n(1+R_j)^{-n(1-a)}$, the same as the discarded atom.
The atom and diffuse tail are mutually singular. Their difference
is exactly $\mu_{u_a,j}-\mu_{u_a}$, proving (RE5).

Set $b=n(1-a)-a\alpha$. On $s<R_j$ the gauge is
$(1-u_{a,j})^\alpha=(1+s)^{a\alpha}$, while at the atom
it is $(1+j)^\alpha=(1+R_j)^{a\alpha}$. Thus
\[
 \int_V(1-u_{a,j})^\alpha\,d\mu_{u_a,j}
   =n a^n(1-a)\int_0^{R_j}(1+s)^{-b-1}\,ds
          +a^n(1+R_j)^{-b}.
 \tag{RE9}
\]
If $b>0$, the right side increases to the quantity in (RE7).
If $b=0$, its first term is
$n a^n(1-a)\log(1+R_j)$ and diverges. If $b<0$, its
first term diverges as a positive power of $1+R_j$.
The latter two divergences already occur on every compact closed
polydisc of radius $e^{-s_0}<1$, since removing the interval
$(0,s_0)$ changes the integral by a bounded quantity. Hence they
are failures of the local criterion at $o$, not only a failure
of a global estimate near the boundary. This proves (RE6)--(RE7).
\end{proof}

\begin{corollary}[The example satisfies every partial-degree energy test]
\label{cor:r11-radial-all-partial-energy}
Assume $0<a<n/(n+1)$, choose
\[
                    1<s<\frac{n(1-a)}a,
\]
and take the $n$-tuple consisting of $u_a$ in every slot.
For each $0\leq k\leq n$ and every original compact smooth
positive complementary form $\Theta$ of degree $(n-k,n-k)$,
\[
 \sup_{j\geq1}
 \int_V (1-nu_{a,j})^s\,
                    \Theta\wedge(Du_{a,j})^k<\infty.
 \tag{RE9a}
\]
The statement includes every finite actual refinement of the
test, not only tests expressed in the radial coordinate.
Thus this genuinely unbounded family satisfies the all-partial
superlinear energy condition used for the recursive unbounded
current construction in the present round.
\end{corollary}

\noindent\textbf{Status: \texttt{DIRECT\_COROLLARY}} of
Proposition~\ref{prop:r11-radial-gauge-threshold} and the
original-test argument below.

\begin{proof}
Every positive-degree original smooth CLD form vanishes on
some ordinary neighborhood of the rigid $K$-rational origin.
Indeed, choose one finite original moment presentation near
that point. For each analytic unit $f$ in that presentation,
$f(o)\in K^\times$ and the ordinary open neighborhood
\[
                 |f-f(o)|<|f(o)|
\]
has $|f|=|f(o)|$ by the ultrametric inequality. Intersecting
the finitely many such neighborhoods makes the whole moment
map constant, so its pulled-back forms of positive degree
vanish. The argument works for any further original moment
presentation that a test requires.

If $k<n$, apply this observation to $\Theta$. Its compact
support avoids an ordinary neighborhood of $o$, and the
continuous function $\max_i|T_i|$ has a strictly positive
minimum on that support. Therefore $u_a$ is bounded on an
ordinary neighborhood of the support. All sufficiently deep
truncations and their products agree there with $u_a$ and its
bounded local products. The finitely many remaining truncations
also have finite weighted mass. This proves (RE9a) for $k<n$.

For $k=n$, $\Theta$ is a bounded compact scalar test. Since
\[
              1-nu_{a,j}\leq n(1-u_{a,j}),
\]
its weighted integral is bounded by
$\|\Theta\|_\infty n^s$ times (RE7) with $\alpha=s$.
The printed condition on $s$ is exactly its convergence
condition. This proves the remaining degree and all test
quantifiers.
\end{proof}

\begin{proposition}[Unbounded comparison fails without loss control]
\label{prop:r11-logarithmic-comparison-failure}
For every $n\geq1$, on the same original unit polydisc, put
\[
 u(x)=\log\max_i|T_i(x)|=-t(x),\qquad
 v(x)=c\max\{u(x),-1\},\qquad 0<c<1.
 \tag{RE10}
\]
Then $u$ is negative extended classical psh, $v$ is bounded
continuous classical psh, and $\mathrm{BG}(u,v)$ holds.
Nevertheless
\[
 \mu_u=0,\qquad
 \mu_v=c^n\delta_{\sigma(1)},\qquad
 \mu_v(\{u<v\})=c^n>0=\mu_u(\{u<v\}).
 \tag{RE11}
\]
Moreover
\[
          \nu_{u,j}=\delta_{\sigma(j)}.
 \tag{RE12}
\]
Thus the assertion that (UC10) holds for every unbounded
classical pair merely by using the truncation measure (UC3) is
\texttt{FALSE}, already for the ambient background in every
dimension.
\end{proposition}

\noindent\textbf{Status: \texttt{NEW\_PROVED}} counterexample
to the precisely stated unrestricted extension.

\begin{proof}
The bounded shell formula gives
$\mu_{u,j}=\delta_{\sigma(j)}$. Its entire mass lies on
$\{u=-j\}$, so the restriction in (UC3) is zero for each
$j$. This proves $\mu_u=0$ and (RE12). Homogeneity gives
$\mu_v=c^n\delta_{\sigma(1)}$.

The gap function $v-u$ equals $(1-c)t$ for $0<t\leq1$,
and $t-c$ for $t\geq1$. It is positive and strictly
increasing in $t$, and tends to zero as $t\downarrow0$.
Every positive-margin gap therefore has closure contained in a
closed polydisc of radius strictly smaller than one. This proves
$\mathrm{BG}(u,v)$. At $\sigma(1)$, $u=-1<-c=v$,
which proves (RE11). The failure is exactly the persistent
discarded mass in (RE12); it involves no change of field, no
special test restriction and no boundary mismatch.
\end{proof}

\paragraph{Relation to a Dirichlet problem.}
The function $u$ in (RE10) and the zero function have the same
zero truncation-limit Monge--Ampere measure and the same ordinary
boundary limit. Thus uniqueness in the unrestricted unbounded
class is also false. The new no-loss comparison theorem is one
way to select a class on which the relative measure can support
a uniqueness statement modulo its own capacity. In the ambient
setting, any further actual pointwise uniqueness uses the precise
ambient gap-detection result separately; no such conclusion is
asserted for arbitrary $S$ here.

% END REVIEWED MODULE research_r11/comparison/radial_unbounded_energy_and_loss.tex


% BEGIN REVIEWED MODULE research_r11/comparison/power_energy_sharpness.tex
\section{Sharpness of the power-energy thresholds and capacity exponent}
\label{sec:r11-power-energy-sharpness}

These examples use $V$ equal to the original ordinary open unit
polydisc in $(\mathbf A_K^q)^{\mathrm{an}}$ and $S=[V]$.
The field and its absolute value remain fixed. They test the
generality of the preceding arbitrary-background theorems;
they are not substituted for those theorems.

\begin{proposition}[Exact power energies of the unbounded radial family]
\label{prop:r11-exact-radial-power-energy}
For $0<a<1$, let $u_a$ be (RE3) in dimension $q$, and set
$\beta=q(1-a)/a$. For every $s>0$,
\[
 u_a\in\mathcal E^s_{[V],\partial}(V)
       \quad\Longleftrightarrow\quad
                 s<\beta
       \quad\Longleftrightarrow\quad a<{q\over q+s}.
 \tag{PS1}
\]
When these conditions hold, its exact energy is
\[
 e_s(u_a)
      =a^q\frac{\Gamma(s+1)\Gamma(\beta-s)}{\Gamma(\beta)}.
 \tag{PS2}
\]
For the canonical truncations one has, without assuming $s<\beta$,
\[
 e_s(u_{a,j})=s a^q\int_0^j y^{s-1}(1+y)^{-\beta}\,dy.
 \tag{PS3}
\]
These increase to (PS2) when finite, diverge logarithmically at
$s=\beta$, and diverge at a positive power rate when $s>\beta$.
Thus the finite-power classes are strictly nested inside the
unbounded no-loss class, in every dimension.
\end{proposition}

\noindent\textbf{Status: \texttt{NEW\_PROVED}} as an exact
energy-class calculation from the new radial measures and the
frozen original shell calibration.

\begin{proof}
Use $y=-u_a=(1+t)^a-1$ in the diffuse density (RE4).
Its pushforward becomes
\[
                   a^q\beta(1+y)^{-\beta-1}\,dy.
\]
The energy is therefore
$a^q\beta\int_0^\infty y^s(1+y)^{-\beta-1}\,dy$.
This beta integral is finite exactly when $s<\beta$ and has
value (PS2). At a truncation level $j$, integrate the same
density on $[0,j)$ and add the shell atom of mass
$a^q(1+j)^{-\beta}$ with weight $j^s$. Integration by parts
cancels the boundary term and gives (PS3). Its integrand behaves
as $y^{s-\beta-1}$ at infinity, proving the endpoint and
supercritical divergence claims. Every $u_a$ already belongs
to the no-loss boundary class by (RE4)--(RE5), so (PS1) follows
from the definition, with no implication from finite non-polar
energy outside that class.
\end{proof}

\begin{proposition}[The capacity exponent is optimal under a power-energy bound]
\label{prop:r11-power-capacity-exponent-sharp}
Fix $s>0$. There is a family of bounded continuous classical
$u_T\in\mathcal E^s_{[V],\partial}(V)$, $T\geq1$, with
$e_s(u_T)=1$, and compact singleton sets $E_T\subset V$, such that
\[
 \operatorname{Cap}_{[V]}(E_T;V)
  =\operatorname{Cap}_{[V]}^*(E_T;V)=T^{-q},\qquad
 \mu_{u_T}(E_T)=T^{-qs/(q+s)}
                =\operatorname{Cap}_{[V]}(E_T;V)^{s/(q+s)}.
 \tag{PS4}
\]
Consequently replacing $s/(q+s)$ in (PE14) by any larger
exponent, with a constant uniform for $e_s\leq1$ on this fixed
domain and background, is \texttt{FALSE}.
\end{proposition}

\noindent\textbf{Status: \texttt{NEW\_PROVED}} sharpness
example. The exact radial compact capacity used below is frozen
\texttt{INTERNAL\_INPUT}: R6,
\path{capacity_\allowbreak laminar_\allowbreak graphs.tex}, equation
\path{eq:r6-arbitrary-\allowbreak radius-\allowbreak capacity},
already proved in R3.

\begin{proof}
Write $t=-\log\max_i|T_i|$, $g_T=-\min(t,T)$, and
$c_T=T^{-s/(q+s)}$. Set
\[
             u_T=c_Tg_T,\qquad E_T=\{\sigma(T)\}.
 \tag{PS5}
\]
The frozen shell calibration gives
$\mu_{u_T}=c_T^q\delta_{\sigma(T)}$.
Since $-u_T(\sigma(T))=c_TT$, its energy is
$c_T^{q+s}T^s=1$. Its ordinary boundary limit is zero, and
it is bounded, so it belongs to the stated energy class.

The frozen exact capacity computation gives raw capacity
$T^{-q}$ for both $E_T$ and the closed radius-$e^{-T}$
polydisc. For clarity, outer capacity is also exact: the ordinary
open polydisc $\{t>T-\varepsilon\}$ contains $E_T$ and has
capacity at most $(T-\varepsilon)^{-q}$ by containment in the
corresponding closed polydisc. Let $\varepsilon\downarrow0$
and use raw capacity as the opposite lower bound. This proves
(PS4), with all radii interpreted over the original field.
If a larger exponent $\gamma$ and one uniform constant were
valid, (PS4) would imply
$T^{q(\gamma-s/(q+s))}\leq C$ for every $T\geq1$, a
contradiction.
\end{proof}

\begin{proposition}[Same-power energy need not be continuous in capacity]
\label{prop:r11-power-energy-continuity-failure}
The family (PS5) converges to the actual classical function zero
in the full truncated local original capacity of (PE16), while
\[
                         e_s(u_T)=1,
                         \qquad e_s(0)=0.
 \tag{PS6}
\]
Thus the lower-semicontinuity assertion (PE18) cannot be upgraded
to energy convergence using only the same uniform power bound.
More strongly, the critical weighted currents satisfy
\[
       (-u_T)^s\mu_{u_T}=\delta_{\sigma(T)}
                  \longrightarrow\delta_o
                  \ne0=(-0)^s\mu_0
 \tag{PS7}
\]
on every original compact smooth scalar test. Thus critical
weighted-current continuity fails even on compact test stages,
not merely through escape of total energy toward the boundary.
The actual pointwise limit of $u_T$ is zero off the rigid origin
and $-\infty$ at that origin; it is not upper semicontinuous.
Consequently capacity convergence also cannot be replaced by
an assertion that this particular pointwise limit is classical
without a separate regularity argument.
\end{proposition}

\noindent\textbf{Status: \texttt{NEW\_PROVED}} counterexamples
to these two stronger assertions, which are \texttt{FALSE}.

\begin{proof}
Away from the origin, $t$ is finite and $u_T=-c_Tt\to0$.
At the origin, $u_T=-c_TT=-T^{q/(q+s)}\to-\infty$.
This proves the pointwise statement and its failure of upper
semicontinuity.

We check the stronger local capacity quantifier, not just the
capacity measured on $V$. Fix an admissible ordinary domain
$\Omega$, a compact $L\Subset\Omega$, and $k,\delta>0$.
If the origin does not belong to $L$, the function $t$ is bounded
on $L$, so the truncated gap set is eventually empty.
Otherwise $\Omega$ contains a smaller original open polydisc
$D=\{t>T_0\}$ about the origin. One may choose
$T_0=N(-\log|\pi|)$ for an integer $N$, where $\pi\in K$ is
an original uniformizer, so this polydisc is identified with the
unit polydisc by the original coordinates $T_i/\pi^N$.
For large $T$, its truncated
gap set is empty if $\delta\geq k$, and otherwise is contained
in
\[
                        \{t>\delta/c_T\}\Subset D.
\]
Domain restriction bounds its $\Omega$-capacity by its
$D$-capacity. The same frozen radial capacity formula on $D$,
with the radial coordinate shifted by $T_0$, bounds that capacity
by $(\delta/c_T-T_0)^{-q}\to0$.
The calculation is independent of the moment presentation of
$\Omega$ or any of its refinements, so verifies (PE16) in full.
Equation (PS6) follows from (PS4), proving the failure of energy
continuity. At the sole atom, the weighted mass is
$(c_TT)^s c_T^q=1$, so the equality in (PS7) is exact.
The original Gauss points $\sigma(T)$ tend to $o$; every original
smooth scalar is locally constant at that rational point, and
therefore evaluates to its value at $o$ for all sufficiently
large $T$. This proves the asserted all-original-current limit
and its disagreement with the weighted current of the zero
potential. It leaves the strict subcritical weighted
continuity in (PE20) intact.
\end{proof}

% END REVIEWED MODULE research_r11/comparison/power_energy_sharpness.tex


% BEGIN REVIEWED MODULE research_r11/comparison/critical_recursive_residue.tex
\section{The critical recursive identity: an exact obstruction and a conditional repair}
\label{sec:r11-critical-recursive-residue}

The distinction between an original distributional derivative and
a non-polar truncation product persists even when the multiplier
is coefficient-integrable. The following example retains every
original test and makes the missing polar term explicit.

\begin{proposition}[A logarithmic polar residue in every dimension]
\label{prop:r11-logarithmic-recursive-residue}
Let $q\geq1$, let $V$ be the original open unit polydisc in
$(\mathbf A_K^q)^{\mathrm{an}}$, let $S=[V]$, and put
\[
                u=\log\max_i|T_i|,\qquad u(o)=-\infty,
 \tag{CR1}
\]
where $o$ is the $K$-rational origin. Then
$\mathrm{BG}(u,0)$ holds. For every $0\leq k<q$, the lower
relative non-polar product $N_k$ of $k$ copies of $u$ exists,
is original positive and separately closed, and is locally finite
on every original coefficient stage. The multiplier $uN_{q-1}$
is coefficient-defined and locally coefficient-integrable.
Nevertheless
\[
                    D(uN_{q-1})=\delta_o,
                    \qquad \mu_u^{\mathrm{np}}=0.
 \tag{CR2}
\]
In fact the uncut bounded top products converge on all original
scalar tests to $\delta_o$, whereas their consistent non-polar
restrictions are zero at every level.
Thus the assertion that local coefficient integrability alone
identifies the recursive derivative with the non-polar top product
is \texttt{FALSE}, even with the compact-sublevel boundary condition.
\end{proposition}

\noindent\textbf{Status: \texttt{NEW\_PROVED}} counterexample,
using only the frozen original unit-slope shell calibration and
bounded recursion. It does not contradict the new no-loss calculus,
because this $u$ has persistent truncation loss.

\begin{proof}
First record an original-test fact at the rational point $o$.
Any original smooth form of positive total degree has zero germ
there. Indeed choose its finite local unit-moment presentation
$f_1,\ldots,f_N$. Each $f_i(o)\in K^\times$, and on an
ordinary neighborhood satisfying
$|f_i-f_i(o)|<|f_i(o)|$ one has $|f_i|=|f_i(o)|$.
Intersect these finitely many neighborhoods. The moment map
is constant there, so the pullback of every positive-degree
superform is zero. This applies to any actual refinement.
In particular every original smooth scalar function is locally
constant at $o$. On a fixed original test stage the neighborhood
can be chosen for its fixed finite moment presentation, so this
observation is compatible with current continuity on that stage.

On $V\setminus\{o\}$ the function $u$ is locally bounded
classical, and all its frozen bounded local products exist.
If $k<q$, every compact complementary test has positive degree
and vanishes on an ordinary neighborhood of $o$. Its support
is therefore compact in the punctured domain, where $u$ is
bounded on a neighborhood of that support. The local bounded
product thus defines a positive coefficient-Radon current
$N_k$ on every original stage of $V$. For sufficiently large
$j$, the original test pairing of $(Du_j)^k$ agrees exactly
with it, by ordinary locality. The same holds for the cut
products, since the support eventually lies in $\{u>-j\}$.
This identifies $N_k$ with the non-polar truncation product.

Every complementary test used for $d'N_k$ or $d''N_k$ still
has positive total degree when $k<q$. It too vanishes near
$o$, and separate closedness reduces to that of the bounded
local product on the punctured domain. Thus $N_k$ is separately
closed. The same compact-support argument with one extra factor
$u$ proves coefficient integrability of $uN_{q-1}$ and gives
\[
                  u_j(Du_j)^{q-1}\longrightarrow uN_{q-1}
 \tag{CR3}
\]
with eventual equality on each fixed original complementary
coefficient test. In particular the convergence holds
distributionally.

The original unit-slope shell formula is
\[
                        (Du_j)^q=\delta_{\sigma(j)},
 \tag{CR4}
\]
where $\sigma(j)$ is the diagonal radius-$e^{-j}$ Gauss point.
Its entire mass is on $\{u=-j\}$, so its restriction to
$\{u>-j\}$ is zero. This proves the second assertion of
(CR2). The Gauss points tend to $o$. More precisely, each
original smooth scalar test is constant on some ordinary
neighborhood of $o$, and therefore has
$\phi(\sigma(j))=\phi(o)$ for all sufficiently large $j$.
Hence (CR4) converges on the full original scalar test family
to $\delta_o$. Apply $D$ to (CR3) and use the frozen bounded
identity $D(u_j(Du_j)^{q-1})=(Du_j)^q$.
This proves the first assertion of (CR2), with the original
normalization and no restricted tropical test family.
\end{proof}

\begin{corollary}[Critical recursion under no loss and explicit integrability]
\label{cor:r11-critical-recursion-if-integrable}
Assume the one-affinoid domain hypothesis (DG1). Let a finite
family $u_1,\ldots,u_m$ belong to
$\mathcal F_{S,\partial}(V)$, let $w=\sum_i u_i$, and let
$N_\alpha$ be the all-degree products of
Theorem~\ref{thm:r11-boundary-all-degree-products}, with
$|\alpha|=q-1$. Fix a family index $i$. Suppose that
$u_i$ is locally coefficient-integrable against $N_\alpha$:
for every original compact positive complementary $(1,1)$
control $\Theta$,
\[
                   \int |u_i|\,\Theta\wedge N_\alpha<\infty.
 \tag{CR5}
\]
Then, on every original compact coefficient stage,
\[
             u_{i,j}T_{\alpha,j}\longrightarrow u_iN_\alpha
                          \quad\hbox{in total variation},
 \qquad D(u_iN_\alpha)=N_{\alpha+e_i}.
 \tag{CR6}
\]
Thus Dini capacity or positive power energy is a sufficient route
to the critical integrability; it is not a necessary hypothesis
for the recursive passage once (CR5) is separately verified.
\end{corollary}

\noindent\textbf{Status: \texttt{CONDITIONAL}} on the explicitly
printed integrability premise. Its automatic validity throughout
the full no-loss boundary cone is not proved here.

\begin{proof}
The global coefficient compensator of (DG3), with $k=q-1$,
gives, for any fixed compact original coefficient test $\eta$,
\[
 \left|\mu_{u_{i,j}T_{\alpha,j},\eta}\right|
                         (\{w\leq-j\})
 \leq C_\eta j^q\operatorname{Cap}_S(\{w<-j/2\};V)
                                             \longrightarrow0.
 \tag{CR7}
\]
The last step is the no-loss capacity-tail criterion for $w$.
The integrability premise and finite original coefficient
domination imply
\[
             \int_{\{w\leq-j\}}|u_i|\,d|\mu_{N_\alpha,\eta}|
                                                   \longrightarrow0.
\]
Indeed the tails decrease to the full family pole union, which
all coefficients of $N_\alpha$ avoid by the all-degree theorem.
On $\{w>-j\}$ strict locality and equality of the actual
bounded slot values identify the two weighted coefficient
measures exactly. The two tail estimates therefore prove
the total-variation assertion in (CR6). Apply $D$ and use
the bounded recursion; the uncut top products converge to
$N_{\alpha+e_i}$ by the no-loss boundary calculus. This proves
the original current identity. The logarithmic example lies
outside that no-loss cone, so (CR2) does not contradict this
conditional repair.
\end{proof}

% END REVIEWED MODULE research_r11/comparison/critical_recursive_residue.tex


% BEGIN REVIEWED MODULE research_r11/capacity/nonlinear_level_sweeping.tex
\section{Nonlinear level sweeping and a scalar truncation-loss residue}
\label{sec:r11-nonlinear-level-sweeping}

Keep the original analytic category, arbitrary allowed background $S$,
and $q=n-p\geq1$. The first results are ordinary-local and require
no affinoid-domain hypothesis. The final Dini identity explicitly
uses (DG1). Every lower-degree equality is tested against all original
compact smooth complementary forms, and every construction respects
all actual moment refinements.

\begin{lemma}[One-sided Stieltjes passage]
\label{lem:r11-one-sided-Stieltjes}
Let $b$ be a bounded nondecreasing left-continuous function on
$\mathbf R$, constant outside a compact interval, and let
$b_\epsilon$ be smooth nondecreasing functions with the same
constant endpoint values such that $b_\epsilon(t)\to b(t)$ at
every real $t$. If $g$ is a compactly supported right-continuous
BV function, then
\[
                  \int g\,db_\epsilon\longrightarrow\int g\,db.
 \tag{NS1}
\]
The measures $db_\epsilon,db$ are distributional derivatives;
in particular the right side includes the actual value $g(t)$
at each atom of $db$.
\end{lemma}

\noindent\textbf{Status: \texttt{DIRECT\_COROLLARY}} of real
measure integration, with the one-sided convention proved here.

\begin{proof}
If $\sigma=db$ and $b_-$ is the left endpoint value, then
$b(t)=b_-+\sigma((-\infty,t))$. The right-continuous BV function
has a finite signed derivative $dg$ satisfying
$dg((s,\infty))=-g(s)$. Fubini therefore gives
\[
                       \int g\,d\sigma=-\int b\,dg.
\]
The same identity holds for smooth $b_\epsilon$.
Their values lie between the same finite endpoint constants,
so dominated convergence against $|dg|$ proves (NS1).
\end{proof}

\begin{theorem}[A nonlinear scalar curvature formula]
\label{thm:r11-nonlinear-Stieltjes}
Let $u$ be locally bounded classical psh, let $1\leq k\leq q$,
and let $F$ be finite convex and nondecreasing on an open
interval containing the local range of $u$. Put
\[
 T_k=T_{S,k}(u,\ldots,u),\qquad
 Q_{k-1}=T_{S,k-1}(u,\ldots,u),\qquad a_-=F'_-,
 \qquad \sigma_{F,k}=d(a_-^k).
\]
Then $F(u)$ is locally bounded classical psh and
\[
 \boxed{\quad
 T_{S,k}(F(u),\ldots,F(u))
  =a_-(u)^kT_k+
       \int\mathsf L_u^{-,Q_{k-1}}(t)\,d\sigma_{F,k}(t).
 \quad}
 \tag{NS2}
\]
Equivalently use $F'_+(u)^k$ and the plus level current of
(SC1) in place of the two displayed terms. The measure
$\sigma_{F,k}$ is positive, so the level integral is an
original positive coefficient-Radon current.
\end{theorem}

\noindent\textbf{Status: \texttt{NEW\_PROVED}} for all original
backgrounds and degrees.

\begin{proof}
For smooth $u$ and smooth $F$, the ordinary chain rule and
$(d'u\wedge d''u)^2=0$ give
\[
 T_{S,k}(F(u))
   =F'(u)^kT_k+
       kF'(u)^{k-1}F''(u)\Gamma_{Q_{k-1}}(u).
 \tag{NS3}
\]
The formula remains valid for bounded classical $u$ and smooth
convex nondecreasing $F$. Indeed take a common local smooth
decreasing row $u_l\downarrow u$. All bounded products and
weights in (NS3) pass by the frozen weighted joint theorem.
For the energy term expand
\[
 \Gamma_{Q_{k-1,l}}(u_l)
       =\tfrac12D(u_l^2Q_{k-1,l})
                          -u_lT_{S,k}(u_l).
\]
Every smooth scalar coefficient of $u_l$, including those
multiplying this identity, is a bounded difference of convex
nondecreasing compositions on its compact range. The same
finite weighted passages apply, as in (LC2). This proves (NS3)
on every original test.

On a slightly larger compact range interval extend $F$ affinely
to the whole line, preserving convexity and nonnegative bounded
slopes. Choose a nonnegative smooth probability density
$\rho$ supported in $(0,1)$ and set
\[
 F_\epsilon(s)=\int_0^1F(s-\epsilon r)\rho(r)\,dr,\qquad
 a_\epsilon=F'_\epsilon,\qquad b_\epsilon=a_\epsilon^k.
\]
These functions are smooth, convex and nondecreasing, and
$F_\epsilon\to F$ uniformly. Moreover
$a_\epsilon(t)\to a_-(t)$ at every $t$, including slope jumps.
The functions $b_\epsilon$ are nondecreasing with the same
constant endpoint values as $b=a_-^k$.

By the energy coarea identity, (NS3) for $F_\epsilon$ becomes
\[
 T_{S,k}(F_\epsilon(u))
     =a_\epsilon(u)^k T_k+
          \int\mathsf L_u^{-,Q_{k-1}}(t)\,db_\epsilon(t).
\]
For each original compact test, the level pairing is a compactly
supported right-continuous BV function by
Proposition~\ref{prop:r11-two-level-currents}.
Lemma~\ref{lem:r11-one-sided-Stieltjes} passes the last integral,
retaining every atomic level term. Fixed-measure bounded
convergence passes the first term.

Finally uniform convergence of bounded psh potentials implies
current convergence of their bounded products here directly:
the bounded telescoping formula writes the difference as
$D$ of their uniformly small difference times finitely many
uniformly CLN-bounded mixed tails. Pairing with a fixed original
test therefore tends to zero. This identifies the left limit
with $T_{S,k}(F(u))$ and proves (NS2).
The alternative slope convention follows from (SC2), because
the atom of $\sigma_{F,k}$ at $t$ is exactly
$F'_+(t)^k-F'_-(t)^k$ and
$D(uQ_{k-1})=T_k$. This also verifies the atomic cancellation
without assuming that $T_k$ avoids the corner levels.
\end{proof}

\begin{corollary}[Exact nonlinear truncation sweeping]
\label{cor:r11-exact-truncation-sweeping}
For every real $t$ and bounded classical $u$ one has
\[
 \begin{split}
 T_{S,k}(\max(u,t),\ldots,\max(u,t))
   &=\mathbf1_{\{u>t\}}T_k+
                           \mathsf L_u^{-,Q_{k-1}}(t)\\
   &=D(\max(u,t)\,Q_{k-1}).
 \end{split}
 \tag{NS4}
\]
In particular the positive excess over the strict finite-level
restriction is carried by the exact Borel level $\{u=t\}$.
The excess need not be separately closed when $k<q$.
\end{corollary}

\noindent\textbf{Status: \texttt{NEW\_PROVED}} as a nonlinear
identity; positivity of the linear level current alone did not
identify the repeated truncation product.

\begin{proof}
Apply (NS2) to $F(s)=\max(s,t)$. Its lower slope is
$\mathbf1_{\{s>t\}}$, whose $k$-th power has derivative
$\delta_t$. The second equality is exactly (SC1).
The carrier and positivity are already proved for that level
current; its differential defect is (LC5).
\end{proof}

\begin{theorem}[A level measure and the total loss residue]
\label{thm:r11-level-loss-residue}
Let $u\leq0$ be negative extended classical psh on an ordinary
domain $V$, and suppose $\mathrm{BG}(u,0)$.
For every $t<0$ define
\[
 \Lambda_u(t)
    :=T_{S,q}(\max(u,t),\ldots,\max(u,t))
                                  -\mathbf1_{\{u>t\}}\mu_u.
 \tag{NS5}
\]
Then $\Lambda_u(t)$ is a positive finite Radon measure carried
by the actual level $\{u=t\}$. For integers $j>-t$,
\[
 \Lambda_u(t)(V)
    =\mu_{u,j}(\{u\leq t\})
    =\mu_u(\{-j<u\leq t\})+\nu_{u,j}(V).
 \tag{NS6}
\]
Consequently the nonnegative numbers $\nu_{u,j}(V)$ are
decreasing for all sufficiently large $j$, their finite limit
$\ell_S(u)$ exists, and
\[
 \Lambda_u(t)(V)=\mu_u(\{u\leq t\})+\ell_S(u)
                         \qquad(t<0).
 \tag{NS7}
\]
The number $\ell_S(u)$ is independent of the choice of level.
It vanishes exactly when $u\in\mathcal F_{S,\partial}(V)$.
No Radon measure concentrated on the possibly nonclosed pole
set is inferred from this scalar residue.
\end{theorem}

\noindent\textbf{Status: \texttt{NEW\_PROVED}}.

\begin{proof}
Local finiteness of $\mu_u$ follows from compact truncation
conservation, without a no-loss assumption. For $j>-t$,
$\max(u_j,t)=\max(u,t)$ and $\{u_j>t\}=\{u>t\}$.
On the latter carrier strict truncation consistency gives
$\mu_{u,j}=\mu_u$. Thus (NS4), applied to $u_j$, identifies
(NS5) with its positive level measure, carried by
$\{u_j=t\}=\{u=t\}$. Its carrier lies in the fixed compact
$\overline{\{u<t/2\}}$, so its total mass is finite.

Write $Q_j=T_{S,q-1}(u_j,\ldots,u_j)$. Equation (NS4) gives
\[
 T_{S,q}(\max(u,t))-\mu_{u,j}
                   =D((\max(u,t)-u_j)Q_j).
\]
The bounded multiplier vanishes outside the compact closure of
$\{u<t\}$. Choose a compact original smooth cutoff equal to one
near the larger compact $\overline{\{u<t/2\}}$; this contains
both the multiplier carrier and all the weak-level measures
appearing below. Its pairing with the displayed double derivative
is zero by Stokes.
Subtract the common measure on $\{u>t\}$ to obtain the first
equality in (NS6). Splitting the remaining carrier at $-j$
and using consistency gives its second equality.

For fixed $t$, the first measure on the right of (NS6)
increases as $j$ increases. Hence the remaining number
$\nu_{u,j}(V)$ decreases and has a nonnegative limit. The
increasing sets exhaust the finite-value part of $\{u\leq t\}$,
and $\mu_u$ vanishes on the omitted pole set. This proves
(NS7). The limit is the same global sequence for every $t$,
so it is independent of the level. Finally the boundary-class
criterion identifies $\nu_{u,j}(V)\to0$ exactly with no loss.
The proof only constructs a number and exact level measures;
it makes no assertion about a polar-supported current limit.
\end{proof}

\begin{corollary}[Comparison with a single finite loss defect]
\label{cor:r11-comparison-scalar-loss}
Under the preceding hypotheses on $u$, every negative extended
classical psh $v$ satisfies
\[
          \mu_v(\{u<v\})
                    \leq\mu_u(\{u<v\})+\ell_S(u).
 \tag{NS8}
\]
Extended infinite values of the right side are allowed.
\end{corollary}

\noindent\textbf{Status: \texttt{DIRECT\_COROLLARY}} of
(NS7) and Theorem~\ref{thm:r11-unbounded-defect-comparison}.

\begin{proof}
Since $v\leq0$, all positive-margin gaps have compact closure
inside the corresponding strict negative sublevel of $u$.
In (UC9), bound each discarded mass on $H_\epsilon$ by its
total mass $\nu_{u,j}(V)$ and use its limit $\ell_S(u)$.
Then let the positive margin decrease to zero; the two
finite-value gap carriers increase to $\{u<v\}$. This gives
(NS8) by continuity from below.
\end{proof}

\begin{proposition}[A Dini Jensen identity on the original domain]
\label{prop:r11-Dini-Jensen}
Assume in addition (DG1), let
$u\in\mathcal F_{S,\partial}(V)$, and suppose
\[
                \int_1^\infty t^{q-1}
                              C_S(\{u<-t\};V)\,dt<\infty.
\]
Set $Q=N_{q-1}(u,\ldots,u)$, with $Q=S$ when $q=1$.
Then for every $t<0$,
\[
 \Lambda_u(t)-\mathbf1_{\{u\leq t\}}\mu_u
                         =D((t-u)^+Q).
 \tag{NS9}
\]
The multiplier on the right is coefficient-integrable and
carried by a compact strict-sublevel closure. Consequently every
original smooth scalar $\varphi$ defined near the larger compact
$\overline{\{u<t/2\}}$ satisfies
\[
 \int\varphi\,d\Lambda_u(t)
       -\int_{\{u\leq t\}}\varphi\,d\mu_u
                         =\langle(t-u)^+Q,D\varphi\rangle.
 \tag{NS10}
\]
If $\varphi$ is also classical psh there, the right side is
nonnegative. Thus (NS10) is a genuine relative Jensen inequality
for original smooth psh tests, retaining its integrability premise.
\end{proposition}

\noindent\textbf{Status: \texttt{NEW\_PROVED}}.

\begin{proof}
The Dini endpoint of Corollary~\ref{cor:r11-boundary-recursion-dini}
gives coefficient integrability of $uQ$ and $D(uQ)=\mu_u$.
The lower-degree total-variation convergence in (DG7) gives
$Q_j\to Q$. Multiplication by the fixed bounded Borel function
$\max(u,t)$ preserves this convergence. Since
$D(\max(u,t)Q_j)=T_{S,q}(\max(u,t))$ for all sufficiently large
$j$ by (NS4), passage to original distributions yields
\[
                    D(\max(u,t)Q)=T_{S,q}(\max(u,t)).
\]
Subtract $D(uQ)=\mu_u$ and use (NS5) to obtain (NS9).
The multiplier is zero outside $\{u<t\}$ and is integrable by
the Dini estimate. Insert an original smooth compact cutoff
equal to one near $\overline{\{u<t/2\}}$ before testing with
$\varphi$; all derivatives of that cutoff vanish on the carrier.
This proves (NS10) on the complete original test sheaf.
For smooth classical $\varphi$, $D\varphi$ is a positive
$(1,1)$-form, hence strongly positive in this degree; pairing
with the middle-positive $Q$ and the nonnegative multiplier
proves the last assertion.
\end{proof}

\paragraph{Scope.}
The nonlinear level formula handles repeated scalar composition of
one actual classical potential. It is not a general quotient for
two analytic PS/PL presentations. The loss residue measures a
specific failure of uncut truncations to retain mass; it is not an
automatic non-Archimedean Lelong number or a canonically supported
polar measure. Formula (NS9) retains an explicit sufficient Dini
endpoint assumption. The later no-loss endpoint argument
(CR6)--(CR7) permits replacing it by separately verified
coefficient integrability of $uQ$, with the same proof above.
Automatic integrability for every boundary-class potential is
not proved. Existence of the closed top measure alone does not
justify that multiplier or its recursive identification.


% END REVIEWED MODULE research_r11/capacity/nonlinear_level_sweeping.tex


% BEGIN REVIEWED MODULE research_r11/capacity/weak_maximum_subsolutions.tex
\section{A weak maximum inequality and finite subsolution maxima}
\label{sec:r11-weak-maximum-subsolutions}

Throughout this section $S$ is the original arbitrary symmetric
middle-positive separately closed $(p,p)$-current and $q=n-p\geq1$.
For one bounded classical psh function write
$T_k(u)=T_{S,k}(u,\ldots,u)$, $0\leq k\leq q$.
All restrictions below are restrictions of the original coefficient
Radon measures to the printed actual Borel sets. In particular the
inequalities are inequalities in the original middle-positive cone,
and are tested against the original positive complementary forms.
Strongly positive forms may be used for coefficient trace control.
The identities themselves hold against every original compact
smooth complementary form and every actual moment refinement.

\begin{theorem}[The equality-level defect of a bounded maximum]
\label{thm:r11-bounded-weak-maximum}
Let $u,v$ be locally bounded classical psh functions on an ordinary
open set, let $w=\max(u,v)$, and fix $1\leq k\leq q$. Put
\[
 f=u-v,\qquad
 Q_\ell=T_{S,k-1}(
       \underbrace{w,\ldots,w}_{\ell},
       \underbrace{u,\ldots,u}_{k-1-\ell}),
       \quad 0\leq\ell\leq k-1.
 \tag{WM1}
\]
Then one has the exact original coefficient-current identity
\[
 \boxed{\quad
 T_k(w)-\mathbf1_{\{u\geq v\}}T_k(u)
             -\mathbf1_{\{u<v\}}T_k(v)
       =\sum_{\ell=0}^{k-1}
                 \mathsf L_{f}^{+,Q_\ell}(0)\geq0.
 \quad}
 \tag{WM2}
\]
The superscript $+$ is the upper-slope convention (SC1), and every
summand is positive and carried by $\{u=v\}$. Swapping $u$ and
$v$ gives the alternative inequality with the weak comparison
assigned to the other branch. No closedness of the equality-level
defect is asserted below top degree.
\end{theorem}

\noindent\textbf{Status: \texttt{NEW\_PROVED}} for the
explicit positive level-current decomposition in every degree.
The bounded top-degree weak maximum inequality itself is
\texttt{INTERNAL\_INPUT}: frozen R3
\path{prop:r3-relative-contact-inequality} already proves it
by decreasing maxima and compact Portmanteau. It is not counted
again here. The new content is the exact sum in (WM2), with
its lower-degree coefficient-current meaning and slope convention.

\begin{proof}
Every $Q_\ell$ is an allowed original symmetric middle-positive
separately closed current by the frozen mixed bounded calculus.
Apply Proposition~\ref{prop:r11-two-level-currents} to $f=u-v$
and each such background. It gives
\[
 \mathsf L_{f}^{+,Q_\ell}(0)
     =\mathsf L_{-f}^{-,Q_\ell}(0)
     =D((w-u)Q_\ell)
                    -\mathbf1_{\{u<v\}}D((v-u)Q_\ell).
 \tag{WM3}
\]
Here $w-u=(-f)_+$, and the replacement of $\max(-f,0)$ by
$(-f)_+$ does not add a constant term. The minus convention
in the second expression is essential: its curvature multiplier
uses the strict set $\{u<v\}$.

The bounded telescoping identity, including all repeated slots,
is
\[
       T_k(w)-T_k(u)
          =\sum_{\ell=0}^{k-1}D((w-u)Q_\ell).
 \tag{WM4}
\]
On the actual strict Borel set $\{u<v\}$, frozen maximum
locality replaces every $w$ slot in $Q_\ell$ by $v$. The
same telescoping formula there consequently gives
\[
 \mathbf1_{\{u<v\}}
       \sum_{\ell=0}^{k-1}D((v-u)Q_\ell)
      =\mathbf1_{\{u<v\}}(T_k(v)-T_k(u)).
 \tag{WM5}
\]
This use of locality concerns products of bounded classical
slots: expand $D((v-u)Q_\ell)$ as the difference of the two
corresponding original mixed products before restricting.
It does not differentiate a Borel indicator or infer locality
from equality on an arbitrary Borel set.

Sum (WM3), insert (WM4)--(WM5), and combine the restrictions
of $T_k(u)$. This is (WM2). Positivity and the exact equality
carrier follow from the two-level theorem before the currents
are added. The bounded construction is local, so finite compact
localizations cover the locally bounded case. Every operation
is intrinsic to the original current sheaf, proving ordinary
restriction and actual-refinement compatibility.
\end{proof}

\begin{theorem}[Weak maxima in the no-loss boundary class]
\label{thm:r11-unbounded-weak-maximum}
Let $V$ be an ordinary open set and let
$u,v\in\mathcal F_{S,\partial}(V)$, with their actual values
and the compact-sublevel boundary condition (BC3). Put
$w=\max(u,v)$, which belongs to the same class.
For top degree, with no further domain hypothesis,
\[
 \mu_w\geq\mathbf1_{\{u\geq v\}}\mu_u
                        +\mathbf1_{\{u<v\}}\mu_v.
 \tag{WM6}
\]
For every $1\leq k<q$ the analogous inequality
\[
 N_k(w)\geq\mathbf1_{\{u\geq v\}}N_k(u)
                        +\mathbf1_{\{u<v\}}N_k(v)
 \tag{WM7}
\]
holds under the one-affinoid compact-containment hypothesis (DG1).
Here $N_k$ are the all-degree positive closed products of
Theorem~\ref{thm:r11-boundary-all-degree-products}.

In either case the difference is a positive original
coefficient-Radon current carried by
$\{u=v> -\infty\}$. It is the increasing limit of the
bounded equality-level defects restricted to the common
finite-value carriers. All involved currents give zero
coefficient variation to $\{u=-\infty\}\cup\{v=-\infty\}$.
Thus the printed weak inequality on the common pole is
unambiguous and contributes zero. Swapping the potentials
again gives the other weak-branch convention.
\end{theorem}

\noindent\textbf{Status: \texttt{NEW\_PROVED}} on the
printed no-loss class, with the lower-degree domain hypothesis
retained.

\begin{proof}
The solid-cone theorem puts $u+v+w$ and $w$ in
$\mathcal F_{S,\partial}(V)$. For any positive degree under
consideration, the positive mixed expansion of $u+v+w$
dominates each diagonal current $N_k(u),N_k(v),N_k(w)$.
The product of $u+v+w$ gives zero coefficient mass to its
own pole set, which is exactly
\[
                       P=\{u=-\infty\}\cup\{v=-\infty\}.
 \tag{WM8}
\]
Positive trace domination therefore proves the same statement
for each of the three diagonal currents. This argument uses
the top mixed theorem (BC7)--(BC9) if $k=q$, and the all-degree
mixed theorem under (DG1) if $k<q$. Avoiding only the potential's
own pole set would not suffice for this step.

For $M>0$ set
\[
                E_M=\{u>-M\}\cap\{v>-M\}.
 \tag{WM9}
\]
For any integer $j>M$, the canonical bounded functions
$u_j,v_j,w_j$ satisfy $w_j=\max(u_j,v_j)$. On $E_M$,
strict maximum locality identifies each of their products
with the corresponding unbounded product, and the comparisons
$u_j\geq v_j$ and $u_j<v_j$ are exactly the printed
comparisons of $u,v$. Restrict the positive defect in (WM2)
for $u_j,v_j$ to $E_M$. It is independent of $j>M$ by this
identity; denote it by $\Delta_M$.

For $M'<M$, the restriction of $\Delta_M$ to $E_{M'}$
is $\Delta_{M'}$. The family is therefore increasing as a
family of positive coefficient currents. It is bounded above,
in the positive cone, by $\mathbf1_{E_M}N_k(w)$, because
the other two terms in (WM2) are positive restrictions.
Finite original trace controls show that the increasing limit
$\Delta$ is a locally finite positive coefficient-Radon current.
Each $\Delta_M$ is carried by
$E_M\cap\{u=v\}$, so $\Delta$ is carried by
$\{u=v> -\infty\}$.

Now $E_M\uparrow V\setminus P$. On this increasing union,
the identity defining $\Delta_M$ gives exactly the difference
in (WM6) or (WM7). All three currents annihilate $P$ by (WM8),
so the identity holds on the whole domain. This proves the
inequality and the claimed defect description. The proof never
subtracts two infinite potential values and never restricts an
uncontrolled weak limit to an equality set. Every assertion
uses the original coefficient measures and hence survives all
actual moment refinements.
\end{proof}

\begin{corollary}[Finite maxima of relative subsolutions]
\label{cor:r11-finite-subsolution-maxima}
Let $\lambda$ be a positive Radon measure on $V$. If
$u_1,\ldots,u_m\in\mathcal F_{S,\partial}(V)$ and
\[
                              \mu_{u_i}\geq\lambda
                                   \quad(1\leq i\leq m),
 \tag{WM10}
\]
then their actual finite maximum $w=\max_i u_i$ belongs
to the same class and satisfies $\mu_w\geq\lambda$.
If the potentials belong to a fixed power class
$\mathcal E^\sigma_{S,\partial}(V)$, $\sigma>0$, then
so does $w$. For bounded classical subsolutions the same
finite-maximum assertion holds on every ordinary open set,
without the boundary condition. Under (DG1) there is also
the analogous assertion in every positive degree for a common
positive coefficient-Radon current dominated by each $N_k(u_i)$.
\end{corollary}

\noindent\textbf{Status: \texttt{DIRECT\_COROLLARY}} of
the new weak maximum inequality and the established solid cones.

\begin{proof}
For two potentials, restrict the two inequalities
$\mu_u\geq\lambda$ and $\mu_v\geq\lambda$ to the
disjoint Borel sets $\{u\geq v\}$ and $\{u<v\}$.
Their sum is $\lambda$, and (WM6) dominates that sum.
If a common pole occurs, $\lambda$ gives it zero mass because
it is dominated by each original product; the same conclusion
therefore holds with the actual extended-value comparisons.
Induction proves the finite assertion. The respective solid-cone
results give the claimed class membership. The bounded and
lower-degree versions use (WM2) and (WM7) with the identical
restriction argument and positive coefficient domination.
\end{proof}

\paragraph{Dependency and scope audit.}
The frozen R1 \texttt{maximum\_locality.tex} proves strict Borel
maximum locality and compact-gap comparison. Frozen R3
\path{relative_\allowbreak contact_\allowbreak inequality.tex}, Proposition
\path{prop:r3-\allowbreak relative-\allowbreak contact-\allowbreak inequality}, already deduces
the bounded top-degree weak maximum and contact inequalities.
The present proof adds the explicit positive level-current
decomposition (WM2) and its all-degree coefficient meaning;
the following theorem extends it to the new unbounded no-loss
boundary class. None of those frozen inequalities is being
renamed as a new result.
All closed backgrounds used in (WM1) are bounded BT products
already furnished by the frozen arbitrary-$S$ calculus. The
nonclosed equality-level defects are never inserted as a new
background for BT iteration.

Finite subsolution maximum closure is one concrete ingredient
for a future envelope construction. It does not establish
arbitrary envelope compactness, membership of an ordinary upper
semicontinuous regularization in the relative class, a contact
measure theorem, or nonlinear Monge--Ampere Dirichlet existence.
The frozen relative-support envelope counterexample continues
to prohibit an unrestricted identification of relative
representatives with ambient upper regularizations.

% END REVIEWED MODULE research_r11/capacity/weak_maximum_subsolutions.tex


% BEGIN REVIEWED MODULE research_r11/capacity/logarithmic_pole_representatives.tex
\section{Preserving compact-sublevel poles in every power class}
\label{sec:r11-logarithmic-pole-representatives}

Keep the arbitrary original background, the strict analytic
Laurent box $V$, its closed box $C$, and its barrier $\rho$
from (TC1)--(TC4). In particular $\overline V$ is contained
in the interior of one original affinoid, and $S$ extends to
an ordinary neighborhood of $C$. Put $q=n-p\geq1$.
All assertions below concern the actual given representatives;
they change a potential by an explicit formula, rather than
changing it on an exceptional set.

\begin{lemma}[Logarithmic attenuation and exponential capacity tails]
\label{lem:r11-logarithmic-attenuation}
Let $u:V\to[-\infty,0]$ be extended classical psh with
$\mathrm{BG}(u,0)$, without any no-loss assumption. Set
\[
                         v=-\log(1-u),
 \tag{LP1}
\]
with the value $-\infty$ at each actual pole of $u$.
Then $v$ is negative extended classical psh, has the same
actual pole set, satisfies $\mathrm{BG}(v,0)$, and belongs
to $\mathcal F_{S,\partial}(V)$. For some finite $C_0,t_0$
depending on the stated data, one has
\[
               C_S(\{v<-t\};V)\leq C_0e^{-qt}
                                     \qquad(t\geq t_0).
 \tag{LP2}
\]
For every compact $H\Subset V$ and every $0<c<q$,
\[
                \int_H e^{-cv}\,d\mu_v<\infty.
 \tag{LP3}
\]
The integral omits the $\mu_v$-null pole set. No global
finite mass at the ordinary boundary is inferred at this stage.
\end{lemma}

\noindent\textbf{Status: \texttt{NEW\_PROVED}} for this
explicit exponential-tail attenuation, using (BQ1) as input.

\begin{proof}
The scalar map $F(r)=-\log(1-r)$ on $r\leq0$ has
$F'(r)=(1-r)^{-1}>0$ and $F''(r)=(1-r)^{-2}>0$.
Apply frozen convex increasing composition to the bounded
canonical truncations of $u$ and pass to their decreasing
limit. This proves the asserted classical property and actual
values. The identity
\[
                    \{v<-t\}=\{u<-(e^t-1)\}
 \tag{LP4}
\]
proves the compact-sublevel condition, the equality of poles,
and, by (BQ1), the bound $C(e^t-1)^{-q}$.
For sufficiently large $t$ this is (LP2).
The capacity-tail criterion (BC3)--(BC5) now puts $v$ in
the no-loss boundary class.

To prove (LP3), partition the sufficiently deep finite-value
locus into the Borel strips
\[
                     A_m=\{m\leq-v<m+1\},
                           \qquad m\in\mathbf N.
 \tag{LP5}
\]
On $A_m$, strict maximum locality identifies $\mu_v$ with
the bounded measure of $v\vee(-m-2)$. Its normalized capacity
bound gives
\[
 \mu_v(A_m)
    \leq(m+2)^q C_S(A_m;V)
    \leq(m+2)^q C_S(\{v<-m+1\};V)
    \leq C_1(m+2)^q e^{-qm}.
 \tag{LP6}
\]
The harmless shift by one keeps the last carrier strictly
open and includes any equality-level mass. Multiplication by
$e^{c(m+1)}$ and summation converge if $c<q$.
On the remaining bounded-value portion of the compact $H$,
local finiteness of $\mu_v$ completes the proof.
\end{proof}

\begin{theorem}[One actual pole representative in every finite power class]
\label{thm:r11-all-power-pole-representative}
Let $u$ satisfy Lemma~\ref{lem:r11-logarithmic-attenuation}.
For every prescribed compact $H\Subset V$ and every
$\varepsilon>0$ there exists a negative extended classical
psh function $g$ on $V$ such that
\[
 \begin{gathered}
  g=-\log(1-u)-\varepsilon
       \quad\hbox{on an ordinary neighborhood of }H
                      \hbox{ and of the pole set},\\
  \{g=-\infty\}=\{u=-\infty\},\qquad
  \mathrm{BG}(g,0),\qquad \mu_g(V)<\infty,\\
  g\in\bigcap_{\sigma>0}\mathcal E^\sigma_{S,\partial}(V).
 \end{gathered}
 \tag{LP7}
\]
More precisely, $g=A\rho$ on an ordinary boundary collar,
for a finite $A>0$, and there are finite constants such that
for all sufficiently large real $t$ and integers $j$,
\[
 \begin{split}
 C_S(\{g<-t\};V)&\leq C_2e^{-qt},\\
 \mu_g\{g\leq-t\}&\leq C_3(1+t)^q e^{-qt},\\
 \|\mu_{g,j}-\mu_g\|_{\mathrm{TV}(V)}
                  &\leq C_4(1+j)^q e^{-qj}.
 \end{split}
 \tag{LP8}
\]
For every $0<c<q$, including the actual canonical truncations,
\[
       \int_V e^{-cg}\,d\mu_g<\infty,
       \qquad
       \sup_{j\geq1}\int_Ve^{-cg_j}\,d\mu_{g,j}<\infty.
 \tag{LP9}
\]
In fact these truncation exponential moments increase to the
displayed limiting moment. The function $g$ itself is chosen
using only $u,H,\varepsilon,V,\rho$: it is independent of the
background. The same $g$ therefore satisfies (LP7)--(LP9)
simultaneously for every original symmetric middle-positive
separately closed background of every degree $p<n$ extending
to a neighborhood of $C$, with the corresponding $q=n-p$.
Constants may depend on $S,u,V,H,\varepsilon,c$; no uniformity
over backgrounds or prescribed potentials is asserted.
\end{theorem}

\noindent\textbf{Status: \texttt{NEW\_PROVED}} with the
explicit strict-box and background-extension hypotheses.

\begin{proof}
Write $v=-\log(1-u)$ and choose $M>0$. The compact set
\[
                      K=H\cup\overline{\{v<-M\}}
                              \Subset V
\]
contains the pole set. Choose an ordinary relatively compact
neighborhood $W$ of $K$ inside $V$. On the compact
$\overline W$, continuity and strict negativity of $\rho$
give an $A>0$ so large that
$A\rho<-M-\varepsilon-1$ everywhere there.
Define $g$ on the following two ordinary open sets:
\[
 \begin{cases}
 g=v-\varepsilon&\text{on }W,\\
 g=\max(v-\varepsilon,A\rho)&\text{on }V\setminus K.
 \end{cases}
 \tag{LP10}
\]
On their overlap $W\setminus K$ one has $v\geq-M$,
so the first branch strictly exceeds $A\rho$.
The definitions therefore agree on the entire ordinary overlap.
The local classical gluing property, equivalently the frozen
curvewise definition after every complete extension, proves
that $g$ is extended classical psh on $V$.
It has the same actual poles as $v$, since those lie in $K$
and the function is bounded below on $V\setminus K$.
It also gives the required equality near $H$ and all poles.

Near the ordinary boundary, the condition
$\rho>-\varepsilon/(2A)$ makes the second branch in (LP10)
strictly larger than $v-\varepsilon\leq-\varepsilon$.
Thus $g=A\rho$ on this boundary collar. Its strict negative
sublevels are consequently contained in a compact subset of
$V$, the union of $\overline W$ and an appropriate compact
negative sublevel of $\rho$.
For $t>M+\varepsilon$, the portion outside $K$ cannot belong
to $\{g<-t\}$, and the exact deep-sublevel identity is
\[
                    \{g<-t\}=\{v<-(t-\varepsilon)\}.
 \tag{LP11}
\]
Equation (LP2) gives the first estimate in (LP8).
The capacity-tail criterion therefore puts $g$ in
$\mathcal F_{S,\partial}(V)$.

Choose a compact subset of $V$ outside which $g=A\rho$.
On that compact subset $\mu_g$ is finite by the no-loss
class theorem. Outside it ordinary locality identifies the
measure with $A^q\mu_\rho$, whose total mass is finite by
(TC4). Thus $\mu_g(V)<\infty$.
In particular this argument does not assert local boundedness
of $g$ near its interior poles or attempt to apply bounded
CLN across them.

Repeat the integer-strip estimate (LP6), now with $g$ and
the global finite mass on its bounded-value part. Summing
from $m=\lfloor t\rfloor$ gives the second estimate in
(LP8): for the decreasing exponential series,
\[
 \sum_{m\geq N}(m+2)^q e^{-qm}
                         \leq C_q(N+2)^q e^{-qN}.
\]
The same strip estimate with $e^{c(m+1)}$ proves the first
assertion of (LP9) for every $c<q$. All finite powers are
bounded by a constant times this exponential, so (LP7)
follows from the definition of the power classes.

For the uncut canonical product, capacity normalization gives
\[
 \nu_{g,j}(V)
       \leq j^q C_S(\{g\leq-j\};V)
       \leq j^q C_S(\{g<-j+1\};V)
       \leq C_5j^q e^{-qj}.
 \tag{LP12}
\]
The current $\mu_{g,j}$ agrees with $\mu_g$ on $\{g>-j\}$.
The variation on the complement is at most the sum of the
discarded mass (LP12) and the second bound in (LP8). This
proves the third bound there. For any fixed $j$, the bounded
function $g_j$ equals $A\rho$ on a sufficiently narrow
ordinary collar and extends as that bounded classical function
across $\partial V$. Bounded CLN and the extension of $S$
give $\mu_{g,j}(V)<\infty$.

Finally split the weighted truncation integral at $\{g>-j\}$:
\[
 \int_V e^{-cg_j}\,d\mu_{g,j}
    =\int_{\{g>-j\}}e^{-cg}\,d\mu_g
                            +e^{cj}\nu_{g,j}(V).
 \tag{LP13}
\]
The first term is bounded by the established finite integral.
The second is $O(j^q e^{-(q-c)j})$ for large $j$, and the
finitely many remaining terms are finite. This proves (LP9).
There is also an exact distributional strengthening. Since
$g$ has no loss, (NS7) at level $-j$ gives
$\nu_{g,j}(V)=\mu_g\{g\leq-j\}$. Consequently
\[
 (-g_j)_*\mu_{g,j}
       =(y\mapsto\min(y,j))_*\bigl((-g)_*\mu_g\bigr).
 \tag{LP18}
\]
Indeed the two sides agree on the finite-value carrier $g>-j$,
and both remaining parts have the same mass at the scalar value
$j$. The omitted pole is $\mu_g$-null. Apply (LP18) to the
increasing weights $e^{c\min(y,j)}$ and use monotone convergence
to obtain the claimed monotone convergence of exponential
moments. The construction (LP10) contains no background datum,
so repeating the proof for each allowed $S$ proves the simultaneous
quantifier stated above.
Every locality or capacity normalization used the full original
tester class on the same $V$.
\end{proof}

\begin{corollary}[An exact family of complete relative pole sets]
\label{cor:r11-all-power-complete-poles}
On the stated domain, the following two families of subsets
are equal:
\[
 \begin{split}
 &\bigl\{\{u=-\infty\}:u\leq0\text{ extended classical psh},
                                  \ \mathrm{BG}(u,0)\bigr\},\\
 &\bigl\{\{g=-\infty\}:g\in\bigcap_{\sigma>0}
       \mathcal E^\sigma_{S,\partial}(V),\ \mu_g(V)<\infty\bigr\}.
 \end{split}
 \tag{LP14}
\]
Each such set, and every subset of it, has zero ordinary-open
outer $S$-capacity. The same set is avoided by every original
coefficient of $S$ and every finite product of power-class
potentials, and by every graded energy current already defined.
\end{corollary}

\noindent\textbf{Status: \texttt{NEW\_PROVED}} for the exact
pole-family equality. The nullity assertions have status
\texttt{DIRECT\_COROLLARY}; their inputs are (BQ2), (DG3), and (TC9).

\begin{proof}
Theorem~\ref{thm:r11-all-power-pole-representative} proves the
first inclusion in (LP14); the reverse is part of the printed
definition of every power class. Outer nullity follows from
(BQ2), and is inherited by arbitrary subsets. For the complete
pole set, which is Borel, (DG3) and
Corollary~\ref{cor:r11-power-class-null-ideal} give all the
stated coefficient-null assertions. Their positive variation
measures then also vanish on all measurable subsets.
\end{proof}

\begin{corollary}[Exponential coefficient and energy calculus for the representative]
\label{cor:r11-logarithmic-all-degree-exponential}
Let $g$ be the function constructed in (LP10), let
$N_k=N_k(g)$, and let $T_{k,j}=T_k(g_j)$.
For every original compact smooth complementary test $\eta$,
$0\leq k\leq q$, and $0<c<q$, the coefficient multiplier
$e^{-cg}N_k$ is locally defined and
\[
 \bigl\|\mu_{e^{-cg_j}T_{k,j}-e^{-cg}N_k,\eta}\bigr\|_{\rm TV}
       \leq C_{\eta,c}(1+j)^k e^{-(q-c)j}
 \tag{LP15}
\]
for all sufficiently large $j$. For $0\leq k\leq q-1$,
the analogous statement for the energy currents is
\[
 \bigl\|\mu_{e^{-cg_j}\Gamma_{T_{k,j}}(g_j)
                     -e^{-cg}\Gamma_{N_k}(g),\eta}\bigr\|_{\rm TV}
       \leq C'_{\eta,c}(1+j)^{k+2}e^{-(q-c)j}.
 \tag{LP16}
\]
Here the test degree in (LP16) is complementary to the energy
degree. In particular, throughout $0\leq k\leq q-1$,
\[
 D(e^{-cg}N_k)
      =-c e^{-cg}N_{k+1}
                          +c^2e^{-cg}\Gamma_{N_k}(g)
 \tag{LP17}
\]
on all original compact smooth tests. The energy currents need
not be separately closed. The signed right side of (LP17) is
separately closed as a whole by that identity, but is not
asserted to be positive. The $N_k$ themselves are the already established
separately closed middle-positive currents, with $N_0=S$.
\end{corollary}

\noindent\textbf{Status: \texttt{NEW\_PROVED}} for the
exponential weights; existence, closedness and the finite-power
energy construction are inputs from the all-degree theorems.

\begin{proof}
Since $g$ belongs to every positive power class, the original
all-degree products, energy currents and their coherent finite
level restrictions exist by
Theorems~\ref{thm:r11-power-full-recursion} and
\ref{thm:r11-power-graded-energy}. No exponential assertion
is used for their initial construction.

Fix a positive original trace control for $\eta$ supplied by
(DG2). On the strip $A_m=\{m\leq-g<m+1\}$, strict
bounded locality allows every slot to be replaced by its
truncation at $-m-2$. Equations (DG3) and (LP8) bound the
resulting trace by
\[
                         C_\eta(m+2)^k e^{-qm}.
\]
The same statement for the energy trace follows from (PG8),
with $k$ replaced by $k+2$. Energy locality here is the
already proved shifted-square/four-branch locality, not an
inference from equality of arbitrary Borel representatives.
Multiplication by $e^{c(m+1)}$ and summation prove local
exponential integrability and the limiting tail bounds.
The pole set gives zero trace mass, including at $k=0$,
by (DG3) and its outer capacity nullity.

On $\{g\leq-j\}$, the entire uncut curvature trace is
bounded by $C_\eta j^k C_S\{g\leq-j\}$, and its weighted
version by $e^{cj}$ times that quantity. The corresponding
energy bound is $C_\eta j^{k+2}e^{cj}C_S\{g\leq-j\}$.
Using the strict enlargement in (LP12) bounds both by the
displayed exponential rates. On $\{g>-j\}$, the respective
weighted measures agree exactly by the established locality.
Adding the uncut and limiting tails proves (LP15)--(LP16),
first for positive traces and then for every original coefficient.

The bounded smooth scalar chain identity is
\[
 D(e^{-cg_j}T_{k,j})
     =-c e^{-cg_j}T_{k+1,j}
                  +c^2e^{-cg_j}\Gamma_{T_{k,j}}(g_j).
\]
Pass the three terms through (LP15)--(LP16) and use
distributional continuity. This proves (LP17).
No differentiation of a restricted current is performed,
and the construction and all bounds remain valid after
each actual moment refinement of a fixed test.
\end{proof}

\paragraph{Exact advance and remaining polar gate.}
The frozen R2 trace-null witness theorem supplies a
measure-dependent attenuation for a specified nonconcentrating
trace. This section instead starts with the explicit compact
sublevel hypothesis, uses the new arbitrary-$S$ capacity tail,
and constructs one actual representative in every finite power
class, with a global boundary collar and quantitative exponential
control. It does not audit or replace that frozen theorem.

The unrestricted converse that every outer-$S$-capacity-null set
lies in a pole set in (LP14) is \texttt{FALSE} without a
nondegeneracy qualification. For example, on an original open
unit disc take $S=0$ and $E=V$. Its relative capacity is zero.
A potential with $E\subset\{u=-\infty\}$ would be identically
$-\infty$ and its strict negative sublevel would be all of $V$,
whose closure meets the ordinary boundary; it cannot satisfy
$\mathrm{BG}(u,0)$. Frozen proper-support examples give the
same obstruction for nonzero backgrounds away from their support.
An appropriate converse with explicit local nondegeneracy and
domain hypotheses remains \texttt{OPEN}. Nor is a general
ambient pluripolar set asserted to admit a global compact-sublevel
defining potential. The frozen relative-support counterexamples
continue to prevent an unrestricted identification of relative
nullity with ambient negligible or pluripolar sets.
The strict-box and extension hypotheses are used to control
the ordinary boundary; removing them is not part of this proof.

% END REVIEWED MODULE research_r11/capacity/logarithmic_pole_representatives.tex


% BEGIN REVIEWED MODULE research_r11/capacity/power_core_envelope_obstruction.tex
\section{A power-core obstruction to ambient Perron regularization}
\label{sec:r11-power-core-envelope-obstruction}

The frozen R2 \texttt{support\_envelopes.tex} already proves
that an ordinary ambient upper regularization can change a
relative capacitary representative on a set of positive relative
capacity. The point of the present extension is to retain all
the new boundary, finite-mass and power-energy hypotheses.
It therefore tests precisely the additional inference one might
try after the weak maximum theorem, rather than recounting the
old bounded counterexample. Frozen R6
\path{capacity_upper_envelope_obstruction.tex}, Proposition
\path{prop:r6-upper-envelope-relative-mass-loss}, already
gives a fixed nonzero measure lost by upper regularization.
The additional assertion here is its realization in the new
finite-mass zero-boundary core, with a common lower barrier
and exactly controlled energy in every positive power class.

\begin{proposition}[Increasing core subsolutions can lose their source under upper regularization]
\label{prop:r11-power-core-envelope-obstruction}
For every $p,q\geq1$, with $n=p+q$, there exist a strict original analytic
Laurent-box domain $V$, a nonzero original positive separately
closed background $S$, a nonzero positive measure $\mu$, and
an increasing sequence $u_j\in\mathscr T_0(S,V)$ such that
\[
 \begin{gathered}
     a\leq u_j\leq0,\qquad
     a\in\mathscr T_0(S,V),\qquad
     u_j,a\in\bigcap_{\sigma>0}\mathcal E^\sigma_{S,\partial}(V),\\
     \mu_{u_j}=\mu_a=\mu,\qquad
     \mu(V)=1,\qquad
     e_\sigma(u_j)=e_\sigma(a)=1
                  \quad(\sigma>0, j\geq1).
 \end{gathered}
 \tag{PO1}
\]
Nevertheless their actual increasing supremum $H=\sup_j u_j$
has ordinary upper semicontinuous regularization $H^*=0$, and
\[
          C_S(\{H<H^*\};V)\geq1,
          \qquad \mu_{H^*}=0\ne\mu.
 \tag{PO2}
\]
Thus ordinary upper regularization does not preserve even the
subsolution inequality in this uniformly controlled power core.
Every measure identity in (PO1)--(PO2) holds on all original
compact smooth scalar tests and all actual moment refinements.
\end{proposition}

\noindent\textbf{Status: \texttt{NEW\_PROVED}} as the
boundary-core and all-power strengthening of the frozen proper-support
obstruction. The corresponding unconditional Perron inference is
\texttt{FALSE}.

\begin{proof}
Work over the unchanged original field on
$X=(\mathbf A_K^n)^{\rm an}$, with coordinates
$T_1,\ldots,T_n$, and set
\[
 \begin{split}
 V&=\{|T_i|<1\ (1\leq i\leq n)\},\qquad
 Y=\{T_1=\cdots=T_p=0\},\qquad S=[Y]|_V,\\
 a&=\max\{-1,\log\max_{1\leq i\leq n}|T_i|\},\\
 u_j&=\max\{a,j^{-1}\log\max_{1\leq l\leq p}|T_l|\}.
 \end{split}
 \tag{PO3}
\]
The logarithm is $-\infty$ on its zero set, as usual.
The closed unit polydisc is compact in the interior of an
original closed polydisc of radius $|\pi|^{-1}>1$.
The current $[Y]$ extends there. Thus the exact domain and
extension hypotheses (TC1)--(TC2) hold, with no auxiliary
field or alteration of the absolute value.

The bounded function $a$ is the strict-box barrier (TC3).
The $u_j$ are bounded continuous classical psh functions,
and increase because $\log\max_{1\leq l\leq p}|T_l|\leq0$ on $V$.
They satisfy $-1\leq a\leq u_j\leq0$ and
$\mathrm{BG}(u_j,0)$ because their strict negative
sublevels lie in those of $a$.

On the ordinary tube $\{\max_{1\leq l\leq p}|T_l|<e^{-j}\}$
containing $Y\cap V$, one has
\[
                 j^{-1}\log\max_{1\leq l\leq p}|T_l|<-1\leq a.
\]
Thus $u_j=a$ there.
Every bounded product over $S=[Y]$ is carried by $Y$:
this follows directly by recursion and ordinary restriction,
since the background vanishes off $Y$.
Ordinary locality on the indicated tube therefore gives
\[
                             \mu_{u_j}=\mu_a
                              \quad\hbox{on the whole }V.
 \tag{PO4}
\]
This is equality on an ordinary neighborhood of the current's
support, not a false locality assertion from equality only
on an arbitrary Borel set.

The frozen analytic integration-current compatibility restricts
the bounded product of $a$ to the original $q$-dimensional
unit polydisc in $Y$. Its restricted potential is
\[
                 a|_Y=\max\{-1,\log\max_{p+1\leq i\leq n}|T_i|\}.
\]
The exact frozen original radial calibration, valid at arbitrary
positive real radii, gives
\[
                \mu_a=\delta_y,\qquad
     y=\text{the joint Gauss point in $Y$ with }|T_i(y)|=e^{-1}
                           \ (p+1\leq i\leq n).
 \tag{PO5}
\]
At this point $a(y)=u_j(y)=-1$. Hence the total masses and
all power energies in (PO1) are exactly one. Finite mass and
the compact boundary condition put each function in
$\mathscr T_0$, and hence in every positive power class.
The measure (PO5) is an equality of original Radon measures,
not merely of tropical test pairings, so it already includes
every original scalar test and every actual refinement.

For a point off $Y$, the finite number
$j^{-1}\log\max_{1\leq l\leq p}|T_l|$ increases to zero. On $Y$ the second
branch is always $-\infty$. Thus
\[
                  H=a\quad\hbox{on }Y,\qquad
                  H=0\quad\hbox{on }V\setminus Y.
 \tag{PO6}
\]
The complement of this original coordinate subspace is
ordinary dense in $V$, as in the frozen support example.
Since $H\leq0$, (PO6) gives $H^*=0$ everywhere.
The normalized tester $a+1$ has measure $\delta_y$, so
$C_S(\{y\};V)\geq1$. Since $H(y)=-1<0=H^*(y)$,
monotonicity proves the first part of (PO2). The second
follows from $D0=0$ in every differentiated slot.
This proves the full claimed strengthening.
\end{proof}

\begin{corollary}[Failure of the ordinary upper Perron operation]
\label{cor:r11-ordinary-power-Perron-false}
For the data in the preceding proposition let
\[
 \mathscr U=\{v\in\mathscr T_0(S,V):\mu_v\geq\mu\},
              \qquad P=\bigl(\sup_{v\in\mathscr U}v\bigr)^*.
 \tag{PO7}
\]
The class is nonempty, contains an actual classical solution
$a$ of $\mu_a=\mu$, and is stable under finite maxima by
Corollary~\ref{cor:r11-finite-subsolution-maxima}.
Nevertheless $P=0$ and $\mu_P\not\geq\mu$.
The same failure occurs if the class is additionally restricted
by the common lower barrier $v\geq a$ and by
$e_\sigma(v)\leq1$ for every $\sigma>0$.
\end{corollary}

\noindent\textbf{Status: \texttt{NEW\_PROVED}} as the
precise failed Perron operation; it does not assert nonexistence
of a solution for these data.

\begin{proof}
The functions $u_j$ belong to each printed class, whose
members are all at most zero. Their upper regularized supremum
therefore lies between $H^*=0$ and zero, so it is zero.
Its top measure is zero, while $\mu=\delta_y\ne0$.
The first class is stable under finite maxima by (WM6) and
the following direct total-mass check. If $w=\max(u,v)$
and $c>1$, then $cu<w$ on $\{w<0\}$, because $u\leq w$.
Moreover $\mu_w\{w=0\}=0$: apply (SC2) to the bounded
nonpositive $w$ and $Q=T_{S,q-1}(w,\ldots,w)$ to obtain
the positive measure
\[
       \mathsf L_w^{+,Q}(0)
              =D(\max(w,0)Q)-\mathbf1_{\{w\geq0\}}D(wQ)
              =-\mathbf1_{\{w=0\}}\mu_w\geq0.
\]
It must be zero. Compact-gap comparison for $cu,w$ therefore
gives $\mu_w(V)\leq c^q\mu_u(V)$; every positive-margin
gap is contained in a compact negative sublevel of $cu$.
Let $c\downarrow1$ to prove finite total mass of $w$.
The boundary-gap condition follows from $w\geq u$.
This argument does not mistake the weak maximum lower bound
for a total-mass upper bound. The restricted family
need not be asserted stable under maxima to prove the stated
failure: it already contains the explicit increasing $u_j$.
\end{proof}

\paragraph{Boundary of the conclusion.}
The background $[Y]$ is used here only as an allowed nonzero
counterexample to a universal assertion. It does not replace
the arbitrary background in any positive theorem.
The actual supremum (PO6) is compatible with the frozen
relative capacitary completion. The obstruction concerns its
replacement by the ordinary ambient upper regularization.
It persists even with one common all-power lower barrier,
finite total masses, compact negative sublevels and a
nonempty classical solution set. A general nonlinear
Dirichlet construction must therefore retain an appropriate
relative representative or impose and prove additional
representability hypotheses.

% END REVIEWED MODULE research_r11/capacity/power_core_envelope_obstruction.tex


% BEGIN REVIEWED MODULE research_r11/comparison/global_first_energy_calculus.tex
\section{Global first-energy Green calculus and relative Aubin--Mabuchi geometry}
\label{sec:r11-global-first-energy}

Throughout this section the domain and background satisfy precisely
(TC1)--(TC2): $V$ is the strict original analytic Laurent box compactly
contained in the interior of one original affinoid domain, and the
arbitrary original symmetric middle-positive separately closed
$(p,p)$-current $S$ is defined on a neighborhood of its closed box.
Put $q=n-p\geq1$, retain the negative bounded barrier $\rho$ of
Lemma~\ref{lem:r11-strict-Laurent-barrier}, and abbreviate
$\mathcal E^1=\mathcal E^1_{S,\partial}(V)$.
The background is not assumed strongly positive. The positive energy
matrices used below come from rank-one original $(1,1)$ squares
wedged with the same middle-positive background.

The compact bounded Green identity is the frozen R4 input
\texttt{eq:relative-green-energy} in
\path{relative_energy_agent.tex}, with the analytic local
scope made explicit in Section~\ref{sec:r11-analytic-Laurent-collars}.
The bounded common-collar Aubin--Mabuchi variation is already frozen
in R3 \path{energy_convex_blocks.tex},
\texttt{lem:r3-energy-midpoint}. The new assertions here remove that
fixed compact collar by the proved unbounded zero-boundary energy
class. The complete argument keeps track of the global boundary
term; it does not assume that it vanishes merely because a potential
has pointwise zero boundary values.

\begin{lemma}[Global bounded Green exchange from finite curvature mass]
\label{lem:r11-global-bounded-Green}
Let $Q$ be any original symmetric middle-positive separately closed
$(n-1,n-1)$-current on $V$. Let $a,b\leq0$ be bounded classical
psh functions satisfying $\mathrm{BG}(a,0)$ and
$\mathrm{BG}(b,0)$. Suppose the two positive measures
$D(aQ)$ and $D(bQ)$ have finite total mass on $V$. Then
\[
 \begin{split}
  \int_V(-a)D(bQ)&=\int_V(-b)D(aQ),\\
  \int_V\Gamma_Q(a,b)&=\int_V(-a)D(bQ),\qquad
  \int_V\Gamma_Q(a)=\int_V(-a)D(aQ)<\infty.
 \end{split}
 \tag{GE1}
\]
The cross energy has finite total variation. These are identities
of the actual original Borel multipliers and energy measures.
\end{lemma}

\noindent\textbf{Status: \texttt{NEW\_PROVED}} as a global
zero-boundary extension of compact bounded Green exchange.

\begin{proof}
For $\epsilon>0$ put $a^\epsilon=\max(a,-\epsilon)$.
The bounded difference $a-a^\epsilon$ is supported in the compact
set $K_\epsilon=\overline{\{a<-\epsilon\}}\Subset V$.
Compact Green exchange therefore gives
\[
 \int(a-a^\epsilon)D(bQ)
     =\int bD(aQ)-\int bD(a^\epsilon Q).
 \tag{GE2}
\]
The measure $D(a^\epsilon Q)$ has the same finite total mass as
$D(aQ)$. Indeed their difference is $D((a^\epsilon-a)Q)$,
with compactly supported multiplier. A compact original smooth
cutoff equal to one on a neighborhood of that support has zero
pairing with this derivative. Outside the compact support the
two measures agree, by ordinary locality. This proves both
finiteness and equality of their total masses.

Fix $\delta>0$. The set
$L_\delta=\overline{\{b<-\delta\}}$ is compact in $V$.
Relative CLN for the fixed $Q$ and the normalized classical
function $(a^\epsilon+\epsilon)/\epsilon\in[0,1]$ gives
$D(a^\epsilon Q)(L_\delta)\leq C_{\delta,Q}\epsilon$.
Consequently
\[
 \left|\int bD(a^\epsilon Q)\right|
 \leq \|b\|_\infty C_{\delta,Q}\epsilon
                     +\delta D(aQ)(V).
 \tag{GE3}
\]
First let $\epsilon\downarrow0$, then $\delta\downarrow0$.
The last term in (GE2) tends to zero. Bounded convergence on
its left proves the first equality in (GE1).

The original bounded clipping rule, including the zero energy
of an exact scalar level, gives
$\Gamma_Q(a-a^\epsilon,a)
 =\mathbf1_{\{a<-\epsilon\}}\Gamma_Q(a)$.
Apply compact Green to this difference and $a$:
\[
 \int_{\{a<-\epsilon\}}\Gamma_Q(a)
                 =\int(a^\epsilon-a)D(aQ).
 \tag{GE4}
\]
As $\epsilon\downarrow0$ the right side increases to the
finite integral $\int(-a)D(aQ)$. The energy gives no mass
to $\{a=0\}$, so monotone convergence gives the diagonal
energy identity. Apply it also to $b$ and $a+b$; the latter
has the same boundary property and finite curvature mass.
Polarization and the first equality of (GE1) identify the
cross integral. The measure Cauchy--Schwarz inequality bounds
its total variation by the product of the square roots of the
two finite diagonal masses.
\end{proof}

For $u_0,\ldots,u_q\in\mathcal E^1$ define
\[
 \mathcal I_S(u_0,\ldots,u_q)
       =\int_V(-u_0)M(u_1,\ldots,u_q).
 \tag{GE5}
\]
It is finite by the already proved mixed estimate (PE11).
Its interpretation uses the actual potentials; all mixed
measures annihilate their joint pole union.

\begin{theorem}[Full symmetry of the global first mixed energy]
\label{thm:r11-global-first-mixed-symmetry}
The nonnegative functional (GE5) is symmetric in all $q+1$
slots and separately additive and positively homogeneous on
$\mathcal E^1$. In particular the undifferentiated slot may
be exchanged with any differentiated slot without a boundary
term. Every integral is over the whole ordinary domain $V$.
\end{theorem}

\noindent\textbf{Status: \texttt{NEW\_PROVED}} for the
unbounded first-power boundary class and arbitrary original $S$.

\begin{proof}
Use the decreasing finite-mass approximants from
Theorem~\ref{thm:r11-finite-mass-power-approximation}:
$b_{i,j}=\max(u_i,j\rho)$. For a fixed $j$ they all lie
in $[-j,0]$ and satisfy the compact-sublevel boundary condition.
Let $A_j=\sum_{i=0}^q b_{i,j}$. Since
$A_j\geq(q+1)j\rho$ and $\rho<0$, comparison with
$c(q+1)j\rho$ for $c>1$ has strict gap all of $V$ and gives
\[
          \mu_{A_j}(V)
              \leq[(q+1)j]^q\mu_\rho(V)<\infty
 \tag{GE6}
\]
after $c\downarrow1$. The positive mixed expansion of
$\mu_{A_j}$ shows that every mixed top measure formed from
these bounded slots has finite total mass.

Fix the last $q-1$ slots and use their bounded mixed product
as $Q$ in Lemma~\ref{lem:r11-global-bounded-Green}. Its two
curvature measures are finite by (GE6), so the lemma exchanges
the first two slots. Symmetry among the differentiated slots
is already part of bounded BT; hence the bounded mixed energy
is fully symmetric.

The approximants in each slot have the fixed lower barrier
$u_i\in\mathcal E^1$. The global critical mixed-energy
continuity of Theorem~\ref{thm:r11-critical-energy-barrier-continuity}
therefore passes the bounded identity to (GE5), on the whole
domain and with no escaped boundary mass. This proves full
symmetry. Additivity and positive homogeneity in the first
slot are ordinary integration identities; in the other slots
they follow from the unbounded mixed-current calculus. The
same symmetry transfers all of them to every slot.
\end{proof}

\begin{theorem}[Globally finite unbounded energy matrices]
\label{thm:r11-global-first-energy-matrix}
Let $a,b,c_2,\ldots,c_q\in\mathcal E^1$ and put
$Q=N(c_2,\ldots,c_q)$, with $Q=S$ when $q=1$.
The graded energy currents at the endpoint $\sigma=1$ are
finite measures on all of $V$ and satisfy
\[
 \begin{split}
 \int_V\Gamma_Q(a,b)
       &=\mathcal I_S(a,b,c_2,\ldots,c_q),\\
 \left|\int_V\Gamma_Q(a,b)\right|^2
       &\leq\mathcal I_S(a,a,c_2,\ldots,c_q)
                   \mathcal I_S(b,b,c_2,\ldots,c_q).
 \end{split}
 \tag{GE7}
\]
Every cross energy has finite total variation, and the same
matrix positivity holds for finite real linear combinations
of the displayed potentials. Thus it applies to differences
even though a difference need not be psh.
\end{theorem}

\noindent\textbf{Status: \texttt{NEW\_PROVED}} as global
finiteness and Green identification; local endpoint existence
is the separate proved theorem (PG10)--(PG12).

\begin{proof}
Fix $\epsilon>0$ and put $a^\epsilon=\max(a,-\epsilon)$.
This bounded function belongs to $\mathcal E^1$ by solidity.
For integers $j>\epsilon$, write $a_j=\max(a,-j)$ and
let $Q_j$ use the same canonical level $j$ in its other slots.
The difference $a_j-a^\epsilon$ is supported in the one
compact set $K_\epsilon=\overline{\{a<-\epsilon\}}$.
The bounded compact Green formula and clipping locality give
\[
 \int_{\{a<-\epsilon\}}\Gamma_{Q_j}(a_j)
                =\int(a^\epsilon-a_j)D(a_jQ_j).
 \tag{GE8}
\]
Insert a fixed original compact cutoff equal to one near
$K_\epsilon$ when reading this identity. The endpoint graded
energy convergence (PG11) passes its left side in local total
variation. Its indicator is independent of $j>\epsilon$.
The weighted curvature convergence (PG4) with $a=1,k=q$
passes its right side, since $a+k=q+1=q+\sigma$ for
$\sigma=1$. Include the bounded slot $a^\epsilon$ in the
fixed family; its canonical truncation is itself for large $j$.
Thus
\[
 \int_{\{a<-\epsilon\}}\Gamma_Q(a)
           =\int(a^\epsilon-a)M(a,c_2,\ldots,c_q).
 \tag{GE9}
\]
The energy gives zero mass to the joint pole set by (PG12).
It also gives no mass to $\{a=0\}$: the corresponding bounded
energies vanish there, and (PG11) is local total-variation
convergence on the identical set $\{a_j=0\}=\{a=0\}$.
Let $\epsilon\downarrow0$ in (GE9). Both sides increase,
and the right side has finite limit by (PE11). This proves
the global diagonal identity in (GE7).

Apply it to $a+b$ as well. The full symmetry from the
preceding theorem and polarization give the cross identity.
The positive local energy matrix and an increasing compact
exhaustion give the global Cauchy--Schwarz inequality and the
total-variation bound. No convergence of global unweighted
energy masses was assumed in order to obtain their finiteness.
\end{proof}

\begin{theorem}[The constant-one first-energy Holder inequality]
\label{thm:r11-optimal-first-energy-Holder}
For all $u_0,\ldots,u_q\in\mathcal E^1$,
\[
 \boxed{\quad
 \mathcal I_S(u_0,\ldots,u_q)
                    \leq\prod_{i=0}^q e_1(u_i)^{1/(q+1)}.
 \quad}
 \tag{GE10}
\]
In particular
\[
 e_1(u+v)^{1/(q+1)}
           \leq e_1(u)^{1/(q+1)}+e_1(v)^{1/(q+1)}.
 \tag{GE11}
\]
The constant one in (GE10) is optimal whenever a potential
of positive first energy exists. This is an integral inequality;
it makes no pointwise Alexandrov--Fenchel or hyperbolicity
claim for the arbitrary middle-positive background.
\end{theorem}

\noindent\textbf{Status: \texttt{NEW\_PROVED}}. The earlier
(PE11) supplied a finite constant for every power and is used
to justify integrability before this sharper first-power proof.

\begin{proof}
The zero-energy case already follows from (PE11), including
its proved zero-factor limiting argument. Otherwise rescale
each potential so that $e_1(u_i)=1$. Put $d=q+1$.
Among the finitely many mixed values
\[
 B_\alpha=\mathcal I_S(
             \underbrace{u_0,\ldots,u_0}_{\alpha_0},\ldots,
             \underbrace{u_q,\ldots,u_q}_{\alpha_q}),
       \qquad \alpha_i\in\mathbf N,\quad\sum_i\alpha_i=d,
\]
choose a largest value $M$, and among its indices choose one
maximizing $\sum_i\alpha_i^2$. If this index has two positive
entries $\alpha_i,\alpha_l$, the two-slot energy
Cauchy--Schwarz inequality (GE7), with all remaining slots
fixed, gives
\[
       M^2=B_\alpha^2
           \leq B_{\alpha+e_i-e_l}B_{\alpha-e_i+e_l}
           \leq M^2.
\]
Since a diagonal value is one, $M\geq1>0$, and both
neighboring values must equal $M$. At least one neighboring
index has a strictly larger sum of squares: the two changes
are $2(\alpha_i-\alpha_l)+2$ and its opposite plus $2$.
This contradicts the choice. Thus a maximizing index is
pure, and $M=e_1(u_i)=1$. Every mixed value is at most one,
proving (GE10) after undoing the normalization.

Expand $e_1(u+v)$ by the fully symmetric multilinearity of
(GE5), and apply (GE10) term by term. The binomial theorem
gives (GE11). If all slots equal one potential with positive
energy, equality holds in (GE10), proving optimality.
\end{proof}

\begin{proposition}[The relative capacity quotient determines every original product]
\label{prop:r11-capacity-representative-quotient}
Let $u_i,v_i\in\mathcal F_{S,\partial}(V)$ for a finite
family of indices, and suppose
\[
       C_S(\{u_i\ne v_i\};V)=0\quad\hbox{for every }i.
 \tag{GE12}
\]
Then every mixed original product of the $u_i$ is equal to
the corresponding product of the $v_i$, in every degree
$0\leq k\leq q$, against every original complementary test
and after every actual moment refinement. For one pair,
\[
 \mu_u=\mu_v
  \quad\Longleftrightarrow\quad C_S(\{u\ne v\};V)=0
  \quad\Longleftrightarrow\quad C_S^*(\{u\ne v\};V)=0.
 \tag{GE13}
\]
If one member of a pair in (GE13) belongs to
$\mathcal E^\sigma_{S,\partial}$, so does the other, with
the same power energy. All defined graded energy currents
and integrable weighted products agree after these replacements.
\end{proposition}

\noindent\textbf{Status: \texttt{NEW\_PROVED}} as a
representative quotient within the specified actual-potential
classes. It is not the analytic PS/PL relation quotient GQ.

\begin{proof}
Let $E$ be the finite union of the exceptional Borel sets in
(GE12); it has zero raw capacity. At each canonical level $j$,
the bounded differences $u_{i,j}-v_{i,j}$ vanish outside $E$.
Every bounded mixed background $Q_j$ over $S$ gives $E$ zero
variation in each original coefficient by the global
coefficient bound (DG3), including degree zero. Therefore
$(u_{i,j}-v_{i,j})Q_j=0$ as an original current. The bounded
telescoping formula
\[
 \prod_{i=1}^k Du_{i,j}\wedge S-
 \prod_{i=1}^k Dv_{i,j}\wedge S
       =\sum_{i=1}^k D\bigl((u_{i,j}-v_{i,j})Q_{i,j}\bigr)
 \tag{GE14}
\]
is zero; $Q_{i,j}$ uses the already replaced slots before
$i$ and the unreplaced slots after it. Pass the uncut local
coefficient total-variation limits from the all-degree theorem.
This proves product equality without assuming a locality
principle on an arbitrary equality Borel set.

For a single pair, equality of the top measures gives both
one-sided outer-capacity domination statements from
Corollary~\ref{cor:r11-boundary-outer-domination}. Their
boundary-gap hypotheses hold because both potentials are
negative and have compact negative sublevels. Countable
subadditivity gives outer capacity zero of their inequality
set. Outer zero implies raw zero, and the proved product
equality gives the remaining implication in (GE13).

All products annihilate $E$ in their coefficient variations,
again by (DG3) and the increasing cut construction. Thus
the actual multiplier values of the two families agree
almost everywhere for each such coefficient. This preserves
integrability, power energies and weighted products.
The graded energy formula (PG10) then proves energy-current
equality wherever its degree hypotheses are satisfied.
\end{proof}

\begin{theorem}[Global Aubin--Mabuchi variation and energy rigidity]
\label{thm:r11-global-Aubin-Mabuchi}
For $u\in\mathcal E^1$ put
\[
             \mathscr A_S(u)=\frac1{q+1}\int_V u\,d\mu_u
                              =-\frac{e_1(u)}{q+1}.
 \tag{GE15}
\]
For $u,v\in\mathcal E^1$, let $h=v-u$ and
$u_t=(1-t)u+tv$, $0\leq t\leq1$. Then $\mathscr A_S(u_t)$
is a polynomial and
\[
 \begin{split}
 \frac d{dt}\mathscr A_S(u_t)&=\int_V h\,d\mu_{u_t},\\
 \frac {d^2}{dt^2}\mathscr A_S(u_t)
              &=-q\int_V\Gamma_{N(u_t^{q-1})}(h)\leq0.
 \end{split}
 \tag{GE16}
\]
All displayed signed integrals are absolutely convergent.
The functional is increasing under pointwise order, and
\[
 \int(v-u)\,d\mu_v
    \leq\mathscr A_S(v)-\mathscr A_S(u)
    \leq\int(v-u)\,d\mu_u.
 \tag{GE17}
\]
The symmetric discrepancy
\[
 \begin{split}
 \mathscr I_S(u,v)
    &:=\int_V(u-v)(d\mu_v-d\mu_u)\\
    &=\sum_{l=0}^{q-1}
       \int_V\Gamma_{N(u^l,v^{q-1-l})}(u-v)\geq0
 \end{split}
 \tag{GE18}
\]
vanishes exactly when the equivalent conditions (GE13) hold.
Consequently $\mathscr A_S$ is strictly concave on the
relative capacity quotient: equality with its chord at one
strict interior parameter holds exactly for equivalent
potentials. No pointwise separation or positivity on every
ordinary open set is assumed for the background.
\end{theorem}

\noindent\textbf{Status: \texttt{NEW\_PROVED}} as the
global unbounded energy extension and quotient rigidity.
The bounded common-collar variation formula itself is an
\texttt{INTERNAL\_INPUT}, as specified above.

\begin{proof}
Full symmetric multilinearity of (GE5) makes the energy a
polynomial along the segment. Differentiating its $q+1$
equal slots gives the first formula in (GE16). Differentiating
once more gives
$q\int hD(hN(u_t^{q-1}))$. The global energy matrix identity
(GE7), extended by bilinearity to the difference $h$, is
exactly its negative in (GE16). Finiteness follows from
(GE10) in each expanded mixed term. Concavity follows from
matrix positivity, and the tangent inequalities (GE17) follow
by integrating the decreasing first derivative. If $u\leq v$
the first derivative is nonnegative, proving monotonicity.

The telescoping identity for $\mu_v-\mu_u$ and the same
global Green identity give (GE18). Suppose its sum is zero.
Every summand is nonnegative, so for each
$Q_l=N(u^l,v^{q-1-l})$ one has
$\Gamma_{Q_l}(h)(V)=0$.
For an original compact smooth scalar test $\varphi$, choose
an original compact cutoff $\chi$ equal to one on its
differential support, and use canonical bounded approximants
$h_j=v_j-u_j$ and $Q_{l,j}$. Compact integration by parts
and energy Cauchy--Schwarz give
\[
 |\langle D(h_jQ_{l,j}),\varphi\rangle|^2
       \leq\left(\int\chi\Gamma_{Q_{l,j}}(h_j)\right)
             \left(\int\chi\Gamma_{Q_{l,j}}(\varphi)\right).
 \tag{GE19}
\]
The first factor tends to zero by the endpoint local
energy total-variation convergence (PG11). The second is
bounded: its scalar differentials are fixed original smooth
forms, and $Q_{l,j}$ has local coefficient total-variation
convergence. The weighted convergence (PG4) also gives
$h_jQ_{l,j}\to hQ_l$, so (GE19) proves $D(hQ_l)=0$
against every original test. Summing the recursive telescoping
formula yields $\mu_u=\mu_v$. Conversely this equality
makes the defining expression in (GE18) zero. Proposition
\ref{prop:r11-capacity-representative-quotient} gives precisely
the other equivalences.

If the energy agrees with its chord at one interior point,
concavity forces it to be affine on the whole segment.
Its endpoint derivatives are equal, so (GE16) gives
$\mathscr I_S(u,v)=\mathscr A_S'(0)-\mathscr A_S'(1)=0$.
Conversely equivalent potentials have identical mixed products
and energies by the preceding proposition, and the energy is
constant on their segment. This proves the claimed strict
concavity on the quotient.
\end{proof}

\begin{corollary}[The variational necessary condition retains its domain gate]
\label{cor:r11-variational-necessary-condition}
Let $\mu$ be a positive Radon datum such that
$\int(-v)\,d\mu<\infty$ for every $v\in\mathcal E^1$.
If $u\in\mathcal E^1$ solves $\mu_u=\mu$, then it maximizes
\[
                   v\longmapsto\mathscr A_S(v)-\int v\,d\mu
 \tag{GE20}
\]
on $\mathcal E^1$. Conversely, a maximizer satisfies the
current equation provided that for every original compact
smooth scalar $\varphi$ there is $\epsilon_\varphi>0$
such that $u+t\varphi\in\mathcal E^1$ for all
$|t|<\epsilon_\varphi$. This last variation hypothesis is
explicit and is not asserted for a general classical psh
potential or arbitrary datum.
\end{corollary}

\noindent\textbf{Status: \texttt{DIRECT\_COROLLARY}} for
necessity and \texttt{CONDITIONAL} for the stated converse.

\begin{proof}
The upper tangent inequality in (GE17) gives necessity.
Under the printed two-sided variation hypothesis, apply
the first derivative identity to the admissible segment
$u+t\varphi$. Maximality for both signs gives
$\int\varphi\,d\mu_u=\int\varphi\,d\mu$ for every
original compact smooth scalar test. This is exactly the
original current equation. No envelope differentiability or
unstated admissibility of arbitrary smooth perturbations is
used.
\end{proof}

\paragraph{Scope and rejected shortcut.}
The tempting inference from first-energy dual boundedness to
an unrestricted ambient classical Dirichlet solution remains
\texttt{FALSE}, by Corollary~\ref{cor:r11-dual-energy-insufficient}.
Concavity alone supplies no compactness or representability
theorem for a maximizing family. The two-sided variation gate
in the last corollary cannot be discarded by analogy with a
smooth strictly elliptic problem. The relative capacity quotient
in (GE12)--(GE14) identifies actual potentials in the proved
boundary classes; it does not identify analytic PS/PL
presentations, does not enlarge the class of delta tests, and
does not close GQ or low-real-degree CTO.

One possible alternate approximation is
$\max(u_i-\epsilon_j,j\rho)$, which imposes a common ordinary
collar. That route requires separate uniform control of the
small-value energy near the moving collar. The proof above
instead uses (GE3) and the already proved decreasing finite-mass
approximation; no claim about the alternate route is needed.

% END REVIEWED MODULE research_r11/comparison/global_first_energy_calculus.tex


% BEGIN REVIEWED MODULE research_r11/comparison/first_energy_capacity_topology.tex
\section{Quantitative first-energy convergence in every local capacity}
\label{sec:r11-first-energy-capacity-topology}

Retain the strict analytic box, background extension and
$\mathcal E^1$ hypotheses of
Section~\ref{sec:r11-global-first-energy}.
All functions below are actual negative classical members
of that class; all capacities use the full original tester
class. Put $d=q+1$ and $\alpha_q=2^{-q}$.
The symmetric discrepancy $\mathscr I_S$ is (GE18).
At a joint pole, a difference is used only as a Borel
coefficient representative; all involved products and
capacities annihilate that pole union.

The bounded common-collar predecessors are the frozen R3
\path{energy_convex_blocks.tex},
\texttt{thm:r3-I-capacity}, and
\path{strong_relative_topology.tex}.
They already contain the positive discrepancy, the mixed
energy transfer recurrence, and bounded relative energy
topologies. The new scope here is the global unbounded
first-energy class, uniform control over its energy testers,
the analytic all-local domain extension, and critical
endpoint continuity without a prescribed common collar.

\begin{theorem}[A uniform first-energy tester estimate]
\label{thm:r11-first-energy-uniform-tester}
For $u,v,\psi\in\mathcal E^1$, put $w=u+v$ and
$M=\max\{e_1(w),e_1(\psi)\}$. Then
\[
 \boxed{\quad
 \int_V|u-v|\,d\mu_\psi
   \leq 2^{(q-1)\alpha_q}
          \mathscr I_S(u,v)^{\alpha_q}M^{1-\alpha_q}.
 \quad}
 \tag{EV1}
\]
The conclusion is zero if $M=0$. For a mixed top measure
in place of $\mu_\psi$, the positive expansion of the sum
of its $q$ potentials gives the same estimate divided by
$q!$, with that sum as $\psi$.
\end{theorem}

\noindent\textbf{Status: \texttt{NEW\_PROVED}} in the
unbounded global first-energy class.

\begin{proof}
Write $h=u-v$. For $0\leq k\leq q-1$ set
\[
 J_k=\int_V\Gamma_{N(\psi^{[k]},w^{[q-1-k]})}(h).
 \tag{EV2}
\]
These are finite nonnegative numbers by (GE7). The positive
expansion of $Dw=Du+Dv$ and (GE18) give
$J_0\leq2^{q-1}\mathscr I_S(u,v)$.
For $k\leq q-2$, full mixed energy symmetry exchanges one
$\psi$ slot with one $h$ slot. Expanding that latter
$h=u-v$ in the background gives
\[
 \begin{split}
 J_{k+1}
 &=\int\Gamma_{N(u,\psi^{[k]},w^{[q-2-k]})}(h,\psi)
   -\int\Gamma_{N(v,\psi^{[k]},w^{[q-2-k]})}(h,\psi)\\
 &\leq\left[J_k\,
       \mathcal I_S(\psi^{[k+2]},w^{[q-1-k]})\right]^{1/2}
 \leq(J_kM)^{1/2}.
 \end{split}
 \tag{EV3}
\]
For the middle inequality apply energy Cauchy--Schwarz in
each of the two positive backgrounds and then Cauchy--Schwarz
to the resulting two nonnegative summands. Their first
diagonal masses add to $J_k$, and their second diagonal
masses add to the displayed mixed energy. The last inequality
is the constant-one Holder inequality (GE10). Thus
\[
 J_{q-1}\leq J_0^{2^{-(q-1)}}M^{1-2^{-(q-1)}}.
 \tag{EV4}
\]
When $q=1$ this is just the identity $J_{q-1}=J_0$.

The real-span element
$|h|=2\max(u,v)-u-v$ has the exact energy identity
$\Gamma_Q(|h|)=\Gamma_Q(h)$ for every background $Q$
in (EV2). To justify it in the present class, use canonical
bounded truncations of $u,v,\max(u,v)$ and the background
slots. At level $j$ one has
$\max(u_j,v_j)=\max(\max(u,v),-j)$.
The bounded absolute-value chain and the zero energy of
the level $\{u_j=v_j\}$ prove the identity at that stage.
Endpoint local total-variation convergence (PG11) passes
it to the unbounded currents. Both are already globally
finite by (GE7).

Global Green, followed by Cauchy--Schwarz, now yields
\[
 \int|h|\,d\mu_\psi
  =-\int\Gamma_{N(\psi^{[q-1]})}(|h|,\psi)
       \leq J_{q-1}^{1/2}e_1(\psi)^{1/2}.
 \tag{EV5}
\]
Combine this with (EV4), the bound on $J_0$, and
$e_1(\psi)\leq M$. If $M=0$, the mixed Holder estimate
already makes the integral zero. The mixed-measure
assertion follows from
$q!M(\psi_1,\ldots,\psi_q)\leq\mu_{\sum_i\psi_i}$.
\end{proof}



\begin{theorem}[Quantitative capacity control on every compact/domain pair]
\label{thm:r11-first-energy-all-local-capacity}
Fix a compact $K\Subset V$ and choose $A_K>0$ so that
$A_K\rho<-2$ on a neighborhood of $K$. Set
\[
 B_K=2A_K^q\mu_\rho(V),\qquad
 \Lambda_K(u,v)=\max\{e_1(u+v),B_K\}.
 \tag{EV6}
\]
For $\delta>0$ put
$G_\delta(u,v)=\{u<v-\delta\}\cup\{v<u-\delta\}$,
using actual extended comparisons. Then
\[
 C_S(K\cap G_\delta(u,v);V)
    \leq\frac{2^{(q-1)\alpha_q}}\delta
          \Lambda_K(u,v)^{1-\alpha_q}
                      \mathscr I_S(u,v)^{\alpha_q}.
 \tag{EV7}
\]
If $K\subset\operatorname{Int}_X K'\Subset V$, the same
bound for $K'$ bounds the outer capacity of
$K\cap G_\delta(u,v)$.

Consequently for every ordinary $\Omega\subset V$ and
compact $K\Subset\Omega$, and every $M_0<\infty$, there
is a finite constant $C_{K,\Omega,M_0,S}$ such that
$e_1(u),e_1(v)\leq M_0$ implies
\[
 C_{S|\Omega}^*(K\cap G_\delta(u,v);\Omega)
       \leq\frac{C_{K,\Omega,M_0,S}}\delta
                            \mathscr I_S(u,v)^{\alpha_q}.
 \tag{EV8}
\]
This is the complete all-local capacity topology, by the
proved analytic affinoid domain comparison, not an assumed
identification of two admissibility domains.
\end{theorem}

\noindent\textbf{Status: \texttt{NEW\_PROVED}} as a
global-energy estimate for all original local testers.

\begin{proof}
For a normalized tester $0\leq a\leq1$ on $V$ put
$\psi=\max(a-2,A_K\rho)$.
It is bounded in $[-2,0]$, has compact negative sublevels,
and agrees with $a-2$ near $K$.
Comparison with $cA_K\rho$ for $c>1$ gives
$\mu_\psi(V)\leq A_K^q\mu_\rho(V)$ after $c\downarrow1$.
Thus $\psi\in\mathcal E^1$ and $e_1(\psi)\leq B_K$,
uniformly in the original tester $a$. Ordinary locality
gives $\mathbf1_K\mu_a=\mathbf1_K\mu_\psi$.
The joint pole set has raw capacity zero, so it has zero
mass for $\mu_a$. On its complement, Chebyshev and (EV1)
give
\[
 \mu_a(K\cap G_\delta(u,v))
       \leq\delta^{-1}\int|u-v|\,d\mu_\psi
       \leq\frac{2^{(q-1)\alpha_q}}\delta
             \Lambda_K(u,v)^{1-\alpha_q}
                          \mathscr I_S(u,v)^{\alpha_q}.
\]
Take the full tester supremum to prove (EV7).

To obtain the outer bound, fix $\eta>0$. Outer
quasicontinuity on $K'$ and outer nullity of the poles
give an ordinary open $O$ with $C_S(O;V)<\eta$ such
that both potentials are finite and continuous on
$K'\setminus O$. In the relative space
$\operatorname{Int}_X K'\setminus O$ the strict gap is
open. Hence there is an ordinary open $H\subset
\operatorname{Int}_X K'$ such that
$H\setminus O=G_\delta(u,v)\cap
(\operatorname{Int}_X K'\setminus O)$.
The open set $H\cup O$ covers $K\cap G_\delta$ and is
contained in $(K'\cap G_\delta)\cup O$. Therefore
\[
 C_S^*(K\cap G_\delta;V)
              \leq C_S(K'\cap G_\delta;V)+\eta.
\]
Let $\eta\downarrow0$ and use (EV7) on $K'$.

Finally choose $K'$ inside $\Omega$ and an ordinary open
core containing $K$ compactly in $\Omega$. The outer
domain comparison (AL9) multiplies this estimate by a fixed
geometric constant. Minkowski (GE11) gives
$e_1(u+v)\leq2^{q+1}M_0$, proving (EV8).
No general Borel raw/outer equality is used.
\end{proof}

\begin{theorem}[The strong first-energy topology]
\label{thm:r11-strong-first-energy-topology}
For any $u,v\in\mathcal E^1$,
\[
 \left|e_1(u)^{1/d}-e_1(v)^{1/d}\right|^d
                                      \leq\mathscr I_S(u,v).
 \tag{EV9}
\]
For a sequence $u_j\in\mathcal E^1$ and an actual
$u\in\mathcal E^1$, the following are equivalent:
\begin{enumerate}
\item $\mathscr I_S(u_j,u)\to0$.
\item $u_j\to u$ in all-local truncated original capacity
and $e_1(u_j)\to e_1(u)$.
\end{enumerate}
The forward implication gives outer-capacity convergence
of the full actual gaps on every compact/domain pair.
This theorem describes convergence to an already specified
actual class member; it does not assert classical
representability of every Cauchy sequence.
\end{theorem}

\noindent\textbf{Status: \texttt{NEW\_PROVED}} for the
unbounded class. The bounded fixed-collar energy topology
was already frozen in R3.

\begin{proof}
Put $a=e_1(u)^{1/d}$ and $b=e_1(v)^{1/d}$.
Expand the definition (GE18), then use the constant-one
mixed Holder inequality on its two cross terms:
\[
 \begin{split}
 \mathscr I_S(u,v)
 &=e_1(u)+e_1(v)
       -\mathcal I_S(u,v^{[q]})-\mathcal I_S(v,u^{[q]})\\
 &\geq a^d+b^d-ab^{d-1}-ba^{d-1}
       =(a-b)(a^{d-1}-b^{d-1})\geq|a-b|^d.
 \end{split}
 \tag{EV10}
\]
The last inequality follows by the binomial expansion after
ordering $a,b$. Thus (1) implies energy convergence and a
uniform energy bound. The capacity estimate (EV8) gives
all the asserted forward capacity conclusions.

Conversely suppose (2). In particular the energies are
uniformly bounded. For each bounded weight truncation and
each original compact cutoff, the proved local bounded-weight
capacity continuity gives
\[
 \begin{split}
 \liminf_j\mathcal I_S(u_j,u^{[q]})&\geq e_1(u),\\
 \liminf_j\mathcal I_S(u,u_j^{[q]})&\geq e_1(u).
 \end{split}
 \tag{EV11}
\]
Here is the order of limits: first replace the
undifferentiated weight by $-\max(\cdot,-L)$ and insert a
compact nonnegative cutoff; the original bounded-weight
continuity passes $j\to\infty$. Then let the cutoff
increase to one and $L\to\infty$, using positivity.
The derivative slots are controlled by the uniform power
energy hypothesis through the already proved top capacity
stability theorem. Thus neither line presupposes critical
weighted continuity.

The optimal Holder upper bounds are respectively
$e_1(u_j)^{1/d}e_1(u)^{q/d}$ and
$e_1(u)^{1/d}e_1(u_j)^{q/d}$, both tending to $e_1(u)$.
Thus both cross integrals converge to $e_1(u)$.
The first line of (EV10) now proves (1).
\end{proof}

\begin{theorem}[Critical mixed energy and weighted-current continuity]
\label{thm:r11-strong-critical-mixed-continuity}
Let $u_{i,j},u_i\in\mathcal E^1$ for $0\leq i\leq q$,
and assume $\mathscr I_S(u_{i,j},u_i)\to0$ for every slot.
Then
\[
 \mathcal I_S(u_{0,j},\ldots,u_{q,j})
                  \longrightarrow\mathcal I_S(u_0,\ldots,u_q),
 \tag{EV12}
\]
and the finite positive measures
\[
 (-u_{0,j})M(u_{1,j},\ldots,u_{q,j})
                      \longrightarrow(-u_0)M(u_1,\ldots,u_q)
 \tag{EV13}
\]
converge against every bounded original smooth scalar test
on $V$, including tests without compact support. In
particular there is neither a critical pole residue nor
escaped boundary energy. No common pointwise lower barrier
is required in this statement.
\end{theorem}

\noindent\textbf{Status: \texttt{NEW\_PROVED}}. The
strong discrepancy hypothesis is substantive; the critical
failure in (PS7) concerns mere capacity convergence with
a uniform energy bound.

\begin{proof}
By (EV9), every slot has uniformly bounded first energy.
Telescope the fully symmetric multilinear integral
$\mathcal I_S$ one slot at a time. Move each changing
slot to the undifferentiated position by (GE5).
The absolute value of that difference is bounded by
$\int|u_{i,j}-u_i|M(\mathbf z_j)$ for a mixed tuple
of unchanged or already changed slots. Its positive
expansion is at most $1/q!$ times the diagonal product
of their sum. That sum has uniformly bounded energy by
(GE11). Estimate (EV1) makes each of the finitely many
differences tend to zero. This proves (EV12).

Write $M_j=M(u_{1,j},\ldots,u_{q,j})$ and $M$ for its
limit. The same estimate gives
\[
                  \int|u_{0,j}-u_0|\,dM_j\longrightarrow0.
 \tag{EV14}
\]
For $u=u_0$, put $u^R=\max(u,-R)$. Since $u^R-u\geq0$,
the nonnegative discrepancy satisfies
\[
 0\leq\mathscr I_S(u,u^R)
       =\int(u^R-u)(d\mu_u-d\mu_{u^R})
       \leq\int(u^R-u)\,d\mu_u\longrightarrow0.
 \tag{EV15}
\]
The last limit is dominated convergence against the finite
first-energy weight $(-u)\mu_u$.
Apply (EV1) to $u,u^R$ and to the uniformly bounded-energy
sum of the differentiated slots of $M_j$. This proves
$\sup_j\int|u-u^R|\,dM_j\to0$; the same holds for $M$.
The constants are uniform because $e_1(u^R)\leq e_1(u)$
by energy monotonicity and all other energies are bounded.

For fixed $R$, local bounded-weight capacity continuity
gives $(-u^R)M_j\to(-u^R)M$ on every original compact
scalar test. Use the preceding uniform weight tail and
(EV14) to pass $R\to\infty$. This proves local convergence
in (EV13). Their total masses converge by (EV12).
Positive finite Radon measures with local convergence
and convergent full masses are tight: choose a compact
original cutoff whose integral in the limit is arbitrarily
close to the full limit mass, and pass its local integral.
This controls the complements uniformly for all large $j$.
It follows that every bounded original smooth scalar test,
with no compact-support restriction, has the same limit.
\end{proof}

\begin{corollary}[Energy-dual maximizing sequences are strongly Cauchy]
\label{cor:r11-dual-maximizing-Cauchy}
Let $\mu$ be a positive Radon datum with
$D=\mathfrak d_1(\mu)<\infty$ as defined in (DO1).
For $\mathcal F_\mu(u)=\mathscr A_S(u)-\int u\,d\mu$,
\[
 0\leq\sup_{u\in\mathcal E^1}\mathcal F_\mu(u)
               \leq\frac q{q+1}D^{(q+1)/q}<\infty.
 \tag{EV16}
\]
Every maximizing sequence is uniformly bounded in first
energy and is Cauchy for $\mathscr I_S$ and in every local
outer capacity. This does not assert the existence of an
actual classical limit or a classical Dirichlet solution.
\end{corollary}

\noindent\textbf{Status: \texttt{NEW\_PROVED}} as a
compactness estimate in the unbounded nonlinear class;
the representation and equation gates remain explicit.

\begin{proof}
Positive homogeneity and the definition of $D$ give
$\int(-u)\,d\mu\leq D e_1(u)^{1/d}$.
If $e_1(u)=0$, scale $u$ by an arbitrary positive number
in (DO1) to obtain the same zero conclusion.
Thus, with $x=e_1(u)^{1/d}$,
\[
                \mathcal F_\mu(u)\leq-x^d/d+Dx.
\]
Maximizing this real polynomial gives (EV16), and it tends
to $-\infty$ as $x\to\infty$. Since $0$ is admissible,
every tail of a maximizing sequence has
$\mathcal F_\mu(u_j)\geq-1$ and hence
\[
          e_1(u_j)\leq
             \max\{2d,(2dD)^{d/(d-1)}\}.
 \tag{EV17}
\]
Discarding finitely many terms gives uniform energy for
the sequence.

The frozen bounded common-collar midpoint argument now
applies to the proved global variation (GE16), giving
\[
 \mathscr A_S\left(\frac{u+v}2\right)
       -\frac{\mathscr A_S(u)+\mathscr A_S(v)}2
            \geq\frac q{4^{q+1}}\mathscr I_S(u,v).
 \tag{EV18}
\]
For completeness the midpoint Green formula is
$\frac q2\int_0^1\min(t,1-t)
 \int\Gamma_{N(u_t^{q-1})}(u-v)\,dt$.
On $[1/4,3/4]$, positive mixed expansion bounds its
energy integrand below by $4^{-(q-1)}\mathscr I_S(u,v)$.
The interval length and the scalar weight give exactly
the constant in (EV18). All global energy integrals
are finite by (GE7).

The datum integral is affine, so it cancels in the midpoint
gap. If $L=\sup\mathcal F_\mu$, then
\[
 \frac q{4^{q+1}}\mathscr I_S(u_j,u_l)
        \leq L-\frac{\mathcal F_\mu(u_j)+
                              \mathcal F_\mu(u_l)}2\longrightarrow0.
\]
The uniform energy bound and (EV8) give the all-local
outer-capacity Cauchy conclusion. Producing an actual
classical representative of that completed class, and
valid admissible variations that identify its curvature
with the datum, are not proved by this argument. The
compact-source obstruction (DO9) prevents omitting those
gates in a universal classical existence assertion.
\end{proof}

\begin{corollary}[Strong first-energy convergence controls every graded endpoint]
\label{cor:r11-strong-first-energy-graded-endpoints}
Let $u_{i,j},u_i\in\mathcal E^1$ be a finite family with
$\mathscr I_S(u_{i,j},u_i)\to0$ in every slot. For
$w_j=\sum_i u_{i,j}$ one has
\[
 \lim_{R\to\infty}\sup_j\int_R^\infty
                 t^q C_S(\{w_j<-t\};V)\,dt=0.
 \tag{EV19}
\]
Consequently all three convergence assertions (GS3), and
the double-prime first-order assertion, hold with
$\sigma=1$ and weak inequalities in their degree/weight
conditions. They hold on every original compact
complementary test and every actual moment refinement,
without a common pointwise lower barrier. These are
original-current convergence assertions; total-variation
convergence for arbitrary varying tuples is not claimed.
\end{corollary}

\noindent\textbf{Status: \texttt{NEW\_PROVED}} as the
uniform-tail consequence of the strong topology; applying
the established endpoint theorem is a direct consequence.

\begin{proof}
First consider one sequence $u_j\to u$ in discrepancy.
Put $a_j=-u_j$ and $b=-u$ off the joint polar set.
The uniform energy bound follows from (EV9). On
$\{a_j>R\}$ split according as $b\leq a_j/2$ or
$b>a_j/2$. In the first region $a_j\leq2|u_j-u|$;
in the second, $a_j<2b$ and $b>R/2$.
For $b>R/2$ one has
$b\leq2(b-R/4)^+=2|u-u^{R/4}|$.
Thus
\[
 \int_{\{a_j>R\}}a_j\,d\mu_{u_j}
    \leq2\int|u_j-u|\,d\mu_{u_j}
             +4\int|u-u^{R/4}|\,d\mu_{u_j}.
 \tag{EV20}
\]
By (EV1), the first term tends to zero as $j\to\infty$.
By (EV1) and (EV15), the second tends to zero as
$R\to\infty$ uniformly in $j$. To obtain the supremum
over all $j$, first make the first term small for
$j\geq j_0$; the finitely many earlier rows each have
integrable first-energy weight, so their tails tend
to zero separately. Then let $R\to\infty$, and finally
$j_0\to\infty$. This proves uniform vanishing of the
left side of (EV20).

The sublevel capacity bound (PE4), before its full
integration, gives for every row
\[
 \int_R^\infty t^q C_S(\{u_j<-t\};V)\,dt
       \leq2^{q+1}\int(a_j-R/2)^+\,d\mu_{u_j}.
 \tag{EV21}
\]
Its right side is bounded by a constant times the
uniformly vanishing energy tail at $R/2$.
This proves (EV19) for one row.

For $m$ slots, the actual sublevel inclusion
$\{w_j<-t\}\subset\bigcup_{i=1}^m\{u_{i,j}<-t/m\}$
and subadditivity imply that its capacity-tail integral
is at most $m^{q+1}$ times the sum of the individual
tail integrals starting at $R/m$. This proves (EV19)
for the family. The full local capacity convergence is
already supplied by (EV8), so every hypothesis of
Theorem~\ref{thm:r11-nonmonotone-graded-stability}, including
its endpoint condition (GS4), now holds at $\sigma=1$.
Applying that theorem proves the current conclusions.
\end{proof}

% END REVIEWED MODULE research_r11/comparison/first_energy_capacity_topology.tex


% BEGIN REVIEWED MODULE research_r11/energy/sharp_midpoint_modulus.tex
\section{The sharp midpoint modulus of the first energy}
\label{sec:r11-sharp-midpoint-modulus}

Retain exactly the strict-box, background-extension and
$\mathcal E^1_{S,\partial}(V)$ hypotheses of the global
first-energy calculus.  Put
\[
 c_q=\frac{1-2^{-q}}{2(q+1)},\qquad
 \mathscr D_S(u,v)=
       \mathscr A_S\left(\frac{u+v}{2}\right)
            -\frac{\mathscr A_S(u)+\mathscr A_S(v)}2 .
 \tag{SM1}
\]
The simpler positive constant in (EV18) already suffices
for the Cauchy argument there.  This section computes
the optimal uniform coefficient and its equality example.

\begin{lemma}[A finite Bernstein coefficient comparison]
\label{lem:r11-midpoint-Bernstein}
For every integer $q\geq1$ and $0\leq l\leq q-1$,
\[
 a_{q,l}:=
 \frac q2\binom{q-1}{l}
    \int_0^1\min(t,1-t)t^l(1-t)^{q-1-l}\,dt
       \geq c_q.
 \tag{SM2}
\]
Equality holds for $l=0$ and $l=q-1$.
\end{lemma}

\noindent\textbf{Status: \texttt{NEW\_PROVED}}
auxiliary finite polynomial inequality; the following
proof holds for every $q$, independently of any numerical
checks.

\begin{proof}
For $1\leq k\leq q$ write
\[
\begin{split}
 p_i(s)&=\binom qi s^i(1-s)^{q-i},\\
 B_k(s)&=\sum_{i=k}^q p_i(s),\\
 b_k(t)&=q\binom{q-1}{k-1}t^{k-1}(1-t)^{q-k}.
\end{split}
 \tag{SM3}
\]
Termwise differentiation of the finite sum telescopes
to $B_k'=b_k$.  The integral of $b_k$ over $[0,1]$
is one.  Also
$B_k(1-s)=1-B_{q-k+1}(s)$, by the binomial theorem.
Consequently
\[
 P_k(s):=\int_s^{1-s}b_k(t)\,dt
              =1-B_k(s)-B_{q-k+1}(s)
                    \quad(0\leq s\leq1/2).
 \tag{SM4}
\]
For $1\leq k\leq(q+1)/2$ this gives
\[
 P_k(s)-P_1(s)
       =\sum_{i=1}^{k-1}\bigl(p_i(s)-p_{q-i}(s)\bigr)\geq0.
 \tag{SM5}
\]
Indeed $i\leq(q-1)/2$, and for $0<s\leq1/2$
the ratio is
$p_i(s)/p_{q-i}(s)=(s/(1-s))^{2i-q}\geq1$.
At $s=0$ use continuity.
The symmetry $P_k=P_{q+1-k}$ proves the same
bound for the remaining $k$.

Layer cake for the nonnegative tent function now gives
\[
 a_{q,k-1}
   =\frac12\int_0^{1/2}P_k(s)\,ds
   \geq\frac12\int_0^{1/2}P_1(s)\,ds=a_{q,0}.
 \tag{SM6}
\]
The endpoint integral, or its reflected version, is
\[
 \frac q2\left[
      \int_0^{1/2}t^q\,dt
        +\int_{1/2}^1(1-t)t^{q-1}\,dt\right]
       =\frac{1-2^{-q}}{2(q+1)}.
 \tag{SM7}
\]
This proves every assertion, including $q=1$.
\end{proof}

\begin{theorem}[Optimal midpoint control of the relative discrepancy]
\label{thm:r11-sharp-midpoint-modulus}
For every $u,v\in\mathcal E^1_{S,\partial}(V)$,
\[
 \boxed{\qquad
       \mathscr D_S(u,v)\geq
          \frac{1-2^{-q}}{2(q+1)}\,\mathscr I_S(u,v).
       \qquad}
 \tag{SM8}
\]
If the class contains a potential of positive first
energy, the constant cannot be increased for that
same domain and background.
No strict positivity or geometric representation of
the original middle-positive $S$ is required.
\end{theorem}

\noindent\textbf{Status: \texttt{NEW\_PROVED}}
as the sharp quantitative form of the global
arbitrary-background midpoint estimate.

\begin{proof}
Let $h=u-v$ and set
\[
 \gamma_l=
       \int_V\Gamma_{N(u^{[l]},v^{[q-1-l]})}(h)
                      \quad(0\leq l\leq q-1).
 \tag{SM9}
\]
Each $\gamma_l$ is finite and nonnegative by (GE7),
and (GE18) gives
$\mathscr I_S(u,v)=\sum_l\gamma_l$.
The scalar Green formula for a twice differentiable
function on $[0,1]$, applied to
$\mathscr A_S((1-t)u+tv)$ and (GE16), yields
\[
 \mathscr D_S(u,v)=
       \frac q2\int_0^1\min(t,1-t)
          \int_V\Gamma_{N(((1-t)u+tv)^{[q-1]})}(h)\,dt
       =\sum_{l=0}^{q-1}a_{q,l}\gamma_l .
 \tag{SM10}
\]
The possible reversal of $l$ under $t\mapsto1-t$
does not affect the coefficient, since the tent
function is symmetric.
All operations are finite sums of the already
proved globally integrable mixed energies.
Apply (SM2) term by term to prove (SM8).

For sharpness choose any $u$ with $e_1(u)>0$ and
take $v=0$.  The zero potential is in the same
boundary power class.  Homogeneity and the definition
of the discrepancy give
\[
 \mathscr I_S(u,0)=e_1(u),\qquad
 \mathscr D_S(u,0)
    =\left(\frac12-2^{-(q+1)}\right)
                       \frac{e_1(u)}{q+1}
    =c_q e_1(u).
 \tag{SM11}
\]
Thus equality occurs for the same original background
and domain, proving optimality.
\end{proof}

\begin{corollary}[An explicit error bound for dual maximization]
\label{cor:r11-sharp-maximizing-error}
In the setting of (EV16), let
$L=\sup\mathcal F_\mu$.
For any $u,v\in\mathcal E^1$,
\[
 \mathscr I_S(u,v)
       \leq \frac{2(q+1)}{1-2^{-q}}\,
                \left[L-\frac{\mathcal F_\mu(u)+\mathcal F_\mu(v)}2\right].
 \tag{SM12}
\]
Thus prescribed objective gaps give a quantitative
bound in every local outer capacity through (EV8).
This remains a convergence estimate; it supplies
neither an ambient classical representative of
a completion limit nor the missing admissible
variations for an unrestricted source equation.
\end{corollary}

\noindent\textbf{Status: \texttt{DIRECT\_COROLLARY}}.

\begin{proof}
The datum integral is affine along the segment,
so it cancels from the midpoint gap.
The value at the midpoint is at most $L$.
Use (SM8) and divide by its positive constant.
\end{proof}

\paragraph{Auxiliary verification.}
The exact rational script
\path{research_r11/verify_midpoint_coefficients.py}
expands the two polynomial integrals for a finite
range of $q$ and checks their minimum against (SM7).
It tests normalization and signs only.
The all-$q$ argument is (SM3)--(SM7), and the analytic
argument is the full original global energy identity
(SM10).

% END REVIEWED MODULE research_r11/energy/sharp_midpoint_modulus.tex


% BEGIN REVIEWED MODULE research_r11/capacity/barrier_relative_completion.tex
\section{A relative completion bounded below by an actual power-energy barrier}
\label{sec:r11-barrier-relative-completion}

Assume exactly the strict analytic box and extension hypotheses
(TC1)--(TC2), with the fixed arbitrary original symmetric
middle-positive separately closed $(p,p)$-current $S$ and
$q=n-p\geq1$. Fix $\sigma>0$ and put $\kappa=q+\sigma$.
Write $C^*=C_S^*(\,\cdot\,;V)$. The original classical
classes $\mathcal F_{S,\partial}$ and
$\mathcal E^\sigma_{S,\partial}$ are unchanged.
This section defines a larger class of actual Borel potentials
modulo local outer-capacity equality. It does not assert that
its members are ambient upper semicontinuous or classical psh.
The background is never approximated or replaced.

\subsection{The bounded input and the exact completion domain}

\begin{lemma}[Analytic bounded capacitary calculus]
\label{lem:r11-analytic-bounded-completion}
On the present original analytic domain, let
$\mathcal A_S^{[-R,0]}(V)$ denote the actual Borel limits,
locally in ordinary-open outer $S$-capacity, of bounded
classical functions on $V$ with values in $[-R,0]$.
The bounded completion, its quasicontinuity, and its original
recursive weighted BT calculus hold on this class.
The same is true of bounded first-order and energy currents,
with their intrinsic distributional definitions.
All these constructions are independent of defining sequences
and of local outer-capacity-null changes of representatives.

In particular, for a fixed original compact complementary test
$\eta$, a bounded completed curvature product $P_R$ of degree
$p+k$ satisfies
\[
                 |\mu_{P_R,\eta}|(E)
                        \leq C_\eta R^k C^*(E)
                    \quad(E\subset V\text{ Borel}),
 \tag{RC1}
\]
with the evident separate degree-zero bound. A bounded completed
first-order current has the corresponding power $R^{k+1}$,
and a bounded completed energy current has power $R^{k+2}$.
The constants on a fixed stage are independent of $R$; all dependence
on the common truncation bound is in the displayed powers of $R$.
They depend only on the background, degree and that stage as in
(DG3), (FO4), and (PG8).
\end{lemma}

\noindent\textbf{Status: \texttt{DIRECT\_COROLLARY}} of
the frozen R2 bounded completion, with its analytic hypotheses
supplied explicitly by the new affinoid localization and collars.

\begin{proof}
The exact frozen inputs are
\path{prop:r2-\allowbreak completion} and
\path{thm:r2-\allowbreak completed-\allowbreak calculus} in
\path{capacitary_\allowbreak closure.tex}, and the bounded increasing
and envelope results
\path{cor:r2-\allowbreak local-\allowbreak increasing} and
\path{thm:r2-\allowbreak local-\allowbreak capacitary-\allowbreak envelope}.
Their coefficient proof uses local bounded regularization,
quasicontinuity, a countable cofinal family of compact test
stages, and coefficient control by the specified local capacity.
The original analytic setting now has these inputs from
(AE2), the analytic Laurent collar theorem, (DG3), and
(AL8)--(AL9). In particular no algebraic neighborhood or
projective compactification is silently imposed.
The explicit collar operator $v\mapsto\max(v+c,b)$ depends only
on the common scalar bound and the chosen collar, applies uniformly
to an arbitrary bounded competitor family, and commutes with finite
maxima. Thus its use for countable envelope selection assumes no
common smooth regularization of an uncountable family.

For precision, the recursive limit proof proceeds as follows.
At degree zero, (DG3) controls the coefficients of $S$.
Bounded actual weights converging in outer capacity can be
multiplied against these coefficients by the ordinary-open
quasicontinuity argument. Inductively, convergence of a bounded
weighted product of degree $k-1$, followed by distributional
$D$, defines the degree-$k$ product. Positivity, symmetry and
separate closedness pass on every original test.
For an ordinary open set $G$, lower semicontinuity of coefficient
variation gives its bound by $C_\eta R^k C_S(G;V)$.
Individual Radon outer regularity and the infimum over such
$G$ give precisely (RC1). This bound for the limit, together
with the corresponding bound for the approximants, permits
the next bounded multiplier passage. It never replaces
$C^*$ by raw capacity on a general Borel set.

The same induction proves independence, the original recursive
identities, and all bounded continuous tuple weights, exactly
as in the frozen proof. First-order currents are obtained by
distributionally differentiating the bounded weighted currents;
(FO4) gives their Radon and outer-capacity bounds. The bounded
three-term formula (GS2) gives the energy limits, while (PG8)
and polarization give their coefficient bounds and positive
matrices. The bounded multiplier argument then gives their
continuous tuple weights. The summable-exception construction
for function limits and the common-collar increasing and
envelope constructions use only the just listed analytic
inputs, so their frozen proofs apply with the same quantifiers.
Every localization is into finitely many actual original test
stages; arbitrary additional moment refinements remain allowed.
\end{proof}

\begin{definition}[Completion with an actual lower barrier]
\label{def:r11-barrier-completion}
For an actual $a\in\mathcal E^\sigma_{S,\partial}(V)$,
let $\widehat{\mathcal E}^{\sigma}_a(V)$ consist of Borel
functions $u:V\to[-\infty,0]$ such that $a\leq u\leq0$
and there exist actual classical
\[
       a\leq u_l\leq0,\qquad
       u_l\in\mathcal E^\sigma_{S,\partial}(V),
 \tag{RC2}
\]
whose fixed bounded truncations converge locally in outer
$S$-capacity to those of $u$. Equality in the completion is
local outer-capacity equality of the actual Borel functions.
Write $u^{(R)}=\max(u,-R)$. By (AL9), the convergence can
equivalently be tested on every compact core using the one
fixed capacity $C^*$, or in every smaller original analytic
capacity domain. Set
$\widehat{\mathcal E}^{\sigma}_{\rm bar}(V)
=\bigcup_{a\in\mathcal E^\sigma_{S,\partial}}
\widehat{\mathcal E}^{\sigma}_a(V)$, with the same quotient.
\end{definition}

\begin{proposition}[Existence, uniqueness and completeness with a fixed barrier]
\label{prop:r11-barrier-completion-existence}
If a sequence satisfying (RC2) is Cauchy in the stated
truncated local outer capacities, it has a Borel limit in
$\widehat{\mathcal E}^{\sigma}_a(V)$, unique modulo local
outer-capacity zero. Its actual representative can be chosen
as the pointwise limsup of one subsequence. The quotient of
$\widehat{\mathcal E}^{\sigma}_a(V)$ is complete for this
topology. Its members have compact strict negative sublevel
closures and are locally finite and continuous outside ordinary
open sets of arbitrarily small relative outer capacity.
The union class is closed under sums, positive scalar multiples
and finite maxima, with suitable actual lower barriers.
\end{proposition}

\noindent\textbf{Status: \texttt{NEW\_PROVED}} for the
unbounded barrier completion; the bounded summable-exception
argument is the frozen input.

\begin{proof}
Use a countable compact exhaustion and the integer truncation
levels. Choose a diagonal subsequence whose consecutive gaps
at threshold $2^{-l}$ have outer capacity at most $2^{-l}$
on each of the first $l$ compact cores and truncation levels.
Choose ordinary open covers with summable errors. Off their
limsup, every fixed truncation converges pointwise, and its
limits are compatible under further truncation.
The pole of $a$ is outer-capacity null by (BQ2).
At every remaining point $a$ is finite, so all terms exceed
one finite truncation level; the untruncated subsequence
therefore converges there. Define $u$ everywhere as its
actual limsup. Then $a\leq u\leq0$, and the bounded
truncations have the required capacity limits. The Cauchy
condition extends convergence to the whole sequence.
The level-set triangle inequality proves uniqueness.

For a Cauchy sequence already in the completion, choose one
classical approximant for its $l$th member satisfying (RC2)
and approximating the first $l$ compact cores, levels and
thresholds. The resulting classical sequence is Cauchy and
has the same limit. This proves completeness without a
compactness assertion for arbitrary bounded-energy families.

The inclusion $\{u<-t\}\subset\{a<-t\}$ gives compact
negative sublevel closures. On a compact core first remove
a deep ordinary sublevel of $a$ of arbitrarily small capacity.
Outside it $u$ equals one bounded truncation, which belongs
to Lemma~\ref{lem:r11-analytic-bounded-completion} and is
quasicontinuous. Removing its additional small ordinary-open
exception proves the assertion for the unchanged function.
For sums, scales and maxima apply the elementary Lipschitz
inequalities to bounded truncations after removing the deep
sublevels of the finitely many actual barriers. The resulting
classical rows have respectively the barriers $a+b$, $ca$,
and $a+b$, which are in the actual power class. This proves
the final assertions; zero scalar coefficients are omitted
before evaluating an infinite value.
\end{proof}

\subsection{A uniform modulus and construction of the currents}

For a finite tuple as in (RC2), allow separate fixed barriers
$a_i\in\mathcal E^\sigma_{S,\partial}(V)$, and put
\[
 b=\sum_i a_i,\qquad
 J_b(R)=\int_R^\infty t^{\kappa-1}C_S(\{b<-t\};V)\,dt.
 \tag{RC3}
\]
Then $J_b(0)<\infty$ and $J_b(R)\to0$. Every classical
row sum is at least $b$, so its capacity integral is bounded
pointwise by $J_b$.

\begin{lemma}[A common-barrier coefficient modulus]
\label{lem:r11-barrier-coefficient-modulus}
Fix one original compact test stage. Let $d=k$, $k+1$, or
$k+2$ be respectively the curvature, first-order, or energy
degree budget. In the latter two cases $k\leq q-1$;
for energy require $k+2\leq\kappa$. For the unbounded
classical row currents $P_l$ of that type, there is a
nondecreasing modulus $\omega_{d,b}$ with
$\omega_{d,b}(s)\to0$ as $s\downarrow0$ such that
\[
          |\mu_{P_l,\eta}|(E)
                       \leq C_\eta\omega_{d,b}(C^*(E))
                  \quad\hbox{uniformly in }l.
 \tag{RC4}
\]
One may use, up to fixed constants and an inessential
right-continuous enlargement,
\[
 \omega_{d,b}(s)
    =\inf_{L\geq1}\{L^d s+L^{d-\kappa}J_b(L/4)\}.
 \tag{RC5}
\]
For $d<\kappa$ a bound by a constant times
$s+s^{1-d/\kappa}$ is available on bounded $s$ ranges.
At $d=\kappa$ the vanishing comes from the actual tail
$J_b(L/4)$, not merely a bound for its full integral.
Every local weak current limit of this family has the same
outer-capacity modulus. Uniformly bounded changing weights
converging in local outer capacity to a quasicontinuous weight
can consequently be multiplied through such a current limit.
\end{lemma}

\noindent\textbf{Status: \texttt{NEW\_PROVED}} for the
common-barrier unbounded modulus, including the energy endpoint.

\begin{proof}
Compare a row with its common canonical truncation at $L$.
The bounded coefficient bounds (DG3), (FO4), (PG8) give
$C_\eta L^d C_S(E;V)$ for the truncated term.
The uncut/truncated variation estimates (PG3), (FO8), (PG12)
give $C_\eta L^{d-\kappa}J_b(L/4)$ for their full compact
stage difference. Taking the infimum proves (RC4)--(RC5).
The degree-zero case also follows directly from (DG3).
For $d<\kappa$ choose $L$ of order $s^{-1/\kappa}$
when $s$ is small and use $J_b(L/4)\leq J_b(0)$.
At equality first choose $L$ so that its tail is small,
then let $s$ tend to zero. This proves the stated modulus.

On an ordinary open set, variation lower semicontinuity
passes the same estimate to a weak limit; individual Radon
outer regularity gives its outer version on every Borel set.
Equivalently keep each fixed $L$ bound before taking the
outer infimum and finally minimize over $L$, avoiding any
continuity assumption on the infimum function.
For a bounded multiplier split the changing-weight error
on its positive gap set and the complement, using (RC4).
For the fixed limiting weight remove an ordinary open
quasicontinuity exception and extend continuously from the
compact remainder. The modulus bounds the errors for
\emph{both} the approximating and limiting measures.
Smooth scalar density then passes the continuous extension.
This is the proof of (GS1) with a vanishing modulus in place
of a linear capacity bound.
\end{proof}

\begin{theorem}[Canonical original products, derivatives and energies in the completion]
\label{thm:r11-barrier-completed-calculus}
Let $u_i\in\widehat{\mathcal E}^{\sigma}_{a_i}(V)$ be a
finite tuple with classical defining rows $u_{i,l}$ as in
(RC2). There are canonical original currents
\[
 \widehat N_\alpha(\mathbf u),\qquad
 \widehat{\mathcal J}'_{\alpha,i}(\mathbf u),\quad
 \widehat{\mathcal J}''_{\alpha,i}(\mathbf u),\qquad
 \widehat G^\alpha_{ih}(\mathbf u)
 \tag{RC6}
\]
with the following quantified scopes. Curvature products have
$k=|\alpha|\leq q$, first-order currents have $k\leq q-1$,
and energy currents have $k\leq q-1$ and $k+2\leq\kappa$.
They are the limits of the corresponding classical row currents
on every original compact smooth complementary test and every
actual moment refinement. They are independent of defining
sequences, of the chosen actual lower barriers, and of locally
outer-capacity-null changes of representatives.

The curvature products are symmetric, middle-positive,
separately $d'$- and $d''$-closed, and multilinear on the
completed positive cone, with $\widehat N_0=S$.
Every current in (RC6) is coefficient-Radon and gives zero
coefficient variation to the joint pole set. The energy
matrices are positive; neither first-order nor energy
currents are asserted to be closed backgrounds.
All the full recursive identities hold:
\[
 \begin{aligned}
 D(u_i\widehat N_\alpha)&=\widehat N_{\alpha+e_i}
                                      &&(|\alpha|\leq q-1),\\
 d'(u_i\widehat N_\alpha)&=\widehat{\mathcal J}'_{\alpha,i},
       \qquad
 d''(u_i\widehat N_\alpha)=\widehat{\mathcal J}''_{\alpha,i}.
 \end{aligned}
 \tag{RC7}
\]
For every continuous real function $F$ on the nonpositive
orthant with $|F(x)|\leq A(1-\sum_i x_i)^a$, all weighted
row limits are the actual coefficient multipliers
\[
                F(\mathbf u_l)P_l\longrightarrow
                                      F(\mathbf u)\widehat P,
                 \qquad a+d\leq\kappa,
 \tag{RC8}
\]
where $d=k,k+1,k+2$ is the corresponding budget.
The endpoint is included because the actual lower barriers
are fixed. The same continuity holds for completed rows
with those fixed barriers and the same capacity convergence.
\end{theorem}

\noindent\textbf{Status: \texttt{NEW\_PROVED}} for this
explicit enlarged potential class and arbitrary original $S$.

\begin{proof}
For each integer $j$ the bounded tuple
$\mathbf u_l^{(j)}$ converges to $\mathbf u^{(j)}$ in the
bounded analytic completion. Let $P^{(j)}$ be its bounded
completed product, derivative, or energy from
Lemma~\ref{lem:r11-analytic-bounded-completion}.
These currents depend only on the actual bounded truncations
modulo outer capacity, not on their defining rows.
For an original classical row, its unbounded current $P_l$
and its bounded level-$j$ current $P_l^{(j)}$ satisfy
\[
    \|\mu_{P_l-P_l^{(j)},\eta}\|_{\rm TV}
               \leq C_\eta j^{d-\kappa}J_b(j/4)
                         =:\epsilon_{j,\eta}\longrightarrow0
 \tag{RC9}
\]
uniformly in $l$. The relevant estimates are exactly
(PG3), (FO8), and (PG12). Apply the triangle inequality
to two truncation levels, then pass $l\to\infty$ using
lower semicontinuity of coefficient variation. It gives
\[
       \|\mu_{P^{(j)}-P^{(r)},\eta}\|_{\rm TV}
                       \leq\epsilon_{j,\eta}+\epsilon_{r,\eta}.
 \tag{RC10}
\]
Thus the bounded completed stages converge in local coefficient
total variation to one current $\widehat P$. The original
rows converge to it on tests by first fixing $j$ and then
using (RC9). Positive trace controls guarantee current
continuity on each original compact stage; the linear
coefficient and refinement identities pass to the limit.
For curvature, the closed positive bounded currents give
separate closedness and middle positivity. For energies,
each positive finite matrix passes to a positive matrix.
This constructs every term in (RC6), with no new assertion
about Borel-restricted nonpolar products required.

Lemma~\ref{lem:r11-barrier-coefficient-modulus} gives the
outer-capacity modulus of every limit. In particular it
annihilates the poles of the actual barriers, hence those
of the completed tuple, and supplies bounded quasicontinuous
multiplier convergence for the row currents.
Every bounded continuous function of a fixed clipped tuple
therefore passes through the current limit.

We give the endpoint integrability argument explicitly,
because integrating the weaker Holder capacity modulus
would not prove it. Write $z=-\sum_i u_i\geq0$ on the
finite-value locus, and $W=1+z$. Choose a continuous
nondecreasing scalar $\theta_R$ which is zero on $[0,R]$
and one on $[2R,\infty)$. For $L\geq1$ the bounded
continuous tuple weight
\[
             \theta_R(z)\min(W^a,L)
 \tag{RC11}
\]
is a continuous function of one sufficiently large fixed
clipped tuple: once any slot is sufficiently negative its
value is the constant $L$. For $a=0$, use the bounded
weight $\theta_R$ itself. Hence the bounded multiplier
passage applies to (RC11). The classical row tail estimates
and $J_{\sum_i u_{i,l}}\leq J_b$ bound its integral against
every positive original trace, for every sufficiently large $R$, by
\[
                C R^{a+d-\kappa}J_b(R/8).
 \tag{RC12}
\]
For signed first-order coefficients the same conclusion
follows from variation lower semicontinuity; cross energies
are controlled by the positive diagonal traces.
Let $L\uparrow\infty$. Monotone convergence proves the
same bound for the limiting trace on $\{z\geq2R\}$.
Adding the finite bounded-value core proves actual local
integrability of $W^a$ at every allowed endpoint.

For a general $F$ in (RC8), clip every slot at $-R$.
The clipped composition is a bounded quasicontinuous
multiplier and has the already proved row limit.
Its difference from $F$ is zero unless some slot is below
$-R$, and otherwise is bounded by $2A W^a$.
The classical and limiting tails are uniformly controlled
by (RC12). Let $R\to\infty$ after the fixed clipped
limit. This proves (RC8) with its actual coefficient
interpretation, including equality in the budget.

Take $F(x)=x_i$ in (RC8). For curvature degree $k\leq q-1$
one has $1+k\leq q<\kappa$. Differentiating the convergent
coefficient multipliers and using the actual classical row
recursions gives (RC7). The three-term energy identity likewise
passes with budgets $(2,k)$ and $(1,k+1)$, both at most
$\kappa$, and identifies
\[
 \widehat G^\alpha_{ih}
   =\tfrac12D(u_i u_h\widehat N_\alpha)
     -\tfrac12u_i\widehat N_{\alpha+e_h}
     -\tfrac12u_h\widehat N_{\alpha+e_i}.
 \tag{RC13}
\]
All multipliers are defined before differentiation.

The construction through the intrinsic bounded stages proves
independence of defining sequence and barrier. It also proves
representative independence because the bounded stages have it.
Symmetry, addition and positive homogeneity pass from simultaneous
classical defining rows, using the sum barriers already established.
For sequences already completed, their bounded stages enjoy
bounded completed stability and inherit (RC9)--(RC12) from
classical defining approximants with the same actual barriers.
The identical two-limit proof gives all asserted continuity.
No ordinary upper regularization, change of $S$, or equality
on extra GK delta tests occurs in this argument.
\end{proof}

\begin{corollary}[Completed scalar differential calculus]
\label{cor:r11-barrier-completed-chain}
The first-order chain identities (FO11) extend to the completed
tuple under the printed growth condition $r+k\leq\kappa$.
The second-order chain identity (PG15) extends when
$2\leq r\leq\kappa-k$ and its first and second derivatives
have the corresponding growth $W^{r-1}$ and $W^{r-2}$.
All identities hold on every original compact smooth test.
\end{corollary}

\noindent\textbf{Status: \texttt{DIRECT\_COROLLARY}}.

\begin{proof}
Apply the classical row identities. Each term has total
budget at most $r+k\leq\kappa$, so (RC8) passes every
coefficient multiplier. Distributional continuity passes
the differentiated term. For a convex coordinatewise
nondecreasing scalar function the curvature expression is
middle-positive by the positive completed energy matrix,
with the same cone argument as for the actual class.
\end{proof}

\subsection{Global critical energies and a relative Perron envelope}

\begin{theorem}[Global mixed energy continuity inside a fixed barrier completion]
\label{thm:r11-completed-global-critical-energy}
For $0\leq i\leq q$, let
$u_i\in\widehat{\mathcal E}^{\sigma}_{a_i}(V)$, and let
$u_{i,l}$ be classical or completed rows with those same
actual barriers converging in the stated capacity topology.
Then
\[
 \int_V \phi\,(-u_{0,l})^\sigma
             N(u_{1,l},\ldots,u_{q,l})
   \longrightarrow
 \int_V \phi\,(-u_0)^\sigma
             \widehat N(u_1,\ldots,u_q)
 \tag{RC14}
\]
for every bounded original smooth scalar $\phi$ on $V$,
including $\phi=1$. Here $N$ on a completed row denotes its
completed product. All displayed weighted measures have finite
total mass; unweighted top products need not have finite boundary
mass. Consequently
\[
        \widehat e_\sigma(u)
            :=\int_V(-u)^\sigma\widehat N_q(u)<\infty
 \tag{RC15}
\]
is canonical and continuous for sequences with a fixed actual
lower barrier. For $\sigma=1$, full global mixed symmetry and
the constant-one Holder inequality (GE10) extend to these
completed potentials by (RC14).
\end{theorem}

\noindent\textbf{Status: \texttt{NEW\_PROVED}} for the
completed global critical energies; the small-value estimate
is the new actual-class input (BA1).

\begin{proof}
Local convergence in (RC14) is (RC8), at the exact budget
$a=\sigma,k=q$. Put $b=\sum_{i=0}^q a_i$.
For the classical rows the small-value estimate (BA1) and
the positive mixed expansion give, for a compact original
cutoff $\chi_\epsilon$ equal to one near
$\overline{\{b<-\epsilon\}}$,
\[
 \int_V(1-\chi_\epsilon)(-u_{0,l})^\sigma
                       N(u_{1,l},\ldots,u_{q,l})
       \leq C_{q,\sigma}
                   \int_0^\epsilon t^{\kappa-1}
                                 C_S(\{b<-t\};V)\,dt.
 \tag{RC16}
\]
The right side tends to zero. It is independent of $l$.
To inherit it for the limiting positive weighted measure,
first multiply $1-\chi_\epsilon$ by any compact smooth
cutoff, pass the local convergence already proved, and
then exhaust $V$. Thus the same bound holds for the limit.
The full mixed-energy bound for the rows and compact exhaustion
give finite total limiting mass. Split any bounded $\phi$
using $\chi_\epsilon$; local convergence treats its compact
part and (RC16) treats both tails. This proves (RC14).
It also extends to completed rows by their defining classical
approximants and the same inherited bound. Symmetry and the
constant-one inequality at $\sigma=1$ pass through these
global finite integrals, proving the last assertion.
\end{proof}

\begin{theorem}[A barrier-bounded relative Perron envelope]
\label{thm:r11-barrier-relative-Perron}
Let $a\in\mathcal E^\sigma_{S,\partial}(V)$ be actual
classical, and let $\mathscr F$ be a nonempty family of
actual classical members of that power class satisfying
$a\leq v\leq0$. There is an increasing sequence $v_l$
of finite maxima of members of $\mathscr F$ whose actual
Borel supremum
\[
                       H=\sup_l v_l
                          \in\widehat{\mathcal E}^{\sigma}_a(V)
 \tag{RC17}
\]
is the least upper bound of $\mathscr F$ modulo local
ordinary-open outer $S$-capacity zero. More explicitly,
every $v\in\mathscr F$ satisfies $v\leq H$ locally
quasi-everywhere, and every Borel $h$ with $v\leq h$
locally quasi-everywhere for all $v\in\mathscr F$ satisfies
$H\leq h$ locally quasi-everywhere. The exceptional set in
the first assertion may depend on $v$.

If every member has $\mu_v\geq\nu$ for one positive
Radon measure $\nu$, then
\[
                              \widehat N_q(H)\geq\nu.
 \tag{RC18}
\]
If every member satisfies a prescribed Borel obstacle
$v\leq\psi$ locally quasi-everywhere on a Borel set $E$,
then so does $H$. The envelope is canonical as a relative
capacity class. It is not replaced by its ordinary ambient
upper semicontinuous regularization.
\end{theorem}

\noindent\textbf{Status: \texttt{NEW\_PROVED}} as an
unbounded relative completed envelope and subsolution theorem.
The countable selection in the bounded truncations is the
frozen relative envelope construction, with analytic scope
supplied by Lemma~\ref{lem:r11-analytic-bounded-completion}.

\begin{proof}
First consider any increasing sequence $v_l$ of actual
classical potentials between $a$ and zero. For every integer
$R$, the sequence $v_l^{(R)}$ is bounded increasing.
The frozen actual bounded increasing theorem, in its analytic
form above, gives convergence in local outer capacity to its
actual supremum. The common lower barrier controls the deep
exceptions, so $H=\sup_l v_l$ belongs to the completion,
with these compatible bounded truncations. The completed
current and global energy theorems therefore apply to it.

For the arbitrary family, apply the bounded relative
capacitary envelope theorem to
$\{\max(v,-R):v\in\mathscr F\}$ for each positive integer
$R$. It chooses countably many original family members whose
finite truncated maxima have the bounded relative least
upper-bound property. Take the countable union of the chosen
members for all $R$, list them, and let $v_l$ be the finite
maxima of the first $l$. Each is actual classical, belongs
to the power class by solidity, and lies above $a$.
The first paragraph constructs its completed supremum $H$.

For a fixed $v\in\mathscr F$ and fixed $R$, the selection
gives $v^{(R)}\leq H^{(R)}$ locally outside a null set.
Take the countable union over $R$. The compatible truncation
inequalities imply $v\leq H$ on its complement, including
extended values. Conversely every $v_l$ is a finite maximum
of original members, so $v_l\leq h$ outside a finite union
of the given exceptional sets. Their countable union over
$l$ proves $H\leq h$ locally quasi-everywhere. This also
proves uniqueness of the relative envelope class.

Under the source hypothesis, the actual unbounded weak
maximum theorem makes every $v_l$ a subsolution with the
same source $\nu$. Passing on every nonnegative original
compact scalar test proves (RC18), and hence the full Radon
measure inequality. For an obstacle, finite maxima preserve
the inequality outside finite unions of exceptional sets;
the countable union in (RC17) preserves it for $H$.
No contact-set measure identity is used for either passage.
\end{proof}

\paragraph{What has and has not been completed.}
The section supplies a noncircular relative completion:
the admissibility premise is a known actual power-energy
lower barrier and capacity approximation by the original
classical potentials. The products are constructed through
bounded completed stages and uniform tail estimates before
their unbounded recursive identities are asserted.
The only weak-limit set control used is outer-capacity
control, explicitly (RC1) and (RC4).

It does not prove that the completed product measures can
always be realized by a normalized ambient classical tester,
that raw and outer capacities agree on compact or Borel sets,
that arbitrary finite dual-energy data are in this operator
image, or that the envelope in (RC17) has equality in (RC18).
It also does not assert new strict Borel nonpolar locality
for arbitrary completed potentials beyond the limits and
identities proved above. Ordinary upper regularization cannot
be inserted: the explicit power-core obstruction (PO1)--(PO7)
already satisfies the common-barrier hypotheses and loses its
nonzero source after that replacement.

% END REVIEWED MODULE research_r11/capacity/barrier_relative_completion.tex


% BEGIN REVIEWED MODULE research_r11/capacity/completed_comparison_domination.tex
\section{Comparison and domination in the barrier completion}
\label{sec:r11-completed-comparison-domination}

Retain exactly (TC1)--(TC2) and the completion
of Definition~\ref{def:r11-barrier-completion}. Thus $S$ is the
unchanged arbitrary original symmetric middle-positive separately
closed $(p,p)$-current, $q=n-p\geq1$, and $V$ is the stated
strict analytic box with compact closure and background extension.
Fix one $\sigma>0$. Every completed potential below has an actual
lower barrier in $\mathcal E^\sigma_{S,\partial}(V)$, and all
defining rows are actual members of that same power class.
Write $\widehat\mu_u=\widehat N_q(u)$ and
$C=C_S(\,\cdot\,;V)$, $C^*=C_S^*(\,\cdot\,;V)$.
The potentials are their specified actual Borel representatives;
no ambient upper semicontinuous representative is asserted.

The frozen bounded comparison theorem is
\texttt{INTERNAL\_INPUT}. The new actual unbounded comparison
theorem (UC10), the analytic completed calculus (RC6)--(RC14),
and the exact finite-mass tester core (TC6) are the additional
inputs. The proof below passes a moving scalar gap by continuous
weights. It does not assume strict Borel nonpolar locality for
completed potentials.

\begin{lemma}[Moving-gap restriction on every original compact test]
\label{lem:r11-completed-moving-gap}
Let $u_l\to u$ and $v_l\to v$ be classical defining rows, or
completed rows, in fixed actual barrier intervals
$a\leq u_l,u\leq0$ and $b\leq v_l,v\leq0$.
The convergence is the local truncated outer-capacity convergence
of the completion. Let $P_l\to\widehat P$ be a curvature
product, a diagonal energy current, or a finite signed
polarization of these currents in (RC6), formed from a finite
tuple of such rows with fixed actual lower barriers and the
printed degree budget.
For each fixed original compact complementary test stage,
and every $t$ outside a countable set depending on its coefficient
traces, the restricted coefficients satisfy
\[
 \mathbf1_{\{u_l<v_l-t\}}P_l
       \longrightarrow \mathbf1_{\{u<v-t\}}\widehat P
 \quad\hbox{on that original test stage}.
 \tag{CP1}
\]
Every additional actual moment refinement is allowed. For
$P_l=\mu_{u_l}$ or $\mu_{v_l}$ and $t>0$, the total masses
in (CP1) converge outside one countable set of positive levels.
For both choices simultaneously that exceptional set can be chosen
countable. The products give zero coefficient mass to every joint
pole set, so values assigned to a difference there do not enter.
\end{lemma}

\noindent\textbf{Status: \texttt{NEW\_PROVED}} as a moving Borel
restriction result for this enlarged class.

\begin{proof}
On the finite-value locus put $h_l=v_l-u_l$ and $h=v-u$,
and assign zero on the omitted joint poles. The outer-capacity
coefficient moduli (RC4)--(RC5) annihilate those poles for every
row and every limit. For every bounded continuous real $F$,
the composition $F(y-x)$ is a continuous bounded function on
the finite nonpositive orthant. Thus (RC8), with growth exponent
zero, gives
\[
                        F(h_l)P_l\longrightarrow F(h)\widehat P
 \tag{CP2}
\]
on every original compact complementary test. Equivalently one
may first clip $u_l,v_l$ and remove their tails by (RC12);
this is why oscillation of $F$ at infinity causes no problem.

For a positive original coefficient trace, apply (CP2) to
continuous scalar upper and lower approximations to
$\mathbf1_{(t,\infty)}$. At a level with no atom for the
pushforward of the limiting trace, their limiting errors tend
to zero. Signed coefficients are controlled by the corresponding
finite positive traces of those same positive row currents.
Their traces have the positive limiting traces given by (RC8),
so the same two scalar bounds control every signed error.
Finite signed polarizations are handled term by term.
A finite trace measure has only
countably many atoms. This proves the assertion on the specified
stage, with no simultaneous exceptional-set claim for an
unfixed uncountable collection of test presentations.

For a positive level $t$, all row and limit strict-gap sets
are contained in
\[
             \{u_l<v_l-t\}\subset\{a<-t\},\qquad
             \{u<v-t\}\subset\{a<-t\}.
 \tag{CP3}
\]
Their closures lie in a common ordinary compact set, by
$\mathrm{BG}(a,0)$. Choose an original compact smooth cutoff
equal to one on a larger compact sublevel core. Then (CP2)
also passes their total mass. Countably many such cutoffs,
for levels bounded away from zero, yield one countable
exceptional set for the two top measures. At a level itself,
$h=t>0$ implies $a\leq-t$, so it is contained in the same
slightly larger compact core. This verifies that its level
mass is the mass detected by the chosen cutoff.
\end{proof}

\begin{theorem}[Comparison for completed power-energy potentials]
\label{thm:r11-completed-comparison}
For $u\in\widehat{\mathcal E}^\sigma_a(V)$ and
$v\in\widehat{\mathcal E}^\sigma_b(V)$, every $t\geq0$ satisfies
\[
       \boxed{\quad
       \widehat\mu_v\{u<v-t\}
                      \leq\widehat\mu_u\{u<v-t\}.
       \quad}
 \tag{CP4}
\]
For $t>0$ both sides are finite. At $t=0$ either side may be
infinite. These are inequalities of actual Borel measures of
the canonical original currents, independent of outer-null
changes of the completed representatives.
\end{theorem}

\noindent\textbf{Status: \texttt{NEW\_PROVED}} beyond the actual
classical class; the scalar comparison for its defining rows is
the already proved input (UC10).

\begin{proof}
Choose simultaneous classical defining rows $u_l,v_l$ above
the fixed actual barriers. For $t>0$, the pair $u_l,v_l-t$
satisfies the required compact boundary-gap hypothesis:
every positive-margin gap lies in a strict negative sublevel
of $a$. Translation of the actual nonpolar product gives
$\mu_{v_l-t}=\mu_{v_l}$, as proved before (BC12).
The actual comparison theorem therefore gives
\[
                   \mu_{v_l}\{u_l<v_l-t\}
                           \leq\mu_{u_l}\{u_l<v_l-t\}.
 \tag{CP5}
\]
At every positive level outside the countable exceptional set
of Lemma~\ref{lem:r11-completed-moving-gap}, both total masses
pass to the limit. This proves (CP4) at those levels.
For any $t\geq0$, take regular positive levels $t_j\downarrow t$.
The strict-gap sets increase to $\{u<v-t\}$, including their
actual extended-value interpretation. Measure continuity from
below proves (CP4). No convergence of a moving discontinuous
indicator was inferred from weak current convergence alone.
\end{proof}

\begin{lemma}[Relative quasi-openness of completed strict gaps]
\label{lem:r11-completed-gap-outer}
For the same completed potentials and every $t\geq0$, the set
$A_t=\{u<v-t\}$ is relatively quasi-open for the full original
capacity. In particular
\[
                              C^*(A_t)=C(A_t).
 \tag{CP6}
\]
This is an assertion about these specified strict gaps, not
about general Borel analytic sets.
\end{lemma}

\noindent\textbf{Status: \texttt{DIRECT\_COROLLARY}} of the
completed quasicontinuity and the ordinary-open definition of
outer capacity.

\begin{proof}
Choose a countable compact exhaustion whose interiors cover $V$.
On its $j$th compact core, the completed quasicontinuity theorem
gives an ordinary open exceptional set of capacity less than
$2^{-j-2}\varepsilon$ on which the two functions may fail to be
finite and continuous. Include ordinary open covers of their
outer-null pole sets with the same summable cost. The union
$G$ has $C(G)<\varepsilon$, and $u,v$ are finite and continuous
in the relative topology of $V\setminus G$. Hence
$A_t\setminus G$ is relatively open there and $A_t\cup G$
is ordinary open in $V$. Monotonicity and subadditivity give
\[
 C(A_t)\leq C^*(A_t)\leq C(A_t\cup G)
                                      \leq C(A_t)+\varepsilon.
\]
Let $\varepsilon\downarrow0$. Local outer-capacity exceptional
sets are measured in the full fixed domain; the analytic domain
comparison (AL9) supplies this interpretation if a smaller
ordinary neighborhood was used for quasicontinuity.
\end{proof}

\begin{theorem}[Quantitative domination and uniqueness in the completion]
\label{thm:r11-completed-domination}
For $s,t>0$, the same completed potentials satisfy
\[
              t^q C^*\{u<v-s-t\}
                       \leq\widehat\mu_u\{u<v-s\}.
 \tag{CP7}
\]
Put $A=\{u<v\}$. If either
\[
       \widehat\mu_u(A)=0,
       \qquad\hbox{or}\qquad
       \mathbf1_A\widehat\mu_u
                              \leq\mathbf1_A\widehat\mu_v,
 \tag{CP8}
\]
then $C^*(A)=0$. Consequently
\[
    \widehat\mu_u=\widehat\mu_v
           \quad\Longrightarrow\quad C^*\{u\ne v\}=0.
 \tag{CP9}
\]
This is uniqueness of the completed capacity class. It is not
pointwise equality for arbitrary degenerate backgrounds.
\end{theorem}

\noindent\textbf{Status: \texttt{NEW\_PROVED}} for the completed
class and the full normalized original capacity.

\begin{proof}
Let $h$ be an arbitrary tester in the exact core
$\mathscr T_0(V)$ of (TC5), so $-1\leq h\leq0$ and $h$
belongs to every positive power class. The completed potential
$w=v+t h$ lies above the actual power-energy barrier $b+t h$.
The completed positive polynomial expansion gives
\[
                  \widehat\mu_w\geq t^q\mu_h.
\]
Since $v-t\leq w\leq v$, (CP4) implies
\[
 \begin{split}
 t^q\mu_h\{u<v-s-t\}
 &\leq\widehat\mu_w\{u<w-s\}\\
 &\leq\widehat\mu_u\{u<w-s\}
                 \leq\widehat\mu_u\{u<v-s\}.
 \end{split}
 \tag{CP10}
\]
Take the supremum over every core tester and use its exact
raw-capacity identity (TC6). Lemma~\ref{lem:r11-completed-gap-outer}
then changes that particular raw gap capacity to outer capacity,
proving (CP7).

If the first alternative in (CP8) holds, (CP7) makes every
positive-margin gap raw-null. Their increasing union is $A$.
Raw capacity is continuous from below, directly from its
definition as a supremum of measures. Hence $C(A)=0$, and
(CP6) gives $C^*(A)=0$.

For the second alternative, use instead
$w=(1+t)v+t h$, still a completed member above the actual
barrier $(1+t)b+t h$. For $s>0$ put $B_s=\{u<w-s\}$.
Then $B_s\subset A$ and $B_s\subset\{a<-s\}$, so all top
measures below are finite on $B_s$. Positivity, comparison and
the restricted measure hypothesis give
\[
 (1+t)^q\widehat\mu_v(B_s)+t^q\mu_h(B_s)
       \leq\widehat\mu_w(B_s)
       \leq\widehat\mu_u(B_s)
       \leq\widehat\mu_v(B_s).
 \tag{CP11}
\]
Subtracting finite quantities shows $\mu_h(B_s)=0$.
First let $s\downarrow0$, then let $t\downarrow0$.
Because $h$ is bounded and $v$ is finite at every point of
the strict extended gap $A$, the resulting increasing union
of these sets is precisely $A$. Thus $\mu_h(A)=0$ for every
core tester. Again (TC6) and (CP6) give $C^*(A)=0$.
Apply this in both orders when the two top measures agree;
subadditivity proves (CP9).
\end{proof}

\paragraph{Scope and remaining gates.}
All products used here are the canonical original currents of
the independently established completion. Every continuous
weight passage occurs before taking a scalar indicator limit.
The core identity is essential: it computes the complete raw
tester supremum without asserting that an arbitrary normalized
tester has the completed boundary condition. No analytic PS/PL
quotient, signed tensor order transfer, weak-current equality
on all GK delta tests, or arbitrary-source existence is inferred.
In particular the compact-source obstruction for the completion
in (BV13)--(BV14) is compatible with uniqueness (CP9).

% END REVIEWED MODULE research_r11/capacity/completed_comparison_domination.tex


% BEGIN REVIEWED MODULE research_r11/capacity/completed_locality_carriers.tex
\section{Strict locality and intrinsic truncation carriers in the completion}
\label{sec:r11-completed-locality-carriers}

Keep the exact domain, background and fixed $\sigma>0$ of
Section~\ref{sec:r11-completed-comparison-domination}.
Every potential is a Borel member of a barrier completion
$\widehat{\mathcal E}^\sigma_a(V)$ with an actual classical
lower barrier $a$. All products are the canonical original
currents constructed independently in (RC6)--(RC14).
The purpose of this section is to prove their Borel locality;
it was not an input to that construction or to (CP4).

The frozen actual strict maximum locality is
\texttt{INTERNAL\_INPUT}. The new ingredients used in its
passage are the completed continuous tuple-weight convergence
and the positive trace argument (CP1)--(CP2). In particular
the same argument works for a rectangular strict condition on
finitely many scalar values: its boundary is contained in the
union of finitely many level sets, and one chooses margins
outside the atoms of the finitely many limiting positive traces.
All equalities below mean every coefficient on every original
compact smooth complementary test, with every actual refinement.

\begin{theorem}[Completed strict Borel maximum locality]
\label{thm:r11-completed-strict-maximum-locality}
Let $u,v$ and the finite background tuple $\mathbf z$ be completed
potentials with actual power-energy lower barriers. For $c\geq0$
put $w=\max(u,v-c)$. This is a completed potential above the
actual lower barrier of $u$. For any integer $m\geq1$ and
multi-index $\beta$ with $m+|\beta|\leq q$, one has
\[
 \begin{split}
 \mathbf1_{\{u>v-c\}}
       \widehat N(w^{[m]},\mathbf z^\beta)
   &=\mathbf1_{\{u>v-c\}}
       \widehat N(u^{[m]},\mathbf z^\beta),\\
 \mathbf1_{\{u<v-c\}}
       \widehat N(w^{[m]},\mathbf z^\beta)
   &=\mathbf1_{\{u<v-c\}}
       \widehat N(v^{[m]},\mathbf z^\beta).
 \end{split}
 \tag{CL1}
\]
The joint pole set has zero coefficient variation for every
term. The notation on the second line uses the constant-shift
invariance of differentiated slots; it does not assert that
$v-c$ itself has a zero-boundary completed representative.
\end{theorem}

\noindent\textbf{Status: \texttt{NEW\_PROVED}} for the completed
Borel class on the unchanged arbitrary background.

\begin{proof}
Choose simultaneous actual defining rows $u_l,v_l,\mathbf z_l$.
The actual function $w_l=\max(u_l,v_l-c)$ is negative classical
psh and lies above the fixed actual barrier of $u$. The upward
stability of the power class makes every $w_l$ a member of that
class. Fixed bounded truncations of this maximum converge in
outer capacity to those of $w$: explicitly, for $R>0$,
\[
       \max(w_l,-R)
          =\max\bigl(\max(u_l,-R),\max(v_l,-R-c)-c\bigr).
 \tag{CL2}
\]
The extra clipping on the second entry may be made at $-R+c$
instead when $R>c$; the displayed deeper clipping already
gives the same maximum because its value is below $-R$.
Continuity of a finite maximum therefore proves membership of
$w$ in the stated completion, with these defining rows.

Fix an original compact test stage. On
$\{u_l>v_l-c+\varepsilon\}$, actual strict maximum locality
replaces every $w_l$ slot by $u_l$. On
$\{u_l<v_l-c-\varepsilon\}$ it replaces every such slot by
$v_l-c$, whose differentiated current is the same as that of
$v_l$. The completed limits of both positive row products
exist by (RC8). Apply (CP1), equivalently its continuous
scalar-gap proof, at a sequence of positive margins
$\varepsilon_j\downarrow0$ outside the countably many
atoms of all the positive traces used on this finite stage.
The two restricted product equalities pass at these margins.
Their increasing unions are the two strict sets in (CL1),
apart from the already coefficient-null joint poles.
Continuity from below of the finite coefficient measures on
the stage proves (CL1). Each stage and every refinement is
allowed separately; no universal exceptional-level assertion
over an unfixed uncountable family was used.
\end{proof}

\begin{theorem}[Exact completed truncation carriers]
\label{thm:r11-completed-truncation-carriers}
Let $\mathbf u=(u_1,\ldots,u_m)$ be a finite completed tuple
and $|\alpha|=k\leq q$. For $j>0$ put
\[
 u_i^{(j)}=\max(u_i,-j),\qquad
 E_j=\bigcap_{i=1}^m\{u_i>-j\}.
 \tag{CL3}
\]
Each clipped function is in both its fixed barrier completion
and the bounded capacitary completion. For every $j>0$,
\[
       \mathbf1_{E_j}\widehat N_\alpha(\mathbf u)
          =\mathbf1_{E_j}
                     T^{\mathrm{comp}}_{S,\alpha}(\mathbf u^{(j)}).
 \tag{CL4}
\]
Consequently the product has the intrinsic representation
\[
       \widehat N_\alpha(\mathbf u)
        =\lim_{j\to\infty}\mathbf1_{E_j}
                     T^{\mathrm{comp}}_{S,\alpha}(\mathbf u^{(j)}).
 \tag{CL5}
\]
The limit is locally in coefficient total variation, and the
restricted positive currents increase in their actual Borel
measure sense. The uncut bounded completed products also
converge locally in coefficient total variation. Their discarded
coefficients satisfy
\[
       \mathbf1_{V\setminus E_j}
                     T^{\mathrm{comp}}_{S,\alpha}(\mathbf u^{(j)})
                         \longrightarrow0
           \quad\hbox{locally in coefficient total variation}.
 \tag{CL6}
\]
These assertions identify a nonpolar truncation construction
inside the proved bounded completed calculus. They neither
assert that the clipped Borel functions are ambient classical
psh nor invoke a new closedness theorem for Borel restrictions.
\end{theorem}

\noindent\textbf{Status: \texttt{NEW\_PROVED}} as the exact
carrier and no-loss identification for the completed products.

\begin{proof}
Take actual defining rows $u_{i,l}$ above fixed actual barriers.
Their clipped rows remain above those same barriers and converge
in bounded outer capacity. Thus their products converge to the
bounded completed products, which agree with the canonical
products (RC6) for that clipped tuple.

For fixed $j$ and a positive scalar margin $\varepsilon$, actual
strict locality gives the full mixed coefficient identity on
\[
                      E_{j,\varepsilon,l}
                      =\bigcap_i\{u_{i,l}>-j+\varepsilon\}.
\]
Use bounded continuous scalar approximations to these finitely
many indicators. Their error is supported in the union of the
corresponding scalar strips. The positive limiting traces of
the two product sequences have atoms at only countably many
values of each $u_i$. Choose one sequence
$\varepsilon_h\downarrow0$ for which every boundary level
$-j+\varepsilon_h$ has zero mass for all these traces.
The rectangular version of (CP2) passes the identity on
$E_{j,\varepsilon_h}=\bigcap_i\{u_i>-j+\varepsilon_h\}$.
Taking its increasing union proves (CL4).

This argument uses continuous scalar weights of $u_{i,l}$
directly. In particular the constant $-j$ is not treated as
if it belonged to a zero-boundary barrier completion.

The sets $E_j$ increase to the complement of the joint pole
set. That set is coefficient-null by the independent construction
in (RC6). Thus the right side of (CL4) is exactly the restriction
of one fixed coefficient-Radon current to an increasing sequence
of carriers, proving (CL5) with local coefficient total variation.

The uncut products converge in the same local variation by the
intrinsic bounded-stage estimate (RC9)--(RC10): for the common
actual lower sum barrier $b$, their error on a fixed original
stage is bounded by
\[
                  C j^{k-q-\sigma}J_b(j/4)\longrightarrow0.
 \tag{CL7}
\]
Finally their variation on $V\setminus E_j$ is at most that
error plus the variation of $\widehat N_\alpha(\mathbf u)$ on
$V\setminus E_j$. The latter tends to zero by pole nullity
and local finiteness. This proves (CL6). Separate closedness
of the unrestricted limit was already established by (RC7);
no differentiated characteristic function is discarded.
\end{proof}

\begin{theorem}[Completed weak maximum inequality including equality]
\label{thm:r11-completed-weak-maximum}
Let $u,v,\mathbf z$ be as above and put $w=\max(u,v)$.
For $m\geq1$ and $m+|\beta|\leq q$,
\[
 \widehat N(w^{[m]},\mathbf z^\beta)
   \geq\mathbf1_{\{u\geq v\}}
                   \widehat N(u^{[m]},\mathbf z^\beta)
       +\mathbf1_{\{u<v\}}
                   \widehat N(v^{[m]},\mathbf z^\beta).
 \tag{CL8}
\]
This is a middle-positive current inequality: it is tested
against every original compact smooth \emph{positive}
complementary form, not just the strongly positive subcone.
In particular finite maxima of completed top-degree
subsolutions for a fixed positive datum remain subsolutions.
\end{theorem}

\noindent\textbf{Status: \texttt{NEW\_PROVED}} in the completed
class. The bounded actual top inequality was already frozen
in R3; the actual unbounded extension is (WM4)--(WM7).

\begin{proof}
For $\varepsilon>0$ let
$w_\varepsilon=\max(u,v-\varepsilon)$.
The first theorem makes it a completed member above the actual
barrier of $u$. On the joint finite-value locus
$0\leq w-w_\varepsilon\leq\varepsilon$; the common pole is
already coefficient-null. Hence its bounded truncations converge
uniformly to those of $w$, and (RC8) gives convergence of every
mixed product on the left of (CL8).

Apply (CL1) \emph{first}, to the completed currents themselves.
The fixed set $\{u\geq v\}$ is contained in
$\{u>v-\varepsilon\}$ except for the coefficient-null common
pole. The disjoint set $\{u<v-\varepsilon\}$ is the other
strict maximum carrier. Positivity of the remaining restriction
therefore gives, before any limit is taken,
\[
 \begin{split}
 \widehat N(w_\varepsilon^{[m]},\mathbf z^\beta)
 \geq{}&\mathbf1_{\{u\geq v\}}
                   \widehat N(u^{[m]},\mathbf z^\beta)\\
       &+\mathbf1_{\{u<v-\varepsilon\}}
                   \widehat N(v^{[m]},\mathbf z^\beta).
 \end{split}
 \tag{CL9}
\]
The first term on the right is one fixed restricted current;
no moving weak limit is multiplied by its equality indicator.
The second term increases as $\varepsilon\downarrow0$ to the
strict $v$-carrier. On an original positive complementary test,
monotone convergence applies to that positive coefficient
measure, while (RC8) passes the left side. This proves (CL8),
including a possibly charged equality level. The top-degree
subsolution assertion follows by adding the two disjoint
restrictions of the fixed datum and iterating finitely many times.
\end{proof}

\paragraph{Dependency audit.}
The order is completed continuous-weight calculus, scalar-margin
restriction, completed strict locality, exact truncation carriers,
and then the weak maximum inequality. The construction of the
completed products did not assume any of these Borel identities.
The arguments preserve every original complementary test and
all actual moment refinements. They do not identify ordinary
current equality with equality on arbitrary GK delta tests,
nor establish the still open full analytic relation quotient.

% END REVIEWED MODULE research_r11/capacity/completed_locality_carriers.tex


% BEGIN REVIEWED MODULE research_r11/comparison/strong_first_energy_completion.tex
\section{The complete relative strong first-energy space}
\label{sec:r11-strong-first-energy-completion}

Retain precisely the strict analytic box and extension hypotheses
(TC1)--(TC2), the arbitrary original separately closed symmetric
middle-positive background $S$, and $q=n-p\geq1$. Write
$d=q+1$, $\alpha=2^{-q}$, and $\mathcal E^1=
\mathcal E^1_{S,\partial}(V)$. The actual classical class
$\mathcal E^1$ is not enlarged by definition. We instead construct
a distinct space of actual Borel functions modulo local ordinary-open
outer $S$-capacity. Its defining rows have no prescribed common
actual lower barrier. We do not assert that its members admit
ambient classical representatives or compact negative sublevel
closures. Their boundary condition is the strong approximation
condition printed below.

The frozen bounded predecessors are the R2 capacitary completion,
R3 bounded common-collar strong topology, and R4 completed
variational retraction. The new inputs are the unbounded global
first-energy calculus, the tester estimate (EV1), and the analytic
bounded completion in Lemma~\ref{lem:r11-analytic-bounded-completion}.
The present construction does not reinterpret the frozen linear
Hilbert completion as a nonlinear classical potential space.

\subsection{Uniform tails and the boundary of an energy-Cauchy row}

Call an actual row $(u_j)\subset\mathcal E^1$ strongly Cauchy if
\[
               \lim_{N\to\infty}\sup_{j,l\geq N}
                    \mathscr I_S(u_j,u_l)=0.
 \tag{IC1}
\]
No triangle inequality for $\mathscr I_S$ is assumed.

\begin{lemma}[Uniform deep tails and global weighted tightness]
\label{lem:r11-strong-cauchy-uniform-tails}
Every strongly Cauchy actual row has bounded first energies and
\[
 \begin{split}
 \lim_{R\to\infty}\sup_j
     \int_{\{u_j<-R\}}(-u_j)\,d\mu_{u_j}&=0,\\
 \lim_{R\to\infty}\sup_j
     \int_R^\infty t^q C_S(\{u_j<-t\};V)\,dt&=0.
 \end{split}
 \tag{IC2}
\]
For finitely many strongly Cauchy actual rows $u_{i,j}$,
$0\leq i\leq q$, the finite positive Radon measures
\[
            (-u_{0,j})M(u_{1,j},\ldots,u_{q,j})
 \tag{IC3}
\]
are uniformly tight on $V$. The same conclusions hold if the
indices of the finitely many rows tend independently to infinity.
\end{lemma}

\noindent\textbf{Status: \texttt{NEW\_PROVED}}. The first assertion
extends (EV19) from convergence to a specified actual member to
arbitrary strongly Cauchy rows; the boundary tightness is proved
without a common lower barrier.

\begin{proof}
The energy-root estimate (EV9) bounds the energies by some $M$.
The case $M=0$ is immediate from the constant-one mixed Holder
inequality; otherwise enlarge $M$ to be positive. On
$\{-u_j>R\}$ split according to whether
$-u_N\leq(-u_j)/2$. The calculation used for (EV20) gives
\[
 \int_{\{u_j<-R\}}(-u_j)\,d\mu_{u_j}
 \leq 2\int|u_j-u_N|\,d\mu_{u_j}
       +4\int|u_N-u_N^{(R/4)}|\,d\mu_{u_j},
 \qquad u^{(L)}=\max(u,-L).
 \tag{IC4}
\]
By (EV1), the first integral is at most
$C_{q,M}\mathscr I_S(u_j,u_N)^\alpha$.
For fixed $N$, the second is at most
$C_{q,M}\mathscr I_S(u_N,u_N^{(R/4)})^\alpha$, uniformly
in $j$. The latter discrepancy tends to zero by (EV15).
Choose $N$ large using (IC1), then let $R$ grow, and finally
make the Cauchy error tend to zero. The finitely many indices
$j<N$ have vanishing tails because their first energies are
finite. This proves the first line of (IC2). The tail form of
(EV21),
\[
 \int_R^\infty t^q C_S(\{u_j<-t\};V)\,dt
       \leq 2^{q+1}\int(-u_j-R/2)^+\,d\mu_{u_j},
 \tag{IC5}
\]
proves the second. For a finite sum of rows, the sublevel union
bound gives the same uniform tail conclusion, with the factor
$(\#\text{slots})^{q+1}$ and a fixed rescaling of $R$.

For (IC3), all slot energies are uniformly bounded. Write
$M_j=M(u_{1,j},\ldots,u_{q,j})$. Positivity of the expansion
of $\mu_{\sum_{i=1}^q u_{i,j}}$ and (EV1) give
\[
             \int|u_{0,j}-u_{0,N}|\,dM_j
                  \leq C\mathscr I_S(u_{0,j},u_{0,N})^\alpha,
 \tag{IC6}
\]
where $C$ depends on the uniform energies and $q$, but not on
the derivative-row indices. Fix a sufficiently late $N$.
For $\epsilon>0$, the set
$K_\epsilon=\overline{\{u_{0,N}<-\epsilon\}}$ is compact.
On its complement $u_{0,N}=u_{0,N}^{(\epsilon)}$, and hence
\[
 \begin{split}
 \int_{V\setminus K_\epsilon}(-u_{0,j})\,dM_j
 &\leq \int|u_{0,j}-u_{0,N}|\,dM_j
          +\int(-u_{0,N}^{(\epsilon)})\,dM_j\\
 &\leq C\mathscr I_S(u_{0,j},u_{0,N})^\alpha
       +e_1(u_{0,N}^{(\epsilon)})^{1/d}
                          \prod_{i=1}^q e_1(u_{i,j})^{1/d}.
 \end{split}
 \tag{IC7}
\]
The last line is the global constant-one Holder inequality.
The small-value energy vanishes:
\[
 e_1(u_{0,N}^{(\epsilon)})
   \leq C_q\int_0^\epsilon t^q
                    C_S(\{u_{0,N}<-t\};V)\,dt\longrightarrow0,
 \tag{IC8}
\]
by the power-energy/capacity equivalence. First make the
Cauchy error small, then choose $\epsilon$ small. The
finitely many omitted weights $u_{0,j}$, with derivative
rows still varying, are handled by the same argument with
zero first term: outside their own compact negative sublevel,
the Holder bound with $u_{0,j}^{(\epsilon)}$ tends uniformly
to zero. A finite union of these compacts proves uniform
tightness. This observation also proves the assertion for
independent row indices, without taking a diagonal assumption
on any derivative slot.
\end{proof}

\subsection{The Borel class and its intrinsic original currents}

\begin{definition}[Strong relative first-energy completion]
\label{def:r11-strong-first-energy-completion}
Let $\widehat{\mathcal E}^1_{\rm str}(V)$ consist of actual
Borel functions $u:V\to[-\infty,0]$ which are limits of an
actual strongly Cauchy row in every local truncated outer
$S$-capacity. Equality means local outer-capacity equality.
The capacities may equivalently be taken on all compact
cores in the fixed domain $V$, by (AL9). A choice of actual
defining row is not part of the element of this quotient.
\end{definition}

\begin{proposition}[Borel limits of all strongly Cauchy actual rows]
\label{prop:r11-strong-cauchy-borel-limit}
Every row satisfying (IC1) has a limit in
Definition~\ref{def:r11-strong-first-energy-completion}, unique
modulo local outer-capacity zero. Every member of that space is
finite and continuous locally outside ordinary open sets of
arbitrarily small outer capacity. Its pole set is locally
outer-capacity zero.
\end{proposition}

\noindent\textbf{Status: \texttt{NEW\_PROVED}} for unrestricted
strongly Cauchy first-energy rows, with no common actual barrier.

\begin{proof}
Uniform energy and (EV8) make each fixed bounded truncation
of the row Cauchy in every local outer capacity. Apply the
countable compact-core and integer-level diagonal construction
from the bounded completion. A subsequence has compatible
pointwise limits of all bounded truncations off a locally
outer-null set. Its actual pointwise limsup defines a Borel
$u\leq0$ on all of $V$, and the Cauchy property gives the
required capacity convergence for the full row.

There is a uniform deep sublevel estimate. The no-loss
sublevel inequality and the first-energy bound imply
\[
            C_S(\{u_j<-R\};V)\leq C_q M R^{-d}.
 \tag{IC9}
\]
These sets are ordinary open, so the same bound holds for
their outer capacities. On a fixed compact core, the set
$\{u<-2R\}$ is contained in $\{u_j<-R\}$ together with
a positive gap of the fixed truncations at $3R$.
Pass $j\to\infty$ in the outer-capacity subadditivity bound.
It follows that the pole of $u$ has outer capacity zero on
that core. Outside an ordinary open cover of a sufficiently
deep sublevel, $u$ agrees with one bounded truncation.
Its bounded-completion quasicontinuity then proves the stated
finite local quasicontinuity. The level-set triangle inequality
proves uniqueness. Neither this argument nor its conclusion
requires compact negative sublevel closures for the limit.
\end{proof}

\begin{theorem}[Canonical critical calculus in the strong completion]
\label{thm:r11-strong-completed-calculus}
For a finite tuple $\mathbf u$ in
$\widehat{\mathcal E}^1_{\rm str}(V)$, with actual strongly
Cauchy defining rows $\mathbf u_j$, there are intrinsic
coefficient-Radon original currents
\[
 \widehat N_\gamma(\mathbf u),\qquad
 \widehat{\mathcal J}'_{\gamma,i},\quad
 \widehat{\mathcal J}''_{\gamma,i},\qquad
 \widehat G^\gamma_{ih}.
 \tag{IC10}
\]
Curvatures have $k=|\gamma|\leq q$; first-order and energy
currents have $k\leq q-1$. All energy degrees are included,
since $k+2\leq q+1=d$. Curvatures are separately closed,
symmetric and middle-positive; energy matrices are positive.
First-order and energy currents are not asserted to be closed.

For any continuous $F$ on the nonpositive orthant satisfying
$|F(x)|\leq A(1-\sum_i x_i)^a$, the actual multiplier limit
\[
                  F(\mathbf u_j)P_j
                         \longrightarrow F(\mathbf u)\widehat P
       \qquad(a+b\leq d)
 \tag{IC11}
\]
holds on every original compact smooth complementary test and
every actual moment refinement. Here the degree budget $b$ is
$k$, $k+1$, or $k+2$ for curvature, first-order, or energy.
All these multipliers are coefficient-integrable. The currents
and weighted limits are independent of defining rows and of
outer-capacity-null representatives, and they annihilate the
joint pole set. They satisfy all original recursive identities,
including
\[
 \begin{split}
 D(u_i\widehat N_\gamma)&=\widehat N_{\gamma+e_i},\\
 d'(u_i\widehat N_\gamma)&=\widehat{\mathcal J}'_{\gamma,i},\qquad
 d''(u_i\widehat N_\gamma)=\widehat{\mathcal J}''_{\gamma,i}
                         \qquad(k\leq q-1),
 \end{split}
 \tag{IC12}
\]
and the actual three-term energy identity (RC13).

For $q+1$ completed slots, the global mixed energy is finite and
\[
 \mathcal I_S(u_{0,j},\ldots,u_{q,j})
      \longrightarrow
 \widehat B(u_0,\ldots,u_q)
       :=\int_V(-u_0)\widehat N(u_1,\ldots,u_q).
 \tag{IC13}
\]
The corresponding positive weighted measures converge against
every bounded original smooth scalar test, including those
without compact support. The assertion also permits independent
indices tending to infinity in the defining rows.
\end{theorem}

\noindent\textbf{Status: \texttt{NEW\_PROVED}} for the strong
completion without a prescribed actual lower barrier.

\begin{proof}
For a finite tuple define
\[
 \Phi(R)=\sup_j\int_R^\infty t^q
          C_S\!\left(\left\{\sum_i u_{i,j}<-t\right\};V\right)dt.
 \tag{IC14}
\]
By Lemma~\ref{lem:r11-strong-cauchy-uniform-tails},
$\Phi(0)<\infty$ and $\Phi(R)\to0$; the same is true for
the supremum over independently large row indices, with the
finitely many earlier indices included if needed. At every
original compact test stage, the actual truncation estimates
(PG3), (FO8), and (PG12) give
\[
 \|\mu_{P_j-P_j^{(L)},\eta}\|_{\rm TV}
             \leq C_\eta L^{b-d}\Phi(L/4).
 \tag{IC15}
\]
The bounded analytic completion supplies intrinsic currents
$P^{(L)}$ of the completed clipped tuple. Passing $j$ to
infinity in the difference of two bounded stages shows that
$P^{(L)}$ are Cauchy in local coefficient total variation,
exactly as in (RC10). Their limit defines $\widehat P$.
Coefficient/refinement identities, positivity and separate
closedness pass on all original tests.

The limiting coefficients have the common modulus
\[
    \inf_{L\geq1}\{L^b s+L^{b-d}\Phi(L/4)\}
                              \longrightarrow0\quad(s\downarrow0).
 \tag{IC16}
\]
Its weak-limit proof keeps $L$ fixed until after ordinary-open
variation lower semicontinuity and Radon outer regularity.
Thus it controls both the defining and limiting coefficients,
including the critical case $b=d$. It annihilates every
locally outer-null set and, in particular, the joint pole set.
The bounded quasicontinuous multiplier argument therefore
passes every bounded continuous function of a fixed clipped
tuple.

For endpoint integrability apply that argument to
$\theta_R(z)\min((1+z)^a,H)$, where
$z=-\sum_i u_i$ and $\theta_R$ is zero below $R$ and one
above $2R$. For fixed $R,H$ this is a continuous function
of one sufficiently large clipped tuple; once a slot exceeds
that clipping level its value is the constant $H$.
When $a=0$, use $\theta_R$ itself. The actual row estimates
bound the integral by
$C R^{a+b-d}\Phi(R/8)$. Passing to the bounded multiplier
limit and then $H\uparrow\infty$ proves the same tail bound
for the limiting positive traces. Variation lower
semicontinuity handles signed first-order coefficients, and
positive diagonal controls handle cross energies.
Adding the bounded-value part gives actual coefficient
integrability. Clipping a general $F$ and then letting its
clipping level tend to infinity proves (IC11). Consequently
all multipliers in (IC12) and (RC13) are defined before
distributional differentiation, and the original row
identities pass to the limit.

For two different defining tuples with the same Borel
capacity limits, their bounded completed stages are identical.
The union of their finitely many row families has the same
vanishing-tail property. The preceding construction therefore
gives identical unbounded currents and actual multipliers.
This proves approximation independence without presupposing
that the two tuples are mutually strongly Cauchy.

Finally, (IC11) gives local convergence of (IC3).
The uniform global tightness in
Lemma~\ref{lem:r11-strong-cauchy-uniform-tails}, together
with the uniform finite mixed-energy bound (GE10), identifies
a finite global limiting measure and gives its total-mass
convergence. This is precisely (IC13). Tightness also passes
every bounded original scalar test. The independent-index
assertions follow from the same bounds, or by applying the
argument to every sequence of independently diverging indices.
\end{proof}

\subsection{Intrinsic discrepancy, completeness, and the exact scope}

Define on this Borel space
\[
 \begin{split}
 \widehat e(u)&=\widehat B(u^{[q+1]}),\\
 \widehat I(u,v)&=\widehat e(u)+\widehat e(v)
       -\widehat B(u,v^{[q]})-\widehat B(v,u^{[q]}).
 \end{split}
 \tag{IC17}
\]

\begin{lemma}[Uniform discrepancy continuity on energy balls]
\label{lem:r11-discrepancy-uniform-perturbation}
The completed space is closed under finite sums and nonnegative
scalar multiples. Its mixed energy is symmetric, positive and
multilinear on that cone, and satisfies the constant-one Holder
inequality. The formulas (EV1), (EV8), (EV9) and the sharp
midpoint inequality (SM8) hold in the completion, with their
actual Borel multipliers and local capacities. In particular,
$\widehat I\geq0$, and it vanishes exactly on equal relative
classes.

If $\widehat e(a),\widehat e(a'),\widehat e(b)\leq M$,
then
\[
 \big|\widehat I(a,b)-\widehat I(a',b)\big|
       \leq L_q M^{1-\alpha}\widehat I(a,a')^\alpha,
 \qquad
 L_q=2(q+1)2^{(q-1)\alpha}3^{d(1-\alpha)}.
 \tag{IC18}
\]
The assertion at $M=0$ means that both sides are zero.
\end{lemma}

\noindent\textbf{Status: \texttt{NEW\_PROVED}} for the uniform
discrepancy modulus; the global algebraic inequalities are
\texttt{DIRECT\_COROLLARY} extensions through (IC13).

\begin{proof}
For two strongly Cauchy defining rows, the sum rows have the
required Borel capacity limit. Expand the discrepancy of two
independently late sum rows into finitely many global mixed
energies. Equation (IC13) makes each term converge to the
corresponding completed mixed energy; the total limiting
discrepancy is zero. Thus the sum row is strongly Cauchy.
Positive scaling is immediate, and zero scaling means the
zero function even at poles. Passing the actual finite mixed
identities and (GE10) through (IC13) now proves all asserted
cone identities.

The actual inequalities (EV1), (EV9), and (SM8) pass by
global mixed-energy convergence. For (EV1), its left side is
the actual weight $|u-v|$, which has degree-one growth and is
bounded by $(-u)+(-v)$; (IC11) and the tightness of these two
critical positive mixed measures pass its global integral.
For (EV8), use defining rows, their convergent energies,
and the level-set triangle inequality for clipped gaps;
then remove deep sublevels using (IC9) and decrease the
intermediate gap threshold. This also proves that
$\widehat I=0$ gives equality in all local outer capacities.
The converse follows from intrinsic current and multiplier
independence.

It remains to prove (IC18). It suffices first to consider
actual potentials. Write $h=a-a'$. Telescoping the three
terms of $I(a,b)$ that depend on $a$, and using full mixed
symmetry, expresses their difference as $2(q+1)$ signed
integrals of $-h$ against positive mixed top products whose
slots lie among $a,a',b$. Every such product is bounded by
$\mu_\psi$, where $\psi=a+a'+b$, by the positive expansion.
Thus
\[
 |I(a,b)-I(a',b)|
                  \leq2(q+1)\int|a-a'|\,d\mu_\psi.
 \tag{IC19}
\]
Minkowski gives $e_1(\psi)\leq3^dM$ and
$e_1(a+a')\leq2^dM$. Apply (EV1) to (IC19) to obtain
(IC18). For completed potentials pass through actual
defining rows whose energies tend to the completed energies;
first allow $M+\epsilon$ and then let $\epsilon\downarrow0$.
The zero-energy case follows separately from Holder.
\end{proof}

\begin{theorem}[Completeness and the intrinsic strong topology]
\label{thm:r11-complete-strong-first-energy-space}
Every $\widehat I$-Cauchy sequence in
$\widehat{\mathcal E}^1_{\rm str}(V)$ converges to a unique
relative class in that space, with discrepancy tending to zero.
The actual classical class is strongly dense.

For members $u_j,u$ of this space, one has
\[
 \widehat I(u_j,u)\longrightarrow0
 \quad\Longleftrightarrow\quad
 \left\{
 \begin{array}{l}
 u_j\longrightarrow u\text{ in all local truncated outer capacities},\\
 \widehat e(u_j)\longrightarrow\widehat e(u).
 \end{array}\right.
 \tag{IC20}
\]
Strongly convergent finite tuples have every original local
weighted-current limit in (IC11), including all its critical
endpoints, and the global mixed-energy and weighted-measure
limits (IC13). These assertions do not claim local total
variation for arbitrary varying strongly convergent rows.

Every fixed actual lower-barrier completion
$\widehat{\mathcal E}^1_a$ of (RC2) embeds canonically in this
space, preserving all its original currents, actual multipliers
and mixed energies.
\end{theorem}

\noindent\textbf{Status: \texttt{NEW\_PROVED}} for the full relative
strong completion. This does not assert actual classical
completeness.

\begin{proof}
If $(v_l)$ is an actual defining row for a completed $u$, then
for fixed $l$, (IC13) gives
$\widehat I(v_l,u)=\lim_{r\to\infty}I(v_l,v_r)$.
Its Cauchy property therefore implies
$\widehat I(v_l,u)\to0$. This proves strong density.

Now let $(u_j)$ be completed and $\widehat I$-Cauchy.
The energy-root estimate bounds its energies. Choose actual
$v_j$ with $\widehat I(v_j,u_j)<2^{-j}$ and uniformly bounded
energies. Applying (IC18) in its two slots shows
\[
 I(v_j,v_l)\leq\widehat I(u_j,u_l)
                   +C\bigl(2^{-j\alpha}+2^{-l\alpha}\bigr).
 \tag{IC21}
\]
Hence $(v_j)$ is strongly Cauchy and has a completed Borel
limit $u$. Its defining property and strong density imply
$\widehat I(v_j,u)\to0$. A further use of (IC18) gives
$\widehat I(u_j,u)\to0$. Separation proves uniqueness.
This proves completeness using a proved uniform modulus,
rather than an unproved triangle inequality.

The forward implication of (IC20) follows from the completed
energy-root and capacity estimates. For the reverse, the
energies are uniformly bounded. Bounded completed stages and
the strict subcritical truncation bound for top unweighted
currents give local convergence of their top measures. More
explicitly, the stage error has order at most $C M/L$ for
degree $q<d$; first fix $L$, use bounded-completion capacity
continuity, and then let $L$ grow. The same uniform coefficient
modulus permits bounded quasicontinuous weight passage.
Consequently, by first truncating the weight and inserting a
compact nonnegative cutoff, then increasing both,
\[
 \liminf_j\widehat B(u_j,u^{[q]})\geq\widehat e(u),\qquad
 \liminf_j\widehat B(u,u_j^{[q]})\geq\widehat e(u).
 \tag{IC22}
\]
Holder gives upper bounds tending to $\widehat e(u)$ for
both cross energies. Expanding (IC17) proves the converse.
This argument uses no critical weighted continuity to prove
that very continuity.

For a strongly convergent completed tuple choose, at the $j$th
row, classical approximants with discrepancy at most $2^{-j}$,
and sufficiently accurate values on the first $j$ compact
test stages and bounded clipped weights. The defining-row
convergence allows these finitely many requirements
simultaneously. By (IC18), these actual rows are strongly
Cauchy, with the same Borel limits and energies. The preceding
construction and (IC11)--(IC13) give the asserted limits.
Equivalently, any failure on a fixed test would furnish a
subsequence and actual approximants contradicting those
already proved actual-row limits. This last argument also
covers bounded original scalar tests in the global statement.

Finally, classical defining rows of a fixed-barrier member
have joint critical mixed-energy convergence by (RC14), even
when two row indices tend independently. Expanding their
discrepancy therefore makes them strongly Cauchy. The Borel
representative is unchanged, and the intrinsic bounded stages
agree in the two constructions. Hence all local currents,
actual weighted multipliers and global mixed energies agree.
\end{proof}

\paragraph{The remaining source and representative gates.}
The completed space solves the strong Cauchy problem for the
relative first-energy geometry. It does not turn a capacity
class into an actual ambient classical psh function. It also
does not, by itself, extend an arbitrary datum integral
$u\mapsto\int u\,d\mu$ continuously to this Borel space.
Such a datum statement requires its own null-set and uniform
integrability hypotheses; finite energy-dual size alone is
not used as a substitute for them. The barrier-constrained
problem (BV7)--(BV8) retains its printed integrable barrier
and BT-diffuse datum assumptions. No unrestricted nonlinear
source equation, GQ, or low-real-degree CTO is inferred.

% END REVIEWED MODULE research_r11/comparison/strong_first_energy_completion.tex


% BEGIN REVIEWED MODULE research_r11/capacity/strong_completed_comparison.tex
\section{Comparison, locality and finite-mass approximation in the full strong completion}
\label{sec:r11-strong-completed-comparison}

Assume exactly the strict analytic box and background extension
(TC1)--(TC2), with the original arbitrary symmetric middle-positive
separately closed $S$ and $q=n-p\geq1$.
Use the distinct Borel space
$\widehat{\mathcal E}^1_{\rm str}(V)$ of
Definition~\ref{def:r11-strong-first-energy-completion}.
Its boundary condition is approximation by actual first-energy
rows Cauchy in the proved discrepancy. No common actual lower
barrier and no compact negative sublevel closure of a limiting
representative is imposed here. The original classical class
remains unchanged.

Write $\widehat\mu_u=\widehat N_q(u)$,
$\widehat e(u)=\widehat B(u^{[q+1]})$ and
$\widehat I$ as in (IC17). All quantities below use those
independently constructed original currents and actual Borel
coefficient multipliers. Put $d=q+1$ and $\alpha=2^{-q}$.
The global finite energy and tightness theorems (IC2)--(IC13)
are essential inputs. The present argument does not deduce
tightness from local weak convergence alone.

\begin{lemma}[Strong continuity of maxima and hard truncations]
\label{lem:r11-strong-max-clipping}
The strong completion is closed under
\[
                 (u,v)\longmapsto\max(u,v-c)\quad(c\geq0),
                 \qquad u\longmapsto\max(u,-L)\quad(L>0).
 \tag{CC1}
\]
These operations are strongly continuous on bounded-energy sets.
For fixed $u,v$, the maximum is also strongly continuous in
the scalar shift $c\geq0$.
More precisely, for $u,u',v,v'$ of energy at most $M$, put
$x=\max(u,v-c)$ and $x'=\max(u',v'-c)$. Then
\[
 \widehat I(x,x')
  \leq K_q M^{1-\alpha}
       \bigl(\widehat I(u,u')^\alpha+
                         \widehat I(v,v')^\alpha\bigr),
 \qquad K_q=2^{1+(q-1)\alpha+d(1-\alpha)}.
 \tag{CC2}
\]
For the two hard truncations at the same depth $L$, the same
bound holds with only the first discrepancy on the right.
The zero-energy case means zero on both sides. In addition,
\[
                  \widehat I(u,\max(u,-L))\longrightarrow0
                                   \quad(L\to\infty).
 \tag{CC3}
\]
\end{lemma}

\noindent\textbf{Status: \texttt{NEW\_PROVED}} for the strong
completion without prescribed actual lower barriers.

\begin{proof}
First consider actual members of $\mathcal E^1$.
Their maxima and hard truncations remain actual members by
upward stability. Since $x\geq u$ and $x'\geq u'$, global
energy monotonicity (GE16)--(GE17) bounds their energies by $M$.
The actual pointwise inequality
\[
                    |x-x'|\leq|u-u'|+|v-v'|
\]
holds off the joint poles, which every involved measure ignores.
Expanding the discrepancy and using positivity gives
\[
 I(x,x')\leq\int|x-x'|(d\mu_x+d\mu_{x'}).
\]
Apply (EV1) to each difference on the right with tester $x$
or $x'$. Minkowski bounds $e_1(u+u')$ and $e_1(v+v')$ by
$2^dM$. The two tester terms give exactly the constant $K_q$
in (CC2). For hard truncations use instead
$|\max(u,-L)-\max(u',-L)|\leq|u-u'|$.

Now choose actual strongly Cauchy defining rows of the
completed potentials. The just proved estimate makes the
rows of maxima, or fixed hard truncations, strongly Cauchy.
Their local bounded truncations have the prescribed Borel
capacity limits, by the scalar Lipschitz maximum and clipping
identities. Thus they define precisely the members in (CC1).
The global mixed-energy convergence (IC13) passes (CC2) to
completed members; energies may first be bounded by $M+\epsilon$
and then $\epsilon\downarrow0$.

We verify the varying-shift assertion separately. It does not
follow merely from a uniform scalar bound when total curvature
mass near the boundary is infinite. Put
$u^\delta=\max(u,-\delta)$ for $\delta>0$. Positivity and
full mixed-energy symmetry make $\widehat B$ decreasing in
each potential with respect to pointwise order. Hence
\[
 \widehat e(u^\delta)
    \leq\widehat B(u^\delta,u^{[q]})
    =\int\min(\delta,-u)\,d\widehat\mu_u
                                      \longrightarrow0.
 \tag{CC2a}
\]
The last limit is dominated convergence by the finite first
energy. For $w_c=\max(u,v-c)$ and $\delta=|c-c'|>0$,
both maxima lie between $u$ and zero, and
$|w_c-w_{c'}|\leq\min(\delta,-u)=-u^\delta$.
The discrepancy formula and completed Holder therefore give
\[
 \widehat I(w_c,w_{c'})
   \leq\widehat B(u^\delta,w_c^{[q]})
          +\widehat B(u^\delta,w_{c'}^{[q]})
   \leq2\widehat e(u^\delta)^{1/d}\widehat e(u)^{q/d}
                                      \longrightarrow0.
 \tag{CC2b}
\]
Their energies are bounded by $\widehat e(u)$, by the already
proved maximum construction and actual energy monotonicity.
This proves the required strong passage as the shift varies,
including $c\downarrow0$.

For $u^{(L)}=\max(u,-L)$, both terms in the actual integrated
discrepancy are defined by (IC13). Positivity gives
\[
 0\leq\widehat I(u,u^{(L)})
   =\int(u^{(L)}-u)(d\widehat\mu_u-d\widehat\mu_{u^{(L)}})
   \leq\int(u^{(L)}-u)\,d\widehat\mu_u\longrightarrow0.
 \tag{CC4}
\]
The last limit is dominated convergence by the finite weight
$(-u)\widehat\mu_u$. Joint poles have zero mass for both
measures before the displayed difference is integrated.
\end{proof}

\begin{theorem}[Strong completed comparison with an energy boundary condition]
\label{thm:r11-strong-completed-comparison}
For arbitrary $u,v\in\widehat{\mathcal E}^1_{\rm str}(V)$ and
every $t\geq0$,
\[
          \boxed{\quad
          \widehat\mu_v\{u<v-t\}
                          \leq\widehat\mu_u\{u<v-t\}.
          \quad}
 \tag{CC5}
\]
For $t>0$ both sides are finite; specifically
\[
 \widehat\mu_v\{u<v-t\}
          \leq t^{-1}\widehat B(u,v^{[q]}),\qquad
 \widehat\mu_u\{u<v-t\}\leq t^{-1}\widehat e(u).
 \tag{CC6}
\]
At $t=0$ the comparison uses extended nonnegative masses.
Every current is the original current, and every gap uses
the completed Borel representative. No ordinary boundary-gap
condition on that representative is being silently assumed.
\end{theorem}

\noindent\textbf{Status: \texttt{NEW\_PROVED}} beyond the
fixed actual lower-barrier comparison (CP4).

\begin{proof}
Choose actual strongly Cauchy rows $u_j,v_j$.
For every bounded continuous scalar $F$, (IC11) passes
$F(v_j-u_j)\mu_{u_j}$ and $F(v_j-u_j)\mu_{v_j}$ on every
original compact scalar test. The poles are null for all
these measures. Exactly the positive-trace scalar strip
argument in (CP1)--(CP2) therefore passes the moving-gap
restrictions locally at every positive level outside a
countable exceptional set. A countable compact exhaustion
allows that set to work for both top measures on all of $V$.

We now justify the global passage, which cannot use a common
compact carrier in this completion. On
$A_{j,t}=\{u_j<v_j-t\}$, one has $-u_j>t$, and hence
\[
 \mathbf1_{A_{j,t}}\mu_{v_j}
             \leq t^{-1}(-u_j)\mu_{v_j},\qquad
 \mathbf1_{A_{j,t}}\mu_{u_j}
             \leq t^{-1}(-u_j)\mu_{u_j}.
 \tag{CC7}
\]
Both positive measures on the right are uniformly globally
tight by (IC3)--(IC7), and their full masses are uniformly
bounded by the mixed-energy inequality. The same inequalities
for the limiting measures follow directly from their actual
Borel coefficients and finite completed mixed energies;
these are (CC6).

For a regular positive $t$, first pass a compact smooth cutoff
in the locally restricted measures. Then exhaust compact sets:
(CC7) bounds the discarded mass uniformly in the row index,
and its limiting counterpart bounds the limiting discarded
mass. Thus the total moving-gap masses converge globally.
This step uses the independently proved boundary tightness,
not a supposed compactness of limiting negative sublevels.

Each actual row satisfies the shifted comparison
$\mu_{v_j}\{u_j<v_j-t\}\leq\mu_{u_j}\{u_j<v_j-t\}$
by applying (UC10) to $(u_j,v_j-t)$; constant-shift invariance
identifies the latter top measure with $\mu_{v_j}$.
Its individual boundary-gap hypothesis follows
from $u_j\in\mathcal E^1$; uniformity of those compact
carriers is unnecessary. Passing total masses proves (CC5)
at regular levels. For arbitrary $t\geq0$, choose regular
positive levels decreasing to $t$ and use measure continuity
from below. This proves the theorem.
\end{proof}

\begin{theorem}[Strong completed outer-capacity domination]
\label{thm:r11-strong-completed-domination}
For the same $u,v$, every strict gap $A_t=\{u<v-t\}$ is
relatively quasi-open and satisfies $C_S^*(A_t;V)=C_S(A_t;V)$.
For $s,t>0$,
\[
        t^q C_S^*(\{u<v-s-t\};V)
                         \leq\widehat\mu_u\{u<v-s\}.
 \tag{CC8}
\]
If $A=\{u<v\}$ satisfies either
$\widehat\mu_u(A)=0$ or
$\mathbf1_A\widehat\mu_u\leq\mathbf1_A\widehat\mu_v$,
then $C_S^*(A;V)=0$. Consequently
\[
      \widehat\mu_u=\widehat\mu_v
             \quad\Longrightarrow\quad C_S^*(\{u\ne v\};V)=0.
 \tag{CC9}
\]
\end{theorem}

\noindent\textbf{Status: \texttt{NEW\_PROVED}} in the full
strong completion, with exactly its encoded energy boundary class.

\begin{proof}
The strong completion has finite local outer quasicontinuity
by Proposition~\ref{prop:r11-strong-cauchy-borel-limit}.
Summing exceptional capacities over a countable compact
exhaustion makes $u,v$ finite and continuous on the complement
of one ordinary open set of arbitrarily small full fixed-domain
capacity. The proof of (CP6) then applies word for word; it
used only this quasicontinuity and subadditivity, not compact
negative sublevels. This proves the stated gap regularity.

For $h\in\mathscr T_0(V)$, the sum $w=v+t h$ is a member
of the strong completion by (IC18), and its positive product
expansion gives $\widehat\mu_w\geq t^q\mu_h$.
Because $v-t\leq w\leq v$, comparison (CC5) gives exactly
(CP10), now for these strong completed functions. Taking the
complete core tester supremum (TC6), and then the proved gap
raw/outer equality, proves (CC8).

The gap-zero-mass conclusion follows by increasing the
positive-margin gaps, as in (CP7)--(CP8). Under the restricted
measure order use $w=(1+t)v+t h$ and $B_s=\{u<w-s\}$.
It is contained in $A$ and in $\{u<-s\}$.
All measures which occur in the subtraction are finite there:
\[
 \widehat\mu_v(B_s)\leq s^{-1}\widehat B(u,v^{[q]}),\qquad
 \widehat\mu_u(B_s)\leq s^{-1}\widehat e(u).
\]
These replace the compact-support argument used in the barrier
case. Positivity and (CC5) give
\[
 (1+t)^q\widehat\mu_v(B_s)+t^q\mu_h(B_s)
        \leq\widehat\mu_u(B_s)\leq\widehat\mu_v(B_s).
 \tag{CC10}
\]
Subtracting the finite masses yields $\mu_h(B_s)=0$.
First let $s\downarrow0$ and then $t\downarrow0$; since $h$
is bounded, the increasing union is the actual extended strict
gap $A$. The complete core tester identity and gap regularity
give $C_S^*(A;V)=0$. Apply this in both orders to obtain (CC9).
\end{proof}

\begin{corollary}[Intrinsic locality and carriers in the strong completion]
\label{cor:r11-strong-completed-locality-carriers}
All strict maximum identities (CL1), exact multituple truncation
identities (CL4)--(CL6), and weak maximum inequalities (CL8)
hold for members of $\widehat{\mathcal E}^1_{\rm str}(V)$,
with every original compact complementary test and every
actual moment refinement. The uncut and restricted canonical
bounded stages converge locally in coefficient total variation.
Thus finite maxima of strong completed subsolutions remain
subsolutions. No common actual lower barrier is required.
\end{corollary}

\noindent\textbf{Status: \texttt{NEW\_PROVED}} as the extension
of the carrier calculus to the full strong completion.

\begin{proof}
Lemma~\ref{lem:r11-strong-max-clipping} supplies the strong
membership of every maximum and every hard truncation used
in the proof of (CL1)--(CL9). In particular their actual
defining rows are strongly Cauchy, so (IC11) supplies every
continuous tuple-weight passage on each fixed original stage.
Positive limiting traces make the scalar or rectangular strip
errors vanish at regular margins. Passing actual strict
maximum locality at such margins and then decreasing the
margins proves (CL1) and (CL4), exactly as in their printed
proofs. Constants are again handled by scalar values and are
not asserted to have the strong boundary condition.

The coefficient-null pole union gives the restricted intrinsic
representation (CL5). The intrinsic bounded stages converge
in local coefficient total variation by (IC15), whose error
is $C L^{k-q-1}\Phi(L/4)\to0$ for curvature degree $k\leq q$.
Adding the limiting current's pole-tail mass proves (CL6).
Finally use the already proved exact strict locality on
$\max(u,v-\varepsilon)$ before taking any limit. This puts
the fixed equality-side restriction on the right of (CL9).
Its other restriction increases; the varying-shift estimate
(CC2a)--(CC2b) proves strong convergence, so (IC11) passes
the left side.
Thus (CL8) holds even on a charged equality locus, and is
tested against the full original positive complementary cone.
\end{proof}

\begin{corollary}[Canonical bounded finite-mass approximation]
\label{cor:r11-strong-finite-mass-barrier-approximation}
Let $\rho$ be the strict negative finite-mass barrier (TC3).
For every $u\in\widehat{\mathcal E}^1_{\rm str}(V)$, put
\[
                             v_j=\max(u,j\rho).
\]
Then $v_j$ lies in the actual-barrier completion
$\widehat{\mathcal E}^1_{j\rho}$, is bounded in $[-j,0]$, and
\[
 \begin{split}
 \widehat\mu_{v_j}(V)&\leq j^q\mu_\rho(V)<\infty,\\
 v_j&\downarrow u,\qquad
 \widehat I(v_j,u)\longrightarrow0,
 \qquad \widehat e(v_j)\longrightarrow\widehat e(u).
 \end{split}
 \tag{CC11}
\]
All original critical local current and global weighted-measure
limits of (IC11)--(IC13) apply to this sequence.
The approximants are completed Borel potentials; actual
classical representability is not part of this assertion.
\end{corollary}

\noindent\textbf{Status: \texttt{NEW\_PROVED}} for the full
strong completion. The actual-class finite-mass approximation
was already proved in (BA8)--(BA9).

\begin{proof}
The maximum lemma proves strong membership. Applying it to
an actual defining row of $u$ gives actual defining functions
between the same actual barrier $j\rho$ and zero. This also
proves membership in the fixed-barrier completion, with the
same intrinsic products by (IC20).

For $c>1$, strict negativity of $\rho$ implies
$cj\rho<j\rho\leq v_j$ everywhere on $V$.
Comparison (CC5) at level zero therefore gives
$\widehat\mu_{v_j}(V)\leq c^qj^q\mu_\rho(V)$.
Let $c\downarrow1$ to prove the finite-mass bound. It is not
obtained by asserting a fixed ordinary collar for $v_j$.

Pointwise $v_j\downarrow u$ because $\rho<0$ everywhere.
The discrepancy identity and positivity yield
\[
 0\leq\widehat I(u,v_j)
      \leq\int(v_j-u)\,d\widehat\mu_u\longrightarrow0,
\]
by domination with $(-u)\widehat\mu_u$. The completed energy-root
inequality gives the energy convergence. The strong continuity
theorem then supplies all the printed current and weighted-measure
limits. The defining strong boundary class, not ordinary
pointwise boundary values, is retained throughout.
\end{proof}

\paragraph{Scope.}
This section removes prescribed actual lower barriers from
comparison, domination, locality and finite-mass approximation
in the first-energy completion. It retains the precise strong
approximation condition and the strict analytic domain.
Uniqueness does not imply source existence: the strong completed
compact-source obstruction (SO1)--(SO6) remains in force.
Neither this completion nor its inequalities identify actual
analytic PS/PL presentations or close low-real-degree CTO.

% END REVIEWED MODULE research_r11/capacity/strong_completed_comparison.tex


% BEGIN REVIEWED MODULE research_r11/capacity/completed_capacity_layers.tex
\section{Intrinsic capacity layers of completed potentials}
\label{sec:r11-completed-capacity-layers}

Retain the exact hypotheses of the barrier completion and
Sections~\ref{sec:r11-completed-comparison-domination}--
\ref{sec:r11-completed-locality-carriers}, or the full strong
completion of Section~\ref{sec:r11-strong-completed-comparison}.
Fix an already constructed potential of either kind:
$u\in\widehat{\mathcal E}^{\sigma}_a(V)$ for some $\sigma>0$
and an actual lower barrier $a$, or
$u\in\widehat{\mathcal E}^1_{\rm str}(V)$ without a prescribed
actual lower barrier. The statements below concern its canonical top measure
$\widehat\mu_u$ and its actual Borel values. They do not assume
or assert a classical representative.

For every real $s>0$, define the extended nonnegative quantities
\[
 \widehat e_s(u)=\int_V(-u)^s\,d\widehat\mu_u,
 \qquad
 \widehat J_s(u)=\int_0^\infty
                   t^{q+s-1}C_S^*(\{u<-t\};V)\,dt.
 \tag{CY1}
\]
The pole is null for the measure and outer capacity. Its value
there is omitted from the integral without changing that
measure. The scalar capacity factor is nonincreasing and
therefore Borel measurable. By (CP6), or by the strong-gap
regularity preceding (CC8), raw capacity gives
exactly the same $\widehat J_s$, although these completed
sublevels need not be ordinary open.

\begin{lemma}[Exact top capacity bound for bounded completed stages]
\label{lem:r11-completed-bounded-top-capacity}
If $z$ belongs to the bounded capacitary completion with
$-L\leq z\leq0$, then
\[
                    T^{\mathrm{comp}}_{S,q}(z)(E)
                                  \leq L^q C_S^*(E;V)
                      \quad(E\subset V\text{ Borel}).
 \tag{CY2}
\]
The constant is one; no compact-test dependent coefficient
constant occurs in this top-degree assertion.
\end{lemma}

\noindent\textbf{Status: \texttt{DIRECT\_COROLLARY}} of the
bounded completed construction and the original full tester
definition, with the ordinary outer step made explicit.

\begin{proof}
Take actual bounded defining rows $z_l\in[-L,0]$.
The functions $1+z_l/L$ are normalized original testers, so
$\mu_{z_l}(E)\leq L^q C_S(E;V)$ on every Borel set.
For an ordinary open $G$, weak convergence of the positive
Radon measures, tested first on compact smooth cutoffs in
$G$, yields
$T^{\mathrm{comp}}_{S,q}(z)(G)\leq L^q C_S(G;V)$.
Individual Radon outer regularity of the limiting measure and
the infimum over such open covers of $E$ give (CY2).
This last step proves an outer-capacity bound, not a raw
bound on arbitrary Borel sets.
\end{proof}

\begin{theorem}[Completed power-energy and capacity-layer equivalence]
\label{thm:r11-completed-energy-capacity-layers}
For every $s>0$, including values different from the exponent
used to construct $u$, one has
\[
       \boxed{\qquad
       \frac{s}{2^{q+s}}\widehat J_s(u)
             \leq\widehat e_s(u)
             \leq\frac{s2^q}{1-2^{-s}}\widehat J_s(u).
       \qquad}
 \tag{CY3}
\]
Both inequalities hold for extended nonnegative integrals.
In particular the finite $s$th moment of the completed top
measure is equivalent to the printed intrinsic capacity-layer
integrability. At $s=\sigma$ in the barrier case, and at $s=1$
in the strong case, both quantities are finite by the already
established completion theorem.
\end{theorem}

\noindent\textbf{Status: \texttt{NEW\_PROVED}} as an extension
of the actual-class theorem (PE3) to completed Borel potentials,
with the same constants and the full original capacity.

\begin{proof}
Use (CP7) in the barrier case or (CC8) in the strong case.
Put $v=0$ and make its two positive scalar parameters equal
to $r$. This gives
\[
                    r^q C_S^*(\{u<-2r\};V)
                                 \leq\widehat\mu_u\{u<-r\}.
 \tag{CY4}
\]
Multiply by $r^{s-1}$ and integrate. Tonelli and the ordinary
layer-cake identity give
$2^{-q-s}\widehat J_s(u)\leq\widehat e_s(u)/s$.
This is the first inequality of (CY3), without a finiteness
assumption on either integral.

For the other inequality fix $r>0$ and $m\geq0$, and put
\[
                         A_{m,r}
                 =\{2^m r<-u\leq2^{m+1}r\}.
\]
Clip $u$ at depth $L=2^{m+1}r+\varepsilon$.
The strict carrier identity (CL4), or its strong extension in
Corollary~\ref{cor:r11-strong-completed-locality-carriers}, identifies its top measure
with $\widehat\mu_u$ on $A_{m,r}$.
By (CY2) and monotonicity of outer capacity,
\[
 \widehat\mu_u(A_{m,r})
    \leq (2^{m+1}r+\varepsilon)^q
                           C_S^*(\{u<-2^m r\};V).
\]
Let $\varepsilon\downarrow0$. The positive extra clipping
margin ensures that atoms at the upper endpoint of the
annulus have not been lost. Summing the disjoint annuli and
using pole nullity yields
\[
 \widehat\mu_u\{u<-r\}
    \leq\sum_{m=0}^\infty(2^{m+1}r)^q
                           C_S^*(\{u<-2^m r\};V).
 \tag{CY5}
\]
Multiply by $s r^{s-1}$ and integrate over $r>0$.
Tonelli justifies every exchange. The substitution $t=2^m r$
makes the $m$th summand equal to
$s2^q2^{-ms}\widehat J_s(u)$.
The geometric series gives the second inequality of (CY3).
\end{proof}

\paragraph{What this characterizes.}
The theorem characterizes the moment of the canonical measure
of an already defined completed potential. It does not say
that an arbitrary Borel function with finite $\widehat J_s$
belongs to the completion, or that it has an actual
$\mathcal E^s$ lower barrier. Nor does it turn raw and outer
capacities into the same set function on every analytic Borel
set. The equality used here is the proved one for completed
strict sublevels. All products and measures retain the original
$S$, original absolute value and original test sheaf.

% END REVIEWED MODULE research_r11/capacity/completed_capacity_layers.tex


% BEGIN REVIEWED MODULE research_r11/capacity/intrinsic_first_energy_characterization.tex
\section{An intrinsic capacity characterization of the strong first-energy class}
\label{sec:r11-intrinsic-first-energy-characterization}

Retain the strict analytic box, original background extension and
arbitrary original $S$ of (TC1)--(TC2). Let
$\mathcal A_S^{[-R,0]}(V)$ denote the global bounded capacitary
completion of Lemma~\ref{lem:r11-analytic-bounded-completion}:
its members are Borel limits, in every local outer capacity,
of actual bounded classical functions on the same whole domain
$V$, with values in $[-R,0]$.
The word ``global'' specifies that admissibility domain; it
does not assert a uniform convergence on all of $V$.

The theorem below replaces a strongly Cauchy approximating row
by an intrinsic condition on all bounded stages and one
capacity integral. The bounded-stage admissibility condition
is retained explicitly. An arbitrary Borel function with a
finite numerical capacity integral is not asserted admissible.

\begin{lemma}[Ordered discrepancy is bounded by energy loss]
\label{lem:r11-ordered-completed-discrepancy}
If $a,b\in\widehat{\mathcal E}^1_{\rm str}(V)$ and
$a\leq b$ outside a locally outer-capacity-null set, then
\[
       0\leq\widehat I(a,b)
                            \leq\widehat e(a)-\widehat e(b).
 \tag{CI1}
\]
\end{lemma}

\noindent\textbf{Status: \texttt{DIRECT\_COROLLARY}} of the
proved completed symmetric mixed energy and its positivity.

\begin{proof}
The completed mixed energy is decreasing in each potential
with respect to pointwise order: in the undifferentiated
slot this is positivity of the product measure, and symmetry
gives it in every slot. All products ignore the exceptional
outer-null set. Therefore
\[
       \widehat B(a,b^{[q]})\geq\widehat e(b),\qquad
       \widehat B(b,a^{[q]})\geq\widehat e(b).
\]
Insert these two inequalities in the definition (IC17) of
$\widehat I$. Its nonnegativity is already proved. This gives
(CI1), with finite global integrals throughout.
\end{proof}

\begin{theorem}[Bounded-stage admissibility and one capacity integral]
\label{thm:r11-intrinsic-strong-first-energy}
For an actual Borel function $u:V\to[-\infty,0]$, the following
are equivalent:
\begin{enumerate}
\item $u\in\widehat{\mathcal E}^1_{\rm str}(V)$.
\item For every positive integer $m$,
$u^{(m)}=\max(u,-m)$ belongs to
$\mathcal A_S^{[-m,0]}(V)$, and
\[
                J_1^*(u):=\int_0^\infty
                       t^q C_S^*(\{u<-t\};V)\,dt<\infty.
 \tag{CI2}
\]
\end{enumerate}
In the second condition no current of the unbounded function
is assumed in advance. Once either condition holds, the
canonical original currents and all critical first-energy
identities are those of (IC10)--(IC20), and
\[
       2^{-q-1}J_1^*(u)\leq\widehat e(u)
                                  \leq2^{q+1}J_1^*(u).
 \tag{CI3}
\]
The capacity integral is unchanged if raw capacity is used.
This last equality is a conclusion for the resulting completed
strict sublevels, not a premise of the second condition.
\end{theorem}

\noindent\textbf{Status: \texttt{NEW\_PROVED}} as a converse
capacity characterization of the full strong Borel class,
with its bounded-stage admissibility gate retained.

\begin{proof}
Suppose first that $u$ belongs to the strong completion.
An actual strongly Cauchy defining row, clipped at $-m$,
is a sequence of actual global bounded classical functions
in $[-m,0]$ converging to $u^{(m)}$ in the required local
outer capacities. This proves every bounded-stage condition.
The completed capacity-layer theorem (CY3) at $s=1$ gives
(CI2) and the lower inequality of (CI3).

Conversely assume precisely the second condition. Let $\rho$
be the actual barrier (TC3), so $-1\leq\rho<0$ on $V$ and
$\rho$ has finite total top mass. For each integer $j\geq1$
define the Borel function
\[
                     v_j=\max(u,j\rho)
                         =\max(u^{(j)},j\rho).
 \tag{CI4}
\]
The equality follows from $j\rho\geq-j$.
Choose global actual bounded defining rows for $u^{(j)}$.
Taking their maximum with $j\rho$ gives actual classical
functions between $j\rho$ and zero, converging to $v_j$
in the bounded local capacity topology. Each such function
belongs to $\mathcal E^1$ by upward stability above the
actual first-energy barrier $j\rho$. Thus
$v_j\in\widehat{\mathcal E}^1_{j\rho}$ by the definition
of the barrier completion, and hence belongs to the strong
completion by its proved embedding. This step defines
products only for $v_j$, not yet for the given $u$.

Since $v_j\geq u$, monotonicity of outer capacity gives
$\widehat J_1(v_j)\leq J_1^*(u)$.
The already proved theorem (CY3), applied to the valid
completed member $v_j$, therefore gives
\[
                    \widehat e(v_j)\leq2^{q+1}J_1^*(u).
 \tag{CI5}
\]
Pointwise $v_{j+1}\leq v_j$. Hence their energies increase
and are bounded by (CI5). For $l\geq j$, (CI1) gives
\[
                  \widehat I(v_l,v_j)
                           \leq\widehat e(v_l)-\widehat e(v_j).
 \tag{CI6}
\]
The increasing bounded energy sequence has a finite limit,
so $(v_j)$ is strongly Cauchy. Completeness (IC20) supplies
a strong completed limit $v$.

It remains to identify its actual capacity class with $u$;
pointwise convergence alone is not used for that purpose.
Fix a compact $H\Subset V$ and a clipping depth $R>0$.
The continuous strict negativity of $\rho$ makes
$\max_H\rho<0$. For all sufficiently large $j$, one has
$j\rho<-R$ on $H$, and therefore the exact scalar identity
\[
                              v_j^{(R)}=u^{(R)}
                                      \quad\hbox{on }H
 \tag{CI7}
\]
holds. Strong convergence gives local clipped capacity
convergence to $v$. Equation (CI7) identifies that limit
with $u^{(R)}$ on every such compact core. The countable
core and integer-depth definition of the completion thus
identifies $u$ with the strong completed limit, using its
specified Borel representative. This proves membership.

The full completed current and energy calculus now applies
to $u$. Applying (CY3) at $s=1$ gives both inequalities in
(CI3); the strict-gap regularity preceding (CC8) gives the
raw/outer equality for its sublevels. No unbounded current
of $u$ was used before it was constructed through completion.
\end{proof}

\begin{corollary}[Monotone closure with the exact energy hypothesis]
\label{cor:r11-strong-monotone-closure}
Let $u_j$ be members of
$\widehat{\mathcal E}^1_{\rm str}(V)$, monotone outside one
locally outer-null set, and let $u$ be their actual pointwise
monotone limit there, with any nonpositive extended Borel values
on that null set.
If $u_j$ increases, then $u$ belongs to the strong completion
and $\widehat I(u_j,u)\to0$.
If $u_j$ decreases, the same conclusion holds provided
\[
                            \sup_j\widehat e(u_j)<\infty.
 \tag{CI8}
\]
Every original critical local current and global mixed
weighted-measure convergence assertion (IC11)--(IC13)
then holds for finite tuples satisfying these hypotheses.
Ordinary upper semicontinuous regularization of the limit
is not part of this conclusion.
\end{corollary}

\noindent\textbf{Status: \texttt{NEW\_PROVED}} for monotone
sequences in the full strong completion; the increasing
bounded relative precursor was already frozen in R2.

\begin{proof}
For an increasing sequence, the nonnegative energies decrease
and are bounded above by the first energy. For a decreasing
sequence they increase, with the explicit bound (CI8).
In either case they converge. Equation (CI1), applied in
the appropriate order to two terms, shows that the sequence
is strongly Cauchy. Completeness gives a strong limit $v$.

Choose a subsequence whose fixed bounded truncations converge
pointwise off a locally outer-null set, using the usual
summable-capacity diagonal construction over compact cores
and integer depths. Monotonicity makes that pointwise limit
equal to the specified $u$. Thus $u=v$ in the completed
capacity class, and the entire sequence converges strongly.
The strong continuity theorem gives all the stated current
and weighted-measure limits. No current equality outside its
original test sheaf is asserted.
\end{proof}

\paragraph{Exact scope.}
The intrinsic condition (CI2) still requires every bounded
truncation to be an admissible global capacitary limit of
actual bounded classical functions. It therefore does not
turn a numerical capacity bound into ambient psh geometry.
It removes a prescribed actual lower barrier and a prescribed
strong approximating row from the characterization of the
completed first-energy class, while retaining the original
background, absolute value and all original tests. It does
not prove arbitrary-source Dirichlet existence or general
raw-capacity outer regularity.

% END REVIEWED MODULE research_r11/capacity/intrinsic_first_energy_characterization.tex


% BEGIN REVIEWED MODULE research_r11/capacity/strong_completed_perron.tex
\section{Relative Perron envelopes of arbitrary strong completed families}
\label{sec:r11-strong-completed-Perron}

Retain exactly the strict analytic box, background extension
and full strong first-energy class of
Section~\ref{sec:r11-strong-completed-comparison}.
The background is the unchanged arbitrary original symmetric
middle-positive separately closed $S$.
Write $f\preceq g$ if $f\leq g$ locally outside an
ordinary-open outer $S$-capacity-null set. This order concerns
the completed Borel representatives. No ordinary ambient
upper regularization is used.

The frozen R2 relative capacitary envelope theorem is
\texttt{INTERNAL\_INPUT} in its bounded scope. Theorem
\ref{thm:r11-barrier-relative-Perron} extends an actual family
above one prescribed actual power-energy barrier. The theorem
below instead permits an arbitrary nonempty family of strong
completed first-energy potentials, with no common actual lower
barrier. Its countable selection is proved directly by global
energy minimization over finite maxima.

\begin{theorem}[The full strong completed relative envelope]
\label{thm:r11-strong-completed-Perron}
Let $\mathscr F\subset\widehat{\mathcal E}^1_{\rm str}(V)$
be any nonempty family, possibly uncountable. Let
$\mathscr F^\vee$ consist of all nonempty finite maxima of
its members, and put
\[
                    L=\inf_{w\in\mathscr F^\vee}\widehat e(w).
 \tag{SE1}
\]
There is a unique relative class
$H\in\widehat{\mathcal E}^1_{\rm str}(V)$ with the following
least-upper-bound properties:
\[
 \begin{split}
 &f\preceq H\qquad(f\in\mathscr F),\\
 &\bigl[f\preceq g\text{ for every }f\in\mathscr F,
                      \ g\text{ Borel}\bigr]
                                      \Longrightarrow H\preceq g.
 \end{split}
 \tag{SE2}
\]
The exceptional set in the first line may depend on $f$;
there is no uncountable-union nullity assertion.

There are countably many selected members of $\mathscr F$
whose cumulative finite maxima $v_j$ satisfy
\[
                  v_j\uparrow H,\qquad
                  \widehat I(v_j,H)\longrightarrow0,
                  \widehat e(H)=L.
 \tag{SE3}
\]
Here the actual increasing Borel limit may be used as the
representative of $H$. More generally, every sequence
$w_j\in\mathscr F^\vee$ with $\widehat e(w_j)\to L$
converges strongly to $H$ and therefore has all the critical
current and mixed-energy limits (IC11)--(IC13).
\end{theorem}

\noindent\textbf{Status: \texttt{NEW\_PROVED}} for arbitrary
completed families and their full energy boundary class.

\begin{proof}
Finite maxima are members of the strong completion by (CC1).
Their energies are nonnegative, and one fixed family member
has finite energy. Thus $0\leq L<\infty$.
Choose $w_j\in\mathscr F^\vee$ with
$\widehat e(w_j)\leq L+1/j$, and put
\[
                            v_j=\max(w_1,\ldots,w_j).
 \tag{SE4}
\]
Each is still a finite maximum of members of the given family.
It increases with $j$, and energy monotonicity gives
$L\leq\widehat e(v_j)\leq\widehat e(w_j)$.
The increasing-closure theorem
\ref{cor:r11-strong-monotone-closure} therefore gives its
actual Borel limit $H$ in the strong completion, with strong
convergence and $\widehat e(H)=L$.
Only countably many members of the given family occur in all the
selected $w_j$; this proves the countable assertion in (SE3).

Fix an arbitrary $f\in\mathscr F$ and set $z_j=\max(v_j,f)$.
Then $z_j\in\mathscr F^\vee$ and
\[
                       L\leq\widehat e(z_j)
                                      \leq\widehat e(v_j).
\]
The strong maximum continuity (CC2) gives
$z_j\to z:=\max(H,f)$ strongly. Hence
$\widehat e(z)=L=\widehat e(H)$.
Since $H\leq z$, the ordered discrepancy estimate (CI1)
gives
\[
               0\leq\widehat I(H,z)
                      \leq\widehat e(H)-\widehat e(z)=0.
 \tag{SE5}
\]
Separation of the completed classes implies $z=H$
locally outside an outer-null set. Thus $f\preceq H$.
This argument fixes one $f$ at a time, as required by (SE2).

Now let $g$ be any Borel upper bound in the sense of the
second line of (SE2). It bounds each of the countably many
selected family members outside its corresponding null set.
Their countable union is locally outer-null. Off that union,
$g$ bounds every $v_j$ and hence their actual pointwise
limit $H$. This proves leastness among Borel upper bounds,
without requiring $g$ to be psh or completed.
Two such least upper bounds bound each other, so their
relative classes agree.

Finally, every $w\in\mathscr F^\vee$ satisfies $w\preceq H$,
by the first line of (SE2) and finite subadditivity.
For any sequence whose energies tend to $L$, (CI1) gives
\[
             0\leq\widehat I(w_j,H)
                    \leq\widehat e(w_j)-\widehat e(H)
                                      \longrightarrow0.
 \tag{SE6}
\]
The strong continuity theorem supplies all the current and
critical weighted-measure limits claimed in the statement.
\end{proof}

\begin{corollary}[Subsolution envelopes and Borel upper constraints]
\label{cor:r11-strong-Perron-subsolution}
Let $\nu$ be a fixed positive Radon measure on $V$ and
$\psi$ a Borel upper constraint. If the preceding nonempty
family satisfies
\[
                    \widehat\mu_f\geq\nu,
                    \qquad f\preceq\psi
                                  \quad(f\in\mathscr F),
 \tag{SE7}
\]
then its canonical relative envelope satisfies
\[
                             \widehat\mu_H\geq\nu,
                             \qquad H\preceq\psi.
 \tag{SE8}
\]
The current inequality holds on every original compact smooth
scalar test and every actual moment refinement. A finite total
mass for $\nu$ is not needed; the printed Radon hypothesis is
local. No equality $\widehat\mu_H=\nu$ or contact support
assertion is inferred.
\end{corollary}

\noindent\textbf{Status: \texttt{NEW\_PROVED}} for full strong
completed subsolution families; it is a consequence of the
new finite-maximum and envelope theorems.

\begin{proof}
The completed weak maximum inequality, proved with its full
middle-positive cone and equality-side mass, shows that every
finite maximum of members of $\mathscr F$ has top measure at
least $\nu$. In particular $\widehat\mu_{v_j}\geq\nu$ for
the selected increasing sequence (SE4).
Its strong convergence passes every original compact scalar
test, so positivity gives $\widehat\mu_H\geq\nu$.
All measures are locally Radon; no global mass subtraction
or unproved boundary tightness is needed for this inequality.
The least-upper-bound property (SE2), or the countable selected
members directly, gives $H\preceq\psi$.
\end{proof}

\paragraph{The distinction from an ambient Perron construction.}
This theorem supplies a canonical relative least upper bound
and a retained source lower bound inside the complete strong
class. The ordinary upper semicontinuous envelope can still
lose the relative mass, as in the all-power counterexample
(PO1)--(PO7). Moreover the exact strong completed compact-source
obstruction (SO1)--(SO6) prevents promoting (SE8) to a universal
source equation. These are separate issues from the proved
existence of this Borel relative envelope. General raw-capacity
outer regularity, unrestricted ordinary-boundary Dirichlet
existence, GQ and low-real-degree CTO remain outside the claim.

% END REVIEWED MODULE research_r11/capacity/strong_completed_perron.tex


% BEGIN REVIEWED MODULE research_r11/energy/completed_barrier_variational_problem.tex
\section{A complete barrier-constrained nonlinear variational problem}
\label{sec:r11-completed-barrier-variational}

Retain the strict analytic box and background extension
(TC1)--(TC2), with the original arbitrary symmetric
middle-positive separately closed $(p,p)$-current $S$.
Set $q=n-p\geq1$, $d=q+1$, and fix an actual
$a\in\mathcal E^1_{S,\partial}(V)$.
Write $\widehat{\mathcal E}^1_a$ for the Borel relative
completion of Definition~\ref{def:r11-barrier-completion}.
The background remains $S$; only the potential class
is enlarged by the stated capacity completion.

The frozen predecessor is R4
\path{completed_variational_retraction.tex}.
It treats a bounded completed problem with a fixed
collar, and already distinguishes its constrained
variational condition from an arbitrary source equation.
The result here is an unbounded global first-energy
problem with one actual power-energy lower barrier,
using the new critical continuity and sharp midpoint
modulus.  It does not claim a new bounded Hilbert
projection theorem.

\subsection{Global geometry in the completed order interval}

For completed potentials with actual first-energy
lower barriers, define
\[
\begin{split}
 \widehat B(u_0,\ldots,u_q)
    &=\int_V(-u_0)\widehat N(u_1,\ldots,u_q),\\
 \widehat e(u)&=\widehat B(u^{[q+1]}),\qquad
 \widehat A(u)=-\widehat e(u)/d,\\
 \widehat I(u,v)
    &=\widehat e(u)+\widehat e(v)
       -\widehat B(u,v^{[q]})-\widehat B(v,u^{[q]}).
\end{split}
 \tag{BV1}
\]
The integrals are finite by (RC14).  They use the
actual Borel multipliers, with the joint pole set
annihilated by every coefficient.

\begin{lemma}[Completed energy geometry and separation]
\label{lem:r11-completed-first-energy-geometry}
The form $\widehat B$ is symmetric, nonnegative and
multilinear on the completed positive cone.
The function $\widehat A$ is concave along every
segment, and
\[
 \widehat A\left(\frac{u+v}{2}\right)
      -\frac{\widehat A(u)+\widehat A(v)}2
                    \geq c_q\widehat I(u,v)\geq0,
 \qquad c_q=\frac{1-2^{-q}}{2(q+1)}.
 \tag{BV2}
\]
Moreover
\[
 \widehat I(u,v)=0
       \quad\Longleftrightarrow\quad
 u=v\text{ locally outside a set of outer }S
                         \text{-capacity zero}.
 \tag{BV3}
\]
On $\widehat{\mathcal E}^1_a$ one has
$0\leq\widehat e(u)\leq e_1(a)$.
For a segment $u_t=(1-t)u+tv$,
\[
             \frac{d}{dt}\widehat A(u_t)
                    =\int_V(v-u)\,d\widehat N_q(u_t).
 \tag{BV4}
\]
All signed integrals in this formula are absolutely
convergent.
\end{lemma}

\noindent\textbf{Status: \texttt{DIRECT\_COROLLARY}}
of the new global critical completion theorem,
the actual first-energy geometry and its sharp
midpoint estimate.  The separation assertion is
proved explicitly to avoid assuming an ambient
classical representative.

\begin{proof}
Choose classical defining rows for all the
finitely many completed slots, retaining their
fixed actual barriers.  Formula (RC14) passes
each global mixed energy and every diagonal
energy to its completed value.  It therefore
passes symmetry, positive multilinearity,
the constant-one Holder inequality, concavity,
and (SM8).  For $u\geq a$, choose its defining
rows above $a$.  Actual energy monotonicity
gives $e_1(u_l)\leq e_1(a)$, and (RC14)
gives the stated bound.

If $\widehat I(u,v)=0$, the same convergence
gives $I(u_l,v_l)\to0$ for their classical
defining rows.  These rows have uniformly
bounded energy.  The all-local outer-capacity
estimate (EV8) makes every actual gap between
$u_l$ and $v_l$ tend to zero in capacity.
The defining convergence of each clipped row,
and the capacity-small deep sublevels of the
fixed barriers, give the same conclusion for
the completed gap.  The level-set triangle
inequality then makes each positive gap of
$u-v$ outer-null.  A countable union proves
the forward implication of (BV3).
Conversely, outer-null changes preserve every
completed current and every actual integrable
multiplier by (RC4)--(RC8), so the expression
in (BV1) is zero.

Finally positive multilinearity makes
$\widehat A(u_t)$ a finite polynomial in $t$.
Differentiate its $d$ symmetric slots.
This is (BV4), with integrability from the
completed mixed-energy bound.  No distributional
product of two unbounded differences is being
introduced by this scalar calculation.
\end{proof}

\begin{lemma}[Finite diffuse data integrate barrier-bounded capacity limits]
\label{lem:r11-barrier-diffuse-integration}
Let $\mu$ be a finite positive Radon measure on $V$
which vanishes on every Borel set of zero ordinary-open
outer $S$-capacity.  Suppose
\[
                              \int_V(-a)\,d\mu<\infty.
 \tag{BV5}
\]
If $u_j,u\in\widehat{\mathcal E}^1_a$ converge in the
defining local truncated outer capacities, then
\[
                           \int_V|u_j-u|\,d\mu\longrightarrow0.
 \tag{BV6}
\]
In particular this integral is independent of the
chosen relative representatives.
\end{lemma}

\noindent\textbf{Status: \texttt{NEW\_PROVED}}
in the unbounded barrier completion; the underlying
finite-measure capacity argument is detailed below.

\begin{proof}
A finite measure that vanishes on outer-null Borel
sets is sequentially absolutely continuous for this
outer capacity.  Otherwise choose Borel $E_j$ with
capacities tending to zero but $\mu(E_j)\geq\epsilon$,
and pass to a subsequence with summable capacities.
The limsup has outer capacity zero by countable
subadditivity.  On the other hand, continuity from
above of the finite measure applied to the decreasing
tail unions gives its mass at least $\epsilon$.
This is a contradiction.

The actual pole set of $a$ is outer-null and hence
$\mu$-null.  Off it the inequalities
$a\leq u_j,u\leq0$ imply
$|u_j-u|\leq-a$.
On a fixed compact set remove a deep sublevel of $a$;
its capacity and hence its $\mu$ mass can be made
small.  On the complement full values agree with
one fixed truncation, whose convergence in capacity
and the preceding sequential absolute continuity
give convergence in $\mu$ measure.
Exhaust $V$ by compact sets and use finiteness of
$\mu$ to obtain convergence in measure on all of $V$.
The integrable bound $-a$ in (BV5) then gives
$L^1(\mu)$ convergence: every subsequence has an
almost-everywhere convergent further subsequence,
and dominated convergence applies to it.
This proves (BV6) for the whole sequence.
The same null-set argument proves representative
independence.
\end{proof}

\begin{theorem}[A unique completed barrier-constrained maximizer]
\label{thm:r11-completed-barrier-maximizer}
For every datum satisfying Lemma
\ref{lem:r11-barrier-diffuse-integration}, the functional
\[
      \widehat{\mathcal F}_\mu(u)
            =\widehat A(u)-\int_Vu\,d\mu,
              \qquad u\in\widehat{\mathcal E}^1_a,
 \tag{BV7}
\]
has a maximizer, unique modulo local outer
$S$-capacity zero.  Its value lies between zero
and $\int(-a)\,d\mu$.
Denote its relative class by $\Pi_a(\mu)$.
Every classical maximizing sequence in the
actual interval $\{a\leq u\leq0\}$ is Cauchy
in all local outer capacities and converges to
this completed class.
The maximizer is characterized by the precise
variational inequality
\[
 \int_V(v-u)\,d\bigl(\widehat N_q(u)-\mu\bigr)
               \leq0
       \quad\bigl(v\in\widehat{\mathcal E}^1_a,\ u=\Pi_a(\mu)\bigr).
 \tag{BV8}
\]
The inequality is equivalent to maximality, and
all its integrals are absolutely convergent.
It is not asserted to be the unrestricted
current equation $\widehat N_q(u)=\mu$.
\end{theorem}

\noindent\textbf{Status: \texttt{NEW\_PROVED}}
for the unbounded completed order interval and
arbitrary original $S$.

\begin{proof}
The zero potential is admissible and has functional
value zero.  Since $\widehat e(u)\geq0$ and
$-u\leq-a$, every value in (BV7) is at most
$\int(-a)\,d\mu<\infty$.
For any completed $u$ its classical defining
rows converge in energy by (RC14) and in
$L^1(\mu)$ by (BV6).
Thus the supremum over the completed interval
equals the supremum over its actual classical
members.

Choose a classical sequence approaching this
finite supremum $L$.
Every term has first energy at most $e_1(a)$.
The midpoint remains in the same actual interval.
Its linear datum term cancels from the midpoint
gap, so (SM8) gives
\[
 c_q I(u_j,u_l)
       \leq L-\frac{\mathcal F_\mu(u_j)+
                              \mathcal F_\mu(u_l)}2\longrightarrow0.
 \tag{BV9}
\]
The local capacity estimate (EV8) makes the
sequence Cauchy in the defining topology.
The fixed-barrier completeness theorem gives
a limit $u\in\widehat{\mathcal E}^1_a$.
Global critical energy continuity (RC14) and
(BV6) give
$\widehat{\mathcal F}_\mu(u)=L$.
This proves existence without an ambient upper
regularization or a classical compactness claim.

If $u,v$ are both maximizers, (BV2) and the
affinity of the datum integral give
$c_q\widehat I(u,v)\leq0$.
By (BV3) they are the same relative class.
The same Cauchy argument and uniqueness apply
to every classical maximizing sequence.

For any completed competitor $v$, the segment
$u_t=(1-t)u+tv$ stays in the interval.
The derivative at $t=0$ is the left side of
(BV8), by (BV4) and (BV6).
Maximality gives its nonpositive sign.
Conversely concavity gives the upper tangent
inequality for (BV7); (BV8) therefore makes
every competitor's value at most the value
at $u$.  The data integrals are controlled
by (BV5), and all curvature integrals by
the completed mixed-energy bound.
This proves both the characterization and
the claimed integrability.
\end{proof}

\begin{theorem}[Data stability and retraction on the finite-mass image]
\label{thm:r11-completed-barrier-data-stability}
Let $\mu,\nu$ satisfy (BV5) and the same
finite BT-diffuse hypotheses, and define
\[
                    \delta_a(\mu,\nu)
                         =\int_V(-a)\,d|\mu-\nu|.
 \tag{BV10}
\]
Then
\[
 \widehat I(\Pi_a(\mu),\Pi_a(\nu))
       \leq\delta_a(\mu,\nu).
 \tag{BV11}
\]
In particular, if $\delta_a(\mu_j,\mu)\to0$,
the completed maximizers converge in every
local outer capacity, every original mixed
current, and all the fixed-barrier weighted
limits of (RC8) and (RC14).

If $b\in\widehat{\mathcal E}^1_a$ and its top
measure $\mu_b=\widehat N_q(b)$ has finite
total mass, then it satisfies the stated
datum hypotheses and
\[
                                \Pi_a(\mu_b)=b
 \tag{BV12}
\]
in the relative capacity quotient.
If the actual barrier has finite total top mass
and positive first energy, the constant one
in (BV11) is optimal for that same barrier,
domain and background.
\end{theorem}

\noindent\textbf{Status: \texttt{NEW\_PROVED}}
for the quantitative completed variational map.
The term retraction means exactly (BV12).

\begin{proof}
Write $u=\Pi_a(\mu)$ and $w=\Pi_a(\nu)$.
Use $w$ as a competitor in (BV8) for $u$,
and $u$ as a competitor for $w$.
Adding the two resulting inequalities gives
\[
 0\leq\widehat I(u,w)
       \leq\int_V(u-w)\,d(\nu-\mu)
       \leq\int_V(-a)\,d|\nu-\mu|.
\]
The last inequality uses $|u-w|\leq-a$.
This proves (BV11).  The two-competitor argument
is the frozen bounded retraction mechanism;
the newly justified unbounded integrals and
completed interval make it applicable here.

For clarity, the passage from $\widehat I\to0$
to local capacity convergence is also valid
for completed rows.  Approximate each pair
by classical rows with their fixed barrier.
Their energies and discrepancies converge
by (RC14).  Apply (EV8) to the classical
rows, use the level-set triangle inequality
and the defining capacity convergence, and
then decrease its intermediate gap threshold
to the requested one.  The resulting bound
is the same fixed-interval capacity modulus.
Thus (BV11) gives the stated convergence,
and (RC8), (RC14) give all the current
and weighted conclusions.

For (BV12), outer-capacity control of the
completed current on each compact stage
shows that $\mu_b$ is BT diffuse.
Also
$\int(-a)\,d\mu_b=\widehat B(a,b^{[q]})<\infty$.
Its finite total mass is the printed
additional assumption.
The completed tangent inequality for
$\widehat A$ gives
\[
 \widehat A(v)-\widehat A(b)
                    \leq\int(v-b)\,d\mu_b .
\]
Thus $b$ maximizes (BV7) for datum $\mu_b$.
Uniqueness proves (BV12).
Finally, if $\mu_a$ has finite total mass and
$e_1(a)>0$, take $\mu=\mu_a$ and $\nu=0$.
The retraction identity gives $\Pi_a(\mu_a)=a$,
and $\Pi_a(0)=0$ because the objective is
$-\widehat e/d\leq0$.
Then
$\widehat I(a,0)=e_1(a)=\delta_a(\mu_a,0)$,
proving the stated optimality.
\end{proof}

\begin{proposition}[A completed variational maximizer need not solve the source equation]
\label{prop:r11-completed-maximizer-source-obstruction}
For the precise smooth moment-Laplacian background
(DO6), with $n\geq2$ and $q=1$, take the
nonzero compact smooth source (DO9).
For every actual
$a\in\mathcal E^1_{S,\partial}(V)$, the
maximizer $\Pi_a(\mu)$ exists, but
\[
                            D(\Pi_a(\mu)\,S)\ne\mu .
 \tag{BV13}
\]
More generally no member of
$\widehat{\mathcal E}^1_a$ has that source.
Thus the implication from completed variational
maximality to an unrestricted original source
equation is \texttt{FALSE}, even after the
classical representative requirement is relaxed
to this relative completion.
\end{proposition}

\noindent\textbf{Status: \texttt{NEW\_PROVED}}
as the extension of the frozen compact-source
obstruction to the new barrier completion.
It refines the earlier actual-class obstruction
without replacing its proof.

\begin{proof}
The datum (DO9) is finite and dominated by a
fixed multiple of the full normalized BT
capacity, so it is BT diffuse.
Equation (DO4), with $s=1$ and $v=a$,
gives $\int(-a)\,d\mu<\infty$.
The existence theorem therefore applies.

Suppose a completed member $u$ had
$D(uS)=\mu$, and choose classical defining
rows $a\leq u_j\leq0$.
Their Gauss traces $g_j=u_j\circ\sigma$
are finite convex functions on the real
box $Q$, by the frozen trace theorem.
The actual barrier has a finite convex
Gauss trace on all of $Q$, by the
finite-trace argument in the proof of
Theorem~\ref{thm:r11-power-compact-source-obstruction}.
Hence on every compact subbox the $g_j$
have uniform upper and lower bounds.
Elementary convex slope bounds give
equicontinuity on each smaller compact
subbox.  A diagonal finite-net argument
therefore supplies a subsequence converging
locally uniformly to a finite convex $g$.

The completed weighted current convergence
in degree zero gives $u_jS\to uS$ on
every original compact test.
For each real $\phi\in C_c^\infty(Q)$,
the exact original calibration now yields
\[
 \langle t_*D(uS),\phi\rangle
     =\lim_j\int_Q g_j\Delta\phi\,dy
     =\int_Q g\Delta\phi\,dy
     =\langle\Delta g,\phi\rangle.
 \tag{BV14}
\]
The purported source has compact tropical
support.  The frozen R10 convex compact-source
rigidity makes $g$ affine and $\Delta g=0$.
Equation (BV14) would therefore give
$t_*\mu=0$, contradicting the nonzero
positive source (DO9).
This proves (BV13) throughout the completed
interval.
It uses the full-original degree-zero
convergence before testing the pushforward;
no toric-only definition of the currents
or auxiliary-field quotient is substituted.
\end{proof}

\paragraph{The remaining equation gate is concrete.}
The theorem proves existence, uniqueness, source
stability and a retraction for a nonlinear
constrained relative problem.
Its Euler condition is exactly (BV8).
Promoting that inequality to equality against
every original test requires additional
admissible variations or a different potential
class.  Proposition
\ref{prop:r11-completed-maximizer-source-obstruction}
shows that this promotion is not a technical
consequence of the completion already built.
The finite signed diffuse relative linear
Dirichlet problem uses the larger specified
energy class and remains compatible with
this obstruction.

% END REVIEWED MODULE research_r11/energy/completed_barrier_variational_problem.tex


% BEGIN REVIEWED MODULE research_r11/dirichlet/diffuse_measure_transposition_dirichlet.tex
\section{Finite diffuse data and an intrinsic transposition Dirichlet problem}
\label{sec:r11-diffuse-dirichlet}

Retain the arbitrary original background and the analytic domains
$(B,\Omega,V)$ of Section~\ref{sec:r11-affinoid-energy}.  Thus
$S$ has degree $(n-1,n-1)$, $F=V\cap\operatorname{supp}S$,
$\lambda=D(hS)$ is faithful, and
$H=\mathscr D_0(V)$ is a regular strongly local coercive part form.
The reference measure $\lambda|_F$ is finite because $V\Subset\Omega$.
Write $c_V$ for the transient $0$-capacity of this part form.
Its usual outer definition agrees, on quasi-open sets, with
\[
 c_V(A)=\inf\{\mathcal E(v,v):v\in H,\ 
                              \widetilde v\geq1\text{ quasi-everywhere on }A\}.
 \tag{MD1}
\]
The tilde denotes the energy-quasicontinuous representative.
Since $\mathcal E\leq\mathcal E_1\leq3\mathcal E$,
\[
                    c_V\leq\operatorname{cap}_{H,1}\leq3c_V.
 \tag{MD2}
\]
The two capacities therefore have identical polar sets.
They are capacities on the original support space $F$, with
its ordinary subspace topology; they are not identified with
relative BT capacity.

\begin{definition}[Finite energy-diffuse data]
\label{def:r11-energy-diffuse-data}
Let $\mathcal M_{\mathrm{diff}}(F)$ be the finite signed Radon
measures $\mu$ on $F$ for which $|\mu|$ vanishes on every
$c_V$-polar Borel set.  This is a condition on the specified
background and domain, not on an unknown solution.
\end{definition}

\noindent\textbf{Status: definition.}

\begin{lemma}[The finite-energy nest input]
\label{lem:r11-finite-energy-nest}
For every $\mu\in\mathcal M_{\mathrm{diff}}(F)$ there are increasing
Borel sets $F_j$ such that
$|\mu|(F\setminus\bigcup_jF_j)=0$ and
$\tau_j={\bf1}_{F_j}|\mu|$ has finite energy integral:
\[
             \int_F|\widetilde v|\,d\tau_j
                         \leq C_j\,\mathcal E(v,v)^{1/2}
                                  \qquad(v\in H).
 \tag{MD3}
\]
In particular $\mu_j={\bf1}_{F_j}\mu$ converges to $\mu$ in
total variation and has a finite-energy variation.
\end{lemma}

\noindent\textbf{Status: \texttt{EXTERNAL\_KNOWN}}.
The precise input is Klimsiak--Rozkosz,
\emph{On the structure of bounded smooth measures associated with a
quasi-regular Dirichlet form}, arXiv:1410.4927v2,
Lemma~2.1 (pp.~2--3).  It asserts a finite $0$-energy nest
for every smooth measure of a transient quasi-regular form.
The original form here is regular, hence quasi-regular, and
coercivity proves transience.  A finite Radon measure charging
no polar sets is smooth: any compact form nest has finite
measure on its members and has polar complement.  Apply the
lemma to $|\mu|$.  The claimed total-variation convergence
then follows by monotone convergence.
No existence theorem for the original current equation is
included in this external input.

\subsection{Quantitative estimates before passing to measure data}

For a finite signed Radon measure $\sigma$ whose variation has
finite energy integral, let $G\sigma\in H$ be its Riesz potential:
\[
               \mathcal E(G\sigma,v)=\int_F\widetilde v\,d\sigma
                                 \qquad(v\in H).
 \tag{MD4}
\]
This is the R6 finite-energy result applied on the analytic
domains now covered by Section~\ref{sec:r11-affinoid-energy}.
For $k>0$ set $T_k(s)=\max(-k,\min(k,s))$.

\begin{lemma}[Mass controls integrability and all truncations]
\label{lem:r11-finite-energy-mass-estimates}
Under the preceding assumptions on $\sigma$, put $w=G\sigma$.
Then
\[
 \|w\|_{L^1(\lambda)}\leq2\|\sigma\|_{\mathrm{TV}},\qquad
 \mathcal E(T_k w,T_k w)\leq k\|\sigma\|_{\mathrm{TV}},\qquad
 c_V(\{|\widetilde w|>k\})\leq\frac{\|\sigma\|_{\mathrm{TV}}}{k}.
 \tag{MD5}
\]
The constants are independent of the finite-energy norm of $\sigma$.
\end{lemma}

\noindent\textbf{Status: \texttt{NEW\_PROVED}} as the uniform
original-background mass estimate needed for the larger solution class.
Only the frozen positivity and coercive barrier are used.

\begin{proof}
Positive finite-energy measures integrate canonical energy
representatives.  To recall the justification, their core norm
bound and the compact-core capacity formula show that they
vanish on polar compact sets, hence on polar Borel sets by Radon
regularity.  For a core approximation $v_j\to v$ in $H$,
the same norm bound applied to smooth approximations of
$|v_j-v_l|$ makes it Cauchy in $L^1$ of the measure.
A subsequence converges quasi-everywhere, so its $L^1$ limit is
the canonical representative $\widetilde v$.
This proves the integration interpretation in (MD4).

Order preservation of $G$ gives
$|G\sigma|\leq G|\sigma|$ quasi-everywhere.  With the torsion
vector $z=G1$, whose representative satisfies $0\leq z\leq2$,
symmetry yields
\[
 \int_FG|\sigma|\,d\lambda
       =\mathcal E(G|\sigma|,z)
       =\int_F\widetilde z\,d|\sigma|
       \leq2\|\sigma\|_{\mathrm{TV}}.
 \tag{MD6}
\]
This proves the first estimate.

The normal contraction $P(s)=s-2T_k(s)$ has Lipschitz constant
one and $P(0)=0$.  The Markov contraction inequality therefore
gives
$\mathcal E(w-2T_kw,w-2T_kw)\leq\mathcal E(w,w)$.
Expanding proves
$\mathcal E(T_kw,T_kw)\leq\mathcal E(w,T_kw)$.
Equation (MD4), tested with $T_kw$, bounds its right side by
$k\|\sigma\|_{\mathrm{TV}}$.  The vector $|T_kw|/k$ is
admissible for the outer-capacity estimate of the quasi-open
set $\{|\widetilde w|>k\}$, giving the last inequality.
\end{proof}

\begin{definition}[The transposition boundary condition]
\label{def:r11-transposition-boundary}
A function $u\in L^1(F,\lambda)$ satisfies the zero transposition
Dirichlet problem for $\mu\in\mathcal M_{\mathrm{diff}}(F)$ if
\[
             \int_Fu f\,d\lambda
                   =-\int_F\widetilde{Gf}\,d\mu
                        \qquad(f\in L^\infty(F,\lambda)).
 \tag{MD7}
\]
Here $Gf=G(f\lambda)\in H$, and its canonical representative is
bounded by $2\|f\|_\infty$.  The right side is therefore defined
for every datum under consideration.
\end{definition}

\noindent\textbf{Status: definition.}
This is an intrinsic energy boundary condition through the
part resolvent.  It does not assert an ordinary pointwise
boundary limit.

\begin{theorem}[All-original-current Dirichlet solutions for finite diffuse data]
\label{thm:r11-diffuse-transposition-dirichlet}
Every $\mu\in\mathcal M_{\mathrm{diff}}(F)$ has a unique
$L^1(F,\lambda)$ solution $u_\mu$ of (MD7).  The solution has
a canonical finite energy-quasicontinuous representative
outside a polar set, and
\[
       T_k(u_\mu)\in H,\qquad
       \mathcal E(T_ku_\mu,T_ku_\mu)\leq k\|\mu\|_{\mathrm{TV}}.
 \tag{MD8}
\]
Its coefficient-defined original current satisfies
\[
                         D(u_\mu S)=\mu\quad\hbox{on }V
 \tag{MD9}
\]
against every original compact smooth scalar test, including
all actual moment refinements.
The solution map is linear, is order reversing on positive data,
and for $\mu,\nu\in\mathcal M_{\mathrm{diff}}(F)$ satisfies
\[
 \begin{split}
 \|u_\mu-u_\nu\|_{L^1(\lambda)}
                   &\leq2\|\mu-\nu\|_{\mathrm{TV}},\\
 \mathcal E\bigl(T_k(u_\mu-u_\nu),T_k(u_\mu-u_\nu)\bigr)
                   &\leq k\|\mu-\nu\|_{\mathrm{TV}},\\
 c_V(\{|\widetilde u_\mu-\widetilde u_\nu|>k\})
                   &\leq\|\mu-\nu\|_{\mathrm{TV}}/k.
 \end{split}
 \tag{MD10}
\]
For positive $\mu$, one has $u_\mu\leq0$.
When $\mu$ already has finite energy, this is exactly the
frozen R6 energy solution.
\end{theorem}

\noindent\textbf{Status: \texttt{NEW\_PROVED}} as the
all-original-current extension to finite diffuse data without a
finite-energy norm assumption.  The abstract exhaustion lemma
above is the stated \texttt{EXTERNAL\_KNOWN} input.
This theorem does not assert ambient classical psh membership
or a nonlinear diagonal Monge--Ampere equation.

\begin{proof}
Choose $\mu_j$ from Lemma~\ref{lem:r11-finite-energy-nest} and
put $u_j=-G\mu_j$.  Differences have finite-energy variation,
so (MD5) gives
\[
 \|u_j-u_l\|_1\leq2\|\mu_j-\mu_l\|_{\mathrm{TV}},\qquad
 c_V(\{|\widetilde u_j-\widetilde u_l|>k\})
       \leq\|\mu_j-\mu_l\|_{\mathrm{TV}}/k.
 \tag{MD11}
\]
Hence $u_j$ converges in $L^1(\lambda)$ to some $u$ and is
Cauchy in energy capacity.  A subsequence with summable
successive capacity errors has a finite limit outside a polar
set: choose summable level thresholds and cover the corresponding
successive gap sets by open sets with summable capacities.
Away from their tails the sequence is uniformly Cauchy.
The quasicontinuity of each $\widetilde u_j$, with another
summable family of exceptional open sets, makes its limit
energy-quasicontinuous.  This limit agrees almost everywhere
with the $L^1$ limit.

For fixed $k$, the vectors $T_ku_j$ are bounded in $H$ by
(MD5).  Since
\[
 \|T_ku_j-T_ku\|_{L^2(\lambda)}^2
                         \leq2k\|u_j-u\|_{L^1(\lambda)},
\]
their $L^2$ limit is $T_ku$.  Weak compactness in $H$ and
coercivity identify every weak cluster point with this vector.
Lower semicontinuity proves (MD8).
The truncation of the capacity-limit representative is itself
quasicontinuous and agrees $\lambda$-almost everywhere with
this $H$ vector.  Uniqueness of quasicontinuous representatives
therefore identifies it with the canonical representative of
$T_ku$; the capacity estimate consequently concerns the same
representative throughout.
The bound on the capacity of its level sets follows by applying
(MD1) to $|T_ku|/k$ and its canonical representative.
It also shows directly that the possible infinite-value set
has capacity zero.

For $f\in L^\infty(F,\lambda)$, finite-energy symmetry gives
\[
                  \int_Fu_j f\,d\lambda
                      =-\int_F\widetilde{Gf}\,d\mu_j .
\]
The left side converges by $L^1$ convergence and the right
side by total-variation convergence and the bound
$|\widetilde{Gf}|\leq2\|f\|_\infty$.
Thus $u$ satisfies (MD7).  Since bounded measurable functions
separate $L^1$ classes, (MD7) proves uniqueness.
It also proves linearity.  Positivity and order follow by
applying (MD11) to common finite-energy exhaustions of the
positive data and passing to the capacity limit.

We verify the current statement without substituting any
restricted test family.  For original
$\phi\in A_c^{0,0}(V;\mathbf R)$, (AE10) gives
$D(\phi S)=b_\phi\lambda$ with $b_\phi$ bounded and compactly
supported.  Coefficient domination makes $uS$ well defined
on every original complementary form, because $u\in L^1(\lambda)$.
Since $\phi\in H$, the Green identity implies
$G(b_\phi)=-\phi$.  Applying (MD7) with $f=b_\phi$ gives
\[
 \langle D(uS),\phi\rangle
        =\int_Fu\,b_\phi\,d\lambda
        =\int_F\phi\,d\mu .
 \tag{MD12}
\]
This proves (MD9) for all original tests.  Alternatively it
follows by $L^1(\lambda)$ convergence in the same coefficient
pairing from the finite-energy equations.

For two data, choose a common finite-energy exhaustion of
$|\mu|+|\nu|$.  Apply (MD5) to the differences of the
approximations, then pass as above.  This proves all of
(MD10) and independence of the initial exhaustion in both
$L^1$ and the canonical capacity class.
If $\mu$ has finite energy, its Riesz solution already
satisfies (MD7), so uniqueness identifies the two solutions.
\end{proof}

\subsection{The relative BT-capacity consequence}

\begin{proposition}[Independence of the faithful reference]
\label{prop:r11-transposition-reference-independence}
Fix $V$ and the original restriction of $S$ to an ordinary
neighborhood of $\overline V$.  The energy capacity $c_V$,
the class $\mathcal M_{\mathrm{diff}}(F)$, and the canonical
solution and original current in
Theorem~\ref{thm:r11-diffuse-transposition-dirichlet} do not depend
on the choices of compensators, faithful potential $h$, or ambient
affinoid neighborhoods for which the construction is available.
\end{proposition}

\noindent\textbf{Status: \texttt{NEW\_PROVED}} for the larger
measure-data solution class.

\begin{proof}
The original core $A_c^{0,0}(V;\mathbf R)$ and its energy
$\mathcal E(f,g)=\int d'f\wedge d''g\wedge S$ are independent
of all these choices.  Coercivity identifies each part space
with the Hilbert completion of this same core modulo zero
energy, by the same energy norm.  The compact-core formula
for transient capacity therefore gives the same $c_V$ on
compact sets and hence, by outer and Choquet regularity, on
all Borel sets.  In particular all canonical core-limit
quasicontinuous representatives and all polar sets coincide.
The two faithful measures have the same null ideal by (AE7).

Finite-energy Riesz potentials are the same abstract Hilbert
vectors and therefore have the same canonical representatives.
Choose one finite-energy nest for $|\mu|$ relative to this
common energy.  The canonical solution is the limit in the
common capacity of the same finite-energy potential sequence.
Equation (MD11) makes this limit independent of the reference;
its two $L^1$ realizations are representatives of that one
class.  Finally the coefficient multiplication $uS$ and its
distributional derivative use only the original $S$ and
the common almost-everywhere class, proving the current assertion.
\end{proof}

\begin{proposition}[Local BT-capacity control and diffuse BT data]
\label{prop:r11-bt-capacity-diffuse-dirichlet}
For each compact $K_0\Subset V$ there is a finite constant
$C_{K_0}$ depending only on the fixed background, domains and
coefficient controls such that
\[
 \operatorname{Cap}_S^*
   \bigl(K_0\cap\{|\widetilde u_\mu-\widetilde u_\nu|>k\};\Omega\bigr)
       \leq C_{K_0}
           \left(\frac{3\|\mu-\nu\|_{\mathrm{TV}}}{k}\right)^{1/2}.
 \tag{MD13}
\]
The representatives may be set to zero off $\operatorname{supp}S$.
They are locally quasicontinuous for relative BT outer capacity.

In particular every finite signed Radon measure on $V$ whose
variation vanishes on all locally relative BT outer-null sets
is an allowed datum for Theorem~\ref{thm:r11-diffuse-transposition-dirichlet}.
This includes any finite datum with a capacity domination
condition whose right side vanishes at capacity zero; no Dini
growth condition is required for this $L^1$ transposition class.
\end{proposition}

\noindent\textbf{Status: \texttt{DIRECT\_COROLLARY}} of the
new theorem and the frozen one-way energy-to-BT-capacity bridge,
now valid on the analytic affinoid domains of
Theorem~\ref{thm:r11-general-affinoid-faithful-form}.

\begin{proof}
The R4 bridge uses only the faithful coefficient bounds,
local smooth regularization, fixed-cutoff energy estimates,
and the regular compact-core capacity formula.  All these
inputs hold on the present analytic domains by (AE5)--(AE13)
and the frozen local bounded BT calculus; its proof is
unchanged.  Hence
\[
 \operatorname{Cap}_S^*(A;\Omega)
       \leq C_{K_0}\operatorname{cap}_{\mathcal E,1}
                  (A\cap\operatorname{supp}S)^{1/2}
                         \quad(A\subset K_0).
 \tag{MD14}
\]
The full-space energy capacity is no larger than the part
capacity, and (MD2) bounds the latter by $3c_V$.
Equation (MD10) gives (MD13).

Energy quasicontinuity passes through (MD14) on every compact
stage.  The open complement of $\operatorname{supp}S$ has
zero relative capacity, since every relative BT measure is
supported on $\operatorname{supp}S$; it may be included in
each exceptional open set.  This justifies the zero extension
without claiming continuity across its boundary.

An energy-polar set is locally BT outer-null by (MD14).
Thus a datum satisfying the displayed null-set hypothesis
vanishes on it.  The same hypothesis makes it vanish off
$\operatorname{supp}S$.  It belongs to
$\mathcal M_{\mathrm{diff}}(F)$, proving the final assertion.
This argument uses only the null ideals, not equality of the
two capacities or outer regularity of raw BT capacity.
\end{proof}

\subsection{A strict extension beyond finite-energy data}

\begin{proposition}[An original finite diffuse datum of infinite energy norm]
\label{prop:r11-infinite-energy-diffuse-example}
For every $n\geq3$ there are a smooth original positive closed
background $S$ on an ordinary torus box, a finite positive Radon
measure $\mu$ carried by a compact subset, and a part domain $V$,
such that $\mu$ is energy diffuse but
\[
                              \|\mu\|_{-1,V}=\infty.
 \tag{MD15}
\]
Theorem~\ref{thm:r11-diffuse-transposition-dirichlet} supplies
its canonical $L^1$ relative solution; that solution is not in $H$.
\end{proposition}

\noindent\textbf{Status: \texttt{NEW\_PROVED}} as a full-dimensional
strictness example.  It is not used to replace the arbitrary-$S$
theorem by a toric theorem.

\begin{proof}
Use original torus coordinates $t_i=\log|T_i|$,
$\beta=\sum_i d't_i\wedge d''t_i$, and
$S=\beta^{n-1}[U]/(n-1)!$ on a bounded ordinary torus box.
Take a smaller Euclidean ball centered at the origin in its
tropical image and a nonnegative compact smooth cutoff $\chi$
which equals one on a still smaller ball.  Put
\[
 \alpha=(n+2)/2,\qquad
                 \mu=\sigma_*(\chi(t)|t|^{-\alpha}\,dt),
 \tag{MD16}
\]
where $\sigma$ is the original joint Gauss section.
Since $\alpha<n$, this is finite and compactly supported.
The smooth coefficient measure $\beta\wedge S=n\,\sigma_*dt$
is locally dominated by the faithful reference measure.
Consequently $\mu\ll\lambda$, so it charges no energy-polar set.

Let $a=(n-2)/2$, and choose an original compact smooth scalar
$\phi_\varepsilon(t)$ which equals
$(|t|^2+\varepsilon^2)^{-a/2}$ on the smaller ball and has
support in the fixed cutoff ball.  Direct radial integration
with the original calibration gives, for constants independent
of $\varepsilon$,
\[
 \int\phi_\varepsilon\,d\mu\geq c\log(1/\varepsilon)-C,\qquad
             \mathcal E(\phi_\varepsilon,\phi_\varepsilon)
                      \leq C\bigl(1+\log(1/\varepsilon)\bigr).
 \tag{MD17}
\]
Indeed the first integrand has radial exponent
$n-1-a-\alpha=-1$ on $\varepsilon\leq|t|\leq r$,
and the squared-gradient energy has exponent
$n-1-2a-2=-1$ there.  The inner ball and the fixed cutoff
annulus contribute bounded amounts to the energy.
The ratio in the negative-energy core norm tends to infinity,
proving (MD15).  If the resulting solution belonged to $H$,
(AE15) and Cauchy--Schwarz would make that norm finite,
a contradiction.
\end{proof}

\paragraph{The larger class and the remaining nonlinear problem.}
Positive data give positive relative curvature $D(u_\mu S)$,
but $u_\mu$ is not asserted ambient classical psh.  The R10
compact-source obstruction therefore does not conflict with
this theorem.  For an arbitrary original background $S_0$
of degree $(p,p)$ and fixed bounded classical slots
$v_2,\ldots,v_{n-p}$, the entire construction applies to
$S=T_{S_0,n-p-1}(v_2,\ldots,v_{n-p})$.  It solves one free
relative slot on every original test.  It does not construct
a diagonal product of several moving $L^1$ solutions, prove a
general nonlinear Monge--Ampere existence theorem, or settle
OPEN-GQ/OPEN-CTO.

% END REVIEWED MODULE research_r11/dirichlet/diffuse_measure_transposition_dirichlet.tex


% BEGIN REVIEWED MODULE research_r11/dirichlet/kato_truncation_measure_data.tex
\section{Kato inequalities and curvature convergence for diffuse data}
\label{sec:r11-diffuse-Kato}

Use the ordinary affinoid energy setting (AE1)--(AE11), fix
$V\Subset\Omega$, and write $H=\mathscr D_0(V)$,
$F=V\cap\operatorname{supp}S$.  In particular $\lambda(V)<\infty$.
Here $S$ is still an arbitrary original middle-positive separately
closed current of degree $(n-1,n-1)$, in arbitrary dimension $n$.
Let $\mu$ be a finite signed Radon measure on $F$ charging no
energy-polar set, and let $u=u_\mu$ be the unique transposition
solution from Theorem~\ref{thm:r11-diffuse-transposition-dirichlet}.
The representative $\widetilde u$ is the canonical energy-
quasicontinuous representative.  Products with $S$ and all
current equations below use every original compact smooth test.

\begin{lemma}[Finite-energy chain rule on the original current]
\label{lem:r11-H-original-chain}
For $v\in H$ there is a finite positive energy measure
$\Gamma_S(v)$, with mass $\mathcal E(v,v)$, obtained by closing
original core energy measures in total variation.  It is absolutely
continuous with respect to $\lambda$ and independent of all moment
presentations.  Suppose $\sigma$ and $|\sigma|$ have finite energy
and $v=u_\sigma\in H$.  For a smooth function $\Phi$ with bounded
first and second derivatives,
\[
 D(\Phi(v)S)=\Phi'(\widetilde v)\sigma
                            +\Phi''(\widetilde v)\Gamma_S(v).
 \tag{KT1}
\]
The identity is local if the stated derivative bounds hold only
on the range used in a local calculation.
\end{lemma}

\noindent\textbf{Status: \texttt{DIRECT\_COROLLARY}} of the closed
energy form and the original core chain rule; the extension argument
is recorded rather than importing an identification with a different
notion of current.

\begin{proof}
For original core functions the positive energy matrix and
Cauchy--Schwarz imply
\[
 \|\Gamma_S(v_j)-\Gamma_S(v_l)\|_{\rm TV}
 \le \mathcal E(v_j-v_l)^{1/2}
        \bigl(\mathcal E(v_j)^{1/2}+\mathcal E(v_l)^{1/2}\bigr).
 \tag{KT2}
\]
An $H$-convergent core sequence therefore defines the asserted
measure by total-variation convergence.  Its mass is the energy.
Every core energy measure is absolutely continuous with respect to
$\lambda$ by (AE10), and this property passes to a total-variation
limit.  Polarization defines the mixed measures.

For completeness the scalar chain rule survives this closure.
Form the Hilbert space of finite sums $\sum_i b_i\,dv_i$ with
bounded Borel $b_i$ and core $v_i$, with squared norm
$\sum_{i,l}\int b_i b_l\,\Gamma_S(v_i,v_l)$, quotienting by its
zero seminorm and completing.  Positivity follows from the
pointwise positive energy matrix after taking densities with
respect to $\lambda$.  The symbols $dv_j$ converge for every
$H$-convergent sequence.  Passing to a subsequence gives
$v_j\to\widetilde v$ almost everywhere for $\lambda$; boundedness
and continuity of $\Phi'$ then give
\[
                 \Phi'(v_j)\,dv_j\longrightarrow
                              \Phi'(\widetilde v)\,dv
\]
in this Hilbert space.  Indeed the part containing $dv_j-dv$
is small in norm, and the remaining part tends to zero by
bounded convergence against the fixed measure $\Gamma_S(v)$.
The $L^2$ convergence of $\Phi(v_j)-\Phi(0)$ follows from its
Lipschitz bound.  Closedness identifies the differential of
$\Phi(v)-\Phi(0)$ with this limit.  The same argument and the
core product rule apply to $\phi\Phi'(v)$ for an original
compact smooth $\phi$.

The finite-energy equation extends to these $H$ tests:
$\mathcal E(v,z)=-\int\widetilde z\,d\sigma$.
Consequently
\[
 \begin{split}
 \langle D(\Phi(v)S),\phi\rangle
  &=-\mathcal E(\Phi(v)-\Phi(0),\phi)\\
  &=-\mathcal E(v,\phi\Phi'(v))
                        +\int\phi\Phi''(\widetilde v)\,d\Gamma_S(v)\\
  &=\int\phi\Phi'(\widetilde v)\,d\sigma
                        +\int\phi\Phi''(\widetilde v)\,d\Gamma_S(v).
 \end{split}
\]
The first equality is exactly the extension of (AE11) to $H$,
not a restricted test equation.  The constant term disappears
because $S$ is closed.  This proves (KT1).  The construction uses
only original measures and their norm limits, so every identity
persists after any actual moment refinement.
\end{proof}

\begin{lemma}[Diffuse measures respect capacity convergence]
\label{lem:r11-diffuse-capacity-absolute-continuity}
For every finite energy-diffuse $\nu$ and every sequence of Borel
sets $E_j\subset F$, $c_V(E_j)\to0$ implies $|\nu|(E_j)\to0$.
Consequently convergence in energy capacity implies convergence
in $|\nu|$-measure, and a subsequence converges $|\nu|$-almost
everywhere.  Bounded continuous scalar compositions therefore
converge in $L^1(|\nu|)$.
\end{lemma}

\noindent\textbf{Status: \texttt{DIRECT\_COROLLARY}} of finiteness,
countable subadditivity of the regular form capacity, and diffuseness.

\begin{proof}
If the assertion fails, choose a subsequence with
$c_V(E_j)\le2^{-j}$ and $|\nu|(E_j)\ge\epsilon>0$.
Its limsup has capacity zero, while finiteness and continuity
from above give measure at least $\epsilon$, a contradiction.
Apply this fact to the gap sets for capacity convergence.
The almost-everywhere subsequence follows by choosing summable
measure errors, and bounded convergence gives the last assertion.
\end{proof}

\begin{theorem}[An all-original Kato inequality for finite diffuse data]
\label{thm:r11-diffuse-original-Kato}
For every real $t$, the original current
\[
 L_t(u;S):=D((u-t)^+S)
                         -\mathbf1_{\{\widetilde u>t\}}\mu
 \tag{KT3}
\]
is a positive Radon measure on $V$.  Equivalently,
\[
              D((u-t)^+S)\ge
                           \mathbf1_{\{\widetilde u>t\}}\mu .
 \tag{KT4}
\]
Thus the curvature of each Lipschitz positive-part truncation is
an actual signed Radon measure, with its original distributional
meaning.  This conclusion holds for every finite energy-diffuse
signed $\mu$, without a finite-energy norm bound on $\mu$.
\end{theorem}

\noindent\textbf{Status: \texttt{NEW\_PROVED}} in this
all-original arbitrary-background transposition class.
The theorem makes no assertion that the measure (KT3) is
carried by the exact level set.  Such a carrier identity for
the general signed solution class is a separate locality question;
it is not silently imported from the bounded classical calculus.
The positive-data case is proved independently in
Section~\ref{sec:r11-positive-diffuse-renormalized}; the full
signed case is subsequently proved in
Section~\ref{sec:r11-signed-diffuse-renormalized} by a separate
strong-truncation argument.

\begin{proof}
Let $\mu_j$ be the finite-energy nest restrictions used to construct
$u$, and put $u_j=u_{\mu_j}\in H$.
Choose smooth convex functions $\Phi_\varepsilon$ with
\[
 0\le\Phi_\varepsilon'\le1,\quad
 \Phi_\varepsilon'(s)=0\ (s\le0),\quad
 \Phi_\varepsilon'(s)=1\ (s\ge\varepsilon),\quad
 \sup_s|\Phi_\varepsilon(s)-s^+|\le\varepsilon.
\]
One can normalize $\Phi_\varepsilon(0)=0$.
Lemma~\ref{lem:r11-H-original-chain} gives
\[
 D(\Phi_\varepsilon(u_j-t)S)
             -\Phi_\varepsilon'(\widetilde u_j-t)\mu_j\ge0.
 \tag{KT5}
\]
The translated scalar composition is normalized at zero when
regarded as an $H$ vector; its constant has zero curvature.

The first term in (KT5) converges as an original current because
$u_j\to u$ in $L^1(\lambda)$ and $\Phi_\varepsilon$ is Lipschitz.
The measures in the second term converge in total variation:
$\mu_j\to\mu$ in total variation, while
$u_j\to\widetilde u$ in capacity, so the preceding lemma gives
\[
 \int\bigl|\Phi_\varepsilon'(\widetilde u_j-t)
            -\Phi_\varepsilon'(\widetilde u-t)\bigr|\,d|\mu|
                                                       \longrightarrow0.
\]
The limit of (KT5) is therefore a positive original top-degree
current, hence a positive Radon measure.
Finally let $\varepsilon\downarrow0$.
The potentials converge in $L^1(\lambda)$, since
$\lambda(V)<\infty$, and the bounded multipliers converge
pointwise for $|\mu|$ to $\mathbf1_{\{\widetilde u>t\}}$,
including at $\widetilde u=t$, where they are zero.
The second passage again preserves positivity and proves
(KT3)--(KT4) on every original test.
\end{proof}

\begin{theorem}[Positive curvature truncations converge in local variation]
\label{thm:r11-diffuse-positive-curvature-truncation}
Suppose $\mu\ge0$.  Then $u\le0$ quasi-everywhere, and for every
$k>0$ the function
\[
                       u^{(k)}=\max(u,-k)=T_k u
 \tag{KT6}
\]
lies in $H\cap L^\infty$, has zero energy boundary value, and
has positive original curvature
\[
 \mu^{(k)}:=D(u^{(k)}S)
       =\mathbf1_{\{\widetilde u>-k\}}\mu+L_{-k}(u;S)\ge0.
 \tag{KT7}
\]
It is a positive finite-energy Radon measure, with
\[
       \|\mu^{(k)}\|_{-1,V}\le\mathcal E(u^{(k)})^{1/2}
                                      \le(k\mu(F))^{1/2}.
 \tag{KT8}
\]
For every compact $K\Subset V$,
\[
 \|\mu^{(k)}-\mu\|_{\rm TV(K)}\longrightarrow0,
 \qquad L_{-k}(u;S)(K)\longrightarrow0.
 \tag{KT9}
\]
Here ``finite-energy'' means the bounded $H$-dual norm; no
unproved finiteness of the total mass of $\mu^{(k)}$ near the
ordinary boundary is asserted.
\end{theorem}

\noindent\textbf{Status: \texttt{NEW\_PROVED}}.  This produces
bounded relative $S$-subharmonic energy potentials approximating
every positive finite diffuse datum in local curvature variation.
Membership in ambient classical psh, in the frozen classical
weak completion, or in an unrestricted nonlinear MA solution
class is not asserted.

\begin{proof}
The sign, $H$ membership and energy estimate are (MD8)--(MD10).
Apply (KT4) with $t=-k$ and subtract the constant $k$ to obtain
(KT7).  For every original core $\phi$,
\[
 \left|\int\phi\,d\mu^{(k)}\right|
     =|\mathcal E(u^{(k)},\phi)|
                   \le\mathcal E(u^{(k)})^{1/2}\mathcal E(\phi)^{1/2}.
\]
Positivity upgrades this to the usual positive finite-energy
measure functional, proving (KT8).

One has $u^{(k)}\to u$ in $L^1(\lambda)$ by dominated
convergence, so $\mu^{(k)}\to\mu$ on all original compact
smooth tests.  Moreover $\widetilde u$ is finite outside an
energy-polar set, which $\mu$ does not charge.  Hence
$\mathbf1_{\{\widetilde u>-k\}}\mu\to\mu$ in total variation.
For a nonnegative original compact smooth $\chi$ equal to one
near $K$, (KT7) gives
\[
 0\le L_{-k}(u;S)(K)
 \le\langle\mu^{(k)},\chi\rangle
           -\int\chi\mathbf1_{\{\widetilde u>-k\}}\,d\mu
                                                   \longrightarrow0.
\]
Finally
\[
 \|\mu^{(k)}-\mu\|_{\rm TV(K)}
 \le L_{-k}(u;S)(K)+\mu(K\cap\{\widetilde u\le-k\})
                                                   \longrightarrow0.
\]
This proves (KT9) without assuming a carrier theorem for the
positive defect measure and without interchanging a distributional
limit with an arbitrary Borel restriction.
\end{proof}

\begin{proposition}[A refined band estimate and vanishing radiation]
\label{prop:r11-diffuse-band-radiation}
For signed finite diffuse $\mu$, $0\le a<b$, and
$B_{a,b}(s)=T_b s-T_a s$, one has
\[
 \begin{split}
 \mathcal E(B_{a,b}(u),B_{a,b}(u))
       &\le -\int B_{a,b}(\widetilde u)\,d\mu\\
       &\le(b-a)|\mu|(\{|\widetilde u|>a\}).
 \end{split}
 \tag{KT10}
\]
In particular,
\[
       \frac{1}{k}\mathcal E(T_{2k}u-T_ku)
          \le|\mu|(\{|\widetilde u|>k\})\longrightarrow0.
 \tag{KT11}
\]
These are estimates in the canonical form; they do not by themselves
assert the full renormalized chain identity for every unbounded test.
\end{proposition}

\noindent\textbf{Status: \texttt{NEW\_PROVED}}, as the strengthened
truncation estimate for arbitrary original $S$ and all finite diffuse
data.

\begin{proof}
The scalar map $P(s)=s-2B_{a,b}(s)$ is a normal contraction:
it fixes zero and has slopes $1,-1,1$ on the successive absolute
value regions.  For each finite-energy $u_j$, the Markov property
and polarization give
\[
 \mathcal E(B_{a,b}(u_j))
      \le\mathcal E(u_j,B_{a,b}(u_j))
      =-\int B_{a,b}(\widetilde u_j)\,d\mu_j.
 \tag{KT12}
\]
The right side is at most $(b-a)\|\mu_j\|_{\rm TV}$.
The vectors on the left are therefore bounded in $H$, and
\[
 \|B_{a,b}(u_j)-B_{a,b}(u)\|_2^2
                         \le2(b-a)\|u_j-u\|_1\longrightarrow0.
\]
Every weak $H$ limit is thus $B_{a,b}(u)$.
The right side of (KT12) converges to the middle term of (KT10)
by total-variation convergence and the diffuse-capacity lemma,
since $B_{a,b}$ is bounded and continuous.
Weak lower semicontinuity proves the first inequality in (KT10).
The second uses $|B_{a,b}|\le(b-a)\mathbf1_{\{|s|>a\}}$.
Taking $a=k,b=2k$ proves (KT11), since $\widetilde u$ is finite
$|\mu|$-almost everywhere.
\end{proof}

\paragraph{Progress beyond this first extension and retained gates.}
The Kato argument alone does not identify (KT3) with an exact-level
carrier.  The following two sections prove that identity, the full
renormalized chain and global curvature variation first for positive
data and then for every finite signed diffuse datum.  They do not prove
ambient classical representability, an unrestricted nonlinear
Monge--Ampere Dirichlet problem, or ordinary pointwise boundary
values.  Appropriate hypotheses for the remaining assertions are
still \texttt{OPEN}; unrestricted ambient classical solvability is
already contradicted by the frozen R10 compact-source examples.
Unrestricted equality of energy and BT polar ideals is
\texttt{FALSE}, by the frozen R6
\path{capacity_energy_null_gap.tex}; it is not an open gate.

% END REVIEWED MODULE research_r11/dirichlet/kato_truncation_measure_data.tex


% BEGIN REVIEWED MODULE research_r11/dirichlet/positive_diffuse_renormalized_dirichlet.tex
\section{Positive diffuse data: renormalized equations and exact level carriers}
\label{sec:r11-positive-diffuse-renormalized}

Keep the original affinoid-neighborhood setting of
Sections~\ref{sec:r11-affinoid-energy}--\ref{sec:r11-diffuse-Kato}:
$S$ is an arbitrary original symmetric middle-positive separately
closed $(n-1,n-1)$-current, $V\Subset\Omega$ is ordinary,
$H=\mathscr D_0(V)$ is its faithful coercive form, and
$\lambda(V)<\infty$.  In the main results of this section
$\mu\ge0$ is a finite energy-diffuse Radon measure on
$F=V\cap\operatorname{supp}S$ and $u=u_\mu\le0$ is the
already constructed transposition solution.  No ambient classical
psh membership is assumed or concluded.

\begin{lemma}[Clipping and exact level carriers in the energy domain]
\label{lem:r11-H-level-carrier}
For $v\in H$ the normalized hinge
$h_t(v)=(v-t)^+-(-t)^+$ belongs to $H$ and
\[
 d h_t(v)=\mathbf1_{\{\widetilde v>t\}}\,dv,
 \qquad \Gamma_S(v)\{\widetilde v=t\}=0.
 \tag{RT1}
\]
Here the differential is the closed original gradient module
constructed in Lemma~\ref{lem:r11-H-original-chain}.
In particular the usual clipping and band chain formulas hold.
Suppose $D(vS)=\sigma$ on the original tests, where $|\sigma|$
is a positive Radon measure of finite energy integral.
No finite total mass near the boundary is required.
Then
\[
 L_t(v;S)=D(h_t(v)S)
                    -\mathbf1_{\{\widetilde v>t\}}\sigma
 \tag{RT2}
\]
is a positive finite-energy Radon measure and
\[
                    \mathbf1_{\{\widetilde v\ne t\}}L_t(v;S)=0.
 \tag{RT3}
\]
These are identities for the canonical energy representatives.
\end{lemma}

\noindent\textbf{Status: \texttt{NEW\_PROVED}} for the
all-original-current level carrier; its closed-form scalar
functional calculus is \texttt{DIRECT\_COROLLARY}.

\begin{proof}
Use the one-sided convex approximations
$\Phi_\varepsilon(s)$ from (KT5).  Their normalized compositions
$\Phi_\varepsilon(v-t)-\Phi_\varepsilon(-t)$ converge in $L^2(\lambda)$
to $h_t(v)$, uniformly in scalar values with error at most
$2\varepsilon$.  Their differentials converge in the gradient
module to $\mathbf1_{\{\widetilde v>t\}}dv$ by bounded convergence
against the fixed finite measure $\Gamma_S(v)$; at $v=t$ all
one-sided derivatives are zero.  Closedness proves strong
convergence in $H$ and the first identity in (RT1).
Apply it also to $-v$ and use
$h_t(v)-h_{-t}(-v)=v$.  The difference of the two differential
identities is $dv=\mathbf1_{\{\widetilde v\ne t\}}dv$,
which proves the level-zero statement.
Differences of hinges give the clipping and band formulas.

The finite-energy chain proof (KT1) applies verbatim to $v$:
its use of the datum requires finite energy integral of
$|\sigma|$, not finite total mass, because every $H$ test is
absolutely integrable for that measure.  Hence
\[
 L_{t,\varepsilon}
 :=D(\Phi_\varepsilon(v-t)S)
                  -\Phi_\varepsilon'(\widetilde v-t)\sigma
   =\Phi_\varepsilon''(\widetilde v-t)\Gamma_S(v)\ge0.
 \tag{RT4}
\]
Its $H$-dual norm is at most
$\|v\|_H+\||\sigma|\|_{-1,V}$: the first potential is a
normal contraction of $v$ after normalization, and the second
measure is dominated in variation by $|\sigma|$.
For each fixed $z\in H$, the first functional in (RT4)
converges by strong $H$ convergence, and the second by dominated
convergence against $|z|\,|\sigma|$.  Thus (RT2) is the weak
$H$-dual limit of positive finite-energy measures.  Its
positive original core functional is a positive Radon measure;
the uniform energy bound gives its canonical integration
interpretation on all of $H$.

Take a nonnegative original compact smooth $\psi$ and a
nonnegative smooth bounded scalar $g$ with bounded derivative
and vanishing on a neighborhood of $t$.
The product $z=\psi g(v)$ belongs to $H$ by the closed
product and chain rules.  For small enough $\varepsilon$,
\[
 \int \widetilde z\,dL_{t,\varepsilon}
  =\int\psi g(\widetilde v)
           \Phi_\varepsilon''(\widetilde v-t)\,d\Gamma_S(v)=0.
\]
Pass in the $H$-dual pairing just proved to obtain
$\int\psi g(\widetilde v)\,dL_t(v;S)=0$.
Countably many such $g$ positive away from $t$, and a compact
smooth exhaustion of $F$, prove (RT3).  In particular this
carrier argument does not restrict an arbitrary weak measure
limit to a discontinuous level set.
\end{proof}

\begin{theorem}[Strong convergence of every positive-data truncation]
\label{thm:r11-positive-diffuse-strong-truncations}
Choose the increasing finite-energy nest restrictions
$\mu_j\uparrow\mu$, put $u_j=-G\mu_j\downarrow u$, and
fix $k>0$.  Then
\[
                  T_k u_j\longrightarrow T_k u
                                  \quad\hbox{strongly in }H.
 \tag{RT5}
\]
The canonical monotonicity assertions hold outside one
energy-polar set.  No compact embedding or unproved density
of special currents is used.
\end{theorem}

\noindent\textbf{Status: \texttt{NEW\_PROVED}}.

\begin{proof}
Put $w_j=-u_j=G\mu_j\ge0$ and $v_j=\min(w_j,k)$.
Order preservation of the already constructed Green map gives
$v_j\uparrow\min(-u,k)$ quasi-everywhere; its capacity and
$L^1$ limits identify this monotone limit with the displayed
canonical representative.  The Kato theorem gives
$-D(v_jS)\ge0$.  This is a positive finite-energy measure,
since $v_j\in H$, so
$\mathcal E(v_j,z)\ge0$ for every $z\in H_+$.
For $l\ge j$, put $z=v_l-v_j\ge0$.  Then
\[
 \mathcal E(v_l)=\mathcal E(v_j)+2\mathcal E(v_j,z)
                                      +\mathcal E(z)
                  \ge\mathcal E(v_j)+\mathcal E(z).
 \tag{RT6}
\]
The energies therefore increase and are bounded by $k\mu(F)$
by (MD5).  Formula (RT6) makes $(v_j)$ Cauchy in $H$.
Its $L^2$ limit is the clipping of $-u$, by the previously
proved $L^1$ convergence and boundedness.  This proves (RT5).
\end{proof}

\begin{definition}[The sigma-finite energy measure of a diffuse potential]
\label{def:r11-positive-diffuse-energy-measure}
For $u=u_\mu$ as above set
\[
              \Gamma_S(u)=\lim_{k\to\infty}\Gamma_S(T_ku).
 \tag{RT7}
\]
This is an increasing limit of positive measures.  Indeed for
$l>k$, (RT1) applied to $T_lu$ gives
\[
 \Gamma_S(T_ku)
       =\mathbf1_{\{|\widetilde u|<k\}}\Gamma_S(T_lu).
 \tag{RT8}
\]
It is a sigma-finite Borel measure carried by the finite-value
locus, with finite mass on every bounded value strip, and
\[
 \Gamma_S(T_ku)
        =\mathbf1_{\{|\widetilde u|<k\}}\Gamma_S(u),
 \qquad \Gamma_S(u)\{\widetilde u=t\}=0\quad(t\in\mathbf R).
 \tag{RT9}
\]
It need not be locally finite on every ordinary compact set.
Only its bounded-value restrictions or explicitly integrable
weights will be used as Radon measures.  It is absolutely
continuous with respect to $\lambda$ on its sigma-finite
finite-value exhaustion.
\end{definition}

\begin{theorem}[A complete positive-data renormalized chain identity]
\label{thm:r11-positive-diffuse-renormalized-chain}
Let $\Phi\in C^2(\mathbf R)$ have compactly supported first
derivative.  Then $\Phi(u)-\Phi(0)\in H$ and
\[
 D(\Phi(u)S)
     =\Phi'(\widetilde u)\mu
                        +\Phi''(\widetilde u)\Gamma_S(u)
 \tag{RT10}
\]
on every original compact smooth test.  The right side is a
finite signed Radon measure, because the energy weight has
bounded value support.  This is an actual original current
identity, not only an identity in an abstract form dual.
\end{theorem}

\noindent\textbf{Status: \texttt{NEW\_PROVED}} for all positive
finite diffuse data and arbitrary original $S$.

\begin{proof}
Choose $M>0$ so that $\Phi'$ and $\Phi''$ vanish on
$(-\infty,-M]$.  Since all potentials are nonpositive,
\[
             \Phi(u_j)=\Phi(T_Mu_j),\qquad
             \Phi(u)=\Phi(T_Mu).
\]
By (RT5) and the closed smooth scalar chain rule, the
normalized compositions converge strongly in $H$.
The finite-energy chain identity for $u_j$ says
\[
 D(\Phi(u_j)S)=\Phi'(\widetilde u_j)\mu_j
                   +\Phi''(\widetilde u_j)\Gamma_S(T_Mu_j).
 \tag{RT11}
\]
At and below the truncation level both scalar derivatives
vanish, so this replacement of the energy is exact.

The source terms converge in total variation by the
capacity-diffuse absolute-continuity lemma and
$\mu_j\to\mu$ in total variation.  The energy measures of
$T_Mu_j$ converge in total variation by strong $H$ convergence
and (KT2).  Multipliers on the remaining fixed limiting
energy measure converge almost everywhere after extraction:
capacity convergence gives a quasi-everywhere subsequence,
and this energy measure charges no polar sets.  Bounded
convergence then gives the desired limit; the usual
subsequence criterion removes the extraction.  The limiting
weighted energy is $\Phi''(\widetilde u)\Gamma_S(u)$ by
(RT8)--(RT9).  On the left, strong $H$ convergence implies
$L^1(\lambda)$ convergence, so the faithful coefficient
bounds pass the original current.  This proves (RT10).
The first term has finite mass and the second is bounded by
a constant times $\Gamma_S(T_Mu)(F)$, proving finiteness.
For $C^2$, rather than smooth, $\Phi$, smooth approximation in
$C^2$ on the bounded clipped range passes the same identities.
\end{proof}

\begin{theorem}[Exact level carriers for positive diffuse data]
\label{thm:r11-positive-diffuse-exact-level-carrier}
For every $t\in\mathbf R$ the positive Kato measure in (KT3)
satisfies
\[
                       \mathbf1_{\{\widetilde u\ne t\}}L_t(u;S)=0.
 \tag{RT12}
\]
For $M>0$, write $\mu^{(M)}=D(\max(u,-M)S)$.  Then
\[
 \mathbf1_{\{\widetilde u>-M\}}\mu^{(M)}
           =\mathbf1_{\{\widetilde u>-M\}}\mu,
 \qquad
 \mathbf1_{\{\widetilde u<-M\}}\mu^{(M)}=0.
 \tag{RT13}
\]
Every restriction is a restriction of actual Radon measures
on the original space.  The statement for arbitrary signed
finite diffuse data is not inferred from this proof.
\end{theorem}

\noindent\textbf{Status: \texttt{NEW\_PROVED}}.

\begin{proof}
Set $v^M=\max(u,-M)\in H$.  Its curvature $\mu^{(M)}$ is
positive and of finite energy by (KT7)--(KT8).
Choose a smooth scalar $\Phi$ whose first derivative is zero
on $(-\infty,-M]$, strictly positive on $(-M,0]$, and zero
above a fixed positive number.  Its second derivative also
vanishes at and below $-M$.
One has $\Phi(v^M)=\Phi(u)$, and (RT9) identifies their
weighted energy terms.  Compare (RT10) for $u$ with the
finite-energy chain identity for $v^M$ to obtain
\[
            \Phi'(\widetilde u)\mu^{(M)}
                         =\Phi'(\widetilde u)\mu.
\]
The strictly positive multiplier on $\{u>-M\}$ proves the
first assertion of (RT13).

Next $v^M=\max(v^{M+1},-M)$.
Lemma~\ref{lem:r11-H-level-carrier}, applied to the $H$ vector
$v^{M+1}$ and its positive finite-energy curvature, gives
\[
 \mu^{(M)}=
 \mathbf1_{\{\widetilde v^{M+1}>-M\}}\mu^{(M+1)}
                          +L_{-M}(v^{M+1};S),
\]
where the last measure is carried by
$\{\widetilde v^{M+1}=-M\}$.
On $\{\widetilde u<-M\}$ one has
$\widetilde v^{M+1}<-M$, so both terms vanish.  This proves
the second assertion of (RT13).  The positive difference
$\mu^{(M)}-\mathbf1_{\{u>-M\}}\mu$ can therefore be carried
only by $\{u=-M\}$, proving (RT12) for $t<0$.
For $t\ge0$ both terms in (KT3) vanish, since $u\le0$.
\end{proof}

\begin{theorem}[An exact energy distribution and global variation control]
\label{thm:r11-positive-diffuse-global-variation}
For $0\le a<b$ one has the exact band identity
\[
 \mathcal E(T_bu-T_au)
       =-\int(T_b\widetilde u-T_a\widetilde u)\,d\mu.
 \tag{RT14}
\]
Consequently the sigma-finite scalar energy image is
\[
       \widetilde u_*\Gamma_S(u)
          =\mathbf1_{\{t<0\}}\mu\{\widetilde u<t\}\,dt.
 \tag{RT15}
\]
For every $t<0$ and $k>0$,
\[
 \begin{gathered}
     L_t(u;S)(F)\le\mu\{\widetilde u\le t\},\qquad
     \mu^{(k)}(F)\le\mu(F),\\
 \|\mu^{(k)}-\mu\|_{\rm TV(F)}
       \le2\mu\{\widetilde u\le-k\}\longrightarrow0.
 \end{gathered}
 \tag{RT16}
\]
Thus hard truncations yield positive finite-total-mass,
finite-energy original curvature measures, converging in
\emph{global} total variation on the part domain.
\end{theorem}

\noindent\textbf{Status: \texttt{NEW\_PROVED}}.  The global
conclusion strengthens (KT9); it does not follow merely from
local vague convergence or the earlier energy norm bound.

\begin{proof}
For the finite-energy $u_j$, (RT1) gives
\[
 \mathcal E(T_bu_j-T_au_j)
  =\mathcal E(u_j,T_bu_j-T_au_j)
  =-\int(T_b\widetilde u_j-T_a\widetilde u_j)\,d\mu_j.
\]
The first equality holds because the differential of the band
map is the indicator of $a<|u_j|<b$, and hence its square
is the same indicator.  Strong convergence of both
truncations proves convergence on the left.  Bounded scalar
continuity, diffuseness and total-variation convergence prove
convergence on the right, giving (RT14).
Writing $w=-u\ge0$, layer cake gives
\[
 \Gamma_S(u)\{a<w<b\}
       =\int_a^b\mu\{w>r\}\,dr.
\]
The zero-energy level assertion (RT9) handles endpoints.
These intervals determine the scalar measure on $(0,\infty)$;
there is no energy mass at $w=0$.  This proves (RT15).

For a fixed $t<0$, apply (RT10) to the one-sided hinge
approximations $\Phi_\varepsilon(u-t)$, modifying their
scalar derivative above the range $u\le0$ to have compact
support.  Their positive energy defects are
\[
       L_{t,\varepsilon}
         =\Phi_\varepsilon''(\widetilde u-t)\Gamma_S(u).
\]
The second derivative is nonnegative, supported in $(0,\varepsilon)$,
and has real integral one.  Formula (RT15) therefore gives
\[
 L_{t,\varepsilon}(F)
   =\int_t^{t+\varepsilon}
        \Phi_\varepsilon''(s-t)\mu\{u<s\}\,ds
   \le\mu\{u<t+\varepsilon\}
\]
for $0<\varepsilon<-t$.
These measures converge on the original compact tests to
$L_t(u;S)$, by the same potential and source convergence as
in the Kato proof.  Positivity and compact smooth exhaustion
of $F$ imply
$L_t(u;S)(F)\le\mu\{u\le t\}$ upon letting
$\varepsilon\downarrow0$.
The Kato decomposition now gives
\[
 \mu^{(k)}(F)
  \le\mu\{u>-k\}+\mu\{u\le-k\}=\mu(F),
\]
and
\[
 \|\mu^{(k)}-\mu\|_{\rm TV(F)}
   \le L_{-k}(u;S)(F)+\mu\{u\le-k\}
   \le2\mu\{u\le-k\}.
\]
The last quantity tends to zero because $u$ is finite outside
an energy-polar set and $\mu$ is diffuse.  This proves (RT16).
\end{proof}

\begin{corollary}[Exact local coarea and the finite-energy threshold]
\label{cor:r11-positive-diffuse-coarea-threshold}
For every bounded Borel scalar $\theta$ with compact support
in $\mathbf R$ one has
\[
           \int_{\mathbf R}\theta(t)L_t(u;S)\,dt
                     =\theta(\widetilde u)\Gamma_S(u).
 \tag{RT17}
\]
For signed $\theta$ this is an equality of finite signed
Radon measures.  Moreover
\[
       u\in H\quad\Longleftrightarrow\quad
                     \int_F|\widetilde u|\,d\mu<\infty,
 \qquad
       \mathcal E(u,u)=\int_F|\widetilde u|\,d\mu
                        \quad\hbox{when finite}.
 \tag{RT18}
\]
\end{corollary}

\noindent\textbf{Status: \texttt{DIRECT\_COROLLARY}} of the
renormalized chain, exact level measures and band identity.

\begin{proof}
First take smooth compactly supported $\theta$ and put
$F(s)=\int\theta(t)(s-t)^+dt$.
On the nonpositive range, $F'$ vanishes sufficiently far to
the left; modify it above that range to have compact support.
The renormalized chain, $F''=\theta$, and Fubini for the
finite source measure prove (RT17) on each original test.
The global bound $L_t(F)\le\mu(F)$ for $t<0$ makes the kernel
locally integrable in the scalar variable.  Its original-test
pairings are Borel by (KT3), and positive Radon kernel
extension gives (RT17) for nonnegative Borel $\theta$;
subtract positive and negative parts for the signed case.
The right side is finite because its value support is bounded.

Set $a=0,b=k$ in (RT14).  The energies
$\mathcal E(T_ku)=\int\min(|u|,k)\,d\mu$ increase to the
right side of (RT18).  If this is finite, the ordered
superharmonic truncations of $-u$ satisfy (RT6), so they
are Cauchy in $H$ and have limit $-u$.  Conversely if $u\in H$,
(RT9) and monotone convergence give
$\mathcal E(T_ku)\uparrow\mathcal E(u)$, proving (RT18).
\end{proof}

\begin{definition}[The positive-data renormalized zero-boundary class]
\label{def:r11-renormalized-positive-boundary}
A candidate $v$ is a nonpositive $L^1(F,\lambda)$ function with
a finite energy-quasicontinuous representative, all $T_kv\in H$,
and their compatible sigma-finite energy measure defined as
in (RT7)--(RT9).  For the specified finite positive diffuse datum
$\mu$, require
\[
 \begin{split}
 D(\Phi(v)S)&=\Phi'(\widetilde v)\mu
                         +\Phi''(\widetilde v)\Gamma_S(v)
,\\
 \frac{1}{k}\mathcal E(T_{2k}v-T_kv)&\longrightarrow0.
 \end{split}
 \tag{RT19}
\]
The first line is required for every $\Phi\in C^2(\mathbf R)$
with compactly supported first derivative and includes every
original compact smooth test.  The second is the vanishing
radiation condition.  The boundary
condition is the explicit membership $T_kv\in H$; the definition
does not use the transposition identity or the unknown Green
potential of the datum.
\end{definition}

\begin{theorem}[Existence and uniqueness in the renormalized class]
\label{thm:r11-positive-renormalized-Dirichlet}
Each finite positive energy-diffuse datum on $F$ has exactly one
candidate in Definition
\ref{def:r11-renormalized-positive-boundary}.  It is $u_\mu$.
Its original current satisfies $D(u_\mu S)=\mu$, its level
carriers and curvature truncations satisfy (RT12)--(RT16), and
its solution class is independent of the faithful reference
measure constructed from the affinoid compensators.
\end{theorem}

\noindent\textbf{Status: \texttt{NEW\_PROVED}} as an intrinsic
renormalized relative Dirichlet theorem, with a concrete
zero-boundary class rather than uniqueness imposed by the
transposition formula itself.

\begin{proof}
Existence follows from the positive-data renormalized chain
and the radiation estimate (KT11).  For uniqueness choose
smooth scalar functions $\Phi_k$ with $\Phi_k(0)=0$,
$0\le\Phi_k'\le1$ on $(-\infty,0]$,
$\Phi_k'=0$ below $-2k$, $\Phi_k'=1$ above $-k$ on this
range, and
$0\le\Phi_k''\le C/k$ on $(-2k,-k)$.
Modify the first derivative above $1$ to have compact support.
Then $|\Phi_k(s)|\le|s|$ for $s\le0$ and
$\Phi_k(s)\to s$ there.

For any candidate $v$ one has
$\Phi_k(v)=\Phi_k(T_{2k}v)\in H$.
Its prescribed chain gives a positive measure
\[
 \nu_k:=D(\Phi_k(v)S)
       =\Phi_k'(\widetilde v)\mu
                           +\Phi_k''(\widetilde v)\Gamma_S(v).
 \tag{RT20}
\]
It has finite total mass, and (RT9) for the candidate's compatible
energy gives
\[
 \|\nu_k-\mu\|_{\rm TV}
 \le\mu\{|\widetilde v|>k\}
       +\frac{C}{k}\mathcal E(T_{2k}v-T_kv)
                                         \longrightarrow0.
 \tag{RT21}
\]
The first term tends to zero because the candidate is finite
quasi-everywhere and the datum is diffuse; the second is its
radiation condition.  Since $\Phi_k(v)\in H$, the positive
measure $\nu_k$ has finite energy integral and its Riesz
solution is exactly $\Phi_k(v)$ with the stated sign convention.
Therefore the established finite-data stability gives
\[
                 \|\Phi_k(v)-u_\mu\|_{L^1(\lambda)}
                                  \le2\|\nu_k-\mu\|_{\rm TV}\to0.
\]
On the other hand dominated convergence gives
$\Phi_k(v)\to v$ in $L^1(\lambda)$.  Thus $v=u_\mu$ in $L^1(\lambda)$.  Each pair of clipped
$H$ vectors is equal almost everywhere and hence quasi-everywhere;
a countable set of clipping levels therefore identifies the
canonical representatives quasi-everywhere.
Reference independence follows either from the proved intrinsic
core-energy independence and the definition (RT19), or from
the already proved reference independence of $u_\mu$.
All current identities use the faithful coefficient definition,
so they include every original test and actual refinement.
\end{proof}

\paragraph{Signed extension and remaining gates.}
The full signed diffuse renormalized chain and exact-level carrier
are not consequences of the positive monotone-nest argument.
They are proved separately in
Section~\ref{sec:r11-signed-diffuse-renormalized} by the positive
$|\mu|$ barrier and localized gradient estimates.
Ambient classical representability, ordinary pointwise boundary
regularity, an unrestricted diagonal nonlinear MA Dirichlet
problem, general GQ and low-real-degree CTO require further
hypotheses or proofs.  The unrestricted ambient classical claim is
\texttt{FALSE} by frozen R10, and unrestricted equality of energy
and BT polar ideals is \texttt{FALSE} by frozen R6
\path{capacity_energy_null_gap.tex}.  Neither failure is relisted
as an open theorem.  Smooth or geometric backgrounds are not
assumptions anywhere in these positive results.

% END REVIEWED MODULE research_r11/dirichlet/positive_diffuse_renormalized_dirichlet.tex


% BEGIN REVIEWED MODULE research_r11/dirichlet/signed_diffuse_renormalized_dirichlet.tex
\section{Signed diffuse data: strong truncations and the full renormalized problem}
\label{sec:r11-signed-diffuse-renormalized}

Keep the same original arbitrary $(n-1,n-1)$ background, affinoid
neighborhood and part form as above.  Let $\mu$ be any finite
signed energy-diffuse Radon measure on $F$, and put
$B=\|\mu\|_{\rm TV}$.  All currents, coefficients and tests are
original CLD objects.  The new argument uses the already proved
positive-data theorem for the \emph{datum} $|\mu|$, with the same
original background $S$.  It does not replace that background.

\begin{lemma}[Bounded energy tests and signed Radon functionals]
\label{lem:r11-bounded-H-tests-diffuse}
Every bounded $z\in H$ has original compact smooth core
approximations converging strongly in $H$, uniformly bounded
in absolute value, and along a subsequence quasi-everywhere.
If a finite signed Radon measure $\nu$ on $F$ has a continuous
core functional for the energy norm, then $|\nu|$ charges no
energy-polar set.  If $h\in H$ and $D(hS)=\nu$ on every original
compact smooth test, then
\[
       \mathcal E(h,z)=-\int\widetilde z\,d\nu
                     \qquad(z\in H\cap L^\infty),
       \qquad h=u_\nu .
 \tag{SD1}
\]
The last equality is the already defined transposition solution.
Neither conclusion assumes that $|\nu|$ has a finite-energy norm.
\end{lemma}

\noindent\textbf{Status: \texttt{NEW\_PROVED}} in the stated
original-core realization; the compact-core capacity formula is
an already specified \texttt{EXTERNAL\_KNOWN} form input.

\begin{proof}
For bounded $z$, start with any original core sequence $z_j\to z$
in $H$ and compose it with one fixed smooth normal clamp that
fixes a neighborhood of the range of $z$ and is bounded globally.
The gradient-module chain rule proves strong convergence of these
compositions to $z$.  They are still original compact smooth
functions because the clamp fixes zero.  A summable-capacity
subsequence converges quasi-everywhere.

Let $K$ be an energy-polar compact set and choose an ordinary
relatively open neighborhood $G$ with compact closure such that
$|\nu|(G\setminus K)<\epsilon$.
The regular compact-core capacity formula, coercivity, a smooth
scalar clamp, and a fixed original cutoff supported in $G$
give core functions $0\le a_j\le1$, equal to one near $K$,
supported in $G$, with $\mathcal E(a_j)\to0$.
Multiplication by that cutoff preserves the energy convergence
by (AE12).  Hence $\int a_j\,d\nu\to0$, while
\[
       \left|\nu(K)-\int a_j\,d\nu\right|
                            \le |\nu|(G\setminus K)<\epsilon.
\]
Let $\epsilon\downarrow0$.  The same argument applies to every
polar compact subset of the positive and negative Hahn sets.
Radon inner regularity of the two variations therefore proves
that $|\nu|$ vanishes on every polar Borel set.

For bounded $z\in H$ use the uniformly bounded core
approximation just constructed.  Energy continuity passes its
left side in the original identity; diffuseness and dominated
convergence pass the right side.  This proves (SD1)'s first
assertion.  For $f\in L^\infty(\lambda)$ take $z=Gf$, which
is bounded by the torsion estimate.  Symmetry gives
\[
 \int h f\,d\lambda=\mathcal E(h,Gf)
                         =-\int\widetilde{Gf}\,d\nu.
\]
Since $h\in L^2(\lambda)\subset L^1(\lambda)$, transposition
uniqueness gives $h=u_\nu$.
\end{proof}

The usual product rule extends to bounded members of $H$:
$d(z_1z_2)=z_1dz_2+z_2dz_1$.  Indeed the preceding uniformly bounded
core approximations and the original gradient module pass this
identity by dominated convergence against the fixed limiting
energy measures.  This fact will be used for a product of two
bounded energy functions, not only a smooth cutoff.

\begin{theorem}[Strong clipping convergence for signed nest data]
\label{thm:r11-signed-strong-truncations}
Let $\mu_j=\mathbf1_{F_j}\mu$ be the finite-energy nest
restrictions from (MD3), and put $u_j=-G\mu_j$.
For every $k>0$,
\[
                          T_ku_j\longrightarrow T_ku_\mu
                                      \quad\hbox{strongly in }H.
 \tag{SD2}
\]
The conclusion does not require monotonicity of the signed
potentials.
\end{theorem}

\noindent\textbf{Status: \texttt{NEW\_PROVED}}.

\begin{proof}
Let $w=-u_{|\mu|}=G|\mu|\ge0$ be the positive-data Green
potential, with its canonical representative.  Positivity and
order of the transposition map give
$|u_j|,|u_\mu|\le w$ quasi-everywhere for all $j$.
Put $a_M=|\mu|\{w>M\}$, so $a_M\to0$, and define
\[
          \zeta_M=\frac{T_{2M}w-T_Mw}{M}\in H,\qquad
          0\le\zeta_M\le1.
\]
It equals one on $\{w\ge2M\}$.  The exact positive-data
band identity and the closed clipping rule give
\[
 \mathcal E(\zeta_M)\le\frac{a_M}{M},\qquad
 d\zeta_M=\frac1M\mathbf1_{\{M<w<2M\}}\,dw.
 \tag{SD3}
\]
Here $dw$ on the printed bounded band means the compatible
clipped gradient $d(T_{2M+1}w)$ there; no globally Radon
first-order current of the unbounded $w$ is required.
Test the finite-energy equation for $u_j$ with
$\zeta_M^2T_ku_j\in H\cap L^\infty$.
The original closed gradient product and clipping rules give
\[
 \int\zeta_M^2\,d\Gamma_S(T_ku_j)
  =-\int\zeta_M^2 T_k\widetilde u_j\,d\mu_j
                  -2\int\zeta_M T_k\widetilde u_j\,
                                      d\Gamma_S(u_j,\zeta_M).
\]
The first term in absolute value is at most $k a_M$.
The cross energy is carried where $M<w<2M$, and there
$|u_j|<2M$, so its $u_j$ factor may be clipped at $2M$.
Cauchy--Schwarz and (MD5) therefore prove
\[
 \sup_j\Gamma_S(T_ku_j)\{w\ge2M\}
   \le k a_M+2k\sqrt{\mathcal E(\zeta_M)
                                 \sup_j\mathcal E(T_{2M}u_j)}
   \le k a_M+2k\sqrt{2B a_M}\longrightarrow0.
 \tag{SD4}
\]
This is the required uniform control of energy in the region
where positive and negative potentials could be large and cancel.

On $A_M=\{w\le M\}$ one has $|u_j-u_l|\le2M$.
The first-order clipping rule and (MD5), applied to the
finite-energy datum $\mu_j-\mu_l$, give
\[
 \|\mathbf1_{A_M}(du_j-du_l)\|^2
      \le\mathcal E(T_{2M+1}(u_j-u_l))
      \le(2M+1)\|\mu_j-\mu_l\|_{\rm TV}\longrightarrow0.
 \tag{SD5}
\]
Thus the localized gradients have a fixed limit in the
original Hilbert gradient module.  Its energy measure is
absolutely continuous with respect to $\lambda$.
For a smooth scalar $P$ with $P(0)=0$ and compactly supported
first derivative, bounded convergence against this fixed
limiting module measure proves that
$\mathbf1_{A_M}d(P(u_j))$ converges strongly.
Here $u_j\to u$ in $L^1(\lambda)$, and hence in measure;
continuity and boundedness of $P'$ supply the required
convergence of its scalar multipliers.
The energy of $P(u_j)$ outside $A_M$ is bounded by a fixed
multiple of the tail of a clipping $T_Lu_j$, where $L$
contains the support of $P'$.  Formula (SD4) makes it small
uniformly in $j$ as $M\to\infty$.
It follows that $P(u_j)$ is Cauchy in $H$; its $L^2$ limit
is $P(u)$, by boundedness and the $L^1$ convergence of $u_j$.

Finally approximate $T_k$ by smooth normal clamps
$P_\epsilon$ with uniform scalar error at most $\epsilon$,
whose derivatives differ from that of $T_k$ only where
$k-\epsilon<|s|<k+\epsilon$.  For $0<\epsilon<k/2$, the
finite-energy band formula and the source equation give
\[
 \begin{split}
 \mathcal E(T_ku_j-P_\epsilon(u_j))
  &\le\mathcal E(T_{k+\epsilon}u_j-T_{k-\epsilon}u_j)\\
  &=-\int(T_{k+\epsilon}\widetilde u_j
                 -T_{k-\epsilon}\widetilde u_j)\,d\mu_j
       \le2\epsilon B.
 \end{split}
 \tag{SD6}
\]
The smooth compositions are strongly convergent by the
preceding paragraph.  First let $j,l\to\infty$ and then
$\epsilon\downarrow0$ in the triangle inequality.
This makes $T_ku_j$ Cauchy in $H$.  Its already known
$L^2$ limit is $T_ku$, proving (SD2).
\end{proof}

\begin{theorem}[Full signed renormalized chain and energy distribution]
\label{thm:r11-signed-renormalized-chain}
For $u=u_\mu$, the increasing compatible energy measure
$\Gamma_S(u)=\lim_k\Gamma_S(T_ku)$ is sigma-finite, finite
on every bounded value strip, and zero on every exact scalar
level.  If $\Phi\in C^2(\mathbf R)$ has bounded first derivative
and compactly supported second derivative, then
\[
 D(\Phi(u)S)=\Phi'(\widetilde u)\mu
                           +\Phi''(\widetilde u)\Gamma_S(u)
 \tag{SD7}
\]
on all original compact smooth tests.  Every term is coefficient-
defined, and the right side is a finite signed Radon measure.
If in addition $\Phi'$ has compact support, then
$\Phi(u)-\Phi(0)\in H$.
For every $0\le a<b$,
\[
 \mathcal E(T_bu-T_au)
                   =-\int(T_b\widetilde u-T_a\widetilde u)\,d\mu.
 \tag{SD8}
\]
The scalar energy image has density
\[
 \widetilde u_*\Gamma_S(u)=k_\mu(t)\,dt,\qquad
 k_\mu(t)=
 \begin{cases}
       \mu\{\widetilde u<t\},&t<0,\\
       -\mu\{\widetilde u>t\},&t>0.
 \end{cases}
 \tag{SD9}
\]
This density is nonnegative and bounded by $B$ almost everywhere;
its value at zero is immaterial.  In particular all bounded-value
weighted energy measures used here are finite Radon measures.
\end{theorem}

\noindent\textbf{Status: \texttt{NEW\_PROVED}} for signed
finite diffuse data.  Unweighted energy need not be an ordinary
locally finite current, as the positive-data counterexample already
shows.

\begin{proof}
Compatibility and level zero follow from the $H$ clipping
identities exactly as in (RT7)--(RT9), without any sign
hypothesis on $u$.  Choose $M$ containing the support of
$\Phi''$ in its interior.  The finite-energy chain for $u_j$
reads
\[
 D(\Phi(u_j)S)=\Phi'(\widetilde u_j)\mu_j
                         +\Phi''(\widetilde u_j)\Gamma_S(T_Mu_j).
\]
The left side passes because $\Phi$ is Lipschitz and
$u_j\to u$ in $L^1(\lambda)$.  The source term passes in
variation by the already proved capacity convergence of $u_j$
to $u$, diffuseness of $|\mu|$, bounded continuity of $\Phi'$, and
total-variation convergence of $\mu_j$ to $\mu$.
Strong convergence in (SD2), the energy variation estimate
(KT2), and bounded convergence against the fixed limiting
energy measure pass the last term.  This proves (SD7).
For merely $C^2$ scalar functions, ordinary smoothing gives
the same limit, uniformly in the first two derivatives on
the bounded active range and uniformly in scalar values
up to the original Lipschitz error.
If $\Phi'$ is compactly supported, its normalized composition
is a smooth function of one bounded clipping and lies in $H$.

The same strong clipping convergence passes the finite-energy
band identity, proving (SD8).  Apply the argument separately
to the bounded one-sided bands
\[
 B^+_{a,b}(s)=\min((s-a)^+,b-a),\qquad
 B^-_{a,b}(s)=-B^+_{a,b}(-s).
\]
Their derivatives are the indicators of $(a,b)$ and
$(-b,-a)$, respectively.  Composing a strongly convergent
clipping with these finite hinge combinations is strongly
continuous: the closed gradient formula applies, and the
limiting $H$ energy is zero on their finitely many corner
levels by (RT1).  Hence
\[
 \begin{split}
 \Gamma_S(u)\{a<u<b\}&=-\int_a^b\mu\{u>t\}\,dt,\\
 \Gamma_S(u)\{-b<u<-a\}&=\int_a^b\mu\{u<-t\}\,dt.
 \end{split}
\]
There is no energy mass at zero or at the interval endpoints.
These formulas prove (SD9).  Nonnegativity follows because
it is the density of a positive measure; the bound follows
from total variation of $\mu$.
\end{proof}

\begin{proposition}[Finite global level measures and clipping curvatures]
\label{prop:r11-signed-global-curvature-truncation}
The Kato measures $L_t(u;S)$ are finite for every real $t$,
with
\[
 \begin{aligned}
 L_t(u;S)(F)&\le |\mu|\{u\le t\}&& (t<0),\\
 L_t(u;S)(F)&\le |\mu|\{u>t\}&& (t>0),\\
 L_0(u;S)(F)&\le B.
 \end{aligned}
 \tag{SD10}
\]
For $M>0$, the actual signed Radon curvature
$\nu_M=D(T_Mu\,S)$ satisfies
\[
 \begin{split}
 \nu_M&=\mathbf1_{\{-M<u\le M\}}\mu+L_{-M}(u;S)-L_M(u;S),\\
 \|\nu_M\|_{\rm TV}&\le B,\qquad
 \|\nu_M-\mu\|_{\rm TV}
                    \le2|\mu|\{|u|\ge M\}\longrightarrow0.
 \end{split}
 \tag{SD11}
\]
Both $\nu_M$ and every $L_t(u;S)$ are energy diffuse.
No finite-energy norm of $|\nu_M|$ is claimed.
\end{proposition}

\noindent\textbf{Status: \texttt{NEW\_PROVED}}.

\begin{proof}
Apply (SD7) to the smooth one-sided hinges.  Their positive
defects are
$\Phi_\epsilon''(u-t)\Gamma_S(u)$, supported in the scalar
strip $t<u<t+\epsilon$.  Formula (SD9), the nonnegative
unit-integral second derivative, and its bounded density
show that their masses are at most $B$.
For $t>0$ the sharper bound is $|\mu|\{u>t\}$; for
$t<0$ and sufficiently small $\epsilon$ it is
$|\mu|\{u<t+\epsilon\}$.
They converge on original compact tests to $L_t(u;S)$.
Positivity, compact smooth exhaustion and then
$\epsilon\downarrow0$ prove (SD10).

The scalar identity $T_Ms=h_{-M}(s)-h_M(s)$ gives the
first formula of (SD11).  The three source regions
$\{u\le-M\}$, $\{-M<u\le M\}$ and $\{u>M\}$
partition the finite-value locus.  Triangle inequalities
and (SD10) therefore give both variation estimates in
(SD11).  Since $T_Mu\in H$, its finite Radon curvature
is diffuse by Lemma~\ref{lem:r11-bounded-H-tests-diffuse}.

For any fixed $t$ and $M>|t|$, the exact scalar identity
\[
                    h_t(T_Mu)=h_t(u)-h_M(u)
 \tag{SD12}
\]
holds on the finite-value locus and hence in the coefficient-defined
$L^1$ sense.  Its left side belongs to $H$.  Its curvature is finite Radon
by (SD10) and the Kato decompositions, hence diffuse by
the same lemma.  Thus
\[
 L_t(u;S)-L_M(u;S)
       =D(h_t(T_Mu)S)-\mathbf1_{\{t<u\le M\}}\mu
\]
is a diffuse finite signed measure.  On any polar Borel
set, the mass of $L_t$ is therefore at most $L_M(F)$,
which tends to zero by (SD10).  This proves diffuseness
of $L_t$ without assuming it passes through an arbitrary
weak measure limit.
\end{proof}

\begin{lemma}[Finite signed Radon data in the $H$ level theorem]
\label{lem:r11-H-level-carrier-finite-signed}
Let $v\in H$ and suppose $D(vS)=\sigma$ is a finite signed
Radon measure on the original tests.  Then the conclusion
(RT2)--(RT3) holds, with a finite diffuse positive level
measure, without assuming finite energy integral of
$|\sigma|$.
\end{lemma}

\noindent\textbf{Status: \texttt{NEW\_PROVED}} as the required
signed-data strengthening of the earlier $H$ carrier lemma.

\begin{proof}
The datum is diffuse and the current equation extends to
all bounded $H$ tests by Lemma
\ref{lem:r11-bounded-H-tests-diffuse}.
Consequently the finite-energy scalar chain proof applies:
its test $\psi\Phi'(v)$ is bounded and lies in $H$.
The one-sided band tests are also bounded $H$ vectors and
give (SD9) with $v,\sigma$ in place of $u,\mu$.
Thus the smoothed hinge defects have uniformly bounded
total mass, and their positive limit $L_t(v;S)$ is finite.
The signed measure
$D(h_t(v)S)=\mathbf1_{\{v>t\}}\sigma+L_t(v;S)$ is finite
and comes from an $H$ vector.  It is therefore diffuse by
the preceding lemma, as is $L_t$.

The original current identity for this hinge source now
extends to every bounded $H$ test.  Take
$z=\psi g(v)$ with $\psi$ an original compact nonnegative
smooth function and $g$ smooth bounded, with bounded
derivative, nonnegative and zero near $t$.
Strong convergence of the smooth hinges in $H$ and bounded
convergence against the finite measure $|\sigma|$ pass
\[
 \int z\,dL_{t,\epsilon}
     =-\mathcal E(h_{t,\epsilon}(v),z)
              -\int z\Phi_\epsilon'(v-t)\,d\sigma
\]
to $\int z\,dL_t$.
Its left side is zero for small $\epsilon$ because
$g\Phi_\epsilon''=0$.  The same countable family of
nonnegative tests as in the proof of (RT3) shows that
$L_t$ is carried by $\{\widetilde v=t\}$.
\end{proof}

\begin{theorem}[Exact level carriers and coarea for all signed diffuse data]
\label{thm:r11-signed-diffuse-exact-level}
For every finite signed energy-diffuse datum and every real $t$,
\[
                    \mathbf1_{\{\widetilde u\ne t\}}L_t(u;S)=0.
 \tag{SD13}
\]
For every bounded compactly supported Borel scalar $\theta$,
\[
                \int\theta(t)L_t(u;S)\,dt
                         =\theta(\widetilde u)\Gamma_S(u)
 \tag{SD14}
\]
as finite signed Radon measures.  The hard-curvature
variation bounds (SD11) hold on the complete original
space and every actual moment refinement.
\end{theorem}

\noindent\textbf{Status: \texttt{NEW\_PROVED}}.

\begin{proof}
First identify $\nu_M=\mu$ on $A_M=\{|u|<M\}$.
Choose a smooth $\Phi$ with $\Phi'$ zero outside $[-M,M]$
and strictly positive on $(-M,M)$.
The identities $\Phi(T_Mu)=\Phi(u)$ and (RT9) make their
weighted energies agree.  Compare (SD7) with the bounded
$H$ chain for $T_Mu$ and its finite signed datum $\nu_M$.
The latter chain is valid by the bounded-test extension,
so $\Phi'(u)(\nu_M-\mu)=0$, proving the assertion on $A_M$.

Fix $t$, choose $M>|t|$, and apply the preceding $H$ level
lemma to $v=T_Mu$.  The level measure $L_t(v;S)$ is carried
by $\{v=t\}=\{u=t\}$.
Compare the curvature of the two sides of (SD12), and then
restrict their actual Radon measures to $A_M$.
The just proved identity $\nu_M=\mu$ on $A_M$ cancels the
source terms and gives
\[
     \mathbf1_{A_M}(L_t(u;S)-L_M(u;S))
                         =\mathbf1_{A_M}L_t(T_Mu;S).
\]
It follows that
$L_t(u;S)(A_M\cap\{u\ne t\})\le L_M(u;S)(F)\to0$.
The sets $A_M$ exhaust the finite-value locus, and $L_t$ is
diffuse, so it gives no mass to the omitted polar infinite-
value set.  This proves (SD13) by monotone measure convergence.

For smooth compactly supported $\theta$, take
$\Phi(s)=\int\theta(t)(s-t)^+dt$.
It has bounded first derivative and compactly supported
second derivative $\theta$, so (SD7) and Fubini prove
(SD14) on every original test.
The positive kernels $L_t$, their global bound by $B$,
and their Borel original-test pairings extend the equality
to nonnegative Borel $\theta$; subtract positive and negative
parts.  The right side is finite by (SD9).
No unbounded unweighted energy current or restricted
background closedness has been assumed.
\end{proof}

\begin{definition}[The signed renormalized zero-boundary class]
\label{def:r11-signed-renormalized-boundary}
For the specified finite signed diffuse datum, a candidate is
$v\in L^1(F,\lambda)$ with a finite energy-quasicontinuous
representative and $T_kv\in H$ for every $k>0$.
Use the compatible sigma-finite energy measure of its
clippings.  Require the original current identity
\[
 D(\Phi(v)S)=\Phi'(\widetilde v)\mu
                     +\Phi''(\widetilde v)\Gamma_S(v)
 \tag{SD15}
\]
for every $\Phi\in C^2(\mathbf R)$ with compactly supported
first derivative, and the radiation condition
$k^{-1}\mathcal E(T_{2k}v-T_kv)\to0$.
These conditions specify the boundary and solution class
without using the transposition identity in their definition.
\end{definition}

\begin{theorem}[The complete finite-diffuse relative renormalized Dirichlet problem]
\label{thm:r11-signed-renormalized-Dirichlet}
Every finite signed energy-diffuse Radon datum has exactly one
candidate in Definition~\ref{def:r11-signed-renormalized-boundary}.
It is the transposition solution $u_\mu$, and all its original
current, exact-level, coarea and global curvature-variation
conclusions above hold.  The construction and solution class
are independent of the chosen faithful reference measure.
\end{theorem}

\noindent\textbf{Status: \texttt{NEW\_PROVED}} for the
printed relative linear operator and boundary class, in
arbitrary ambient dimension and arbitrary original $S$ of
the specified degree.

\begin{proof}
Existence follows from (SD7), the truncation membership,
quasicontinuity, and the already proved signed radiation
estimate (KT11).  For uniqueness choose smooth odd scalar
clamps $\Phi_k$ fixing the interval $[-k,k]$, constant
outside $[-2k,2k]$, with $\Phi_k(0)=0$,
$0\le\Phi_k'\le1$ and
$|\Phi_k''|\le C/k$ on the two transition bands.
Then $|\Phi_k(s)|\le|s|$ and $\Phi_k(s)\to s$.
For a candidate $v$, its prescribed chain gives
\[
 \begin{split}
 \nu_k:=D(\Phi_k(v)S)
       &=\Phi_k'(v)\mu+\Phi_k''(v)\Gamma_S(v),\\
 \|\nu_k-\mu\|_{\rm TV}
       &\le |\mu|\{|v|>k\}
                +\frac Ck\mathcal E(T_{2k}v-T_kv)\to0.
 \end{split}
 \tag{SD16}
\]
The measures $\nu_k$ are finite signed Radon measures.
The potential $\Phi_k(v)=\Phi_k(T_{2k}v)$ belongs to $H$,
so Lemma~\ref{lem:r11-bounded-H-tests-diffuse} proves
both diffuseness of $\nu_k$ and
$\Phi_k(v)=u_{\nu_k}$.  This step does not infer a
finite-energy norm for $|\nu_k|$ from its signed core norm.
The established total-variation stability gives
$\|\Phi_k(v)-u_\mu\|_1\le2\|\nu_k-\mu\|_{\rm TV}\to0$.
Dominated convergence also gives $\Phi_k(v)\to v$ in
$L^1(\lambda)$.  Thus $v=u_\mu$ almost everywhere.
Equality of all clipped $H$ vectors identifies the canonical
representatives quasi-everywhere.  Reference independence
follows from the already proved independence of the original
core form and the transposition solution.
\end{proof}

\paragraph{The remaining general problem.}
This section closes the signed diffuse chain and level-carrier
gates left open at the preceding positive-data stage.  It does
not establish ambient classical psh representability, ordinary
pointwise boundary regularity, or an unrestricted diagonal nonlinear Monge--Ampere
Dirichlet theory.  General analytic PS/PL relation descent and
low-real-degree CTO remain \texttt{OPEN}.  The background in
all the proved relative statements remains arbitrary.
Unrestricted equality of energy and BT polar ideals is already
\texttt{FALSE} by frozen R6, and unrestricted ambient classical
solvability is already \texttt{FALSE} by frozen R10; the larger
relative solution class is consistent with both obstructions.

% END REVIEWED MODULE research_r11/dirichlet/signed_diffuse_renormalized_dirichlet.tex


% BEGIN REVIEWED MODULE research_r11/dirichlet/signed_data_strong_clipping_stability.tex
\section{Strong clipped energy under total-variation convergence of signed data}
\label{sec:r11-signed-data-strong-clipping}

Retain the original arbitrary degree-$(n-1,n-1)$ background,
the affinoid neighborhood and the ordinary relatively compact
part domain of the signed diffuse Dirichlet theorem.  All
measures below are finite signed energy-diffuse measures on
$F$, and all potentials use their canonical energy-quasicontinuous
representatives.  The norm of $H$ is $\mathcal E^{1/2}$.

\begin{theorem}[Strong clipping and smooth-curvature stability]
\label{thm:r11-signed-data-strong-clipping}
Suppose $\|\mu_j-\mu\|_{\rm TV}\to0$, and write
$u_j=u_{\mu_j}$, $u=u_\mu$ for the unique intrinsic
renormalized solutions of $D(u_jS)=\mu_j$, $D(uS)=\mu$.
Then, for every fixed $k>0$,
\[
                    \|T_ku_j-T_ku\|_H\longrightarrow0.
 \tag{DS1}
\]
In particular, for any fixed $k,l>0$, the finite original
energy measures converge in global total variation:
\[
 \Gamma_S(T_ku_j,T_lu_j)
               \longrightarrow\Gamma_S(T_ku,T_lu).
 \tag{DS2}
\]
If $\Phi\in C^2(\mathbf R)$ has bounded first derivative
and compactly supported second derivative, then
\[
           D(\Phi(u_j)S)\longrightarrow D(\Phi(u)S)
                       \quad\hbox{in global total variation}.
 \tag{DS3}
\]
If, more restrictively, $\Phi'$ has compact support, then
$\Phi(u_j)-\Phi(0)\to\Phi(u)-\Phi(0)$ strongly in $H$.
For this latter class both original first-order currents
$d'(\Phi(u_j)S)$ and $d''(\Phi(u_j)S)$ converge in local
coefficient total variation, against every original compact
complementary smooth test and every actual moment refinement.
The zero-order multipliers $\Phi(u_j)S$ converge in that same
local coefficient sense under the hypotheses of (DS3).
\end{theorem}

\noindent\textbf{Status: \texttt{DIRECT\_COROLLARY}} of
the completed signed renormalized theorem, especially the exact
band-energy identity (SD8), and the proved finite-datum capacity
stability.  This conclusion is not used to prove (SD8).

\begin{proof}
The finite-datum estimates give $u_j\to u$ in energy capacity
and in $L^1(\lambda)$.  The clipped functions converge in
$L^2(\lambda)$: they are bounded by $k$, and $\lambda(F)<\infty$
because the part domain is relatively compact in the neighborhood
where $\lambda$ is Radon.  Also
$\mathcal E(T_ku_j)\leq k\|\mu_j\|_{\rm TV}$, so these
vectors are bounded in $H$.  Every weakly convergent subsequence
therefore has the unique $L^2$ representative $T_ku$.
Consequently the whole sequence converges weakly in $H$.

Take $a=0,b=k$ in (SD8).  The exact norm identity is
\[
 \mathcal E(T_ku_j)=-\int_F T_k\widetilde u_j\,d\mu_j,
 \qquad
 \mathcal E(T_ku)=-\int_F T_k\widetilde u\,d\mu.
 \tag{DS4}
\]
The difference between the first datum integral and the second
is bounded in absolute value by
\[
 k\|\mu_j-\mu\|_{\rm TV}
       +\int_F|T_k\widetilde u_j-T_k\widetilde u|\,d|\mu|.
 \tag{DS5}
\]
The second term tends to zero by energy-capacity convergence,
finite diffuseness of $|\mu|$, and boundedness of the clippings.
Thus the norms in (DS4) converge.  Weak convergence and convergence
of norms in the Hilbert space prove (DS1).  In particular the
argument does not infer strong clipping merely from a uniform
energy bound and an $L^1$ limit; its decisive input is the exact
identity already proved in the signed renormalized theorem.

The original closed gradient module gives, for $f_j,g_j,f,g\in H$,
\[
 \|\Gamma_S(f_j,g_j)-\Gamma_S(f,g)\|_{\rm TV}
  \leq \|f_j-f\|_H\|g_j\|_H
                         +\|f\|_H\|g_j-g\|_H.
 \tag{DS6}
\]
This is measure Cauchy--Schwarz after adding and subtracting
$\Gamma_S(f,g_j)$.  Applying (DS1) proves (DS2).

For (DS3), choose $M$ larger than the support radius of $\Phi''$.
The all-original chain identity reads
\[
 D(\Phi(u_j)S)=\Phi'(\widetilde u_j)\mu_j
       +\Phi''(\widetilde u_j)\Gamma_S(T_Mu_j).
 \tag{DS7}
\]
The first term converges in variation: its changing-datum error
is at most $\|\Phi'\|_\infty\|\mu_j-\mu\|_{\rm TV}$,
and its fixed-datum multiplier error tends to zero by finite
diffuseness and capacity convergence.  For the second term use
(DS2) and then bounded convergence in measure against the fixed
finite diffuse measure $\Gamma_S(T_Mu)$.  This gives its variation
convergence as well and proves (DS3).  The zero-order multiplier
claim follows from the Lipschitz bound for $\Phi$, $L^1(\lambda)$
stability and the faithful original coefficient domination.

If $\Phi'$ has compact support, its normalized composition is
a smooth function of one sufficiently large clipping.  The
closed gradient formula and (DS1) give strong convergence:
the changing-gradient term is bounded by
$\|\Phi'\|_\infty\|T_Mu_j-T_Mu\|_H$, and the
changing scalar factor tends to zero in the fixed finite
energy measure.  Its $L^2$ representative converges by the
same bounded composition argument.  Finally, on a fixed
original compact stage, write any complementary one-form
test as a finite sum of bounded smooth coefficients times
differentials of original local moment functions.  Multiply
those moment functions by original compact cutoffs equal to
one on the test support.  The coefficient pairing of a
first-order derivative with each such term is an original
mixed gradient measure; (DS6), or its one-variable
Cauchy--Schwarz bound, controls its variation by a fixed
constant times the $H$-norm of the changing potential.
This proves local coefficient-TV convergence of both first-order
currents.  The argument starts with an arbitrary original test
and is valid after any further actual moment refinement.
It is not a claim of one global variation norm for all
lower-degree coefficients simultaneously.
\end{proof}

\begin{proposition}[Hard clipping curvatures need not be TV continuous in the datum]
\label{prop:r11-hard-clipping-source-TV-failure}
For a fixed $k>0$, the implication
\[
 \mu_j\longrightarrow\mu\text{ in total variation}
 \quad\Longrightarrow\quad
 D(T_ku_{\mu_j}S)\longrightarrow D(T_ku_\mu S)
                   \text{ in total variation}
 \tag{DS8}
\]
is \texttt{FALSE}, even for positive smooth original data on an
original annulus.  This does not contradict (DS1)--(DS3) or
the global hard-curvature limit as $k\to\infty$ in (SD11).
\end{proposition}

\noindent\textbf{Status: \texttt{DIRECT\_COROLLARY}} of the
frozen all-original one-dimensional torus calibration, with the
following explicit computation.  The rejected inference (DS8)
is recorded as \texttt{FALSE}.

\begin{proof}
Let $V=\{x\in\mathbf G_{m,K}^{\rm an}:-1<t(x)<1\}$,
where $t=\log|T|$, and take $S=[V]$.  It extends to a larger
ordinary annulus inside one original affinoid.  Write
$\sigma:(-1,1)\to V$ for the original Gauss section.  Put
\[
 u=t^2-1,\qquad u_j=c_j u,\qquad c_j=1+j^{-1},
 \qquad \mu_j=2c_j\sigma_*(dt),\quad\mu=2\sigma_*(dt).
 \tag{DS9}
\]
These potentials belong to the part space $H$: smooth real
compact cutoffs approximate their zero-boundary real Sobolev
profiles in gradient energy, and the exact original calibration
passes that approximation to the full original form.  Their
all-original curvature measures are the displayed data.
They are therefore the canonical solutions by the finite-energy
or signed renormalized uniqueness theorem.  In particular these
finite data are energy diffuse and
$\|\mu_j-\mu\|_{\rm TV}=4(c_j-1)\to0$.

Fix $0<k<1$, and put
\[
                  a_j=\sqrt{1-k/c_j},\qquad a=\sqrt{1-k}.
\]
Since the potentials are nonpositive, $T_ku_j=\max(u_j,-k)$.
Its real slope is zero on $(-a_j,a_j)$ and $2c_jt$ outside.
The two positive slope jumps have size $2c_ja_j$.  The
original calibrated measure is consequently
\[
 D(T_ku_jS)
  =\sigma_*\left(2c_j\mathbf1_{\{|t|>a_j\}}\,dt
                +2c_ja_j(\delta_{-a_j}+\delta_{a_j})\right).
 \tag{DS10}
\]
The same formula with $c_j=1,a_j=a$ gives $D(T_kuS)$.
The Gauss section is injective, $a_j>a$, and all four atoms
in the difference are distinct.  The diffuse parts converge
in total variation, while the total variation of the atomic
difference tends to
\[
                         4c_ja_j+4a\longrightarrow8a>0.
 \tag{DS11}
\]
Thus the total variation in (DS8) does not tend to zero.
For an arbitrary prescribed $k>0$, replace $u$ by
$A(t^2-1)$ with $A>k$, and replace $u_j$ by $c_j u$.
The same calculation has $a_j=\sqrt{1-k/(Ac_j)}$ and
limiting atomic variation $8A\sqrt{1-k/A}>0$.
Each equality is an original Radon measure equality from the
frozen calibration, valid on every original scalar test and
every actual moment refinement.  Strong convergence of the
clipped gradients still holds by (DS1); it is the moving
hard-level curvature atoms, not their energy measures, that
prevent the stronger asserted continuity.
\end{proof}

\paragraph{Dependency boundary.}
The new stability statement is downstream of the signed
renormalized chain and exact band identity.  It does not close
their initial proof by a circular norm argument.  The smooth
scalar curvature limit (DS3) has its full derivative hypotheses;
the explicit failure (DS8) prevents replacing those hypotheses
by an arbitrary hard clamp.  Unweighted energy of an unbounded
solution may still fail to be locally Radon, and no ambient
classical potential or ordinary pointwise boundary value is
asserted.

% END REVIEWED MODULE research_r11/dirichlet/signed_data_strong_clipping_stability.tex


% BEGIN REVIEWED MODULE research_r11/dirichlet/energy_measure_radon_obstruction.tex
\section{A non-Radon energy measure and sharp global variation loss}
\label{sec:r11-energy-Radon-obstruction}

\begin{proposition}[The sigma-finite qualification is necessary]
\label{prop:r11-diffuse-energy-not-Radon}
For every $n\ge3$, within the original field and absolute value,
there are an ordinary torus domain, a smooth original positive
closed background $S$ of degree $(n-1,n-1)$, and a finite positive
compactly supported energy-diffuse datum $\mu$ whose unique
renormalized solution $u$ has the following properties.
Its energy measure $\Gamma_S(u)$ has infinite mass on every
ordinary compact neighborhood of a specified interior Gauss
point, and is therefore not an ordinary coefficient-Radon
current.  Its bounded hard-truncation curvatures nevertheless
satisfy the exact formula
\[
 \|D(\max(u,-k)S)-\mu\|_{\rm TV(F)}
                  =2\mu\{\widetilde u\le-k\}>0
                         \qquad(k>0).
 \tag{RO1}
\]
Thus the factor $2$ in (RT16) is sharp.
\end{proposition}

\noindent\textbf{Status: \texttt{NEW\_PROVED}} counterexample
and sharpness calculation.  The assertion that every finite
diffuse datum has an ordinary locally finite energy current
is \texttt{FALSE}.  The general Dirichlet theorem continues
to allow arbitrary original $S$; this special background is
used only to disprove that additional assertion.

\begin{proof}
Take original coordinates $T_1,\ldots,T_n$ on the split torus,
write $t_i=\log|T_i|$, $r=|t|$, and let
\[
 V=\{r<R\},\qquad
 \beta=\sum_i d't_i\wedge d''t_i,\qquad
 S=\frac{\beta^{n-1}}{(n-1)!}[V].
 \tag{RO2}
\]
The same smooth background is first defined on a larger ordinary
polyannulus $\Omega$, and then restricted to $V$.  Choose it so
that $V\Subset\Omega$ and $\overline\Omega$ lies in the interior
of a still larger compact polyannulus $B$.  Thus the affinoid
domain hypothesis holds exactly as required for the part form.
Let $\sigma$ be the original joint Gauss section.
Original torus calibration gives
\[
 D(g(t)S)=\sigma_*((\Delta g)\,dt),\qquad
 \Gamma_S(g(t))=\sigma_*(|\nabla g|^2\,dt)
 \tag{RO3}
\]
for original smooth scalar $g(t)$, on all original compact
scalar tests, including their actual refinements.
These are identities of original top-degree measures; they
are not assertions obtained by testing only radial functions.

Put $a=(n-2)/2>0$, so $\alpha=(n+2)/2=a+2<n$.
Choose a nonnegative smooth radial cutoff $\chi$ equal to
one for $0\le r\le r_0$, and zero for $r\ge r_1<R$.
Set
\[
 f(r)=\chi(r)r^{-\alpha},\quad
 m(r)=\int_0^r \tau^{n-1}f(\tau)\,d\tau,\quad
 U(r)=-\int_r^R s^{1-n}m(s)\,ds,\quad
 \mu=\sigma_*(f(|t|)\,dt).
 \tag{RO4}
\]
The datum is finite since $\alpha<n$, and compactly supported.
It is energy diffuse: $\mu\ll\sigma_*dt$, while the original
smooth coefficient $\beta\wedge S=n\sigma_*dt$ is locally
dominated by the faithful $\lambda$; energy-polar sets are
$\lambda$-null.  The radial identity
$U''+(n-1)U'/r=f$ follows by differentiation.
The function $U$ is strictly increasing, is zero at $R$, and
tends to $-\infty$ at zero.

For each $k>0$ let $r_k$ be the unique solution of $U(r_k)=-k$,
and let $U_k=\max(U,-k)$.  It is constant on $r\le r_k$,
is locally Lipschitz on the real ball, and is smooth near
its outer boundary, where $U_k=O(R-r)$ and its gradient is
bounded.  Smoothing the radial corner and multiplying by
smooth cutoffs supported in $r<R$ approximates $U_k(t)$ by
original compact smooth core functions in $H$.
Indeed the boundary collar energy error is $O(\varepsilon)$:
$|U_k|\le C\varepsilon$ and the cutoff gradient is
$O(\varepsilon^{-1})$ on a collar of volume $O(\varepsilon)$.
The remaining energy error tends to zero by ordinary smooth
approximation of the bounded radial Lipschitz profile.
The $L^2(\lambda)$ error tends to zero because $U_k$ is bounded,
the scalar smoothing can be uniform, and $\lambda(V)<\infty$.
Consequently $U_k(t)\in H$ with the printed zero energy boundary.

Real radial integration by parts and (RO3) give
\[
 \begin{split}
 \mu_k:=D(U_k(t)S)
 =\sigma_*\bigl(&\mathbf1_{\{r>r_k\}}f(r)\,dt\\
          &+U'(r_k)\,\mathcal H^{n-1}|_{\{r=r_k\}}\bigr)\ge0.
 \end{split}
 \tag{RO5}
\]
To justify the same formula on arbitrary original tests, smooth
the radial corner uniformly on each compact stage.  Its
Laplacian measures converge to the displayed real measure.
The Gauss section is continuous, so every original smooth test
restricted to that section is a legitimate continuous test for
this convergence.  Uniform convergence of the scalar potentials
also passes $D(U_kS)$ in the original coefficient calculus.
Thus (RO5) has the full original-current meaning.

Writing $\omega_{n-1}$ for the area of the unit sphere,
\[
 U'(r_k)\omega_{n-1}r_k^{n-1}
       =\omega_{n-1}m(r_k)=\mu\{r<r_k\}.
 \tag{RO6}
\]
Therefore $\mu_k(F)=\mu(F)$, and the two nonzero parts of
$\mu_k-\mu$ lie on the disjoint sphere and its interior.
This proves the exact variation formula (RO1), because
$\{U\le-k\}=\{r\le r_k\}$ and $\mu$ gives no mass to the sphere.
The right side tends to zero as $k\to\infty$.
Each $U_k(t)$ is the finite-energy solution for $\mu_k$ by
(RO5), so the proved $L^1$ transposition stability identifies
the pointwise radial limit with $u_\mu$.  This also proves
$U(t)\in L^1(\lambda)$ rather than assuming it in the
unbounded coefficient product.  Each $U_k(t)$ is its continuous
canonical $H$ representative; equality of the clipped $H$ vectors
therefore identifies the radial limit with the canonical solution
outside one energy-polar set.

For $r<r_0$ one has $m(r)=r^a/a$, and hence
\[
            U(r)=-\frac{1}{a^2}r^{-a}+C,
                   \qquad U'(r)=\frac1a r^{-a-1}.
 \tag{RO7}
\]
The compatible energy measures of the hard truncations give
\[
 \Gamma_S(u)=\sigma_*(U'(r)^2\,dt\text{ on }r>0).
\]
For every $0<\varepsilon<r_0$ its mass over the compact
ordinary set $\{r\le\varepsilon\}\Subset V$ is
\[
        \frac{\omega_{n-1}}{a^2}
                    \int_0^\varepsilon\frac{dr}{r}=\infty.
 \tag{RO8}
\]
Every ordinary neighborhood of the interior point $\sigma(0)$
contains $\sigma(t)$ for all sufficiently small $|t|$, by
continuity of the Gauss section.  Thus the measure is not
locally finite at that actual original point.  It is precisely
the sigma-finite energy in (RT7), not an additional Radon
current.  The original current equation $D(uS)=\mu$ and all
bounded-value weighted energy identities remain valid.
\end{proof}

% END REVIEWED MODULE research_r11/dirichlet/energy_measure_radon_obstruction.tex


% BEGIN REVIEWED MODULE research_r11/dirichlet/relative_energy_obstacle_contact.tex
\section{Relative energy obstacles and an exact contact equation}
\label{sec:r11-energy-obstacles}

Retain the original arbitrary background $S$ of degree $(n-1,n-1)$,
the part domain $V$, and its faithful energy Hilbert space $H$.
All inequalities between energy functions mean inequalities of
their canonical energy-quasicontinuous representatives outside
an energy-polar set.  By the proved energy-to-BT-capacity bridge,
this convention also gives a canonical relative BT-capacitary class.
It does not impose ambient classical psh membership.

Let $\sigma$ be a signed Radon measure of finite energy on $F$.
Its core integral has the continuous extension $\ell_\sigma\in H^*$.
For a prescribed obstacle $\psi\in H$, put
\[
 K_\psi=\{v\in H:\widetilde v\geq\widetilde\psi
                         \text{ energy-quasi-everywhere}\},
 \qquad
 J_\sigma(v)=\tfrac12\mathcal E(v,v)-\ell_\sigma(v).
 \tag{EO1}
\]
The zero-boundary condition is membership in $H$.

\begin{theorem}[The original relative obstacle reaction]
\label{thm:r11-relative-energy-obstacle}
There is a unique minimizer $w=w_{\sigma,\psi}$ of $J_\sigma$
on $K_\psi$.  It satisfies
\[
       \mathcal E(w,v-w)\geq\ell_\sigma(v-w)
                                  \qquad(v\in K_\psi).
 \tag{EO2}
\]
There is a positive Radon measure $\rho$ on $F$, of finite
energy integral, such that
\[
             -D(wS)=\sigma+\rho
                 \quad\hbox{on every original compact smooth test in }V,
 \qquad
             \int_F(\widetilde w-\widetilde\psi)\,d\rho=0.
 \tag{EO3}
\]
In particular
\[
                  \rho\{\widetilde w>\widetilde\psi\}=0 .
 \tag{EO4}
\]
The reaction is characterized by
$\ell_\rho(v)=\mathcal E(w,v)-\ell_\sigma(v)$ for every $v\in H$.
Its total mass near the ordinary boundary need not be finite.

The minimizer is the least element of the following class:
\[
 \{v\in H:\widetilde v\geq\widetilde\psi,\
               \mathcal E(v,\phi)\geq\ell_\sigma(\phi)
                              \text{ for every nonnegative }\phi\in H\}.
 \tag{EO5}
\]
Thus the displayed variational problem supplies a true relative
superharmonic envelope and its exact reaction, with no contact
formula assumed in advance.
\end{theorem}

\noindent\textbf{Status: \texttt{NEW\_PROVED}} as an
all-original-current obstacle/contact theorem for the relative
energy class on general affinoid neighborhoods.
Hilbert projection is \texttt{EXTERNAL\_KNOWN};
the original current identification and the precise solution
class are part of the proof below.

\begin{proof}
The cone $H_+=\{v\in H:v\geq0\ \lambda\text{-almost everywhere}\}$
is closed and convex, by coercivity and the $L^2$ realization.
For energy vectors this is equivalent to nonnegativity of the
canonical representative quasi-everywhere: the negative part is
the zero vector and has zero canonical representative.
Consequently $K_\psi=\psi+H_+$ is a nonempty closed convex set.
The Riesz vector $y=G\sigma$ satisfies
\[
          J_\sigma(v)=\tfrac12\|v-y\|_H^2-\tfrac12\|y\|_H^2.
 \tag{EO6}
\]
Hilbert projection gives a unique minimizer, and the elementary
projection inequality is exactly (EO2).

For every nonnegative $\phi\in H$ the vector $w+\phi$ belongs
to $K_\psi$, so
\[
                  \mathcal E(w,\phi)-\ell_\sigma(\phi)\geq0.
 \tag{EO7}
\]
On original compact smooth scalars this positive functional is
a positive Radon measure.  Here is a direct original-test
justification.  On a fixed compact core choose an original
smooth nonnegative cutoff $\chi$ equal to one near that core.
If an original smooth scalar $\phi$ is supported there, then
$\|\phi\|_\infty\chi\pm\phi\geq0$; (EO7) therefore bounds its
functional value by a fixed multiple of $\|\phi\|_\infty$.
The original scalar algebra is uniformly dense in continuous
compact functions by the exact affinoid evaluation argument.
The positive functional extends uniquely to those functions
and hence is represented by a Radon measure $\rho$.
It is supported on $F$, since both the energy and the finite-energy
datum vanish on core functions supported off $F$.

The functional has energy norm at most
$\|w\|_H+\|\sigma\|_{-1,V}$.  Applying it to smooth
approximations of absolute values shows that its positive
Radon measure has finite energy integral.  The polar-null and
canonical integration argument in the proof of
Lemma~\ref{lem:r11-finite-energy-mass-estimates} gives
\[
       \ell_\rho(v)=\int_F\widetilde v\,d\rho
                              \qquad(v\in H),
 \tag{EO8}
\]
including absolute integrability.
Equation (AE15) now identifies (EO7) as the original current
equation $-D(wS)=\sigma+\rho$ on every original scalar test.
No selected moment family has replaced that equation.

Use the admissible competitor $v=\psi$ in (EO2).  It gives
$\ell_\rho(w-\psi)\leq0$.  Since $w-\psi\in H_+$,
positivity gives the reverse inequality.  Equation (EO8)
therefore proves the second identity in (EO3).
The integrand is nonnegative and the measure charges no polar
sets, so (EO4) follows.

Finally let $v$ belong to (EO5), and put $\eta=(w-v)^+\in H$.
The vector $w-\eta=w\wedge v$ still lies above $\psi$,
hence is an admissible competitor.  Formula (EO2) gives
$\mathcal E(w,\eta)\leq\ell_\sigma(\eta)$, whereas (EO5)
gives $\mathcal E(v,\eta)\geq\ell_\sigma(\eta)$.
Thus $\mathcal E(w-v,\eta)\leq0$.
The normal-contraction inequality gives
$\mathcal E(w-v,(w-v)^+)\geq\mathcal E(\eta,\eta)$,
so coercivity implies $\eta=0$.  This proves $w\leq v$
quasi-everywhere and the least-envelope characterization.
\end{proof}

\begin{corollary}[An exact relative energy extremal]
\label{cor:r11-relative-energy-extremal}
For every $a\in H$ there is a greatest element $P_Sa$ of
\[
 \mathscr Q_a=
 \{u\in H:\widetilde u\leq\widetilde a,\
                 D(uS)\geq0\text{ on every original compact smooth scalar}\}.
 \tag{EO9}
\]
It is $P_Sa=-w_{0,-a}$ and satisfies
\[
 \begin{gathered}
       D((P_Sa)S)=\rho_a\geq0,\qquad
       \rho_a\{\widetilde{P_Sa}<\widetilde a\}=0,\\
                    \int_F(\widetilde a-\widetilde{P_Sa})\,d\rho_a=0.
 \end{gathered}
 \tag{EO10}
\]
It is canonical modulo energy-polar sets and hence determines
the canonical relative BT-capacitary class.
All identities are on the complete original test sheaf and
all actual moment refinements.
\end{corollary}

\noindent\textbf{Status: \texttt{DIRECT\_COROLLARY}} of the
preceding theorem.  This is an energy extremal, with its
printed zero-energy-boundary condition.

\begin{proof}
With $\sigma=0$ and $\psi=-a$, put $u=-w$.
The theorem gives $u\leq a$ and $D(uS)=\rho_a\geq0$.
If $v\in\mathscr Q_a$, then $-v\geq-a$ and its core
superharmonic inequality extends to nonnegative vectors of
$H$ by the positive approximation argument of R6.
Thus $-v$ belongs to (EO5), and leastness gives $v\leq u$.
The contact and orthogonality statements are (EO3)--(EO4)
with the signs reversed.
\end{proof}

\begin{proposition}[Order and stable dependence of the relative envelope]
\label{prop:r11-energy-obstacle-stability}
If $\psi\leq\psi'$ and $\ell_\sigma\leq\ell_{\sigma'}$ on $H_+$,
then
$w_{\sigma,\psi}\leq w_{\sigma',\psi'}$ quasi-everywhere.
For arbitrary finite-energy data and obstacles,
\[
 \|w_{\sigma,\psi}-w_{\sigma',\psi'}\|_H
          \leq\|\sigma-\sigma'\|_{-1,V}+\|\psi-\psi'\|_H.
 \tag{EO11}
\]
In particular $a\mapsto P_Sa$ is order preserving and
nonexpansive for the energy norm.  The local original
BT-capacity estimate is
\[
 \operatorname{Cap}_S^*
   (K_0\cap\{|\widetilde{P_Sa}-\widetilde{P_Sb}|>\delta\};\Omega)
       \leq\sqrt3\,C_{K_0}\delta^{-1}\|a-b\|_H
                         \qquad(K_0\Subset V).
 \tag{EO12}
\]
\end{proposition}

\noindent\textbf{Status: \texttt{DIRECT\_COROLLARY}} of the
proved envelope characterization, Hilbert projection and
the frozen one-way capacity bridge.

\begin{proof}
The second obstacle solution is a superharmonic majorant
for the first problem, so leastness proves order.
Let $P$ denote orthogonal projection onto $H_+$.
Formula (EO6) gives
\[
                     w_{\sigma,\psi}
                       =\psi+P(G\sigma-\psi).
 \tag{EO13}
\]
Both $P$ and $I-P$ are nonexpansive, as follows by adding
their two projection inequalities.
At fixed $\psi$, (EO13) is nonexpansive in $G\sigma$.
At fixed $\sigma$, write it as
$G\sigma-(I-P)(G\sigma-\psi)$ to obtain nonexpansiveness
in $\psi$.  The triangle inequality proves (EO11).
The extremal statements follow by changing signs.
Finally $c_V\{|P_Sa-P_Sb|>\delta\}\leq
\delta^{-2}\|P_Sa-P_Sb\|_H^2$.
Apply (MD2), the part/full capacity comparison and (MD14)
to get (EO12).
\end{proof}

\paragraph{Relation to the ambient classical obstruction.}
The contact assertion here concerns the energy class (EO9).
It cannot be inserted into an envelope taken over ambient
classical psh functions.  R10 exhibits such ambient envelopes
with nonzero relative mass outside their contact set.
Likewise the assertion that an energy obstacle reaction is
finite-energy is not an identification of energy-polar and
relative BT-null ideals.  The current $S$ remains arbitrary
throughout; only the solution class and boundary condition
have been specified precisely.
The same results apply to every original fixed bounded BT
mixed tail of an arbitrary allowed background.  Several
unfrozen nonlinear slots and unrestricted ambient
Monge--Ampere Dirichlet existence remain open.

% END REVIEWED MODULE research_r11/dirichlet/relative_energy_obstacle_contact.tex


% BEGIN REVIEWED MODULE research_r11/dirichlet/diffuse_obstacle_contact.tex
\section{Finite signed diffuse data and an exact relative obstacle problem}
\label{sec:r11-diffuse-obstacle-contact}

Keep the original arbitrary symmetric middle-positive separately
closed $(n-1,n-1)$-current $S$, the analytic affinoid neighborhood
$\Omega$, and the ordinary part domain $V\Subset\Omega$ from the
signed diffuse Dirichlet theorem.  Write $(H,\mathcal E,c_V,G,\lambda)$
for its faithful part form and put
\[
                              Az=-D(zS).
 \tag{OB1}
\]
Thus $G\sigma=-u_\sigma$ for a finite signed energy-diffuse measure
$\sigma$, and $A(G\sigma)=\sigma$ on every original test.
All functions below use their canonical energy-quasicontinuous
representatives on $F$.  Inequalities mean quasi-everywhere on $F$;
no pointwise boundary condition or ambient classical psh condition
is substituted for the specified part-form boundary class.

\begin{lemma}[Positive parts are Green potentials]
\label{lem:r11-diffuse-positive-part-Green}
Let $\sigma$ be any finite signed energy-diffuse Radon measure and
let $z=G\sigma$.  Then $z^+$ is the Green potential of its own finite
signed diffuse curvature and
\[
 \begin{split}
 Az^+&=\mathbf1_{\{\widetilde z>0\}}\sigma-L_0(z;S),\\
 z^+&=G(Az^+),\qquad L_0(z;S)\ge0.
 \end{split}
 \tag{OB2}
\]
In particular,
\[
        \mathbf1_{\{\widetilde z>0\}}\sigma\le0
                              \quad\Longrightarrow\quad z\le0.
 \tag{OB3}
\]
These assertions do not require $z$ or $z^+$ to belong to $H$.
\end{lemma}

\noindent\textbf{Status: \texttt{DIRECT\_COROLLARY}} of the full
signed exact-level and global curvature-variation theorems, with the
Green representation proved explicitly below.

\begin{proof}
The signed Kato theorem gives the first line of (OB2), and its level
measure is finite and diffuse by (SD10)--(SD13).  For $M>0$ put
$v_M=(T_Mz)^+\in H$.  The actual scalar identity
$v_M=z^+-(z-M)^+$ and the signed level formula give
\[
 D(v_MS)=\mathbf1_{\{0<\widetilde z\le M\}}D(zS)
                                    +L_0(z;S)-L_M(z;S).
 \tag{OB4}
\]
Consequently
\[
 \|D(v_MS)-D(z^+S)\|_{\rm TV}
                   \le2|\sigma|\{\widetilde z>M\}\longrightarrow0.
\]
Each $v_M$ is an $H$ vector with finite signed Radon curvature, so
Lemma~\ref{lem:r11-bounded-H-tests-diffuse} identifies it with the
transposition solution of that curvature.  Total-variation to
$L^1(\lambda)$ stability then gives
$v_M\to u_{D(z^+S)}$.  Also $v_M\to z^+$ in $L^1(\lambda)$ by
dominated convergence, since $|v_M|\le|z|$.  This proves the second
line of (OB2), including the canonical representative: its bounded
clippings agree as $H$ vectors and hence quasi-everywhere.

Under the hypothesis in (OB3), (OB2) gives $Az^+\le0$.
Order preservation of $G$ therefore gives $z^+=G(Az^+)\le0$.
Thus $z\le0$.  No test of an unbounded vector as if it belonged
to $H$ occurs in this proof.
\end{proof}

\begin{definition}[Obstacles with finite curvature and the diffuse solution class]
\label{def:r11-diffuse-obstacle-class}
An allowed obstacle is a vector $\psi\in H$ such that
$\tau:=A\psi$ is a finite signed Radon measure and $|\tau|$ has
finite energy integral.  In particular every real original compact
smooth function is an allowed obstacle: its curvature has a bounded
density with respect to a faithful reference measure, by the
original compact coefficient control.

Fix a finite signed energy-diffuse datum $\mu$.  A solution pair
$(w,\rho)$ consists of a finite positive energy-diffuse Radon measure
$\rho$ and a signed renormalized zero-boundary solution $w$ of
\[
                               Aw=\mu+\rho
 \tag{OB5}
\]
such that
\[
                     w\ge\psi,\qquad
                     \rho\{\widetilde w>\widetilde\psi\}=0.
 \tag{OB6}
\]
The renormalized condition in (OB5) means the intrinsic class from
Definition~\ref{def:r11-signed-renormalized-boundary}, with sign
reversed: $w\in L^1(\lambda)$ is finite and quasicontinuous,
$T_kw\in H$ for every $k>0$, the original chain identity is
\[
 D(\Phi(w)S)=-\Phi'(\widetilde w)(\mu+\rho)
                         +\Phi''(\widetilde w)\Gamma_S(w)
 \tag{OB7}
\]
for $C^2$ scalar $\Phi$ with compactly supported first derivative,
and $k^{-1}\mathcal E(T_{2k}w-T_kw)\to0$.
Thus the definition does not prescribe $w$ by a Green formula.
\end{definition}

\begin{theorem}[The full finite-diffuse obstacle and contact theorem]
\label{thm:r11-diffuse-obstacle-contact}
For every allowed obstacle $\psi$ and every finite signed
energy-diffuse datum $\mu$, there is exactly one solution pair in
Definition~\ref{def:r11-diffuse-obstacle-class}.  It satisfies
\[
 \begin{split}
 w&=G(\mu+\rho),\qquad Aw=\mu+\rho,\\
 0\le\rho&\le(\tau-\mu)^+\le\tau^++\mu^-,\qquad \tau=A\psi.
 \end{split}
 \tag{OB8}
\]
The second line is the actual measure inequality, including the
contact set and every original test; no density of the arbitrary
background is assumed.  In particular the reaction is finite.

The solution is the least member of the following explicitly
nonempty class:
\[
 \mathscr M_{\mu,\psi}=
 \{G(\mu+\eta):\ \eta\ge0\text{ finite and energy diffuse},\
                         G(\mu+\eta)\ge\psi\}.
 \tag{OB9}
\]
It is therefore an actual relative envelope, with a finite positive
reaction carried by contact.  All equations are original current
equalities on the ordinary domain, for every compact smooth scalar
test and every actual moment refinement.
\end{theorem}

\noindent\textbf{Status: \texttt{NEW\_PROVED}}.  This extends the
finite-energy Hilbert obstacle problem to every finite signed diffuse
datum.  The obstacle hypothesis and part-form boundary class are
retained explicitly; it is not an unrestricted classical obstacle or
nonlinear Monge--Ampere existence assertion.

\begin{proof}
First consider a finite-energy datum $\mu_j$ and the Hilbert obstacle
solution $w_j$ from Theorem~\ref{thm:r11-relative-energy-obstacle}.
Its reaction $\rho_j$ is positive, has finite energy integral, and is
carried by $\{w_j=\psi\}$.  It need not initially be assumed finite
in total mass.  Put $f_j=w_j-\psi\in H_+$.  Its curvature is
\[
                           D(f_jS)=\tau-\mu_j-\rho_j.
\]
The total variation on the right has finite energy integral, so the
$H$ level theorem (RT2) applies even if its total mass near the
boundary is infinite.  Since $f_j^+=f_j$, its zero-level formula is
\[
        \mathbf1_{\{\widetilde f_j=0\}}
                         (\tau-\mu_j-\rho_j)\ge0.
 \tag{OB10}
\]
Combine this with the contact carrier of $\rho_j$ to obtain
\[
                          0\le\rho_j\le(\tau-\mu_j)^+.
 \tag{OB11}
\]
This is the required finite-mass bound; it is derived before being
used in any finite-data approximation argument.

Now let $\mu_j=\mathbf1_{F_j}\mu$ be an increasing finite-energy
nest for $|\mu|$, and solve these finite-energy obstacle problems.
The signed measures converge in total variation to $\mu$ and
$\mu_j^-\le\mu^-$.  Thus all reactions are dominated by the same
finite diffuse measure
\[
                         \pi=\tau^++\mu^-,\qquad 0\le\rho_j\le\pi.
 \tag{OB12}
\]
For two indices put $d=w_j-w_l\in H$.  The source equation extends
to $T_kd\in H$.  On $\{d>0\}$ the measure $\rho_j$ is zero,
because $w_j>w_l\ge\psi$ there.  On $\{d<0\}$ the measure
$\rho_l$ is zero.  Therefore
\[
 \begin{split}
 \mathcal E(T_kd)
   &=\int T_k\widetilde d\,d(\mu_j-\mu_l+\rho_j-\rho_l)\\
   &\le k\|\mu_j-\mu_l\|_{\rm TV},\qquad
 c_V\{|\widetilde d|>k\}
                         \le\frac{\|\mu_j-\mu_l\|_{\rm TV}}{k}.
 \end{split}
 \tag{OB13}
\]
The first equality follows from the closed clipping rule and the
$H$ source identity.  The reaction integral is nonpositive by the
two contact observations.  The capacity estimate follows by the
normal clipped capacity test.  Thus $w_j$ is Cauchy in capacity.

Since $w_j=G(\mu_j+\rho_j)$ by the finite-energy theorem, positivity
of the transposition Green operator and (OB12) give
\[
                  |w_j|\le G(2|\mu|+\tau^+)=:B
                                        \quad\hbox{quasi-everywhere}.
\]
The fixed potential $B$ is finite quasi-everywhere and belongs to
$L^1(\lambda)$.  The usual summable-capacity subsequence construction
from (MD7)--(MD10) gives a finite quasicontinuous limit $w$ and
capacity convergence of the whole sequence.  Coercivity implies
convergence in $\lambda$-measure; the displayed domination then
gives $w_j\to w$ in $L^1(\lambda)$.  The inequalities $w_j\ge\psi$
pass outside a polar set to $w\ge\psi$.
For completeness, the source identity also gives
$\mathcal E(T_kw_j)\le k(\|\mu\|_{\rm TV}+\pi(F))$.
Thus each sequence of bounded clippings has a weak $H$ limit, whose
$L^2$ representative is $T_kw$.  Its canonical representative is the
capacity limit: choose a summable-capacity subsequence, then apply
Mazur's lemma to its tails.  The convex combinations still converge
uniformly outside the same arbitrarily small-capacity exceptional
sets, while their strong $H$ convergence identifies the canonical
limit.  Hence $T_kw\in H$ with the same quasi-everywhere values.

Write $\rho_j=r_j\pi$ with $0\le r_j\le1$.  The finite Radon
space here is second countable, so $L^1(\pi)$ is separable.
Weak-star sequential compactness gives a subsequence with
$r_j\rightharpoonup^*r$ in $L^\infty(\pi)$, where $0\le r\le1$.
Set $\rho=r\pi$.  It is finite, positive and diffuse.  For every
$f\in L^\infty(\lambda)$ the canonical $Gf$ is bounded and hence
belongs to $L^1(\pi)$.  Passing to the limit in the finite-energy
Green identity gives
\[
 \int wf\,d\lambda
       =\lim_j\int\widetilde{Gf}\,d(\mu_j+\rho_j)
       =\int\widetilde{Gf}\,d(\mu+\rho).
 \tag{OB14}
\]
The first passage uses $L^1(\lambda)$ convergence; the two datum
passages use total variation and weak-star convergence, respectively.
Thus $w=G(\mu+\rho)$.  The signed diffuse theorem now supplies its
complete intrinsic renormalized class and every original current
equation, so no current-only uniqueness assumption is needed.
The preceding clipped representative identification makes this an
equality of canonical functions quasi-everywhere as well as in $L^1$.

For contact let $g(s)=\min(1,s^+)$.  The functions
$g(w_j-\psi)$ converge in capacity to $g(w-\psi)$ and are bounded.
Finite diffuseness of $\pi$ gives convergence in $L^1(\pi)$.
The exact finite-energy contact condition and (OB12) therefore imply
\[
 \begin{split}
 \int g(\widetilde w-\widetilde\psi)\,d\rho
 &=\lim_j\int g(\widetilde w-\widetilde\psi)\,d\rho_j\\
 &=\lim_j\int g(\widetilde w_j-\widetilde\psi)\,d\rho_j=0.
 \end{split}
 \tag{OB15}
\]
The first passage is weak-star convergence against the fixed bounded
function.  The second is bounded by the $L^1(\pi)$ error.  Since
$w\ge\psi$, this proves the exact contact carrier.  In particular
the proof does not restrict an uncontrolled weak measure limit to a
discontinuous contact set.

For any solution pair in the stated class, signed renormalized
uniqueness first gives $w=G(\mu+\rho)$.  Since $\psi=G\tau$,
the difference $f=w-\psi=G(\mu+\rho-\tau)$ is a signed diffuse
Green potential and is nonnegative.  Its signed zero-level Kato
formula gives
\[
                  \mathbf1_{\{\widetilde f=0\}}
                              (\tau-\mu-\rho)\ge0.
\]
The contact carrier yields the second line of (OB8), exactly as in
(OB10)--(OB11).

Finally let $v=G(\mu+\eta)\in\mathscr M_{\mu,\psi}$.
Then $z=w-v=G(\rho-\eta)$.  On $\{z>0\}$ one has
$w>v\ge\psi$, so $\rho$ is zero there and
$\mathbf1_{\{z>0\}}(\rho-\eta)\le0$.
Lemma~\ref{lem:r11-diffuse-positive-part-Green} gives $z\le0$.
This proves leastness.  The class is nonempty before solving the
obstacle: take $\eta=\tau^++\mu^-$, since
$G(\mu+\eta)=G(\mu^++\tau^+)\ge G\tau=\psi$.
Two solution pairs are each least, so their potentials coincide;
their original equations then identify their reaction measures.
This proves uniqueness and independence of the extracted subsequence.
\end{proof}

\begin{theorem}[Order and quantitative dependence on the signed source]
\label{thm:r11-diffuse-obstacle-stability}
For a fixed allowed obstacle, write $w_\mu$ for the solution.
For any two finite signed diffuse data $\mu,\nu$ one has
\[
 \begin{split}
 |w_\mu-w_\nu|&\le G|\mu-\nu|,\\
 \|w_\mu-w_\nu\|_{L^1(\lambda)}&\le2\|\mu-\nu\|_{\rm TV},\\
 c_V\{|\widetilde w_\mu-\widetilde w_\nu|>k\}
                              &\le\frac{\|\mu-\nu\|_{\rm TV}}{k}.
 \end{split}
 \tag{OB16}
\]
If $\mu\le\nu$, then $w_\mu\le w_\nu$.
For fixed source the solution also increases when the allowed
obstacle increases.  These are relative energy orders; no pointwise
statement outside the coefficient support is asserted.
\end{theorem}

\noindent\textbf{Status: \texttt{NEW\_PROVED}} for the finite
signed diffuse obstacle map.

\begin{proof}
Put $d=w_\mu-w_\nu$ and $b=G|\mu-\nu|\ge0$.
The signed Green potential $z=d-b$ has source
\[
            Az=\mu-\nu-|\mu-\nu|+\rho_\mu-\rho_\nu.
\]
On $\{z>0\}$ one has $d>b\ge0$, so the contact condition gives
$\rho_\mu=0$ there.  The remaining restricted source is
nonpositive.  Apply (OB3) to obtain $d\le b$; interchange the two
data to obtain $-d\le b$.  This is the first line of (OB16).
The torsion bound and the finite-datum capacity bound for $b$
give the remaining two lines.  If $\mu\le\nu$, apply (OB3)
directly to $d$, using its same contact sign, to get order.
For ordered obstacles, the solution of the larger obstacle belongs
to the class (OB9) for the smaller one, so leastness proves order.
\end{proof}

\begin{corollary}[Positive nest data give decreasing finite reactions]
\label{cor:r11-positive-obstacle-reaction-TV}
If $\mu\ge0$ and $\mu_j\uparrow\mu$ are its finite-energy nest
restrictions, then their finite-energy obstacle solutions satisfy
\[
       w_j\uparrow w_\mu,\qquad
       \rho_j\downarrow\rho_\mu,\qquad
       \|\rho_j-\rho_\mu\|_{\rm TV}\longrightarrow0,\qquad
       0\le\rho_j\le\tau^+.
 \tag{OB17}
\]
In this positive-data case the limiting reaction itself has finite
energy integral.  The potential $w_\mu$ need not belong to $H$.
\end{corollary}

\noindent\textbf{Status: \texttt{NEW\_PROVED}} as a stronger
reaction convergence statement for positive nested data.

\begin{proof}
Order gives $w_l\ge w_j$ for $l\ge j$.  On the contact set
$\{w_l=\psi\}$ one also has $w_j=\psi$.
The nonnegative $H$ vector $f=w_l-w_j$ has curvature
\[
                         D(fS)=-(\mu_l-\mu_j)-(\rho_l-\rho_j).
\]
Its zero-level Kato inequality, restricted to the contact set of
$w_l$, gives $\rho_l\le\rho_j$ there because $\mu_l\ge\mu_j$.
Off that set $\rho_l=0$.  Hence the reactions decrease as measures.
Their common finite bound $\tau^+$ from (OB11) gives total-variation
convergence.  The preceding uniqueness and source convergence identify
the limits with $w_\mu,\rho_\mu$.  Domination by the finite-energy
measure $\tau^+$ gives finite energy integral for the reaction.
This does not put the possibly infinite-energy source $\mu$ or the
whole potential $w_\mu$ in $H$.
\end{proof}

\paragraph{Dependency and scope audit.}
The finite-energy obstacle/contact theorem is used only for the nest
approximants.  Its level-zero Kato consequence proves the uniform
finite reaction bound before extracting measures.  The signed
renormalized theorem supplies both the limiting boundary class and
the positive-part Green representation used in comparison; the
unbounded potentials are never tested as if they were $H$ vectors.
The common finite diffuse domination is what passes contact through
the signed approximation.  No ambient regularization of an arbitrary
upper envelope, no current pullback under an auxiliary field, and no
analytic PS/PL quotient is used.  Removing the finite-curvature
obstacle hypothesis, imposing ordinary pointwise boundary values, or
realizing these envelopes by ambient classical psh functions requires
additional arguments and is not asserted here.

% END REVIEWED MODULE research_r11/dirichlet/diffuse_obstacle_contact.tex


% BEGIN REVIEWED MODULE research_r11/dirichlet/linearized_power_class_dirichlet.tex
\section{Actual power-class variations and the local linearized Dirichlet problem}
\label{sec:r11-linearized-power-Dirichlet}

This section records the precise link between the new unbounded mixed
calculus and the relative linear Dirichlet theory.  It does not identify
every energy tangent with a classical psh variation.  Let the original
arbitrary background $S$ have degree $(p,p)$, put $q=n-p\ge1$, and assume
the fixed ordinary domain $V$ satisfies the one-affinoid hypothesis
(DG1).  Fix $\sigma>0$ and two actual potentials
$u_0,u_1\in\mathcal E^\sigma_{S,\partial}(V)$.  Write
\[
 u_t=(1-t)u_0+tu_1\quad(0\le t\le1),\qquad
 P=\{u_0=-\infty\}\cup\{u_1=-\infty\}.
 \tag{LD1}
\]
On $V\setminus P$ put $h=u_1-u_0$; assign $h=0$ on $P$ only to
choose a Borel coefficient representative.  Every coefficient used
below annihilates $P$, so this assignment has no effect on a current.
It is not an interpretation of the indeterminate difference of two
infinite potential values as an actual classical function.

Let $N_k(a)$ denote the diagonal degree-$k$ product over $S$, and let
mixed notation retain the same background.  The solid power cone puts
every $u_t$ in the same power class.  Set
\[
             Q_t=N_{q-1}(u_t),\qquad
             R_t=N_{q-2}(u_t)\quad(q\ge2).
 \tag{LD2}
\]
Here $N_0=S$.  Thus $Q_t$ is an actual original symmetric
middle-positive separately closed $(n-1,n-1)$-current, without a
smoothness, geometric-origin, or nondegeneracy condition.

\begin{theorem}[Differentiation of an actual admissible segment]
\label{thm:r11-actual-power-segment-derivative}
The map $t\mapsto\mu_{u_t}$ is a polynomial of degree at most $q$
with values in the locally finite signed Radon measures on $V$:
\[
 \mu_{u_t}
  =\sum_{i=0}^q\binom qi(1-t)^{q-i}t^i
                  M(u_0^{[q-i]},u_1^{[i]}).
 \tag{LD3}
\]
On each fixed compact $K\subset V$, differentiation is valid in total
variation.  In particular
\[
                  \frac d{dt}\mu_{u_t}=qD(hQ_t).
 \tag{LD4}
\]
The multiplier $hQ_t$ is coefficient integrable on every compact
stage, before taking its distributional curvature.  Endpoints in
(LD4) mean the corresponding one-sided derivative of the actual
segment.  All identities are on the original CLD current sheaf.

If $q\ge2$, the original positive energy current
\[
 \Gamma_{R_t}(h)
       =\tfrac12D(h^2R_t)-hD(hR_t)
 \tag{LD5}
\]
is coefficient Radon and satisfies
\[
             \frac {d^2}{dt^2}\mu_{u_t}
                       =-q(q-1)D\Gamma_{R_t}(h).
 \tag{LD6}
\]
Every multiplier in (LD5) is defined by the graded moment estimates.
The energy current itself need not be separately closed.  For $q=1$
the second derivative is zero and no $R_t$ is introduced.
\end{theorem}

\noindent\textbf{Status: \texttt{DIRECT\_COROLLARY}} of the new
power-class mixed polynomial identities, graded recursion, and graded
energy theorem.  The derivative is a consequence of an actual segment,
not an assumption that an arbitrary signed energy direction is
classically admissible.

\begin{proof}
The positive mixed polynomial identity gives (LD3), including the
endpoints.  Every coefficient is a locally finite original top
measure.  Its restriction to a fixed compact set is finite, so a finite
polynomial in these measures is differentiable in the total-variation
norm on that compact set.  This proves the asserted sense of
differentiation without a weak-limit interchange.

Set $W=1-u_0-u_1$ and $\kappa=q+\sigma$.  Then $|h|\le W$ off $P$.
The weighted curvature estimate (PG3) with weight degree $1$ and
current degree $q-1$ applies because $1+(q-1)=q<\kappa$.
It proves coefficient integrability of $hQ_t$.  The graded recursion
and mixed additivity give
\[
 D(hQ_t)=M(u_1,u_t^{[q-1]})-M(u_0,u_t^{[q-1]}).
 \tag{LD7}
\]
Differentiating the finite polynomial (LD3) gives exactly $q$ times
this difference, proving (LD4).

For $q\ge2$, the energy theorem (PG10) applies in degree $q-2$ for
every $\sigma>0$, because $(q-2)+2=q<\kappa$.  Polarization for
the coefficient vector $(-1,1)$ gives its positive diagonal energy
for $h=u_1-u_0$, and the intrinsic three-term expression is (LD5).
The same graded bounds with degrees $(2,q-2)$ and $(1,q-1)$ define
the two weighted terms.  Applying $D$ to the already defined current
identity, and using $D^2=0$, gives
\[
                  D\Gamma_{R_t}(h)=-D(hD(hR_t)).
 \tag{LD8}
\]
No curvature of a Borel-restricted current is used here.
Expand $D(hR_t)$ as the difference of the two mixed currents with
first slot $u_1$ and $u_0$.  The weighted terms are integrable by the
same degree calculation; graded recursion then expands
$D(hD(hR_t))$ as
\[
 M(u_1,u_1,u_t^{[q-2]})
 -2M(u_0,u_1,u_t^{[q-2]})
 +M(u_0,u_0,u_t^{[q-2]}).
\]
The second derivative of (LD3) is $q(q-1)$ times this expression.
This proves (LD6).  The compact coefficient controls and every
recursion used in the proof quantify over all original tests and all
actual moment refinements, as do the resulting equalities.
\end{proof}

\begin{corollary}[A local second-order remainder]
\label{cor:r11-power-segment-remainder}
For each compact $K\subset V$ there is a finite $C_K$ such that,
whenever $t,t+s\in[0,1]$,
\[
 \big\|\mu_{u_{t+s}}-\mu_{u_t}-sqD(hQ_t)\big\|_{\mathrm{TV}(K)}
                                                    \le C_Ks^2.
 \tag{LD9}
\]
For $q=1$ the left side is identically zero.  This is local in the
ordinary domain; no finite total mass near its boundary is assumed.
\end{corollary}

\noindent\textbf{Status: \texttt{DIRECT\_COROLLARY}}.

\begin{proof}
The second derivative of the finite polynomial (LD3) has bounded
total variation on $K$, uniformly for $0\le t\le1$, because it is a
finite linear combination of fixed finite measures there with bounded
scalar coefficients.  The integral Taylor remainder gives (LD9).
For $q=1$ the polynomial is affine.
\end{proof}

\subsection{Signed diffuse solvability for the resulting linear operator}

Fix a parameter $t\in[0,1]$.  Choose an ordinary affine-energy
neighborhood $\Omega$ compactly contained in $V$: explicitly, there
is a compact original $K$-affinoid $B'\subset V$ with
$\overline\Omega\Subset\operatorname{Int}_X(B')$.
Choose an ordinary $O\Subset\Omega$ and apply the faithful
construction to the actual current $Q_t|_\Omega$.
Write $\lambda_t$ for its finite faithful reference measure,
$F_t=\operatorname{supp}(\lambda_t)\cap O$, and
$(\mathcal E_t,H_t,c_t,G_t)$ for the original-core part energy,
its zero-boundary domain, its energy capacity, and its Green operator.
The notation records the fixed parameter.  No uniform comparison of
the forms for different parameters is assumed.

\begin{theorem}[The local linearized equation for every signed diffuse datum]
\label{thm:r11-local-linearized-diffuse-Dirichlet}
Let $\nu$ be any finite signed Radon measure on $F_t$ whose variation
charges no $c_t$-polar set.  Extend it by zero to $O$.
Then the original equation
\[
                            qD(vQ_t)=\nu\quad\hbox{on }O
 \tag{LD10}
\]
has exactly one canonical relative renormalized zero-boundary solution.
Its class is the one defined by finite energy-quasicontinuous
$v\in L^1(F_t,\lambda_t)$, $T_kv\in H_t$ for all $k>0$, the scalar
original chain identities with datum $\nu/q$, and vanishing radiation.
The solution is
\[
 \int_{F_t}vf\,d\lambda_t
                 =-\frac1q\int_{F_t}\widetilde{G_tf}\,d\nu
                  \qquad(f\in L^\infty(\lambda_t)).
 \tag{LD11}
\]
In the normalization of the faithful form, it satisfies
\[
 \|v\|_{L^1(\lambda_t)}\le\frac2q\|\nu\|_{\rm TV},\qquad
 \mathcal E_t(T_kv)\le\frac{k}{q}\|\nu\|_{\rm TV},\qquad
 c_t\{|\widetilde v|>k\}\le\frac{\|\nu\|_{\rm TV}}{qk}.
 \tag{LD12}
\]
The hard curvatures $\nu_k=qD(T_kv\,Q_t)$ are finite signed diffuse
measures, and
\[
 \|\nu_k\|_{\rm TV}\le\|\nu\|_{\rm TV},\qquad
 \|\nu_k-\nu\|_{\rm TV}
                  \le2|\nu|\{|\widetilde v|\ge k\}\longrightarrow0.
 \tag{LD13}
\]
There is also the complete original Kato, exact-level, and compact
scalar coarea calculus, with the level measure and energy both
multiplied by $q$.  The solution, current equation, and canonical
class are independent of the faithful compensators and reference
measure used to realize the same part energy.
\end{theorem}

\noindent\textbf{Status: \texttt{DIRECT\_COROLLARY}} of the full
signed diffuse theorem, applied after the unbounded mixed product
$Q_t$ has already been proved to be an original positive closed
current.  It gives linearized solvability in every ambient dimension
for the original arbitrary $S$ and every printed power-class
potential, subject to the explicit local domain and datum conditions.

\begin{proof}
The all-degree theorem and the graded power calculus construct $Q_t$
before any use of this section's energy form.  It is symmetric,
middle-positive, and separately $d'$- and $d''$-closed on the ordinary
neighborhood $\Omega$.  Thus every hypothesis of the analytic
affinoid faithful-form theorem applies to $Q_t$, followed by the
part construction on $O$.

Apply Theorem~\ref{thm:r11-signed-renormalized-Dirichlet} with this
background and with datum $\mu=\nu/q$.  Its existence and uniqueness
give (LD10)--(LD11), in the stated intrinsic class.  The mass estimates
(MD5)--(MD8) give (LD12), and (SD11), multiplied by $q$, gives
(LD13).  The signed original chain, exact-level and coarea theorems
give the remaining identities with the same scaling.  The faithful
reference-independence result proves independence of its realization.
All test functions are original compact smooth scalar functions on
$O$; the construction of $Q_t$ and the coefficient definition of
$vQ_t$ include every actual moment refinement.

If $Q_t=0$, then $F_t$ is empty and the only allowed datum is zero.
The resulting uniqueness is uniqueness of the zero energy class and
its original current, not of arbitrary ambient values outside the
support of $Q_t$.  The same distinction holds on any coefficient-null
region for a nonzero but degenerate background.
\end{proof}

\paragraph{Why the two statements are distinct.}
An actual segment supplies the signed derivative $qD(hQ_t)$ by
(LD4).  The local Dirichlet theorem supplies a canonical relative
energy solution for every allowed signed diffuse source, with its
own zero-boundary class.  It does not assert that every such solution
is the difference of two admissible classical power-class potentials,
or that $u_t+sv$ stays classical psh for either sign of $s$.
Conversely, an arbitrary actual segment difference need not have the
printed zero-boundary condition on the smaller part domain.  No such
condition, and no energy-diffuseness of its derivative datum, is
inferred merely from the polynomial identity.

If an actual direction $h$ does satisfy the printed renormalized
zero-boundary conditions on $O$ for a finite diffuse derivative datum
$\nu=qD(hQ_t)$, uniqueness identifies it with the solution of
(LD10).  This is a comparison of two already specified classes, not
an existence claim for classical variations.  The frozen compact-source
obstructions to unrestricted ambient classical solvability remain in
force.  This section solves neither the nonlinear diagonal
Monge--Ampere Dirichlet problem nor the analytic PS/PL quotient, and
does not use either problem as an input.

% END REVIEWED MODULE research_r11/dirichlet/linearized_power_class_dirichlet.tex


% BEGIN REVIEWED MODULE research_r11/comparison/power_class_source_obstruction.tex
\section{Energy-dual data and a source obstruction in the unbounded power classes}
\label{sec:r11-power-source-obstruction}

The finite-power calculus provides necessary estimates for its
Monge--Ampere image. These estimates do not by themselves remove
the distinction between ambient classical psh potentials and
relative energy solutions. We first give the general estimate,
then extend the precise frozen R10 compact-source obstruction to
the new unbounded power classes.

\begin{definition}[A power-energy dual bound for a positive datum]
\label{def:r11-power-dual-datum}
For a positive Radon measure $\mu$ on $V$, define
\[
 \mathfrak d_s(\mu)
   =\sup_{\substack{v\in\mathcal E^s_{S,\partial}(V)\\e_s(v)\leq1}}
                                  \int_V(-v)^s\,d\mu.
 \tag{DO1}
\]
The integral uses the extended actual value at a pole and is
$+\infty$ if that pole has positive datum mass. Thus finiteness
of (DO1) includes the required pole non-concentration; it is
not imposed by redefining the potential on a null set.
\end{definition}

\begin{proposition}[Necessary dual control and a capacity-dominated sufficient bound]
\label{prop:r11-power-dual-bound}
For every $u\in\mathcal E^s_{S,\partial}(V)$,
\[
                     \mathfrak d_s(\mu_u)
                               \leq C_{q,s}e_s(u)^{q/(q+s)}.
 \tag{DO2}
\]
If $\mu$ has finite total mass $M$ and, for some $A>0$,
\[
                     \mu(E)\leq A\operatorname{Cap}_S(E;V)
                             \quad(E\subset V\text{ Borel}),
 \tag{DO3}
\]
then
\[
 \int(-v)^s\,d\mu
    \leq2^{s+1}A^{s/(q+s)}M^{q/(q+s)}e_s(v)^{s/(q+s)}
           \quad(v\in\mathcal E^s_{S,\partial}(V)),
 \tag{DO4}
\]
and in particular $\mathfrak d_s(\mu)<\infty$.
If $M=0$ the assertion is interpreted as the zero measure.
\end{proposition}

\noindent\textbf{Status: \texttt{NEW\_PROVED}} as a general
datum bound in the original arbitrary-background power class.
Neither direction is asserted to be a solvability criterion.

\begin{proof}
Apply (PE11) with $v$ as the undifferentiated potential and
$q$ copies of $u$ in the differentiated slots. Normalize
$e_s(v)\leq1$ and take the supremum to obtain (DO2).

For (DO4), the pole of $v$ is capacity null and hence $\mu$-null
by (DO3). Let $R>0$. Layer cake, the total mass bound on the
small-value part, and (DO3) on the remaining part give
\[
 \begin{split}
 \int(-v)^s\,d\mu
 &\leq M R^s+A s\int_R^\infty
                        t^{s-1}\operatorname{Cap}_S(\{v<-t\};V)\,dt\\
 &\leq M R^s+A s R^{-q}J_s(v)\\
 &\leq M R^s+A2^{q+s}R^{-q}e_s(v).
 \end{split}
 \tag{DO5}
\]
For $M,e_s(v)>0$, take
$R=(A2^{q+s}e_s(v)/M)^{1/(q+s)}$.
The two terms then agree and their sum is (DO4).
If $e_s(v)=0$, let $R\downarrow0$ in (DO5).
This also proves the zero-energy case without a nondegeneracy
assumption on the background.
\end{proof}

\subsection{The fixed original moment-Laplacian background}

Let $n\geq2$, let $Q=\prod_{i=1}^n(\ell_i,r_i)$ be a bounded
open real box, and set
\[
 U=\{x\in(\mathbf G_{m,K}^n)^{\mathrm{an}}:
                         t(x)=(\log|T_1(x)|,\ldots,\log|T_n(x)|)\in Q\}.
\]
Let $\sigma:Q\to U$ be the original joint Gauss section and put
\[
 \beta=\sum_i d't_i\wedge d''t_i,\qquad
                 S={\beta^{n-1}\over(n-1)!}\wedge[U].
 \tag{DO6}
\]
This is an original smooth positive separately closed background
with remaining degree $q=1$. It extends to a larger ordinary
polyannulus, so the one-affinoid hypotheses of the full power
calculus hold.

The identity matrix is positive definite in the real moment
coordinates, and the transverse calibration measures in the
frozen R10 argument are Lebesgue measures with full support in
their transverse real boxes. This does \emph{not} assert
$\operatorname{supp}S=U$ as an ordinary analytic support:
on an ordinary residue polydisc where all $|T_i|$ are constant,
$\beta$ and therefore $S$ vanish.

The exact \texttt{INTERNAL\_INPUT}s are the R6 lemma
\path{lem:r6-joint-gauss-convexity}, the original smooth
calibration $\beta^n[U]/n!=\sigma_*(dt)$, and the R10 real
convex rigidity lemma
\path{lem:r10-convex-compact-marginal-rigidity} in
\path{full_support_\allowbreak rigidity_of_\allowbreak compact_sources.tex}.
The old locally bounded obstruction is retained as input; its
new extension below supplies the missing unbounded trace and
current-integrability steps.

\begin{theorem}[No compact classical source in the unbounded power classes]
\label{thm:r11-power-compact-source-obstruction}
For the precise background (DO6), let $u$ be negative extended
classical psh with $\mathrm{BG}(u,0)$. Assume that $uS$ is
locally coefficient-defined and that $D(uS)$ is a positive
Radon measure whose tropical pushforward is carried by a compact
subset of $Q$. Then
\[
                             D(uS)=0.
 \tag{DO7}
\]
In particular, for every $s>0$ there is no
$u\in\mathcal E^s_{S,\partial}(U)$ satisfying
$\mu_u=\mu$ for a nonzero positive compactly supported datum
$\mu$ on $U$.
\end{theorem}

\noindent\textbf{Status: \texttt{NEW\_PROVED}} as the extension
to coefficient-integrable compact-sublevel potentials and hence
to the new unbounded power classes. It is a counterexample to
unrestricted ambient classical solvability, not a theorem
restricting the arbitrary background in the positive results.

\begin{proof}
The bounded Gauss traces $g_j=u_j\circ\sigma$ are finite convex
functions on $Q$ by the frozen R6 lemma. Their decreasing actual
limit $g=u\circ\sigma$ is extended convex and at most zero.
It has at least one finite value. Indeed
$K=\overline{\{u<-1\}}$ is compact in $U$, so $t(K)$ is
a compact proper subset of the open box $Q$. If $y\notin t(K)$,
then $u(\sigma(y))\geq-1$.

An extended convex function bounded above by zero on an open
convex box, with one finite value, is finite everywhere. To
verify this explicitly, suppose $g(x)=-\infty$ and $g(y)$ is
finite. Extend the segment from $x$ through $y$ slightly past
$y$ to a point $z\in Q$. Then $y$ is a strict convex
combination of $x,z$, and the convex inequality gives
$g(y)\leq-\infty$, a contradiction. Thus $g$ is an ordinary
finite convex function, locally bounded and continuous on $Q$.

For $\varphi\in C_c^\infty(Q)$, its pullback $\varphi(t)$
is an original compact smooth scalar test on $U$. The exact
calibration gives
\[
              D(\varphi(t))\wedge S
                         =\sigma_*((\Delta\varphi)\,dt).
\]
The assumed original coefficient multiplication therefore gives
\[
 \bigl\langle t_*D(uS),\varphi\bigr\rangle
      =\int_Q g\,\Delta\varphi\,dt
      =\langle\Delta g,\varphi\rangle.
 \tag{DO8}
\]
All integrals are finite on the compact real test core. This
uses the original current equation; no equality on arbitrary
delta tests is substituted for it.

The positive measure $\Delta g$ has compact support in $Q$.
The frozen R10 convex compact-marginal rigidity lemma, with
Lebesgue transverse measures, makes $g$ affine and
$\Delta g=0$. Positivity then forces $D(uS)=0$: a compact
ordinary set maps to a compact real set, and its measure is
bounded by the mass of the full inverse image of that real set,
which is zero by (DO8). This proves (DO7).

For $u\in\mathcal E^s_{S,\partial}(U)$, the full recursive
calculus of Theorem~\ref{thm:r11-power-full-recursion}, in
remaining degree one, supplies precisely the missing hypotheses
$uS$ coefficient-defined and $D(uS)=\mu_u\geq0$.
A compact ordinary source has compact tropical image, so (DO7)
proves the stated nonexistence in every positive power class.
\end{proof}

\begin{corollary}[Finite dual energy does not imply ambient classical existence]
\label{cor:r11-dual-energy-insufficient}
Take $Q=(-1,1)^n$ in (DO6), choose a nonzero
$0\leq\chi\in C_c^\infty(Q)$ with $\chi\leq1$, and set
\[
                       \mu=\sigma_*(\chi\,dt).
 \tag{DO9}
\]
For every $s>0$, this nonzero finite smooth compact source has
$\mathfrak d_s(\mu)<\infty$. It is even dominated by one half
of the curvature measure of one normalized smooth original
classical tester. Nevertheless it is not $\mu_u$ for any
$u\in\mathcal E^s_{S,\partial}(U)$, for any $s>0$.
Thus sufficiency of the energy-dual bound (DO1) for such an
ambient classical Dirichlet solution is \texttt{FALSE}.
\end{corollary}

\noindent\textbf{Status: \texttt{NEW\_PROVED}} for the
unbounded energy-dual inference. The datum and its tester
domination are the frozen R10 example, not a new example of
locally bounded nonexistence.

\begin{proof}
The original smooth tester $v_0=n^{-1}\sum_i t_i^2$ satisfies
$0\leq v_0\leq1$ and
$D(v_0S)=2\sigma_*(dt)$ on the whole same domain, by the
frozen all-original calibration. Therefore
\[
                     \mu\leq\tfrac12D(v_0S),\qquad
                     \mu(E)\leq\tfrac12\operatorname{Cap}_S(E;U).
\]
Its mass is finite and nonzero. Proposition
\ref{prop:r11-power-dual-bound} proves finite $\mathfrak d_s$
for every $s>0$. The preceding theorem excludes a power-class
classical solution because the source has compact ordinary
support. The admissible power classes are not empty: the bounded
Laurent-box barrier belongs to every one of them.

There is no contradiction with the positive diffuse renormalized
Dirichlet theorem. That theorem uses the relative energy class
with its specified part-form boundary condition, and does not
claim ambient classical psh membership. The obstruction here
shows why a growth or energy bound on the datum cannot erase
that difference even in this original smooth moment-Laplacian
example.
\end{proof}

% END REVIEWED MODULE research_r11/comparison/power_class_source_obstruction.tex


% BEGIN REVIEWED MODULE research_r11/energy/strong_completion_source_obstruction.tex
\section{A compact-source obstruction even in the unrestricted strong completion}
\label{sec:r11-strong-completion-source-obstruction}

This section specializes only to the explicit counterexample background
(DO6), with $n\geq2$, $Q=(-1,1)^n$, and remaining degree $q=1$.
The field, norm, analytic space, and original test category are unchanged.
Write $U=t^{-1}(Q)$ and let $\sigma:Q\to U$ be the joint original
Gauss section.  The new potential space is precisely
$\widehat{\mathcal E}^{1}_{\rm str}(U)$ of
Definition~\ref{def:r11-strong-first-energy-completion}.
Its elements need not have an actual common lower barrier or an
ambient classical representative.  Their currents and multipliers
are those already constructed in (IC10)--(IC13).

The frozen inputs are exactly the R6 joint-Gauss convexity and
original calibration, and the R10 compact-source convex rigidity
listed in the power-class obstruction section.  The already proved
actual and fixed-barrier obstructions (DO7) and (BV13) are not
reproved here.  The new issue is whether the energy bound of an
arbitrary strongly Cauchy row still forces a finite convex trace
in its current limit.  The following estimate answers that issue.

\begin{lemma}[Local trace compactness from a first-energy bound]
\label{lem:r11-energy-bounded-convex-trace-compactness}
Let $(u_j)\subset\mathcal E^1_{S,\partial}(U)$ be actual classical
potentials with $\sup_j e_1(u_j)\leq M<\infty$.
Their finite convex Gauss traces $g_j=u_j\circ\sigma\leq0$ have
uniformly bounded $L^1$ norms on every compact real subbox of $Q$.
They are locally uniformly bounded and locally equicontinuous.
Consequently a subsequence converges locally uniformly on $Q$
to a finite convex function $g\leq0$.
No common pointwise lower barrier is assumed.
\end{lemma}

\noindent\textbf{Status: \texttt{NEW\_PROVED}} in this precise
original moment example.  The compactness assertion concerns real
convex traces, not an asserted compactness of ambient classical psh
potentials in their original category.

\begin{proof}
The finite-convex-trace argument preceding (DO8) applies separately
to every actual $u_j$: its compact negative sublevels supply one
finite Gauss value, and extended convexity bounded above by zero
then gives finiteness everywhere.  Fix a compact real subbox
$B\Subset Q$ and put $K=\sigma(B)$, a compact ordinary analytic
set.  The normalized smooth original tester
\[
                v_0=\frac1n\sum_{i=1}^n t_i^2
           \quad\hbox{satisfies}\quad
                D(v_0S)=2\sigma_*(dy).
 \tag{SO1}
\]
Let $\rho$ be the strict-box barrier from (TC2).
Choose $A>0$ so that $A\rho<-2$ near $K$ and set
\[
                     \psi=\max(v_0-2,A\rho).
 \tag{SO2}
\]
The exact core construction (TC7)--(TC8), with its harmless
constant rescaling, shows that $\psi$ is an actual bounded
finite-mass member of every positive power class.  In particular
$e_1(\psi)<\infty$, and it agrees with $v_0-2$ near $K$.
Original locality therefore gives
$\mathbf1_K\mu_\psi=2\sigma_*(\mathbf1_Bdy)$.
The global constant-one mixed energy inequality yields
\[
 \int_B(-g_j)\,dy
   =\frac12\int_K(-u_j)\,d\mu_\psi
   \leq\frac12\int_U(-u_j)\,d\mu_\psi
   \leq\frac12M^{1/2}e_1(\psi)^{1/2}.
 \tag{SO3}
\]
Every multiplier in this equation is the previously justified
actual multiplier.  This proves the local $L^1$ bound, including
the case $M=0$.

For completeness, the convex compactness step is elementary.
Fix a compact real set $C\Subset\operatorname{Int}(B)$ and
choose $\delta>0$ such that each Euclidean ball
$B(x,\delta)$ with $x\in C$ lies in $B$.
The functions $h_j=-g_j$ are nonnegative and concave.
For $z\in B(x,\delta/2)$, write $z=(x+y)/2$ with
$y\in B(x,\delta)$.  Concavity gives
$h_j(z)\geq h_j(x)/2$.  Integration implies
\[
 0\leq -g_j(x)
       \leq\frac{2}{\operatorname{vol}(B(0,\delta/2))}
                       \int_B(-g_j)\,dy.
 \tag{SO4}
\]
Thus the traces have uniform two-sided bounds on every compact
interior set.  Apply this bound on a slightly larger compact
set.  On any line through a point of the smaller set, the
monotone secant slopes of a convex function are then bounded
by its oscillation on the larger set divided by the fixed
distance to its boundary.  This gives a uniform Lipschitz
bound on every smaller compact subbox.  A finite-net
subsequence argument and a countable exhaustion yield local
uniform convergence of a subsequence.  Convexity and the
upper bound zero pass to its finite limit.
\end{proof}

\begin{theorem}[Strong completion does not create a nonzero compact source]
\label{thm:r11-strong-completed-compact-source-obstruction}
For the precise background just specified, let
$u\in\widehat{\mathcal E}^{1}_{\rm str}(U)$.
If the positive original measure $D(uS)=\widehat N_1(u)$
has tropical pushforward carried by a compact subset of $Q$,
then
\[
                              D(uS)=0.
 \tag{SO5}
\]
In particular, the nonzero positive compact source
$\mu=\sigma_*(\chi\,dy)$ of (DO9) is not $D(uS)$ for
any member of this unrestricted strong completion.
Thus passing from actual power-class potentials to the complete
relative strong first-energy space does not make finite-dual or
capacity-controlled data universally solvable.
\end{theorem}

\noindent\textbf{Status: \texttt{NEW\_PROVED}} as an extension
of the frozen source obstruction to the new strong completion
without an actual lower barrier.  The universal source-solvability
inference being rejected is \texttt{FALSE}.

\begin{proof}
Choose an actual strongly Cauchy defining row $(u_j)$ for $u$.
Its first energies are bounded by (EV9), or (IC2).
The preceding lemma supplies a subsequence of its Gauss traces
converging locally uniformly to a finite convex $g$.
The strong completed calculus gives the original coefficient
convergence $u_jS\to uS$, before applying any tropical test.
Consequently, for every $\varphi\in C_c^\infty(Q)$,
\[
 \begin{split}
 \langle t_*D(uS),\varphi\rangle
   &=\lim_j\langle u_jS,D(\varphi(t))\rangle\\
   &=\lim_j\int_Qg_j\,\Delta\varphi\,dy
     =\int_Qg\,\Delta\varphi\,dy
     =\langle\Delta g,\varphi\rangle.
 \end{split}
 \tag{SO6}
\]
The pullback $\varphi(t)$ is an original compact smooth test:
inverse images of compact real subboxes in this original torus
are compact polyannuli.  Each identity uses the already fixed
all-original calibration.  It does not define the completed
current by a selected toric test subfamily.

The positive measure $\Delta g$ has compact support in $Q$.
The exact frozen R10 convex compact-marginal rigidity theorem,
valid here for $n\geq2$, gives $\Delta g=0$.
Hence $t_*D(uS)=0$.  Positivity makes the original measure zero:
for every compact ordinary $L\subset U$, its mass is at most
the mass of $t^{-1}(t(L))$, which is the zero pushforward
mass of the compact set $t(L)$.  Radon inner regularity
then gives (SO5).

The datum (DO9) is finite, nonzero, compactly supported,
dominated by one half of an actual normalized tester measure,
and has finite energy-dual size, as already proved there.
It contradicts (SO5) if it were a completed source.
No claim is made that this special background exhausts arbitrary
backgrounds; one allowed background suffices to disprove the
unqualified universal source assertion.
\end{proof}

\paragraph{The resulting minimal equation gate.}
The strong completion solves a precise nonlinear Cauchy problem
and retains all-original closed curvature currents.  It does
not contain every solution of the larger relative linear
energy problem even when $q=1$.  The surviving convex trace
is an explicit obstruction.  A prospective general nonlinear
Dirichlet theorem must therefore specify an admissible operator
image or a different potential class; neither completion alone
nor a finite datum energy norm supplies the missing equation.
The analytic PS/PL relation quotient and low-real-degree CTO
are separate and are not used in this argument.

% END REVIEWED MODULE research_r11/energy/strong_completion_source_obstruction.tex


% BEGIN REVIEWED MODULE research_r11/dirichlet/strong_space_variational_dominated_data.tex
\section{A nonlinear variational problem for explicit dominated data on the full strong space}
\label{sec:r11-strong-space-dominated-variational}

Retain precisely the strict analytic box and background-extension
hypotheses (TC1)--(TC2).  The background $S$ is the original
arbitrary symmetric middle-positive separately closed $(p,p)$-current,
and $q=n-p\geq1$.  Put $d=q+1$, $\alpha=2^{-q}$, and write
$\mathscr C=\widehat{\mathcal E}^{1}_{\rm str}(V)$ for the
relative strong completion constructed in (IC1)--(IC22).
Its original top current is $\widehat\mu_u=\widehat N_q(u)$,
and its global mixed first energy, diagonal energy and discrepancy
are $\widehat B$, $\widehat e$ and $\widehat I$.

The frozen R4 completed variational retraction treats a bounded
common-collar problem.  The new fixed-barrier problem (BV7) treats
unbounded functions above one actual integrable lower barrier.
Here the admissible potential cone is the entire relative strong
completion, without a prescribed common lower barrier.  The
source hypothesis is instead the explicit domination below.

\begin{lemma}[Literal datum integrals and strong continuity]
\label{lem:r11-dominated-datum-strong-integration}
Let $\psi\in\mathcal E^1_{S,\partial}(V)$ be an actual
classical potential and let $A\geq0$.  Suppose $\mu$ is a
positive Radon measure on $V$ satisfying the actual Borel
measure inequality
\[
                              0\leq\mu\leq A\mu_\psi.
 \tag{VF1}
\]
Total mass on $V$ is not assumed finite.  For every $u\in\mathscr C$,
the actual Borel integral is finite and
\[
 0\leq\int_V(-u)\,d\mu
       \leq A\widehat B(u,\psi^{[q]})
       \leq A\widehat e(u)^{1/d}e_1(\psi)^{q/d}.
 \tag{VF2}
\]
It is independent of the relative representative.  If
$\widehat I(u_j,u)\to0$, then
\[
                         \int_V|u_j-u|\,d\mu\longrightarrow0.
 \tag{VF3}
\]
More precisely, for $u,v\in\mathscr C$ and
$M=\max\{\widehat e(u+v),e_1(\psi)\}$,
\[
 \int_V|u-v|\,d\mu
   \leq A\,2^{(q-1)\alpha}
                   \widehat I(u,v)^\alpha M^{1-\alpha},
 \tag{VF4}
\]
with value zero when $M=0$.
\end{lemma}

\noindent\textbf{Status: \texttt{DIRECT\_COROLLARY}} of
the actual-multiplier construction (IC11), the global limit
(IC13), completed Holder, and the completed extension of (EV1).
The measure inequality (VF1) is a checkable sufficient source
hypothesis, not an assumption that the desired equation is solvable.

\begin{proof}
The original positive measure $\mu_\psi$ gives zero mass to
every locally outer-capacity-null set.  This follows from its
local coefficient-capacity modulus and a countable compact
exhaustion.  In particular it ignores the pole of every member
of $\mathscr C$, by the Borel-limit theorem.  The same is true
of $\mu$ by (VF1).  The global mixed integral in (VF2) is the
actual coefficient multiplier constructed in (IC11)--(IC13),
not a newly postulated functional on an abstract completion.
Domination and the global mixed Holder inequality prove (VF2).
This also proves representative independence, including the
possibly infinite values on a null pole set.

The completed tester estimate (EV1), with the fixed actual
tester potential $\psi$, gives (VF4) after applying (VF1).
If $u_j\to u$ strongly, their energies are bounded by the
energy-root estimate.  Minkowski bounds $\widehat e(u_j+u)$
uniformly.  Formula (VF4) therefore gives (VF3) on the whole
domain, even if $\mu(V)=\infty$.  No global finite-measure
dominated-convergence theorem is used in this passage.
\end{proof}

\begin{theorem}[Unique variational maximizer on the full relative strong space]
\label{thm:r11-full-strong-dominated-maximizer}
For every datum satisfying (VF1), the functional
\[
 \mathcal F_\mu(u)
      =-\frac{\widehat e(u)}{d}-\int_Vu\,d\mu,
                       \qquad u\in\mathscr C,
 \tag{VF5}
\]
has a maximizer, unique modulo local outer $S$-capacity zero.
Write its class as $\Pi_{\rm str}(\mu)$.  Every maximizing
sequence in $\mathscr C$ converges strongly to it.  The same
supremum and limit are obtained from actual classical
$\mathcal E^1$ maximizing sequences.

For $u=\Pi_{\rm str}(\mu)$ the exact variational characterization is
\[
 \int_V(v-u)\,d(\widehat\mu_u-\mu)\leq0
                                      \qquad(v\in\mathscr C).
 \tag{VF6}
\]
All signed integrals in (VF6) are absolutely convergent, even
when either positive measure has infinite total mass.  Equivalently,
\[
 \begin{split}
  \widehat e(u)&=\int_V(-u)\,d\mu,\\
  \int_V(-v)\,d\widehat\mu_u
                  &\geq\int_V(-v)\,d\mu\qquad(v\in\mathscr C).
 \end{split}
 \tag{VF7}
\]
The energy and maximal value satisfy the sharp bounds
\[
 \begin{split}
  \widehat e(u)&\leq A^{d/q}e_1(\psi),\\
  0\leq\sup_{\mathscr C}\mathcal F_\mu
       &=\frac qd\widehat e(u)
                          \leq\frac qd A^{d/q}e_1(\psi).
 \end{split}
 \tag{VF8}
\]
For $\mu=A\mu_\psi$ the maximizer is the actual potential
$A^{1/q}\psi$ and both bounds are attained.  The case $A=0$
has the zero relative class as maximizer.
\end{theorem}

\noindent\textbf{Status: \texttt{NEW\_PROVED}} for the
unconstrained relative strong first-energy cone and the explicit
dominated source class.  The admissible cone still has its
specified strong-approximation boundary condition.  Formula
(VF6) is not asserted to be an unrestricted current equation.

\begin{proof}
Put $C=Ae_1(\psi)^{q/d}$.  For $r=\widehat e(u)^{1/d}$,
(VF2) gives the coercive bound
\[
                    \mathcal F_\mu(u)\leq-r^d/d+Cr.
 \tag{VF9}
\]
The zero potential has value zero.  The supremum $L$ is finite,
and any maximizing sequence can be taken with values at least
$-1$; (VF9) bounds its energies.  A sequence already maximizing
has this property after finitely many terms, so no choice is
lost by this reduction.  The completed cone is closed under
midpoints.  The linear datum term cancels in their midpoint gap,
and the sharp estimate gives
\[
 c_q\widehat I(u_j,u_l)
   \leq L-\frac{\mathcal F_\mu(u_j)+\mathcal F_\mu(u_l)}2
                      \longrightarrow0,
 \qquad c_q=\frac{1-2^{-q}}{2(q+1)}.
 \tag{VF10}
\]
The strong completion theorem supplies a relative Borel limit
$u\in\mathscr C$ with discrepancy tending to zero.  Its energy
converges by the energy-root estimate, and its datum integral
converges by (VF3).  Hence $\mathcal F_\mu(u)=L$.
Two maximizers have zero discrepancy by the same midpoint
estimate and thus are the same relative class.  This proves
existence, uniqueness and convergence of every maximizing
sequence.  Actual classical strong density, (IC13), and (VF3)
show that restricting the supremum to the actual class does
not change it; any actual maximizing sequence has the same limit.

For a competitor $v\in\mathscr C$, positive multilinearity
and symmetry of the global mixed energy make the energy of
$(1-t)u+tv$ a finite polynomial.  Its derivative at zero gives
\[
 \left.\frac d{dt}\right|_{t=0}
       \left[-\frac{\widehat e((1-t)u+tv)}d\right]
                           =\int_V(v-u)\,d\widehat\mu_u.
\]
The datum term is linear and finite by (VF2).  Thus maximality
implies (VF6), and concavity gives the converse tangent inequality.
Absolute integrability follows from (VF2) for the datum and
from the completed mixed first energies for the curvature.

The path $t\mapsto tu$ is admissible for every $t\geq0$.
Differentiating at the interior parameter $t=1$ gives
$\widehat e(u)=\int(-u)\,d\mu$.  Using the admissible paths
$u+tv$, $t\geq0$, gives the second line of (VF7).  Conversely
those two statements imply (VF6) by subtracting the equality
for the $u$ term.  If $\widehat e(u)>0$, (VF2) and the
scaling equality imply
$\widehat e(u)^{q/d}\leq Ae_1(\psi)^{q/d}$; if it is zero,
the same energy bound is immediate.  The scaling equality also
gives $L=q\widehat e(u)/d$.  This proves (VF8).

When $\mu=A\mu_\psi$, the actual potential $b=A^{1/q}\psi$
has original top measure $\mu_b=\mu$, by the original
homogeneity.  Its Euler expression in (VF6) is identically
zero, so it is the maximizer.  Its energy is
$A^{d/q}e_1(\psi)$, giving equality in the bounds.
When $A=0$, the objective is $-\widehat e/d\leq0$;
the zero potential maximizes it, uniquely as a relative class.
\end{proof}

\begin{corollary}[Quantitative data dependence and the dominated operator image]
\label{cor:r11-strong-dominated-data-stability}
Suppose $\mu,\nu$ are positive Radon measures, each dominated
by a finite constant times $\mu_\psi$ for the same actual
$\psi\in\mathcal E^1$.  Define
\[
 \delta_1(\mu,\nu)
   =\sup_{\substack{v\in\mathcal E^1\\ e_1(v)\leq1}}
                          \int_V(-v)\,d|\mu-\nu|<\infty.
 \tag{VF11}
\]
For $u=\Pi_{\rm str}(\mu)$ and $w=\Pi_{\rm str}(\nu)$,
\[
 \widehat I(u,w)
  \leq\int_V(u-w)\,d(\nu-\mu)
  \leq\delta_1(\mu,\nu)
                 \bigl(\widehat e(u)^{1/d}+\widehat e(w)^{1/d}\bigr).
 \tag{VF12}
\]
In particular, under one common bound $\mu_j,\mu\leq A\mu_\psi$,
convergence $\delta_1(\mu_j,\mu)\to0$ gives strong convergence
of the maximizers, all original local current and critical weighted
limits in (IC11), and all global mixed-energy limits in (IC13).

If $b\in\mathscr C$ and $\widehat\mu_b$ satisfies a domination
of the form (VF1), then
$\Pi_{\rm str}(\widehat\mu_b)=b$ as a relative class.  This
retraction does not require finite total mass of $\widehat\mu_b$.
\end{corollary}

\noindent\textbf{Status: \texttt{DIRECT\_COROLLARY}} of
the full variational theorem and the literal datum integration lemma.

\begin{proof}
The variation $|\mu-\nu|$ is dominated by the sum of the two
source domination bounds, so (VF2) proves finiteness in (VF11).
Strong actual approximation and (VF3), applied to this positive
dominated variation, extend its normalized bound to every
completed $v$:
$\int(-v)\,d|\mu-\nu|\leq\delta_1\widehat e(v)^{1/d}$.
For zero energy this follows from relative separation and the
datum's outer diffuseness; otherwise use positive scalar
normalization and approximation.

Use $w$ and $u$ as the two competitors in their respective
Euler inequalities.  Adding them gives the first inequality
in (VF12) with the displayed sign.  The pointwise bound
$|u-w|\leq(-u)+(-w)$ and the extended normalized estimate
give the second.  The common domination bounds the solution
energies by (VF8), so (VF12) proves strong convergence.
The current and weighted conclusions are then exactly the
proved strong-completion continuity statements.  They are
not assertions of local total variation for arbitrary moving
nonlinear currents.  Finally a dominated operator-image datum
makes the Euler expression for $b$ identically zero; uniqueness
proves the stated retraction.
\end{proof}

\begin{corollary}[Dominated-data variational solvability does not imply the equation]
\label{cor:r11-strong-dominated-maximizer-not-equation}
Even within (VF1), the unique full strong-space maximizer need
not satisfy $\widehat\mu_u=\mu$.  Thus the implication from
the all-cone variational characterization (VF6)--(VF7) to
the unrestricted original source equation is \texttt{FALSE}.
\end{corollary}

\noindent\textbf{Status: \texttt{DIRECT\_COROLLARY}} of
the explicit strong-completion compact-source obstruction
(SO1)--(SO6), with the source-class membership checked below.

\begin{proof}
Use the original smooth moment-Laplacian background (DO6),
with $n\geq2$, $q=1$, and the nonzero compact positive source
$\mu=\sigma_*(\chi\,dy)$ in (DO9), where $0\leq\chi\leq1$.
Choose the real compact subbox $B$ in (SO2) to contain the
support of $\chi$.  The corresponding actual bounded
finite-mass potential $\psi$ agrees with $v_0-2$ near
$\sigma(B)$, so original locality and (SO1) give
\[
                              \mu\leq\tfrac12\mu_\psi
                                      \quad\hbox{on the whole domain}.
 \tag{VF13}
\]
Outside that compact source support, the left side is zero
and the right side is positive.  Thus the datum satisfies
the explicit hypothesis with $A=1/2$.  The preceding theorem
produces its unique maximizer on the entire relative strong
completion.  Theorem
\ref{thm:r11-strong-completed-compact-source-obstruction}
excludes that nonzero source from the top-current image of
every member of this same space.  Hence its maximizer cannot
satisfy the source equation.  This verifies the failure within
the proved admissible datum class, rather than by leaving
that class in a counterexample.
\end{proof}

\paragraph{Exact boundary of the conclusion.}
The new theorem solves a coercive nonlinear variational problem
on a complete relative Borel cone, for an explicit source class
that may have infinite boundary mass.  Every datum and mixed
curvature integral is defined by an actual coefficient multiplier.
The background remains arbitrary and all current constructions
retain the full original test sheaf and every actual refinement.
The admissible variations are exactly this cone; they are not
all compact original smooth signed perturbations.  The explicit
counterexample shows why its Euler inequality cannot be promoted
to a general Monge--Ampere equation.  No ordinary pointwise
boundary theorem, ambient classical representative, analytic
PS/PL quotient, or low-degree signed order transfer is supplied
by this variational construction.

% END REVIEWED MODULE research_r11/dirichlet/strong_space_variational_dominated_data.tex


% BEGIN REVIEWED MODULE research_r11/comparison/first_energy_hilbert_source_cone.tex
\section{The exact degree-one source cone of the strong completion}
\label{sec:r11-first-energy-Hilbert-source-cone}

Assume the strict analytic box and background extension
(TC1)--(TC2), and now put $q=1$. Thus $S$ is an arbitrary
original symmetric middle-positive separately closed
$(n-1,n-1)$-current. Choose an ordinary neighborhood $\Omega$
of $\overline V$, contained in the interior of the original
affinoid and within the given extension domain of $S$.
The analytic faithful-form theorem supplies
$\lambda=D(hS)$ with $0\leq h\leq1$ on $\Omega$.
Use its coercive part space
\[
                 H=\mathscr D_0(V),\qquad
                 \|f\|_H^2=\mathcal E(f,f),\qquad
                 F=V\cap\operatorname{supp}S.
 \tag{HC1}
\]
The measure $\lambda|_F$ is finite, and
$\|f\|_{L^2(\lambda)}^2\leq2\|f\|_H^2$.
All currents are tested against the full original compact
smooth CLD scalar sheaf, with every actual moment refinement.

The frozen R4 Proposition \path{prop:r4-classical-in-energy}
identifies bounded local classical BT energy with the closed
form energy. Its proof uses local original bounded
regularization and coefficient domination; (AE2)--(AE10)
and the analytic collars supply these same inputs on the
present analytic domain. The frozen R6
\path{dirichlet_energy_separation.tex} already proves
Hilbert projection and source separation for a bounded
fixed-collar constraint. The new identification here concerns
the full global unbounded strong first-energy cone.

\begin{theorem}[An isometric closed-cone realization]
\label{thm:r11-first-energy-closed-Hilbert-cone}
Every actual $u\in\mathcal E^1_{S,\partial}(V)$ belongs to
$H$ as its $L^2(\lambda)$ representative, and
\[
 \mathcal E(u,v)=\widehat B(u,v),\qquad
 \|u-v\|_H^2=\mathscr I_S(u,v).
 \tag{HC2}
\]
Here the hat may be omitted on actual members. This map
extends uniquely to an injective map
\[
 J:\widehat{\mathcal E}^1_{\rm str}(V)\longrightarrow H,
 \qquad
 \|Ju-Jv\|_H^2=\widehat I(u,v),\quad
 \|Ju\|_H^2=\widehat e(u).
 \tag{HC3}
\]
Its image $\mathcal C$ is precisely the $H$-closure of the
actual classical first-energy cone. It is a closed convex
cone. The $L^2(\lambda)$ representative of $Ju$ equals the
actual Borel $u$ almost everywhere for $\lambda$.

The closed-form energy measures agree with the original
completed energy currents, and
\[
       \mathcal E(Ju,\phi)=-\langle D(uS),\phi\rangle
            \quad\bigl(\phi\text{ original smooth compact on }V\bigr).
 \tag{HC4}
\]
Neither this statement nor injectivity identifies the full
energy-polar and BT-capacity-null ideals on arbitrary Borel sets.
\end{theorem}

\noindent\textbf{Status: \texttt{NEW\_PROVED}} for the global
unbounded embedding and its strong-completion closed cone.
The bounded local energy embedding is \texttt{INTERNAL\_INPUT}.

\begin{proof}
First let $u\in\mathcal E^1$ be bounded. For $\epsilon>0$
put $u^{(\epsilon)}=\max(u,-\epsilon)$ and
$w_\epsilon=u-u^{(\epsilon)}$. Its support is contained in
the compact set $\overline{\{u<-\epsilon\}}$.
Choose an original compact cutoff equal to one on a
neighborhood of that set. The bounded local classical
energy embedding applied to the difference of these two
bounded potentials makes $w_\epsilon$ a member of the part
space $H$. Its closed energy is its original bounded BT
energy, without an exterior cutoff term.

Both $u^{(\epsilon)}$ and $u^{(\eta)}$ are actual members
of $\mathcal E^1$. The global first-energy identity for
$q=1$ therefore gives
\[
 \begin{split}
 \|w_\epsilon-w_\eta\|_H^2
   &=\mathscr I_S(u^{(\epsilon)},u^{(\eta)})\\
   &=e_1(u^{(\epsilon)})+e_1(u^{(\eta)})
           -2\mathcal I_S(u^{(\epsilon)},u^{(\eta)})\\
   &\leq e_1(u^{(\epsilon)})+e_1(u^{(\eta)})\longrightarrow0.
 \end{split}
 \tag{HC5}
\]
The final limit is the near-zero energy estimate (IC8).
Since $\lambda(V)<\infty$ and $u$ is bounded,
$w_\epsilon\to u$ in $L^2(\lambda)$ by bounded convergence.
Completeness and coercivity identify its $H$ limit with $u$.
Moreover
\[
 \|w_\epsilon\|_H^2
   =e_1(u)+e_1(u^{(\epsilon)})
                   -2\mathcal I_S(u,u^{(\epsilon)})
                         \longrightarrow e_1(u),
 \tag{HC6}
\]
by the constant-one Holder bound. Thus $\|u\|_H^2=e_1(u)$.
The same approximation for two bounded potentials and
polarization give (HC2) for them.

For an arbitrary actual $u\in\mathcal E^1$, its bounded
canonical truncations $u^{(R)}$ lie in $H$ by the preceding
argument. Equation (EV15) gives
$I(u,u^{(R)})\to0$. The global positive energy-square
inequality in the fixed background $S$ yields
\[
 I(u^{(R)},u^{(L)})
                  \leq2I(u^{(R)},u)+2I(u,u^{(L)})\longrightarrow0.
 \tag{HC7}
\]
Consequently the truncations are Cauchy in $H$.
Their monotone pointwise limit is $u$. An $L^2(\lambda)$
convergent subsequence has a finite almost-everywhere limit,
which therefore equals $u$. This proves actual membership
and identification without presupposing square integrability.
The global mixed-energy convergence of the truncations gives
the norm and pairing identities (HC2).

Now let $(u_j)$ be an actual strongly Cauchy defining row for
a completed $u$. By (HC2) it is Cauchy in $H$; define $Ju$
as its $H$ limit. The datum $\lambda|_V=D(hS)|_V$ is itself
the measure of a normalized actual degree-one tester on $V$.
Thus $\lambda(E)\leq C_S(E;V)$ on Borel sets. The defining
local outer-capacity convergence identifies the local
convergence in $\lambda$ measure of each bounded truncation.
Comparison with the $L^2$ limit, or its almost-everywhere
subsequence, shows that the actual Borel $u$ is the
$\lambda$-almost-everywhere representative of $Ju$.

Independence of the defining row follows from the intrinsic
global mixed-energy limits (IC13); they also give (HC3).
The separation theorem for $\widehat I$ proves injectivity
on relative classes. Addition and nonnegative scaling are
preserved, so the image is a convex cone. Actual strong
density and completeness (IC21) show that it is exactly
the closed $H$-closure asserted in the statement.

For actual bounded approximations the local energy measure
identity is the frozen bounded one. Strong $H$ convergence
passes these energy measures in global total variation by
measure Cauchy--Schwarz. The already constructed original
completed energy currents have the same local test limits,
so they agree with the closed-form measures. Finally, the
core Green identity and coefficient domination (AE10) pass
to every $H$ limit. Together with $u_jS\to uS$ on original
tests, this proves (HC4). All identifications are made in
the original test category before any special representation
of the background or its support is considered.
\end{proof}

\subsection{The precise image criterion and its projection}

For a finite signed Radon measure $\mu$ on $F$, call the
following a finite $H$-dual condition:
\[
 \left|\int_F\phi\,d\mu\right|
             \leq C\|\phi\|_H
       \quad\bigl(\phi\text{ in the original compact smooth core}\bigr).
 \tag{HC8}
\]
It defines a unique continuous functional $\ell_\mu\in H^*$.
The subsequent statements also apply to an abstract functional
$\ell\in H^*$, but the equation with a measure uses (HC8).
Pairings with unbounded $H$ vectors mean this continuous
functional. When a measure integral interpretation is justified,
it uses the canonical energy-quasicontinuous representative.
It is not automatically the integral of an arbitrary literal
BT-capacity Borel representative. In particular no finite-energy
bound for $|\mu|$ is silently inferred from the signed bound
(HC8).

Let $y_\mu\in H$ be defined by
\[
                  \mathcal E(y_\mu,\phi)=-\ell_\mu(\phi)
                                   \qquad(\phi\in H),
 \tag{HC9}
\]
and define the polar cone
\[
 \mathcal C^\circ=
   \{\phi\in H:\mathcal E(v,\phi)\leq0
                       \text{ for every actual }v\in\mathcal E^1\}.
 \tag{HC10}
\]
By continuity this is also the polar of the whole closed cone.

\begin{theorem}[Exact finite-dual source criterion]
\label{thm:r11-strong-degree-one-source-criterion}
For a datum satisfying (HC8), the following are equivalent:
\begin{enumerate}
\item There is $u\in\widehat{\mathcal E}^1_{\rm str}(V)$
      with $D(uS)=\mu$ on every original compact smooth test.
\item $y_\mu\in\mathcal C$.
\item $\ell_\mu(\phi)\geq0$ for every $\phi\in\mathcal C^\circ$.
\end{enumerate}
The solution, if it exists, is unique in the relative
outer-capacity quotient and equals $J^{-1}y_\mu$.
The third condition quantifies over the specified closed
energy space; it is not replaced by an unproved smaller
family of polar original smooth tests.
\end{theorem}

\noindent\textbf{Status: \texttt{DIRECT\_COROLLARY}} of the new
closed-cone identification, the frozen Riesz construction, and
the elementary Hilbert separation theorem.

\begin{proof}
By (HC4), an original current solution gives
$\mathcal E(Ju,\phi)=-\ell_\mu(\phi)$ on the original core.
Density extends the equality to $H$, and coercivity gives
$Ju=y_\mu$. Conversely membership in $\mathcal C$ supplies
its unique completed preimage and (HC4) proves the full
original current equation. A closed convex cone equals its
bipolar in a Hilbert space. Since
$\mathcal E(y_\mu,\phi)=-\ell_\mu(\phi)$, the bipolar
inequality is exactly the third condition with its printed
sign. Uniqueness follows from (HC3).
\end{proof}

\begin{theorem}[The constrained source projection and its exact defect]
\label{thm:r11-strong-cone-source-projection}
For every $\ell\in H^*$, let $y$ be its signed Riesz vector
as in (HC9). There is a unique completed relative class
$u_\ell$ maximizing
\[
                  -\tfrac12\widehat e(u)-\ell(Ju)
                    \quad\bigl(u\in\widehat{\mathcal E}^1_{\rm str}\bigr).
 \tag{HC11}
\]
It is characterized by $Ju_\ell=P_{\mathcal C}y$ and the
projection variational inequality. For two data,
\[
            \widehat I(u_\ell,u_{\ell'})
                           \leq\|\ell-\ell'\|_{H^*}^2.
 \tag{HC12}
\]
Let $T:H\to H^*$ be $Tz(\phi)=-\mathcal E(z,\phi)$.
The closed source cone is $T(\mathcal C)$, and
\[
 \begin{split}
 \operatorname{dist}_{H^*}(\ell,T(\mathcal C))
   &=\|y-P_{\mathcal C}y\|_H\\
   &=\sup_{\substack{\phi\in\mathcal C^\circ\\\|\phi\|_H\leq1}}
                                      \bigl(-\ell(\phi)\bigr).
 \end{split}
 \tag{HC13}
\]
This is a source-functional projection, not a claim that every
projected source has finite total Radon mass up to the boundary.
The constant one in (HC12) is optimal whenever the admissible
actual class contains a finite-mass member of positive first
energy, by taking its source and the zero source.
\end{theorem}

\noindent\textbf{Status: \texttt{DIRECT\_COROLLARY}}. The Hilbert
projection mechanism is frozen; its new domain is the unbounded
closed cone proved above.

\begin{proof}
For $z=Ju\in\mathcal C$, the objective is
$\|y\|_H^2/2-\|z-y\|_H^2/2$. The Hilbert projection theorem
therefore proves existence, uniqueness and its variational
inequality. Testing the inequalities for two projections
against one another gives
\[
 \|P_{\mathcal C}y-P_{\mathcal C}y'\|_H^2
    \leq\mathcal E(P_{\mathcal C}y-P_{\mathcal C}y',y-y')
    \leq\|P_{\mathcal C}y-P_{\mathcal C}y'\|_H\|y-y'\|_H.
\]
This is (HC12), since $T$ is an isometry onto $H^*$.
If $r=y-P_{\mathcal C}y$, the cone projection inequality
gives $r\in\mathcal C^\circ$ and
$\mathcal E(r,P_{\mathcal C}y)=0$. Hence for every polar
unit vector $\phi$,
$-\ell(\phi)=\mathcal E(y,\phi)\leq\|r\|_H$;
equality is attained at $\phi=r/\|r\|_H$ when $r\ne0$,
and the zero case is immediate. This proves (HC13).
For an actual member $a$ with finite positive energy, take
$\ell=Ta$, $\ell'=0$. Their projections are $a,0$ and
$\widehat I(a,0)=e_1(a)=\|Ta\|_{H^*}^2$, proving sharpness
under the stated measure-domain qualification.
\end{proof}

\begin{corollary}[The cone condition is a real restriction]
\label{cor:r11-strong-Hilbert-cone-proper}
For the precise moment-Laplacian background (DO6), $n\geq2$,
the nonzero smooth compact datum (DO9) has finite $H$-dual
size but its Riesz solution lies outside $\mathcal C$.
Its defect in (HC13) is strictly positive. Thus the exact
source criterion is compatible with the existence of a
relative linear energy solution and the failure of a solution
in the strongly completed classical first-energy cone.
\end{corollary}

\noindent\textbf{Status: \texttt{DIRECT\_COROLLARY}} of the new
strong-completion obstruction and the preceding source criterion.

\begin{proof}
The compact smooth datum is an original smooth coefficient
of $S$ on a compact stage. The faithful domination (AE7)--(AE10)
bounds its density by a constant times $\lambda$ there.
Finite reference mass and coercivity therefore give (HC8).
The source obstruction (SO5) excludes its Riesz vector from
$\mathcal C$. Since this cone is closed, its distance is
strictly positive, and (HC13) identifies that distance.
The Riesz solution itself still exists in $H$ by coercivity;
the assertion concerns exactly its membership in the printed
potential cone, not failure of the larger linear theory.
\end{proof}

\begin{corollary}[The original smooth core is a difference of finite-mass members]
\label{cor:r11-finite-mass-core-differences}
Under (TC1)--(TC2), the potential construction in the next
sentence holds for every remaining degree $q\geq1$.
Every real original smooth compact scalar $\phi$ on $V$ is
an exact actual difference $\phi=a-b$, where $a,b$ are bounded
classical members of every positive power-energy class, have
finite total top mass, and agree with the same $A\rho$ outside
one compact subset of $V$. In particular, in the degree-one
setting of this section,
\[
                    A_c^{0,0}(V)\subset\mathcal C-\mathcal C,
                    \qquad\overline{\mathcal C-\mathcal C}^{\,H}=H.
 \tag{HC14}
\]
This dense span does not imply that a fixed member of
$\mathcal C$ admits both signs of every original smooth
variation.
\end{corollary}

\noindent\textbf{Status: \texttt{DIRECT\_COROLLARY}} of the
analytic compensators and finite-mass barrier construction.
The frozen R6 \path{lem:r6-smooth-direction-dsh} already
proves the local classical-difference assertion, and the
explicit global algebraic compensator is in frozen R4
\path{relative_capacity_next.tex}. The additional conclusion
here is the printed original analytic finite-mass/common-collar
realization and its application to the newly identified cone.

\begin{proof}
Extend $\phi$ by zero across the complement of its compact
support. On finitely many original moment neighborhoods of
that support write $\phi$ as a smooth function of unit
logarithms. The exact analytic replacement (AE2) writes each
unit norm as the quotient of norms of functions on the one
original affinoid. Use the convex-flat functions of these
numerator and denominator logarithms from (AE5).
As in the frozen R4 compensator proof, take a sufficiently
large finite positive sum $g$ of them. On each smaller moment
stage the selected summands have Hessian dominating the
negative of the Hessian of $\phi$. Thus $g+\phi$ is classical
there by convex composition of harmonic unit logarithms;
the other summands are already classical. Off the support
it equals $g$. Ordinary-open gluing proves that both $g$
and $g+\phi$ are global smooth classical functions on the
required neighborhood. This is an explicit classical
construction, not an inference from positivity against one
background current.

Choose $C$ so large that $g-C$ and $g+\phi-C$ are at most
$-1$ on $V$. Choose $A>0$ such that $A\rho$ is below both
on a neighborhood of $\operatorname{supp}\phi$. Put
\[
             a=\max(g+\phi-C,A\rho),\qquad
             b=\max(g-C,A\rho).
 \tag{HC15}
\]
Near the support their difference is $\phi$; outside it the
two branches coincide. Hence the equality holds at every
actual point of $V$. They equal $A\rho$ on the ordinary
collar $\{A\rho>-1\}$, whose complement is compact.
Both are bounded, negative and have compact strict negative
sublevels. Comparing each with $cA\rho$, $c>1$, and then
letting $c\downarrow1$, gives total mass at most
$A\mu_\rho(V)$ in the present degree-one case. For arbitrary
remaining degree the same construction gives the bound
$A^q\mu_\rho(V)$. Boundedness and finite total mass give
every positive power energy. Thus $a,b$ have all the printed
properties.

The original smooth core is dense in $H$ by its definition,
so (HC14) follows. Although a smooth direction is a
difference of two cone members, subtracting one of them
from a specified maximizer need not preserve the cone.
Accordingly the two one-sided Euler inequalities cannot
be subtracted to manufacture an unrestricted equation.
The proper-cone example above makes this distinction
concrete. No approximation of background currents, and
in particular no $W_{\rm loc}$ density assertion, is involved.
\end{proof}

% END REVIEWED MODULE research_r11/comparison/first_energy_hilbert_source_cone.tex


% BEGIN REVIEWED MODULE research_r11/dirichlet/strong_dominated_operator_image.tex
\section{Closedness of the dominated nonlinear source image}
\label{sec:r11-dominated-operator-image}

Retain the full arbitrary-background degree $q=n-p\geq1$,
the strict analytic box and extension hypotheses (TC1)--(TC2),
and the strong completion of the preceding variational theorem.
Put $d=q+1$.
Fix one actual $\psi\in\mathcal E^1_{S,\partial}(V)$ and
$A\geq0$.  The data may have infinite total mass near the
boundary, but are positive Radon measures on $V$.
Use exactly the absolute dual-energy distance $\delta_1$ in
(VF11), with the full actual classical first-energy tester cone.

\begin{corollary}[A precise approximation criterion for dominated sources]
\label{cor:r11-dominated-source-image-closed}
Suppose
\[
 \begin{gathered}
 0\leq\mu_j,\mu\leq A\mu_\psi,
       \qquad\delta_1(\mu_j,\mu)\longrightarrow0,\\
 u_j\in\widehat{\mathcal E}^1_{\rm str}(V),
                         \qquad\widehat\mu_{u_j}=\mu_j.
 \end{gathered}
 \tag{OI1}
\]
Then $u_j$ converges strongly to the relative class
$u=\Pi_{\rm str}(\mu)$, and
\[
                              \widehat\mu_u=\mu
 \tag{OI2}
\]
on every original compact smooth scalar test and every actual
moment refinement.  Thus the source image of the strong
completion inside $\{0\leq\nu\leq A\mu_\psi\}$ is closed
for $\delta_1$.  This is a class-valued approximation criterion;
it does not assert that every dominated datum lies in that image.
\end{corollary}

\noindent\textbf{Status: \texttt{DIRECT\_COROLLARY}} of
the dominated-data variational retraction and stability, the
strong completed current continuity, and the exact original
smooth core decomposition (HC14)--(HC15).  The latter potential
construction is explicitly valid for every $q\geq1$; only its
separate Hilbert dense-span consequence has $q=1$.

\begin{proof}
The dominated operator-image retraction gives
$u_j=\Pi_{\rm str}(\mu_j)$ as relative classes.  The data
stability (VF12), with the uniform energy bound (VF8), proves
$\widehat I(u_j,\Pi_{\rm str}(\mu))\to0$.  Hence the full
strong completed calculus gives
$\widehat\mu_{u_j}\to\widehat\mu_u$ on every original
compact smooth scalar test.

It remains to identify that current limit with the prescribed
measure, rather than only with its potential-dual pairings.
Let $\phi$ be an arbitrary real original smooth compact scalar.
The all-degree classical construction (HC15) writes it as the
exact actual difference $\phi=a-b$, where $a,b$ are bounded
finite-mass members of every positive power class and have a
common actual boundary collar.  In particular $a,b\in\mathcal E^1$.
Normalization in (VF11), with the zero-energy case supplied by
domination and Holder, gives
\[
 \left|\int_V\phi\,d(\mu_j-\mu)\right|
 \leq\bigl(e_1(a)^{1/d}+e_1(b)^{1/d}\bigr)
                              \delta_1(\mu_j,\mu)\longrightarrow0.
 \tag{OI3}
\]
Each integral against the absolute variation is finite by the
printed domination, even if the unweighted variation has
infinite total mass.  Consequently the prescribed source
measures converge on the complete original test sheaf.  Their
other limit is $\widehat\mu_u$, so (OI2) follows there.
Original smooth scalar density and Radon regularity identify
the same original measures.  The argument began with an
arbitrary original test and therefore includes every actual
moment refinement.  No approximation of background currents
or replacement by a special tester subfamily is used.
\end{proof}

The compact moment-Laplacian datum in (VF13) belongs to the
same type of dominated interval but is excluded from the
entire strong completed source image by (SO5).  Therefore it
cannot admit approximations satisfying all of (OI1).  This
retains the explicit source obstruction while identifying
exactly what a valid approximation argument would have to
provide.  It does not infer the missing original current
equation from the cone Euler inequality alone.

% END REVIEWED MODULE research_r11/dirichlet/strong_dominated_operator_image.tex


% BEGIN REVIEWED MODULE research_r11/dirichlet/strict_relative_polar_ideals.tex
\section{Applying the frozen polar-ideal obstruction to the new solution class}
\label{sec:r11-strict-polar-ideals}

This section is a provenance integration and direct application,
not a new proof of the separation of two capacities.
The entire separation theorem, its full-original Gauss-trace
contraction, the uniform classical-tester estimate, and the
Cantor--Newton construction are already frozen in R6,
\path{research_r6/sections/capacity_energy_null_gap.tex}.
The exact source SHA256 is
{\small\nolinkurl{a84378b70e36eaee1c83253630cae036028869d2f8bdd42984dc3c09ab8cee04}}.
The nested R6 archive was checked against the R7 manifest,
and its source bytes agree with the source read here.
Its earlier proofs are not reproduced or modified.

\paragraph{The exact frozen input.}
Take $n\geq3$, $Q=(-2,2)^n$,
$V=t^{-1}(Q)\subset(\mathbf G_{m,K}^n)^{\mathrm{an}}$,
with $t_i=\log|T_i|$, and set
\[
 \beta=\sum_i d't_i\wedge d''t_i,\qquad
 S=\frac{\beta^{n-1}}{(n-1)!}\wedge[V].
 \tag{PI1}
\]
The background extends to a larger polyannulus.
For the original joint Gauss section $\sigma$, R6 constructs
a compact $A\Subset Q$ and a probability $\vartheta$ on $A$
such that, with $K=\sigma(A)$ and $\mu=\sigma_*\vartheta$,
\[
\begin{gathered}
 C_S^*(K;V)=0,\qquad \mu(K)=1,\\
 \left|\int f\,d\mu\right|
       \leq C\,\mathcal E_S(f,f)^{1/2}
       \quad\text{for every original }f\in A_c^{0,0}(V;\mathbf R).
\end{gathered}
 \tag{PI2}
\]
The exact frozen labels are
\path{lem:r6-energy-gauss-restriction},
\path{lem:r6-gauss-lipschitz-capacity},
\texttt{lem:r6-cantor-newton-set}, and
\path{thm:r6-energy-BT-null-gap}.
In particular the source bound in (PI2) is for the entire
original core, not just for toric functions.

\noindent\textbf{Status: \texttt{INTERNAL\_INPUT}}.
The failure of a converse energy/BT polar implication remains
a frozen \texttt{FALSE} assertion, not an R11 discovery.

\begin{corollary}[The new renormalized conclusions apply beyond BT-diffuse data]
\label{cor:r11-energy-source-beyond-BT-diffuse}
For the frozen datum in (PI2), realize the part form on $V$
using a faithful reference measure on a slightly larger
ordinary polyannulus.  Then $\mu$ has finite energy integral,
and the original equation
\[
                         D(uS)=\mu
 \tag{PI3}
\]
has its unique zero-boundary energy solution $u=-G\mu\in H$.
It belongs to the new signed renormalized class, with the
exact-level, coarea, and global hard-curvature variation
conclusions (SD7)--(SD16).
The datum is energy diffuse but is not BT diffuse.
Moreover $\mu$ cannot equal $\mu_w$ for any
$w\in\mathcal E^\sigma_{S,\partial}(V)$ and any $\sigma>0$.
\end{corollary}

\noindent\textbf{Status: \texttt{DIRECT\_COROLLARY}}.
Finite-energy Green solvability itself is also frozen;
the new content of this application is its membership in
the R11 renormalized theorem and its comparison with the
new unbounded power classes.

\begin{proof}
Apply (PI2) to smooth normal approximations of $|f|$.
Their energy is bounded by that of $f$, so the positive
measure $\mu$ has finite energy integral for the full
original core.  The original Green construction and its
canonical representative therefore give $u=-G\mu\in H$
and (PI3), independently of the chosen faithful reference.
The same bound gives energy diffuseness, by the compact-core
capacity argument.  All hypotheses of the signed diffuse
renormalized theorem are now satisfied.
Its original-current, exact-level, coarea, uniqueness and
curvature-variation conclusions apply to this already
specified datum.

The strict polyannulus satisfies (TC1)--(TC2), and the
background extends across its closure.  Every power-class
top measure annihilates the raw-capacity-null Borel set
$K$, by (TC9), whereas $\mu(K)=1$.  This proves the
nonrealization in every positive power class.
Equation (PI2) directly proves that the datum is not
BT diffuse.  No converse comparison of the two capacities
and no ambient classical property of the energy solution
has been assumed.
\end{proof}

\paragraph{Preserved failed novelty attempt.}
An independent R11 derivation initially rediscovered
the frozen R6 separation and its mollifier argument.
The complete draft is preserved in
\path{research_r11/failed_attempts/rediscovered_R6_polar_gap.tex},
with a hash-bound provenance correction.
It is excluded from all R11 new-proof counts and from
this manuscript's proof claims.

% END REVIEWED MODULE research_r11/dirichlet/strict_relative_polar_ideals.tex


% BEGIN REVIEWED MODULE research_r11/capacity/fixed_tail_energy_obstruction.tex
\section{A fixed linear energy tail does not control the full nonlinear capacity}
\label{sec:r11-fixed-tail-energy-obstruction}

The frozen R6 null-ideal counterexample compares the degree-one
BT capacity of a background with the energy capacity of that
\emph{same} background. It gives a BT-null compact set of
positive energy capacity. The assertion tested here is different:
can one fix a smooth $(q-1,q-1)$ tail over a higher-degree
BT problem and use that single linear energy capacity to control
the capacity with all $q$ slots free? The following opposite
null direction disproves this proposed reduction.

\begin{proposition}[An ambient BT atom at an auxiliary energy-polar point]
\label{prop:r11-fixed-tail-energy-obstruction}
For every $n\geq2$, over the unchanged original field, let
\[
 \begin{split}
 V&=t^{-1}((-2,2)^n)\subset(\mathbf G_{m,K}^n)^{\rm an},
       \qquad t_i=\log|T_i|,\\
 S_0&=[V],\qquad
 \beta=\sum_{i=1}^n d't_i\wedge d''t_i,\qquad
 Q=\frac{\beta^{n-1}}{(n-1)!}\wedge S_0.
 \end{split}
 \tag{FT1}
\]
The backgrounds extend to a larger ordinary polyannulus inside
one original affinoid. Let $x=\sigma(0)$ be the original joint
Gauss point and let $c_{Q,V}$ be the faithful part energy capacity
of the original degree-one current $Q$, using the complete
original smooth core. Then
\[
             c_{Q,V}(\{x\})=0,
             \qquad
             C_{S_0}(\{x\};V)
                      \geq\frac{n!}{(2n)^n}>0.
 \tag{FT2}
\]
The first equality holds also for its $1$-capacity, and the
second inequality also holds for the ordinary-open outer BT
capacity. In particular there is no modulus $\omega$ with
$\omega(0)=0$ for which
\[
                  C_{S_0}^*(A;V)\leq\omega(c_{Q,V}(A))
 \tag{FT3}
\]
holds for all compact $A\Subset V$.
\end{proposition}

\noindent\textbf{Status: \texttt{NEW\_PROVED}} as a
counterexample to the fixed-tail reduction (FT3), which is
\texttt{FALSE}. The original target background in the
capacity on the left is $S_0$ throughout; $Q$ is only the
explicit proposed auxiliary energy tail.

\begin{proof}
The original torus calibration, frozen with its complete
test scope, gives
\[
 \mathcal E_Q(\phi(t),\phi(t))
                         =\int_{\mathbf R^n}|\nabla\phi|^2\,dy
       \quad\bigl(\phi\in C_c^\infty((-2,2)^n)\bigr).
 \tag{FT4}
\]
The faithful form construction supplies
$\mathcal E_1\leq3\mathcal E$ on the original core by
coercivity. Thus a sequence with the energy in (FT4) tending
to zero also has full faithful $1$-form norm tending to zero;
no assertion about the support of the reference measure on
the Gauss section is needed.

First suppose $n\geq3$. Choose a real smooth function
$\psi$ supported in the unit ball, equal to two near zero,
and put $\phi_\epsilon(y)=\psi(y/\epsilon)$.
Its pullback is an original smooth compact core function,
equal to two on an ordinary neighborhood of $x$, and
\[
 \mathcal E_Q(\phi_\epsilon(t),\phi_\epsilon(t))
      =\epsilon^{n-2}\int|\nabla\psi|^2\longrightarrow0.
 \tag{FT5}
\]
Here the pullback support is compact because tropicalization
of the original split torus is proper over compact real boxes.

For $n=2$, choose $0<R<1$ and a smooth real function
$\theta$ equal to zero on $(-\infty,0]$, equal to two on
$[1,\infty)$, and with bounded derivative. For $L>1$ put
\[
 \phi_L(y)=\theta\left(\frac{\log(R/|y|)}L\right)
       \quad(y\ne0),\qquad \phi_L(0)=2.
 \tag{FT6}
\]
It is smooth: it is constant near zero and its scalar profile
is flat at both ends of the transition. Its support is contained
in the closed ball of radius $R$, and it equals two on the
ball of radius $Re^{-L}$. Polar-coordinate integration gives
the exact energy
\[
       \int_{\mathbf R^2}|\nabla\phi_L|^2\,dy
                      =\frac{2\pi}{L}
                                 \int_0^1|\theta'(s)|^2\,ds
                      \longrightarrow0.
 \tag{FT7}
\]
Thus the same original compact-core capacity argument applies
also in the two-dimensional case.

The point $x$ lies in the support of $Q$: for every ordinary
neighborhood of $x$, continuity of the Gauss section and the
calibration $\beta\wedge Q=n\sigma_*(dy)$ provide a
nonzero positive original coefficient pairing in that
neighborhood. The admissible functions in (FT5) or (FT6)
are greater than one near this actual point. The regular
compact-core capacity formula and $\mathcal E_1\leq3\mathcal E$
therefore give $c_{Q,V}(\{x\})=0$ and zero $1$-capacity.
The full original core is larger than this explicit test
subfamily, so this upper bound remains valid for that full
core and every refinement. No restriction to toric energy
testers is made in the definition.

For the target nonlinear capacity, use the actual normalized
bounded continuous classical tester
\[
                 h=\frac1{2n}\sum_{i=1}^n\max(t_i,0),
                          \qquad0\leq h\leq1\text{ on }V.
\]
The frozen R3 Gauss-box atom identity
\path{eq:r3-gauss-box-test-atom} gives the exact original
Radon measure
\[
                  (Dh)^n\wedge S_0
                           =\frac{n!}{(2n)^n}\delta_x.
 \tag{FT8}
\]
That identity was proved by smooth one-variable hinge
regularization and the full-original torus calibration:
repeated coordinate squares vanish and the $n!$ distinct
coordinate terms each have unit atom. Its normalization
uses the original logarithmic coordinates and every original
scalar test, including any extra actual moments.
Consequently (FT8) is an admissible measure in the complete
normalized capacity defining $C_{S_0}$. It proves the positive
lower bound in (FT2), and monotonicity of outer regularization
gives the outer version. Taking $A=\{x\}$ disproves (FT3).
\end{proof}

\paragraph{Provenance and failed route.}
The prior audit for this statement was limited to the relevant
frozen capacity files. R6
\path{capacity_energy_null_gap.tex} proves the opposite
null gap for one degree-one background; its all-original
double-mollifier trace argument is retained as frozen input
and is not counted as a new R11 result. R6
\path{capacity_energy_boundary.tex} proves qualified
outer regularity from an energy-polar relative boundary,
again for one remaining slot. Neither statement is being
extended silently to the $n$-slot capacity in (FT2).

The rejected higher-degree route would use the regularity of
one faithful linear energy form to settle compact or Borel
regularity of the entire nonlinear tester family. Equations
(FT2)--(FT3) show that this domination cannot hold even for
the displayed smooth elliptic real moment tail. This is a
\texttt{FAILED\_ATTEMPT} at that reduction, not a
counterexample to compact continuity or ordinary-open outer
regularity of the full BT capacity itself. Estimates uniform
over the moving positive tails, or a different relative
completion mechanism, are still required.

% END REVIEWED MODULE research_r11/capacity/fixed_tail_energy_obstruction.tex


% BEGIN REVIEWED MODULE research_r11/remaining_gates_and_status.tex
\section{Remaining gates, false formulations, and the exact delta}
\label{sec:r11-remaining-gates}

The status of a universal assertion is not changed by
proving a qualified theorem with a specific boundary
class.  Conversely, a counterexample to an unqualified
assertion should not be recorded merely as an open
problem.  The following ledger distinguishes the two.

\subsection{The capacity and unbounded-calculus boundary}

\paragraph{Proved unbounded sector.}
The original arbitrary-background construction now includes
coefficient-finite non-polar closedness on ordinary germs;
compact-sublevel outer quasicontinuity; the exact no-loss
capacity-tail cone; its comparison, domination and
relative uniqueness; all positive power-energy classes;
all-degree curvature recursion and graded first-order
and energy currents; and their precise capacity and
strong-energy continuity statements.
On one original affinoid neighborhood, the new analytic
tester extension compares the full local capacities
on compact cores.
On strict analytic boxes with an extending background,
the finite-mass boundary tester core computes raw
capacity, and the full global $E^1$ Green and energy
geometry hold.
These are \texttt{NEW\_PROVED} results with the
individual assumptions printed in their modules.

\paragraph{General raw and outer capacities: \texttt{OPEN}.}
An unrestricted proof of
$C_S(K;V)=C_S^*(K;V)$ for every compact $K$,
or of raw/outer equality on every Borel set, is not
provided.  The elementary continuity of the outer
regularization on decreasing compact sets was already
frozen and is not the missing assertion.
The remaining weak-limit gate is concrete: for a
locally weakly convergent family of full normalized
original tester measures, one needs control on a
specified compact or Borel set by the raw tester
capacity, rather than merely the outer-capacity
bound inherited on opens.
No infimum over neighborhoods has been interchanged
with a supremum over testers.

The exact raw equality for the finite-mass core
(TC6) concerns two tester families on the same
domain.  The domain comparisons (AL8)--(AL9)
compare two capacities of the same kind.
The strict-gap equality (BC17) and the scalar
coarea estimates concern specified sets.
None of these statements proves the unrestricted
raw/outer assertion.

\paragraph{A single auxiliary energy tail: \texttt{FALSE}.}
The new fixed-tail example (FT2) compares the full
nonlinear capacity of one background with the
linear energy capacity of a particular auxiliary
tail: its point has positive nonlinear capacity
and zero auxiliary energy capacity. The frozen R6
example instead compares the BT and energy capacities
of the same degree-one background and gives a BT-null
compact set of positive energy capacity.
Together they explain why neither identifying
the two polar ideals nor controlling an entire
nonlinear tester family by one fixed linear
form resolves the raw-capacity gate.
They are not counterexamples to raw/outer
regularity itself.

\paragraph{Unrestricted recursion: \texttt{FALSE}.}
The logarithmic example (CR2) has a locally
coefficient-integrable multiplier
$uN_{q-1}$, yet
$D(uN_{q-1})=\delta_o$ while the non-polar
top product is zero.
No-loss and the printed integrability
condition repair that passage; every positive
power class supplies sufficient integrability.
The discarded polar residue cannot be removed
by redefining the measure on its pole after
differentiation.

\paragraph{Unrestricted endpoint continuity: \texttt{FALSE}.}
Uniform $E^\sigma$ energy and capacity convergence
alone do not give critical weighted continuity.
The shell family in (PS7) has a nonzero limiting
weighted atom.
The proved repairs are fixed-tuple tails,
uniform capacity-tail integrability, actual
decreasing lower barriers, or the strong
$E^1$ discrepancy topology with its derived
uniform tails.

\subsection{Poles, envelopes, and admissible representatives}

\paragraph{Complete compact-sublevel pole sets.}
The logarithmic construction gives one actual
representative of each such pole set in every
positive power class, simultaneously for all
allowed extending backgrounds.
This is \texttt{NEW\_PROVED}.
The ambient R2 pole/negligible statements on
their specified exhaustion domains remain
\texttt{INTERNAL\_INPUT}.
The unrestricted claim that a relative
capacity-null set must be an ambient or
compact-sublevel complete pole set is
\texttt{FALSE} without nondegeneracy and
domain qualifications.
An appropriate converse with explicit
nondegeneracy, representation and domain
hypotheses remains \texttt{OPEN}.

\paragraph{Ordinary upper regularization: \texttt{FALSE}.}
The all-power boundary-core example (PO1)--(PO7)
retains a common actual lower barrier, finite
mass, unit energies and even an actual
classical solution of the prescribed source.
Its ordinary upper-regularized Perron
supremum nevertheless loses the source.
This is a strengthening of the frozen
proper-support example, not a new proof
of the original bounded failure.

\paragraph{Relative completed envelopes: \texttt{NEW\_PROVED}.}
The barrier completion constructs original
closed currents, derivatives and allowed
energy currents before asserting their
unbounded recursive identities.
Its unbounded relative Perron envelope
preserves Borel constraints and a common
source lower bound.
The completed barrier-constrained
variational problem has a unique relative
maximizer, a precise Euler inequality,
optimal weighted-data stability, and a
retraction on its finite-mass operator image.
These are actual existence statements in
an explicitly enlarged Borel potential
class.  They do not claim that every
completed potential has an ambient
classical representative.

\paragraph{Full relative strong completion: \texttt{NEW\_PROVED}.}
The first-energy strong space removes the prescribed actual
lower barrier. Uniform deep tails and global weighted tightness
give its complete actual-Borel/current realization. Its full
comparison, domination, strict locality, exact truncation
carriers and finite-mass approximation retain the encoded
strong boundary class. The intrinsic capacity-layer formula
characterizes each power moment of an already constructed
completed member; its strict-gap raw/outer equality still
does not give unrestricted capacity regularity.

The exact intrinsic characterization now requires same-domain
bounded capacitary admissibility of every truncation and a
finite outer first-energy capacity integral. This removes
the need to prescribe a strongly Cauchy defining row, while
retaining a genuine analytic admissibility gate. It also
gives the precise monotone closure theorem. Every nonempty
strong completed family has a relative Perron envelope by
energy minimization over finite maxima, preserving source
lower bounds and Borel upper constraints.

For positive data dominated by a constant times one actual
first-energy top measure, the full strong cone has a unique
variational maximizer, sharp energy bounds and quantitative
data stability. Total source mass at the boundary may be
infinite. Literal datum integrals are proved continuous from
the completed tester estimate before maximization is used.
The source equation remains distinct even in this explicit
dominated class.
Within one fixed dominated interval, the actual completed
operator image is closed for the proved absolute
energy-dual distance. This is an approximation criterion:
it requires approximants already in that operator image
and does not assert that the image fills the interval.

\paragraph{Maximality implies the source equation: \texttt{FALSE}.}
The completed maximizer in (BV13), and even the
unconstrained strong-cone maximizer in (VF13), can
fail the prescribed current equation.
The R10 compact-source rigidity extends through
the new degree-zero current convergence to the
entire strong completion: a uniform first-energy
bound controls local convex Gauss traces by (SO3).
Thus a classical representative is not
the only missing issue: the constrained
admissible variation cone is itself a
substantive equation gate.
The exact proved condition is (BV8).
The current equation follows only under
additional valid variations, as explicitly
stated in the actual-class conditional
corollary.

\subsection{Dirichlet and obstacle problems}

\paragraph{The complete relative linear measure-data problem.}
For an arbitrary original degree-$(n-1,n-1)$
background on the printed affinoid part domain,
every finite signed energy-diffuse Radon datum
has a unique relative renormalized zero-boundary
solution.  Its original current equation,
exact level carriers, bounded-value coarea,
hard-curvature global total-variation convergence,
and reference independence are
\texttt{NEW\_PROVED}.
The signed finite-diffuse obstacle problem
likewise has a unique least majorant and
an exact positive contact reaction, under
the stated finite-curvature obstacle
assumption.
The reaction bound and source stability
are proved before and during the limiting
construction; contact is not inferred
from an uncontrolled weak restriction.

The word zero boundary here denotes the
specified part-space clipping and radiation
conditions.  It does not denote arbitrary
ordinary pointwise boundary values.
Unweighted energy of a measure-data
solution is not always Radon locally;
the explicit positive-data example
proves this assertion \texttt{FALSE}.
The chain and coarea theorems use only
the finite bounded-value energy measures
they actually require.

\paragraph{The nonlinear equation remains appropriately qualified.}
For $q>1$, the new calculus identifies the
derivative of every actual admissible
power-class segment and solves its
relative linearized signed diffuse
problem on a small part domain.
This is a \texttt{DIRECT\_COROLLARY};
it does not realize an arbitrary energy
tangent as a classical variation.

An unrestricted assertion that every
capacity-controlled or finite-dual-energy
datum has an ambient classical solution
for every original $S$ is already
\texttt{FALSE} in the precise remaining-degree-one
moment example, now also in its unbounded
power classes and full strong completion.
The genuinely open nonlinear program is
therefore to specify a valid admissible
data class, potential class and boundary
condition, and prove existence, uniqueness,
comparison and suitable regularity there.
No theorem here removes those three
specifications.

\paragraph{A precise degree-one image criterion.}
For $q=1$ the strong first-energy completion embeds
isometrically as the closed cone $\mathcal C$ in the
original part Hilbert space. The full finite-$H$-dual
source equation is solvable in that cone exactly when
its Riesz vector lies in $\mathcal C$, equivalently
when the stated polar-cone inequalities hold.
The source-functional projection, its sharp stability,
and its exact polar defect are direct Hilbert consequences.
The original smooth core is a difference of actual
finite-mass common-collar potentials, so the cone has
dense linear span. Its properness shows why dense span
does not provide feasible two-sided directions at a
specified maximizer.
This is an exact criterion in the enlarged relative
class, not a constructive classification of all ambient
classical data or an identification of the BT and energy
polar ideals. Pairings with general signed $H$-dual data
remain continuous energy functionals unless their literal
measure interpretation is separately justified.

\subsection{The analytic relation quotient is still separate}

\paragraph{Full GQ: \texttt{OPEN}.}
For every compatible finite actual PS/PL
presentation $e$, the target is still
\[
 e=0\text{ as an actual analytic delta-form}
           \quad\Longrightarrow\quad A_S(e)=0
 \tag{OGQ}
\]
for every fixed original arbitrary $S$,
on every ordinary germ, after all actual
moment refinements, and against all original
complementary tests.
The positive-zero kernel, fixed-moment
top-real-degree signed relations, and the
R8--R10 matroidal, fan and actual analytic
sectors remain frozen inputs.
The new quotient of actual power-class
potentials modulo relative capacity
does not imply (OGQ).
No equality on the full GK delta-test
category has been silently added.

\paragraph{Low-real-degree CTO: \texttt{OPEN}.}
Signed constant-tensor order transfer in
$1<b<r$, including the critical
$b=2<r=3$ range, is not completed by
the new ordinary energy and capacity
identities.  They do not provide an
arbitrary positive lift of an actual
signed analytic relation.
The precise R10 conormal and order
sectors remain as frozen, with their
counterexamples to stronger assertions.

\paragraph{Other frozen open gates.}
$W_{\rm loc}$ density and the Gauss-residue
auxiliary-field quotient remain
\texttt{OPEN}.  They are not inputs to
any R11 new proof.
The old internal R/B/A2 results remain
authorized \texttt{INTERNAL\_INPUT},
and completed A2 geometric adelic
nefness remains completed.

\subsection{Verification, provenance and continuation}

Every module is retained as exact source,
with independent review records binding
the reviewed SHA256.
The standalone manuscript is generated
from those modules with a recorded source
map; its layout substitutions do not
change mathematics.
The exact midpoint script is an auxiliary
rational normalization check only.
The infinite-degree proof is the finite
binomial argument in (SM3)--(SM7), not
the finite range checked by the script.

Two provenance corrections are explicit.
The R6 polar-ideal separation was
rediscovered independently and excluded
from new-proof counts; its complete
draft remains in the failed-attempt record.
The bounded top weak maximum inequality
was already frozen in R3; R11 adds its
exact all-degree positive level-current
decomposition and the unbounded extension.
The R3/R4 bounded midpoint and constrained
retraction mechanisms are also credited;
R11 supplies the global unbounded domain
and the new sharp constant.

The minimal next tasks are mathematical:
raw weak-limit domination for the full
normalized tester family; a valid source
or variation criterion beyond the now
explicit false unrestricted one; completed
nonlinear contact/equation mechanisms with
their true boundary class; and the
independent full GQ/low-degree CTO gates.
The completed barrier calculus and
finite signed diffuse contact theorem
can be used as new inputs, with all their
domain and representative assumptions
carried forward.

\paragraph{External source used for the new measure-data extension.}
T.~Klimsiak and A.~Rozkosz,
\emph{On the structure of bounded smooth measures
associated with a quasi-regular Dirichlet form},
arXiv:1410.4927v2, Lemma 2.1:
\url{https://arxiv.org/abs/1410.4927}.
The exact finite-zero-energy nest input and
its hypotheses are stated in the main
Dirichlet proof.  General form-theoretic
facts already used in the frozen faithful
energy theory retain that earlier dependency
record; this round does not add an
unidentified PDE existence theorem.

% END REVIEWED MODULE research_r11/remaining_gates_and_status.tex

\end{document}
